\chapter{Cálculo de volúmenes}

En este capítulo se presentan varios ejemplos de sólidos con sus respectivos gráficos y diferentes maneras de hallar su volumen. En primera instancia se pueden usar coordenadas rectangulares y plantear los seis posibles órdenes para evaluar la integral triple; otra posibilidad es utilizar un cambio de coordenadas, ya sean cilíndricas o esféricas; también parametrizar el sólido, esta última idea se apoya en la fórmula \ref{math2}.
En algunos casos se ha incluido la función de densidad del sólido, motivo por el cual se agrega el cálculo de la masa del sólido.
Además, se presenta la manera de obtener los gráficos y evaluar estas integrales con los procedimientos respectivos utilizando el programa \verb"Mathematica".

\section{Sólidos acotados por planos}
Presentamos algunos sólidos acotados principalmente por planos y superficies cuadráticas.

\subsection*{Una pirámide}
A continuación se muestra la figura del sólido acotado por las figuras de los planos $x=2y$, $z=4-x-2y$
en el primer octante.
\vskip0.2cm
Se debe tener en cuenta que en el plano $xy$ este se proyecta en la recta $x=4-2y$. La intersección de los planos proyectada en el plano $yz$ se determina por la ecuación $z=4-4y$ y en el plano $xz$ esta región se determina por $z=4-2x$.

\noindent
\begin{minipage}{\textwidth}
	\centering
	\captionof{figure}{Intersección entre los planos $x + 2y + z = 4$ \quad y \quad $x = 2y$.}
	\label{fig:interseccion_planos}
	\begin{minipage}{0.36\linewidth}
		\centering
		\vspace{-3mm}
		\includegraphics[width=\linewidth]{Fig/148.pdf}
	\end{minipage}
	\hspace{1em}
\begin{minipage}{0.35\linewidth}
	   \centering
	   \vspace{-3mm}
		\includegraphics[width=\linewidth]{Fig/149.pdf}
	\end{minipage}
	
	\vspace{1ex}
	\fuentepropia
\end{minipage}

\vskip0.2cm
Sobre los planos coordenados se determinan las siguientes regiones.

\vskip0.2cm
\noindent
\begin{minipage}{\textwidth}
	\centering
	\captionof{figure}{Proyección de los planos $x + 2y + z = 4$ \quad y \quad $x = 2y$.}
	\label{fig:proyeccion_planos}
	\begin{minipage}{0.28\linewidth}
		\centering
		\includegraphics[width=\linewidth]{Fig/150.pdf}
	\end{minipage}
	\hspace{1em}
	\begin{minipage}{0.28\linewidth}
		\centering
		\includegraphics[width=\linewidth]{Fig/151.pdf}
	\end{minipage}
	\hspace{1em}
	\begin{minipage}{0.28\linewidth}
		\centering
		\includegraphics[width=\linewidth]{Fig/152.pdf}
	\end{minipage}
	
	\vspace{1ex}
	\fuentepropia
\end{minipage}

\vskip0.2cm
Estos gráficos nos ayudan a establecer los límites de las integrales triples que determinan su volumen.
\begin{multline*}
\ds \int_{0}^{1}\int_{0}^{4-4y}\int_{0}^{2y}dxdzdy +\int_{0}^{1}\int_{4-4y}^{4-2y}\int_{0}^{4-z-2y} dxdzdy\\
+ \int_{1}^{2}\int_{0}^{4-2y}\int_{0}^{4-z-2y}dxdzdy=
\ds \int_{0}^{2}\int_{x/2}^{\frac{4-x}{2}}\int_{0}^{4-x-2y}dzdydx\\
= \ds \int_{0}^{2}\int_{0}^{4-2x}\int_{x/2}^{\frac{4-x-z}{2}}dydzdx= \int_{0}^{4}\int_{0}^{\frac{4-z}{2}}\int_{x/2}^{\frac{4-x-z}{2}}dydxdz\\
= \ds \int_{0}^{1}\int_{0}^{2y}\int_{0}^{4-x-2y}dzdxdy+\int_{1}^{2}\int_{0}^{4-2y}\int_{0}^{4-x-2y}dzdxdy \\
= \ds \int_{0}^{4}\int_{0}^{\frac{4-z}{4}}\int_{0}^{2y}dxdydz+\int_{0}^{4}\int_{\frac{4-z}{4}}^{\frac{4-z}{2}}\int_{0}^{4-z-2y}dxdydz =\frac{8}{3}.
\end{multline*}
\vskip0.2cm
La primera de estas integrales se evaluó en \verb"Mathematica" ejecutando las siguientes instrucciones:
\enlargethispage{\baselineskip}
\vspace{0.6em}
\noindent
\begin{tcolorbox}[breakable,enhanced,title=\bf Instrucción en \verb"Mathematica", top=2pt,bottom=2pt,boxsep=0pt,before skip=2pt,after skip=0pt]
\begin{Verbatim}[formatcom=\letraverbatim]
Integrate[1,{y,0,1},{z,0,4-4*y},{x,0,2*y}]+
Integrate[1,{y,0,1},{z,4-4y,4-2*y},{x,0,4-z-2*y}]+
Integrate[1,{y,1,2},{z,0,4-2*y},{x,0,4-z-2*y}]
\end{Verbatim}
\end{tcolorbox}
\vspace{0.7em}
%\noindent
\pagebreak[0]
El volumen también se puede hallar de la siguiente manera.

%\begin{tcolorbox}[title=Parametrizando el sólido]
%\small
	Los vértices de la base son los puntos $(0,0,0)$, $(2,1,0)$ y $(0,2,0)$. Este triángulo se puede parametrizar primero, considerando un segmento entre dos puntos de estos y luego \textbf{formar segmentos rectilíneos con el tercer punto}. Es decir $(1-u)(0,0,0) + u(2,1,0) = (2u, u, 0)$ y luego $(1-v)(2u, u, 0) + v(0,2,0) = (2(1-v)u, (1-v)u + 2v, 0)$. Por tal razón la función $\mathbf{r}(u, v) = \langle 2(1-v)u, (1-v)u + 2v, 0 \rangle$, $u \in [0,1]$, $v \in [0,1]$, representa la base del sólido. El sólido completo se recorre involucrando el vértice $(0,0,4)$, es decir, se debe formar la siguiente combinación convexa: $(1-w)(0,0,4) + w(2(1-v)u, (1-v)u + 2v, 0)$ que al simplificarse, determina la función que parametriza el sólido
	
	\[ \mathbf{r}(u,v,w) = \left\langle 2wu(1-v), wu(1-v)+2wv, 4-4w \right\rangle, \]
	con $u \in [0,1]$, $v \in [0,1]$, $w \in [0,1]$. Es sencillo ver que
	\begin{align*}
		\mathbf{r}_u &= \left\langle 2w(1-v), w(1-v), 0 \right\rangle \\
		\mathbf{r}_v &= \left\langle -2wu, w(-u+2), 0 \right\rangle \\
		\mathbf{r}_w &= \left\langle 2(1-v)u, u-uv+2v, -4 \right\rangle,
	\end{align*}
	con lo que se puede establecer $\mathbf{r}_u \times \mathbf{r}_v = \langle 4w(u-2), -8wu, -4wu \rangle;$ por lo tanto, $|\mathbf{r}_u \cdot (\mathbf{r}_v \times \mathbf{r}_w)| = 16|w^2(v-1)|$, el volumen del sólido está determinado por
	\[ \int_{0}^{1}\int_{0}^{1}\int_{0}^{1} 16|w^2(v-1)|\, dudvdw = \int_{0}^{1}\int_{0}^{1}\int_{0}^{1} 16w^2(1-v)\, dudvdw = \frac{8}{3}. \]
%\end{tcolorbox}
%\clearpage

\noindent
Si la densidad del sólido en un punto $(x,y,z)$ está dada por $\delta(x,y,z) = xyz^{2}$, entonces la masa es:
\begin{align*}
\int_{0}^{2} \int_{x/2}^{\frac{4-x}{2}} \int_{0}^{4 - x - 2y} xyz^{2} \, dz \, dy \, dx 
&= \int_{0}^{2} \int_{0}^{4 - 2x} \int_{x/2}^{\frac{4 - x - z}{2}} xyz^{2} \, dy \, dz \, dx \\
&= \frac{256}{315}.
\end{align*}

\vspace{0.5em}

Con la parametrización escogida anteriormente, la densidad está dada por:
\[
\delta(u,v,w) = 32w^{2}u(1-v)(u - uv + 2v)(1 - w)^{2}.
\]

Y la masa del sólido es:
\begin{center} % Mantiene el centrado si es deseado
	\resizebox{1.0\textwidth}{!}{%
		$\begin{aligned}
			\int_{0}^{1} \int_{0}^{1} \int_{0}^{1} 16 \cdot \left(32w^{2}u(1 - v)^{2}(u - uv + 2v)(1 - w)^{2}\right) w^{2} \, du \, dv \, dw = \frac{256}{315}.
		\end{aligned}$
	}
\end{center}
\subsection*{Una pirámide II}
Las figuras de las expresiones $z=y-x,\,\, z=0,\,\, y=1, \,\, x=0,$
corresponden a cuatro planos que determinan un sólido; hallaremos su volumen.

\begin{figure}[H]
	\centering
	\caption{Intersección entre los planos $x - y + z = 0$ \quad y \quad $y = 1$.}
	\label{fig:interseccion_xy_planos}
	\begin{minipage}[t]{0.48\textwidth}
		\centering
		\includegraphics[width=\linewidth]{Fig/153.pdf}
	\end{minipage}
	\hspace{1em}
	\begin{minipage}[t]{0.35\textwidth}
		\centering
		\includegraphics[width=\linewidth]{Fig/154.pdf}
	\end{minipage}
	
	\vspace{1ex}
	\fuentepropia
\end{figure}



\noindent
El sólido proyectado en cada uno de los planos coordenados proporciona las siguientes figuras:

\noindent
\begin{minipage}{\textwidth}
	\centering
	% Usar \captionof es más robusto cuando no estás en un entorno 'figure' flotante.
	\captionof{figure}{Proyección de los planos $x - y + z = 0$, $y = 1$ y su intersección.} 
	\label{fig:proyeccion_interseccion_planos}
	
	% AJUSTE CRÍTICO DE ANCHOS Y ESPACIOS:
	% (3 * 0.31) + (2 * 0.015) = 0.96 < 1.0 (deja un poco de margen)
	\begin{minipage}{0.31\linewidth} % Ancho reducido
		\centering
		\vspace{-5mm}
		\includegraphics[width=\linewidth]{Fig/155.pdf}
	\end{minipage}
	\hfill % Relleno de espacio flexible
	\begin{minipage}{0.31\linewidth} % Ancho reducido
		\centering
		\vspace{-5mm}
		\includegraphics[width=\linewidth]{Fig/156.pdf}
	\end{minipage}
	\hfill % Relleno de espacio flexible
	\begin{minipage}{0.31\linewidth} % Ancho reducido
		\centering
		\vspace{-5mm}
		\includegraphics[width=\linewidth]{Fig/157.pdf}
	\end{minipage}
	
	\vspace{1ex}
	% Nota al pie con fuente pequeña.
	\fuentepropia
\end{minipage}
\vskip0.2cm


Con ayuda de ellos es posible plantear las integrales triples que determinan el volumen del sólido presentado.
\vspace{-2mm}
\begin{multline*}
\int_{0}^{1}\int_{x}^{1}\int_{0}^{y-x}dzdydx= \int_{0}^{1}\int_{0}^{y}\int_{0}^{y-x}dzdxdy \\ =\int_{0}^{1}\int_{z}^{1}\int_{0}^{y-z}dxdydz
\int_{0}^{1}\int_{0}^{y}\int_{0}^{y-z}dxdzdy \\ =\int_{0}^{1}\int_{0}^{1-z}\int_{x+z}^{1}dydxdz= \int_{0}^{1}\int_{0}^{1-x}\int_{x+z}^{1}dydzdx= \frac{1}{6}
\end{multline*}
Estos cálculos se pueden hacer en \verb"Mathematica" ejecutando las siguientes instrucciones:

\vskip0.2cm
\noindent
\begin{tcolorbox}[breakable,enhanced,title=\bf Instrucción en \verb"Mathematica", top=2pt,bottom=2pt,boxsep=0pt,before skip=2pt,after skip=0pt]
\begin{Verbatim}[formatcom=\letraverbatim]
Integrate[1,{x,0,1},{y,x,1},{z,0,y-x}]
Integrate[1,{y,0,1},{x,0,y},{z,0,y-x}]
Integrate[1,{z,0,1},{y,z,1},{x,0,y-z}]
Integrate[1,{y,0,1},{z,0,y},{x,0,y-z}]
Integrate[1,{z,0,1},{x,0,1-z},{y,x+z,1}]
Integrate[1,{x,0,1},{z,0,1-x},{y,x+z,1}]
\end{Verbatim}
\end{tcolorbox}
\vspace{0.7em}
\noindent

El volumen también se puede hallar con el siguiente procedimiento. En el plano $xy$, el segmento que une los puntos $(0,0)$ con $(0,1)$ se parametrizan como $(0,t)$; y el segmento que une cualquiera de estos puntos con $(1,1)$ es de la forma $(u, t - tu + u)$. En el espacio, el sólido se recorre mediante segmentos que unen los puntos $(u, t - tu + u, 0)$ con $(1, 1, 1)$, esto es $(1-w)(u, t - tu + u, 0) + w(1, 1, 1)$. Por tanto, la función $\mathbf{r}(t, u, w) = \left\langle (u + w(1-u)), t + (t-1)(uw - w - u), w \right\rangle$, con $t \in [0,1]$, $u \in [0,1]$, $w \in [0,1]$, parametriza el sólido. Con esta función hallamos los siguientes vectores.
\vspace{-4mm}
	\begin{align*}
		\mathbf{r}_{t} &= \langle 0, 1 - w + uw - u, 0 \rangle \\
		\mathbf{r}_{u} &= \langle 1 - w, t(1 - w) - (1 - w), 0 \rangle \\
		\mathbf{r}_{w} &= \langle 1 - u, (t - 1)(u - 1), 1 \rangle
	\end{align*}
	
	\noindent
	\textbf{El producto vectorial triple} es $|\mathbf{r}_{t} \cdot (\mathbf{r}_{u} \times \mathbf{r}_{w})| = |(w-1)^{2}(u-1)|$, en este caso es $(1-w)^{2}(1-u)$, al integrarlo nos proporciona el volumen del sólido
	
	\[ \int_{0}^{1}\int_{0}^{1}\int_{0}^{1} (1-w)^{2}(1-u) \, dt\,du\,dw = \frac{1}{6}. \]
	
%\end{tcolorbox}


\vskip0.2cm
El cálculo anterior, incluyendo la parametrización, se consigue con las siguientes instrucciones:
\vskip0.2cm
\noindent
\begin{tcolorbox}[breakable,enhanced,title=\bf Instrucción en \verb"Mathematica", top=2pt,bottom=2pt,boxsep=0pt,before skip=2pt,after skip=0pt]
\begin{Verbatim}[formatcom=\letraverbatim]
r=(1-w)*((1-u){0,t,0}+u*{1,1,0})+w*{1,1,1};
A=Abs[Dot[D[r,t],Cross[D[r,u],D[r,w]]]];
Integrate[A,{u,0,1},{v,0,1},{w,0,1}]
\end{Verbatim}
\end{tcolorbox}
\vspace{0.7em}
\noindent

\subsection*{Tres planos y un cilindro parabólico}
El cilindro parabólico con ecuación $z=4-y^{2}$ y los planos $y=x,$ $x=0$ y $z=0$ limitan en el primer octante del sólido que se muestra a conti\-nuación.

\vskip0.3cm
\noindent
\begin{minipage}{\textwidth}
	\centering
	\captionof{figure}{Sólido determinado por $0 \leq z \leq 4 - y^2$, \quad $0 \leq x \leq y$.}
	\label{fig:solido_parabolico_inclinado_2}
	\begin{minipage}{0.23\linewidth}
		\centering
		\includegraphics[width=\linewidth]{Img/16c.pdf}
	\end{minipage}
	\hspace{1em}
	\begin{minipage}{0.23\linewidth}
		\centering
		\includegraphics[width=\linewidth]{Img/16a.pdf}
	\end{minipage}
	\hspace{1em}
	\begin{minipage}{0.23\linewidth}
		\centering
		\includegraphics[width=\linewidth]{Img/16b.pdf}
	\end{minipage}
	
	\vspace{1ex}
	\fuentepropia
\end{minipage}

\vskip0.3cm

Estas figuras se obtienen ejecutando las siguientes instrucciones:

\vskip0.2cm
\noindent
\begin{tcolorbox}[breakable,enhanced,title=\bf Instrucción en \verb"Mathematica", top=2pt,bottom=2pt,boxsep=0pt,before skip=2pt,after skip=0pt]
\begin{Verbatim}[formatcom=\letraverbatim]
ContourPlot3D[{z==4-y^2,x==y,x==0},{x,-1,2},{y,-2,2},
{z,0,4},PlotPoints->60,Mesh->None,
ContourStyle->{Orange,Blend[{Yellow,Pink,Green}]}]
RegionPlot3D[0<z<4-y^2&&x<y<2&&0<x<2,{x,0,2},{y,0,2},
{z,0,4},PlotPoints->90,Mesh->None]
\end{Verbatim}
\end{tcolorbox}
\vspace{0.7em}
\noindent


\noindent
\begin{minipage}{\textwidth}
	\centering
	\captionof{figure}{Proyección del sólido sobre los planos coordenados.}
	\label{fig:proyeccion_sólido}
	\begin{minipage}{0.29\linewidth}
		\centering
		\vspace{-4mm}
		\includegraphics[width=\linewidth]{Fig/158.pdf}
	\end{minipage}
	\hspace{1em}
	\begin{minipage}{0.29\linewidth}
		\centering
		\vspace{-4mm}
		\includegraphics[width=\linewidth]{Fig/159.pdf}
	\end{minipage}
	\hspace{1em}
	\begin{minipage}{0.29\linewidth}
		\centering
		\vspace{-4mm}
		\includegraphics[width=\linewidth]{Fig/160.pdf}
	\end{minipage}
	
	\vspace{1ex}
	\fuentepropia
\end{minipage}
\vskip0.2cm
Estas figuras nos orientan para plantear las siguientes integrales triples.
\begin{multline*}
\int_{0}^{2}\int_{x}^{2}\int_{0}^{4-y^{2}}dzdydx= \int_{0}^{2}\int_{0}^{y}\int_{0}^{4-y^{2}}dzdxdy \\
= \int_{0}^{4}\int_{0}^{\sqrt{4-z}}\int_{x}^{\sqrt{4-z}}dydxdz
\int_{0}^{2}\int_{0}^{4-x^{2}}\int_{x}^{\sqrt{4-z}}dydzdx \\
\int_{0}^{4}\int_{0}^{\sqrt{4-z}}\int_{0}^{y}dxdydz= \int_{0}^{2}\int_{0}^{4-y^{2}}\int_{0}^{y}dxdzdy= 4.
\end{multline*}
Estas integrales se evalúan en \verb"Mathematica" con las siguientes sentencias:

\vskip0.2cm
\noindent
\begin{tcolorbox}[breakable,enhanced,title=\bf Instrucción en \verb"Mathematica", top=2pt,bottom=2pt,boxsep=0pt,before skip=2pt,after skip=0pt]
\begin{Verbatim}[formatcom=\letraverbatim]
Integrate[1,{x,0,2},{y,x,2},{z,0,4-y^2}]
Integrate[1,{y,0,2},{x,0,y},{z,0,4-y^2}]
Integrate[1,{z,0,4},{x,0,Sqrt[4-z]},{y,x,Sqrt[4-z]}]
Integrate[1,{x,0,2},{z,0,4-x^2},{y,x,Sqrt[4-z]}]
Integrate[1,{z,0,4},{y,0,Sqrt[4-z]},{y,0,y}]
Integrate[1,{y,0,2},{z,0,4-y^2},{x,0,y}]
\end{Verbatim}
\end{tcolorbox}
\vspace{0.7em}
\noindent

Otra posibilidad de obtener este valor es la siguiente.
\vskip0.2cm


%\begin{tcolorbox}[title=\bfseries Parametrizando el sólido]
\noindent	
La base de esta región es un triángulo en el plano $xy$, cuyos vértices son los puntos $(0,0),$ $(0,2)$ y $(2,2).$
Primero, se parametriza el segmento que une los puntos $(0,0)$ con $(2,2),$ con una combinación convexa entre dichos puntos, es decir,
$(1-u)(0,0) +u( 2,2) =( 2u,2u) .$
	
\vspace{1ex} 
	
\noindent
	

Luego se une el punto $(0,2) $ con los puntos de la forma $(2u,2u)$, es decir, se forman los segmentos rectilíneos por medio de la expresión
$(1-v)(0,2)+v ( 2u,2u)=\langle 2uv,2-2v+2uv\rangle.$
Algunos de estos segmentos se muestran a continuación.

	
\vskip0.2cm
\noindent
\begin{minipage}{\textwidth}
	\centering
	\captionof{figure}{$(2uv, 2(1 - v + uv)).$}
	\label{fig:curva_parametrica_uv_2}
	\begin{minipage}{0.32\linewidth}
		\centering
		\includegraphics[width=\linewidth]{Fig/161.pdf}
	\end{minipage}
	
	\vspace{0.3em}
	\fuentepropia
\end{minipage}
\vskip0.2cm
	
	\vspace{1ex} 
	\noindent
	
Sobre la superficie del cilindro parabólico, la coordenada en $z$ se puede obtener con la ecuación que lo define, es decir, \[z= 4-y^{2}=4-(2-2v+2uv)^{2} =4v(1-u)(uv+2-v).\]
%\end{tcolorbox}


%\begin{tcolorbox}
%\small
\begin{figure}[H]
	\centering
	\caption{Sólido determinado por $0 \leq z \leq 4 - y^2$, \quad $x \leq y \leq 2$.}
	\label{fig:solido_parabolico_y}
	\vspace{0.3cm} 
	\begin{minipage}[t]{0.33\textwidth}
		\centering
		\includegraphics[width=\linewidth]{Fig/162.pdf}
	\end{minipage}
	\hspace{1em}
	\begin{minipage}[t]{0.33\textwidth}
		\centering
		\includegraphics[width=\linewidth]{Fig/163.pdf}
	\end{minipage}
	
	\vspace{0.4cm} 
	
	\fuentepropia
\end{figure}

\vspace{-4mm}
Es de anotar que la curva sobre el cilindro en el plano $x=0$, se determina por medio de la expresión
$\langle 0,t,4-t^{2}\rangle$ y sobre el plano $y=x$, se usa $\langle t,t,4-t^{2}\rangle$.\\

Para efectos del gráfico, la parte sobre el cilindro se parametriza con
$(1-u)\langle 0,t,4-t^{2}\rangle +u\langle t,t,4-t^{2}\rangle=\langle ut,t,4-t^{2}\rangle$. \\[0.25cm]
Con lo anterior, se puede establecer que un punto $Q$ sobre la superficie del cilindro tiene la forma
$(2uv,2-2v+2uv,4v(1-u)(uv+2-v))$, mientras que un punto $P$ sobre la base del sólido tiene la forma $(2uv,2-2v+2uv,0).$
Ahora se puede parametrizar el sólido formando una combinación convexa de los puntos $P$ y $Q$,
esto es $(1-w)Q+wP$,
%$(1-w) \langle 2uv,2-2v+2uv,0\rangle+w\langle 2uv,2-2v+2uv,4v(1-u)(uv+2-v)\rangle$
que al simplificarse nos permite definir la función
${\bf r}( u,v,w) = \langle 2uv,2-2v+2uv,4wv(1-u)(uv+2-v)\rangle,$
para $u\in [0,1]$, $v\in [0,1]$, $w\in [0,1].$ De ella se obtienen los siguientes vectores:
\vspace{-3mm}
\begin{eqnarray*}
{\bf r}_{u}&=& \langle 2v,2v,8wv( v-uv-1)\rangle \\
{\bf r}_{v}&=& \langle 2u,-2+2u,-8w(u-1) ( uv-v+1) \rangle \\
{\bf r}_{w}&=& \langle 0,0,4v( 1-u) (uv+2-v) \rangle
\end{eqnarray*}
\vspace{-4pt}
Además, ${\bf r}_{v}\times {\bf r}_{w}= \big\langle -8v( u-1)^{2}(uv+2-v),8uv( u-1)(uv+2-v),0 \big \rangle,$
entonces, $ |{\bf r}_{u}\cdot({\bf r}_{v}\times {\bf r}_{w})|=16|v^{2}(u-1)(uv+2-v)|.$ Por último, el volumen del sólido está dado por
{\small
\[\int_{0}^{1}\int_{0}^{1}\int_{0}^{1}16| v^{2}( u-1)(uv+2-v)|dudvdw= 4.\]
}
%\end{tcolorbox}

%\clearpage
Los cálculos anteriores se pueden efectuar en \verb"Mathematica"
definiendo una parametrización $p$ para la base del sólido, $x$ e $y$ son las coordenadas de $p.$ Se determina una parametrización ${\bf r}$ del sólido y se evalúa el triple producto vectorial, que se integra para encontrar el volumen deseado.
\vskip0.2cm
\noindent
\begin{tcolorbox}[breakable,enhanced,title=\bf Instrucción en \verb"Mathematica",top=2pt,bottom=2pt,boxsep=0pt,before skip=2pt,after skip=0pt]
\begin{Verbatim}[formatcom=\letraverbatim]
p={x,y}=(1-v)*{0,2}+v*{2*u,2*u}
r=(1-w)*{x,y,0}+w*{x,y,4-y^2};
A=Factor[Dot[D[r,u],Cross[D[r,v],D[r,w]]]]
Integrate[Abs[A],{u,0,1},{v,0,1},{w,0,1}]
\end{Verbatim}
\end{tcolorbox}
\vspace{0.7em}
\noindent

\subsection*{Cuatro planos y un cilindro parabólico}

La porción de las figuras de las expresiones $z=2-x$, $z^2=x$, $y=0$, $y=3$ y $z=0$ en el primer octante se muestra a continuación; estas acotan una región en el espacio.

\vskip0.2cm
\noindent
\begin{minipage}{\textwidth}
	\centering
	\captionof{figure}{Superficies delimitantes: $z = 2 - x$, $z^2 = x$, $y = 0$, $y = 3$, $z = 0$.}
	\label{fig:superficies_z2x_z2x_y03_2}
	\begin{minipage}{0.34\linewidth}
		\centering
		\vspace{-3mm}
		\includegraphics[width=\linewidth]{Fig/164.pdf}
	\end{minipage}
	\hspace{1em}
	\begin{minipage}{0.29\linewidth}
		\centering
		\vspace{-3mm}
		\includegraphics[width=\linewidth]{Fig/165.pdf}
	\end{minipage}
	
	\vspace{1ex}
	\fuentepropia
\end{minipage}

\vskip0.2cm
Ahora se muestra el sólido en mención y las instrucciones que lo generan.


\vskip0.2cm
\noindent
\begin{minipage}{\textwidth}
	\centering
	\captionof{figure}{Sólido acotado por $z = 2 - x$, $z^2 = x$, $y = 0$, $y = 3$, $z = 0$}
	\label{fig:solido_z2x_z2x_y03}
	\begin{minipage}{0.19\linewidth}
		\centering
		\includegraphics[width=\linewidth]{Img/15a.pdf}
	\end{minipage}
	\begin{minipage}{0.19\linewidth}
		\centering
		\includegraphics[width=\linewidth]{Img/15b.pdf}
	\end{minipage}
	
	\vspace{0.3em}
	\fuentepropia
\end{minipage}
\vskip0.2cm



\noindent
\begin{tcolorbox}[breakable,enhanced,title=\bf Instrucción en \verb"Mathematica", top=2pt,bottom=2pt,boxsep=0pt,before skip=2pt,after skip=0pt]
\begin{Verbatim}[formatcom=\letraverbatim]
RegionPlot3D[z^2<x<2-z&&0<y<3&&0<x<2,{x,0,2},{y,0,3},
{z,0,1},PlotPoints->90,Mesh->None]
\end{Verbatim}
\end{tcolorbox}
\vspace{0.7em}
\noindent


\noindent
\begin{minipage}{\textwidth}
	\centering
	\captionof{figure}{Proyección del sólido en los planos coordenados.}
	\label{fig:proyeccion_sólido_coordenados}
	\begin{minipage}{0.28\linewidth}
		\centering
		\includegraphics[width=\linewidth]{Fig/166.pdf}
	\end{minipage}
	\hspace{1em}
	\begin{minipage}{0.28\linewidth}
		\centering
		\includegraphics[width=\linewidth]{Fig/167.pdf}
	\end{minipage}
	\hspace{1em}
	\begin{minipage}{0.28\linewidth}
		\centering
		\includegraphics[width=\linewidth]{Fig/168.pdf}
	\end{minipage}
	\vspace{0.2cm}
	\vspace{0.3em}
	\fuentepropia
\end{minipage}
	

El volumen del sólido se puede determinar evaluando, según el caso, alguna de las siguientes integrales.
{\small
\[\int_{0}^{1}\int_{0}^{3}\int_{z^{2}}^{2-z}dxdydz=\int_{1}^{2}\int_{0}^{3}\int_{0}^{2-x}dzdydx+\int_{0}^{1}\int_{0}^{3}\int_{0}^{\sqrt{x}}dzdydx =\frac{7}{2}\]
\[\int_{0}^{3}\int_{1}^{2}\int_{0}^{2-x}dzdxdy+\int_{0}^{3}\int_{0}^{1}\int_{0}^{\sqrt{x}}dzdxdy =\int_{0}^{3}\int_{0}^{1}\int_{z^{2}}^{2-z}dxdzdy=\frac{7}{2}\]
\[\int_{0}^{1}\int_{z^{2}}^{2-z}\int_{0}^{3}dydxdz= \int_{0}^{1}\int_{0}^{\sqrt{x}}\int_{0}^{3}dydzdx
+\int_{1}^{2}\int_{0}^{2-x}\int_{0}^{3}dydzdx= \frac{7}{2}.\]
}

\vskip0.2cm
Dos de estas integrales se pueden evaluar con las siguientes sentencias:

\vskip0.2cm
\noindent
\begin{tcolorbox}[breakable,enhanced,title=\bf Instrucción en \verb"Mathematica", top=2pt,bottom=2pt,boxsep=0pt,before skip=2pt,after skip=0pt]
\begin{Verbatim}[formatcom=\letraverbatim]
Integrate[1,{z,0,1},{y,0,3},{x,z^2,2-z}]
Integrate[1,{x,1,2},{y,0,3},{z,0,2-x}]+
Integrate[1,{x,0,1},{y,0,3},{z,0,Sqrt[x]}]
\end{Verbatim}
\end{tcolorbox}
\vspace{0.7em}
\noindent

Este volumen también se puede obtener de la siguiente manera.
\vskip0.2cm
%\begin{tcolorbox}[title=\bfseries Parametrizando el sólido]
	%\small
	\noindent
	En el plano $xz$, para $t \in [0,1]$, los puntos $A$ y $B$ son de la forma $(t^{2},t)$ y $(2-t,t)$ respectivamente; el segmento que los une se parametriza mediante la combinación convexa $(1-u)A+uB,$ que se simplifica como $(t^{2}(1-u)+u(2-t),t),$ para $u\in [0,1].$
	\vskip0.3cm
	%\noindent
	\begin{minipage}{\textwidth}
		\captionof{figure}{Superficies delimitantes: $z = 2 - x$, $z^2 = x$, $y = 0$, $y = 3$, $z = 0$.}
		\label{fig:superficies_z2x_z2x_y03_3}
		\begin{minipage}{0.31\linewidth}
			\centering
			\vspace{-4mm}
			\includegraphics[width=\linewidth]{Fig/206.pdf}
		\end{minipage}
		\hspace{1em}
		\begin{minipage}{0.37\linewidth}
			\centering
			\vspace{-4mm}
			\includegraphics[width=\linewidth]{Fig/207.pdf}
		\end{minipage}
		%\vspace{0.2cm}
		\vspace{0.3em}
		\fuentepropia
	\end{minipage}
	
	\vspace{1ex} 
	El sólido se parametriza formando segmentos rectilíneos entre dos puntos $P(t^{2}(1-u)+(2-t)u,0,t)$ y $Q(t^{2}(1-u)+(2-t)u,3,t).$
	Simplificando $(1-w)P+wQ$ se define
	 \[\mathbf{r}(t,u,w) = \left\langle t^{2}(1-u)-u(t-2), 3w, t \right\rangle,\]
con $t\in [0,1]$, $u\in [0,1]$ y $w\in [0,1]$. Además,
%\vspace{-10mm}
{\small
\begin{align*}
	\mathbf{r}_t &= \frac{\partial}{\partial t}\left\langle t^{2}(1-u)-u(t-2),3w,t\right\rangle = \left\langle 2t(1-u)-u,0,1\right\rangle \\
	\mathbf{r}_u &= \frac{\partial}{\partial u}\left\langle t^{2}(1-u)-u(t-2),3w,t\right\rangle = \left\langle -t^{2}+2-t,0,0\right\rangle \\
	\mathbf{r}_w &= \frac{\partial}{\partial w}\left\langle t^{2}(1-u)-u(t-2),3w,t\right\rangle = \left\langle 0,3,0\right\rangle.
\end{align*}
}
	
	El elemento diferencial de volumen es $\left| \mathbf{r}_t\cdot (\mathbf{r}_u\times \mathbf{r}_w) \right|
	= \left|6-3t-3t^{2}\right| ,$
	al integrarse determina el volumen del sólido
	
	\vspace{-5mm}
	
	\[ \int_{0}^{1}\int_{0}^{1}\int_{0}^{1} |6-3t-3t^{2}| \, dt\,du\,dw = \int_{0}^{1}\int_{0}^{1}\int_{0}^{1}(6-3t-3t^{2}) \, dt\,du\,dw = \frac{7}{2}. \]

%\end{tcolorbox}


\subsection*{Un cilindro y dos planos}
Considere la región en el primer octante, acotada por las figuras de las ecuaciones
$y^2+z^2=9$, $z=x$ y $y=0$. Su gráfico se muestra a conti\-nuación:

\vskip0.2cm
\noindent
\begin{minipage}{\textwidth}
	\centering
	\captionof{figure}{Sólido determinado por $0 \leq y \leq \sqrt{9-z^2}$ \quad y \quad $0 \leq x \leq z$.}
	\label{fig:solido_yz_xz}
	\begin{minipage}{0.35\linewidth}
		\centering
		\vspace{-2mm}
		\includegraphics[width=\linewidth]{Fig/208.pdf}
	\end{minipage}
	\hspace{1em}
	\begin{minipage}{0.35\linewidth}
		\centering
		\vspace{-2mm}
		\includegraphics[width=\linewidth]{Fig/209.pdf}
	\end{minipage}
	\vspace{0.2cm}
	\fuentepropia
\end{minipage}


Enseguida se muestra la región proyectada sobre los planos coordenados.

\vskip0.2cm
\noindent
\begin{minipage}{\textwidth}
	\centering
	\captionof{figure}{Proyecciones de $y^2 + z^2 = 9$, $z = x$, $y = 0$ y sus intersecciones.}
	\label{fig:proyecciones_cilindro_plano}
	\begin{minipage}{0.29\linewidth}
		\centering
		\vspace{-2mm}
		\includegraphics[width=\linewidth]{Fig/210.pdf}
	\end{minipage}
	\hspace{1em}
	\begin{minipage}{0.29\linewidth}
	\centering
	\vspace{-2mm}
		\includegraphics[width=\linewidth]{Fig/211.pdf}
	\end{minipage}
	\hspace{1em}
	\begin{minipage}{0.29\linewidth}
		\centering
		\vspace{-2mm}
		\includegraphics[width=\linewidth]{Fig/212.pdf}
	\end{minipage}
	\vspace{0.2cm}
	
	\fuentepropia
\end{minipage}

\vskip0.2cm
Con ayuda de estas figura se pueden plantear las integrales necesarias para determinar el volumen del mencionado sólido.
\begin{multline*}
\int_{0}^{3}\int_{0}^{\sqrt{9-x^{2}}}\int_{x}^{\sqrt{9-y^{2}}}dzdydx =
\int_{0}^{3}\int_{0}^{\sqrt{9-y^{2}}}\int_{0}^{z}dxdzdy \\
=\int_{0}^{3}\int_{0}^{\sqrt{9-y^{2}}}\int_{x}^{\sqrt{9-y^{2}}}dzdxdy
\int_{0}^{3}\int_{0}^{\sqrt{9-z^{2}}}\int_{0}^{z}dxdydz \\
= \int_{0}^{3}\int_{0}^{z}\int_{0}^{\sqrt{9-z^{2}}}dydxdz = \int_{0}^{3}\int_{x}^{3}\int_{0}^{\sqrt{9-z^{2}}}dydzdx = 9
\end{multline*}
\vskip0.2cm
Tres de estas integrales se evalúan con \verb"Mathematica" usando las siguientes instrucciones:

\vskip0.2cm
\noindent
\begin{tcolorbox}[breakable,enhanced,title=\bf Instrucción en \verb"Mathematica", top=2pt,bottom=2pt,boxsep=0pt,before skip=2pt,after skip=0pt]
\begin{Verbatim}[formatcom=\letraverbatim]
Integrate[1,{z,0,3},{y,0,Sqrt[9-z^2]},{x,0,z}]
Integrate[1,{y,0,3},{x,0,Sqrt[9-y^2]},{z,x,
Sqrt[9-y^2]}]
Integrate[1,{x,0,3},{z,x,3},{y,0,Sqrt[9-z^2]}]
\end{Verbatim}
\end{tcolorbox}
\vspace{0.7em}
\noindent

Con el siguiente proceso también se halla este volumen.
\vskip0.2cm
%\begin{tcolorbox}[title=\bfseries Parametrizando el sólido]
Sobre el plano $yz,$ la región de integración se parametriza como
\[(0,3u\cos v,3u\sin v),\] para $u\in \lbrack 0,3],$ $v\in [0,\frac{\pi}{2}].$ Un punto sobre el plano $z=x,$ tiene la forma
$P(3u\sin v,3u\cos v,3u\sin v) ,$ mientras que sobre el plano $yz$ un punto arbitrario $Q$ tiene la forma $Q(0,3u\cos v,3u\sin v) .$
El sólido se parametriza formando segmentos rectilíneos puntos $P$ y $Q,$
esto define la función ${\bf r}$, simplificando la expresión $(1-w)P + wQ,$
\begin{eqnarray*}
{\bf r}(u,v,w) &=&\langle3 u w \sin (v),3 u \cos (v),3 u \sin (v)\rangle,
\end{eqnarray*}
con $u\in [0,1],$ $v\in [ 0,\frac{\pi }{2}],$ $w\in [0,1].$ Ahora se hallan los siguientes vectores
\begin{eqnarray*}
{\bf r}_{u}&=&\left\langle 3 w \sin (v),3 \cos (v),3 \sin (v)\right\rangle \\
{\bf r}_{v}&=&\left\langle 3 u w \cos (v),-3 u \sin (v),3 u \cos (v) \right\rangle \\
{\bf r}_{w}&=&\left\langle 3 u \sin (v),0,0 \right\rangle .
\end{eqnarray*}
Además $\left| {\bf r}_{u}\cdot \left( {\bf r}_{v}\times {\bf r}_{w}\right) \right| = 27u^{2}\sin(v),$ que determina el elemento diferencial de volumen, al integrarlo se halla el volumen del sólido.
\[\int_{0}^{1}\int_{0}^{\pi /2}\int_{0}^{1} 27u^{2}\sin(v) dudvdw=9.\]
%\end{tcolorbox}

Los cálculos mostrados antes se pueden efectuar brevemente, usando \verb"Mathematica",
mediante la ejecución de las siguientes instrucciones:

\vskip0.2cm
\noindent
\begin{tcolorbox}[breakable,enhanced,title=\bf Instrucción en \verb"Mathematica", top=2pt,bottom=2pt,boxsep=0pt,before skip=2pt,after skip=0pt]
\begin{Verbatim}[formatcom=\letraverbatim]
{y,z}={3*u*Cos[v],3*u*Sin[v]};
r=(1-w)*{0,y,z}+w*{z,y,z};
A=Dot[D[r,u],Cross[D[r,v],D[r,w]]];
Integrate[A,{u,0,1},{v,0,Pi/2},{w,0,1}]
\end{Verbatim}
\end{tcolorbox}
\vspace{0.7em}
\noindent

En coordenadas cilíndricas, escritas en  forma
$y=r \cos \theta$, $z=r \sin \theta$, con $\theta \in[0,\pi/2]$, $r \in [0,3];$
esta integral se evalúa como sigue:

$\int_0^{\pi/2} \int_0^3 \int_0^{r\sin\theta} r dxdr d\theta =9. $
El cálculo anterior se obtuvo en \verb"Mathematica" mediante la ejecución de la siguiente instrucción:

\begin{tcolorbox}[breakable,enhanced]
\begin{Verbatim}[formatcom=\letraverbatim]
Integrate[r,{\theta,0,Pi/2},{r,0,3},{x,0,r*Cos[\theta]}
\end{Verbatim}
\end{tcolorbox}


\subsection*{Un cilindro y dos planos II}
Las expresiones
$z=-y$, $x^2+y^2=x$ y $z=0$ corresponden a dos planos y un cilindro. Para $y<0$, estas expresiones acotan el sólido que se muestra a continuación.

\begin{figure}[H]
\centering
\caption{Sólido determinado por $x^2 + y^2 \leq x$, \quad para \quad $0 \leq z \leq -y$.}
\label{fig:solido_disco_inclinado}
\vspace{-2mm}
\includegraphics[width=4.0cm,height=3.5cm]{Img/12c.pdf}
\vspace{-2mm}
\includegraphics[height=4.5cm]{Img/12a.pdf}
\vspace{-2mm}
\includegraphics[height=4.5cm]{Img/12b.pdf}
\vspace{2mm}
\fuentepropia
\end{figure}


Estas figuras se obtienen ejecutando las siguientes instrucciones:

\vskip0.2cm
\noindent
\begin{tcolorbox}[breakable,enhanced,title=\bf Instrucción en \verb"Mathematica", top=2pt,bottom=2pt,boxsep=0pt,before skip=2pt,after skip=0pt]
\begin{Verbatim}[formatcom=\letraverbatim]
ContourPlot3D[{z==-y,x^2+y^2==x,z==0},{x,0,1},
{y,-1,1},{z,-1,1},Ticks->{{0,1},{0,1},{0,1}},
Mesh->False,ContourStyle->{Orange,LightGreen,Yellow}]

RegionPlot3D[0<z<-y&&x>x^2+y^2,{x,0,1},{y,-1/2,0},
{z,0,1/2},PlotPoints->90,Mesh->False,PlotStyle->Orange]
\end{Verbatim}
\end{tcolorbox}
\vspace{0.7em}
\noindent

Las expresiones $x^2+y^2\leq x,$ $0\leq z \leq -y$ determinan un triángulo en el plano $yz$ y un semicírculo en los planos $xz$ y $xy$.


\vskip0.2cm
\noindent
\begin{minipage}{\textwidth}
	\centering
	\captionof{figure}{Proyecciones del sólido en los planos coordenados}
	\label{fig:proyecciones_solido}
	\begin{minipage}{0.27\linewidth}
		\centering
		\includegraphics[width=\linewidth]{Fig/213.pdf}
	\end{minipage}
	\begin{minipage}{0.27\linewidth}
		\centering
		\includegraphics[width=\linewidth]{Fig/214.pdf}
	\end{minipage}
	\begin{minipage}{0.27\linewidth}
		\centering
		\includegraphics[width=\linewidth]{Fig/215.pdf}
	\end{minipage}
	
	\vspace{0.3em}
	\fuentepropia
\end{minipage}
\vskip0.2cm


El volumen del sólido está dado por cualquiera de las siguientes integrales.

\vspace{1ex} 
\noindent
\makebox[\linewidth][l]{
	\scriptsize % <--- REDUCCIÓN MÁXIMA DE LA FUENTE
	$\begin{aligned}
		\bigintsss_{0}^{1} \bigintsss_{-\sqrt{x-x^{2}}}^{0}\ds \bigintsss_{0}^{-y}dz\,dy\,dx &= \bigintsss_{-1/2}^{0}\bigintsss_{\frac{1}{2}-\sqrt{\frac{1}{4}-y^{2}}}^{\frac{1}{2}+ \sqrt{\frac{1}{4}-y^{2}}}\bigintsss_{0}^{-y}dz\,dx\,dy \\
		&= \bigintsss_{0}^{1}\bigintsss_{0}^{\sqrt{x-x^{2}}} \bigintsss_{-\sqrt{x-x^{2}}}^{-z}dy\,dz\,dx \\
		&= \bigintsss_{0}^{1/2}\bigintsss_{\frac{1}{2}-\sqrt{\frac{1}{4}-z^{2}}}^{\frac{1}{2}+ \sqrt{\frac{1}{4}-z^{2}}}\bigintsss_{-\sqrt{x-x^{2}}}^{-z}dy\,dx\,dz \\
		&= \bigintsss_{0}^{1/2}\bigintsss_{-1/2}^{-z}\bigintsss_{\frac{1}{2}-\sqrt{\frac{1}{4}-y^{2}}}^{\frac{1}{2}+\sqrt{\frac{1}{4}-y^{2}}}dx\,dy\,dz \\
		&= \bigintsss_{-1/2}^{0} \bigintsss_{0}^{-y}\bigintsss_{\frac{1}{2}-\sqrt{\frac{1}{4}-y^{2}}}^{\frac{1}{2}+\sqrt{\frac{1}{4}-y^{2}}}dx\,dz\,dy = \frac{1}{12}
	\end{aligned}$
}
\vspace{1ex}

Estas integrales se pueden evaluar en \verb"Mathematica" mediante el uso de las siguientes instrucciones:

\vskip0.2cm
\noindent
\begin{tcolorbox}[breakable,enhanced,title=\bf Instrucción en \verb"Mathematica", top=2pt,bottom=2pt,boxsep=0pt,before skip=2pt,after skip=0pt]
\begin{Verbatim}[formatcom=\letraverbatim]
Integrate[1,{x,0,1},{y,-Sqrt[x-x^2],0},{z,0,-y}]
Integrate[1,{y,-1/2,0},{x,1/2-Sqrt[1/4-y^2],1/2
+Sqrt[1/4-y^2]},{z,0,-y}]
Integrate[1,{x,0,1},{z,0,Sqrt[x-x^2]},
{y,-Sqrt[x-x^2],-z}]
NIntegrate[1,{z,0,1/2},{x,1/2-Sqrt[1/4-z^2],
1/2+Sqrt[1/4-z^2]},{y,-Sqrt[x-x^2],-z}]
NIntegrate[1,{z,0,1/2},{y,-1/2,-z},
{x,1/2-Sqrt[1/4-y^2],1/2+Sqrt[1/4-y^2]}]
NIntegrate[1,{y,-1/2,0},{z,0,-y},{x,1/2-Sqrt[1/4-y^2],
1/2+Sqrt[1/4-y^2]}]
\end{Verbatim}
\end{tcolorbox}
\vspace{0.7em}
\noindent


Si la densidad del sólido está dada por $\delta(x,y,z) =xy^{2},$ entonces la masa del sólido está dada por
\begin{multline*} \bigintsss_{0}^{1}\bigintsss_{-\sqrt{x-x^{2}}}^{0}\bigintsss_{0}^{-y}xy^{2}dzdydx \\ =\bigintsss_{-1/2}^{0}\bigintsss_{\frac{1}{2}-\sqrt{\frac{1}{4}-y^{2}}}^{\frac{1}{2}+\sqrt{\frac{1}{4}-y^{2}}}\bigintsss_{0}^{-y}xy^{2}dzdxdy= \frac{1}{240}.
\end{multline*}
% $\bigintsss_{0}^{1}\bigintsss_{0}^{\sqrt{x-x^{2}}}\bigintsss_{-\sqrt{x-x^{2}}}^{-z}xz^{2}dydzdx=\bigintsss_{0}^{1/2}\bigintsss_{\frac{1}{2}-\sqrt{\frac{1}{4}-z^{2}}}^{\frac{1}{2}+\sqrt{\frac{1}{4}-z^{2}}}\bigintsss_{-\sqrt{x-x^{2}}}^{-z}xz^{2}dydxdz= \frac{1}{720}$
%$\bigintsss_{0}^{1/2}\bigintsss_{-1/2}^{-z}\bigintsss_{\frac{1}{2}-\sqrt{\frac{1}{4}-y^{2}}}^{\frac{1}{2}+\sqrt{\frac{1}{4}-y^{2}}}xz^{2}dxdydz=\bigintsss_{-1/2}^{0}\bigintsss_{0}^{-y}\bigintsss_{\frac{1}{2}-\sqrt{\frac{1}{4}-y^{2}}}^{\frac{1}{2}+\sqrt{\frac{1}{4}-y^{2}}}xz^{2}dxdzdy= \frac{1}{720}$
El volumen y la masa del sólido también se pueden determinar con el siguiente procedimiento.
%\vskip0.1cm

\vskip0.3cm
%\begin{tcolorbox}[title=\bfseries Parametrizando el sólido]
A partir de la ecuación $x^{2}+y^{2}=x$, que equivale a
$\left( x-\frac{1}{2}\right) ^{2}+y^{2}=\frac{1}{4},$ que representa una circunferencia cuyo interior se puede parametrizar como $x=\frac{1}{2}+u\cos v$, $y=u\sin v$, para concluir que un punto en la base del sólido puede parametrizarse en la forma,
$P=\left( \frac{1}{2}+u\cos v,u\sin v,0\right) $ y un punto del sólido sobre el plano con ecuación $z=-y,$ es la forma
$Q=\left( \frac{1}{2}+u\cos v,u\sin v,-u\sin v\right).$
\vskip0.3cm
Se pueden formar segmentos rectilíneos que unan estos puntos y parame- \\ trizar así el sólido. Esto es $(1-w)Q+wP$
es por esta razón que se usara la función
\[{\bf r}\left( u,v,w\right) =\left\langle \frac{1}{2}+u\cos (v) ,u\sin
(v) , (w-1)u\sin (v) \right\rangle\]
para $u\in [0,1/2],\,v\in [-\pi,0],\,w\in [0,1],\, $
con lo anterior, obtenemos los siguientes vectores
\begin{eqnarray*}
{\bf r}_{u}&=&\langle \cos (v) ,\sin (v),\sin (v) \left( w-1\right) \rangle\\
{\bf r}_{v}&=&\langle -u\sin (v) ,u\cos (v) ,u\cos(v) \left( w-1\right) \rangle\\
{\bf r}_{w}&=& \langle 0,0,u\sin (v) \rangle
\end{eqnarray*}
El elemento diferencial de volumen es
$| {\bf r}_{u}\cdot \left( {\bf r}_{v}\times {\bf r}_{w}\right)| =-u^{2}\sin (v) ,$
entonces, el volumen del sólido esta dado por
\vspace{-4mm}
{\small
\begin{multline*}
\ds \int_{0}^{1}\ds \int_{-\pi }^{0}\ds \int_{0}^{1/2}\left| u^{2}\sin (v)
\right| dudvdw \\ = \int_{0}^{1}\ds \int_{-\pi }^{0} \int_{0}^{1/2} -u^{2}\sin (v)dudvdw=\frac{1}{12}.
\end{multline*}
}
%\end{tcolorbox}
	
Con las siguientes indicaciones se integra el término
$|{\bf r}_{u}\cdot \left( {\bf r}_{v}\times {\bf r}_{w}\right)|$, que con los límites adecuados nos produce su volumen.

\vspace{0.6em}
\noindent
\begin{tcolorbox}[breakable,enhanced,title=\bf Instrucción en \verb"Mathematica", top=2pt,bottom=2pt,boxsep=0pt,before skip=2pt,after skip=0pt]
\begin{Verbatim}[formatcom=\letraverbatim]
{x,y,z}={1/2+u*Cos[v],u*Sin[v],-y}
r=(1-w)*{x,y,0}+w*{x,y,z};
M=Abs[Dot[D[r,u],Cross[D[r,v],D[r,w]]]];
Integrate[M,{u,0,1/2},{v,-Pi,0},{w,0,1}]
\end{Verbatim}
\end{tcolorbox}
\vspace{0.7em}
\noindent

\begin{tcolorbox}
Si la densidad en cualquier punto $(x,y,z)$ de sólido es $xy^{2}$,
entonces, $\delta \left( u,v,w\right) =u^2 \sin ^2(v) \left(u \cos (v)+\frac{1}{2}\right)$. Además
{\small
\[\delta \left( u,v,w\right) | {\bf r}_{u}\cdot \left( {\bf r}_{v}\times {\bf r}_{w}\right)|=\left|-\frac{1}{2}u^4\sin^3(v)(2u\cos(v)+1)\right|\]}

por lo tanto, la masa del sólido esta dada por
$\int_{0}^{1}\int_{-\pi }^{0}\int_{0}^{1/2}\left|-\frac{1}{2}u^4\sin^3(v)(2u\cos(v)+1)\right| dudvdw=\frac{1}{240} $
\end{tcolorbox}


\subsection*{Tres planos y un cilindro parabólico II}
Las siguientes ecuaciones $z=2-y,$ $x= 4-y^2,$ $x=0,$ $y=0 $ definen un sólido en el espacio.

\vskip0.2cm
\noindent
\begin{minipage}{\textwidth}
	\centering
	\captionof{figure}{Sólido determinado por $0 \leq y + z \leq 2$ \quad y \quad $0 \leq x \leq 4 - y^2$.}
	\label{fig:solido_yz_xy}
	\begin{minipage}{0.29\linewidth}
		\centering
		\vspace{-4mm}
		\includegraphics[width=\linewidth]{Fig/216.pdf}
	\end{minipage}
	\hspace{1em}
	\begin{minipage}{0.29\linewidth}
		\centering
		\vspace{-4mm}
		\includegraphics[width=\linewidth]{Fig/217.pdf}
	\end{minipage}
	
	\vspace{1ex}
	\fuentepropia
\end{minipage}
\vskip0.2cm


Este sólido determina a su vez las siguientes regiones.

\vskip0.2cm
\noindent
\begin{minipage}{\textwidth}
	\centering
	\captionof{figure}{Proyección del sólido en los planos coordenados.}
	\label{fig:proyecciones_planes_coordenados}
	\begin{minipage}{0.28\linewidth}
		\centering
		\includegraphics[width=\linewidth]{Fig/218.pdf}
	\end{minipage}
	\hspace{1em}
	\begin{minipage}{0.28\linewidth}
		\centering
		\includegraphics[width=\linewidth]{Fig/219.pdf}
	\end{minipage}
	\hspace{1em}
	\begin{minipage}{0.28\linewidth}
		\centering
		\includegraphics[width=\linewidth]{Fig/220.pdf}
	\end{minipage}
	\vspace{0.2cm}
	\fuentepropia
\end{minipage}
\vskip0.2cm


Con ellas puede ser más sencillo ver los límites de integración al momento de plantear las integrales múltiples en los seis órdenes posibles, que son, respectivamente, las siguientes:
\begin{multline*}
\int_{0}^{4}\int_{0}^{\sqrt{4-x}}\int_{0}^{2-y}dzdydx= \int_{0}^{2}\int_{0}^{4-y^{2}}\int_{0}^{2-y}dzdxdy\\
=\int_{0}^{2}\int_{0}^{2-z}\int_{0}^{4-y^{2}}dxdydz= \int_{0}^{2}\int_{0}^{2-y}\int_{0}^{4-y^{2}}dxdzdy\\
=\int_{0}^{2}\int_{0}^{4z-z^{2}}\int_{0}^{2-z}dydxdz+\int_{0}^{2}%
\int_{4z-z^{2}}^{4}\int_{0}^{\sqrt{4-x}}dydxdz \\
= \int_{0}^{4}\int_{2-\sqrt{4-x}}^{2}\int_{0}^{2-z}dydzdx+\int_{0}^{4}%
\int_{0}^{2-\sqrt{4-x}}\int_{0}^{\sqrt{4-x}}dydzdx= \frac{20}{3}.
\end{multline*}

En \verb"Mathematica", las siguientes instrucciones evalúan estas integrales.

\vspace{0.6em}
\noindent
\begin{tcolorbox}[breakable,enhanced,title=\bf Instrucción en \verb"Mathematica", top=2pt,bottom=2pt,boxsep=0pt,before skip=2pt,after skip=0pt]
\begin{Verbatim}[formatcom=\letraverbatim]
Integrate[1,{x,0,4},{z,2-Sqrt[4-x],2},{y,0,2-z}]+
Integrate[1,{x,0,4},{z,0,2-Sqrt[4-x]},{y,0,Sqrt[4-x]}]
Integrate[1,{z,0,2},{x,0,4*z-z^2},{y,0,2-z}]
+Integrate[1,{z,0,2},{x,4*z-z^2,4},{y,0,Sqrt[4-x]}]
Integrate[1,{z,0,2},{y,0,2-z},{x,0,4-y^2}]
Integrate[1,{y,0,2},{z,0,2-y},{x,0,4-y^2}]
Integrate[1,{x,0,4},{y,0,Sqrt[4-x]},{z,0,2-y}]
Integrate[1,{y,0,2},{x,0,4-y^2},{z,0,2-y}]
\end{Verbatim}
\end{tcolorbox}
\vspace{0.7em}
\noindent


El volumen también puede hallarse mediante el siguiente procedimiento.
\vskip0.2cm


%\begin{tcolorbox}
Tomaremos como base del sólido al triángulo que se forma en el plano $yz.$ Un punto sobre el eje $y$ tiene la forma $(t,0) $ y otro punto sobre la recta $z=2-y$ posee la forma $(t,2-t),$ con $t \in [0,2].$
Uniendo estos puntos con un segmento rectilíneo, se paramétrica este triángulo, esto es $(1-u)(t,0) +u(t,2-t)=(t,u(2-t)).$
Como $y=t$, sobre el cilindro se satisface que $x=4-t^2$.

\vskip0.2cm
Con el fin de parametrizar el sólido, se formará la combinación convexa de la forma
\[(1-w) \langle 0,t,u( 2-t)\rangle +w\langle 4-t^{2},t,u( 2-t) \rangle,\]
con lo cual se define la función
${\bf r}( t,u,w) = \langle w( 4-t^{2}),t,2u-tu\rangle ,$
en donde $t\in \lbrack 0,2],$ $u\in \lbrack 0,1],$ $w\in \lbrack 0,1].$
Al derivar parcialmente se obtienen los siguientes vectores
\[{\bf r}_{t} = \langle -2wt,1,-u\rangle,\]
\[{\bf r}_{u} = \langle 0,0,2-t\rangle\]
\[{\bf r}_{w} = \langle 4-t^{2},0,0\rangle.\]
\vskip0.2cm
El triple producto vectorial es $|{\bf r}_{t}\cdot ({\bf r}_{u}\times {\bf r}_{v}) | = ( t+2)( t-2) ^{2},$
integrándolo nos proporciona el volumen del sólido
\[\int_{0}^{1}\int_{0}^{1}\int_{0}^{2}(t+2)(t-2)^{2}dtdudw= \frac{20}{3}.\]
%\end{tcolorbox}

Este cálculo se puede evaluar con las siguientes instrucciones:

\vspace{0.6em}
\noindent
\begin{tcolorbox}[breakable,enhanced,title=\bf Instrucción en \verb"Mathematica", top=2pt,bottom=2pt,boxsep=0pt,before skip=2pt,after skip=0pt]
\begin{Verbatim}[formatcom=\letraverbatim]
r=(1-w)*{0,t,u*(2-t)}+w*{4-t^2,t,u*(2-t)};
Integrate[Dot[D[r,w],Cross[D[r,t],D[r,u]]],{u,0,1},
{t,0,2},{w,0,1}]
\end{Verbatim}
\end{tcolorbox}
\vspace{0.7em}
\noindent


\subsection*{Tres planos y un cilindro parabólico III}
Evaluar {\small $\ds \iiint_{E}6xydV$,} en donde $E$ es el sólido acotado por las figuras de las ecuaciones $z=1+x+y,\,\, y=\sqrt{x},\, x=1,\,\,z=0,\,\, y=0.$ Los siguientes gráficos muestran la región en mención.

\vskip0.2cm
\noindent
\begin{minipage}{\textwidth}
	\centering
	\captionof{figure}{Sólido determinado por $0 \leq z \leq 1 + x + y$, \quad $y^2 \leq x \leq 1$, \quad $0 \leq y \leq 1$.}
	\label{fig:solido_parabolico_inclinado}
	\begin{minipage}{0.34\linewidth}
		\centering
		\vspace{-6mm}
		\includegraphics[width=\linewidth]{Fig/221.pdf}
	\end{minipage}
	\hspace{1em}
	\begin{minipage}{0.16\linewidth}
		\centering
		\vspace{-6mm}
		\includegraphics[width=\linewidth]{Img/9b.pdf}
	\end{minipage}
	\hspace{1em}
	\begin{minipage}{0.16\linewidth}
		\centering
		\vspace{-6mm}
		\includegraphics[width=\linewidth]{Img/9a.pdf}
	\end{minipage}
	\vspace{0.2cm}
	\fuentepropia
\end{minipage}
	


Las siguientes instrucciones permiten obtener las figuras anteriores:

\vspace{0.6em}
\noindent
\begin{tcolorbox}[breakable,enhanced,title=\bf Instrucción en \verb"Mathematica", top=2pt,bottom=2pt,boxsep=0pt,before skip=2pt,after skip=0pt]
\begin{Verbatim}[formatcom=\letraverbatim]
ContourPlot3D[{z==1+x+y,x==y^2,x==1},{x,0,1},{y,0,1},
{z,0,3},Mesh->False,Ticks->{{1},{0,1},{0,1,2}},
ContourStyle->{Orange,Blend[{Yellow,Pink,Green}]}]
RegionPlot3D[0<z<1+x+y&&y^2<x<1&&0<y<1,{x,0,1},
{y,0,1},{z,0,3},PlotStyle->Blend[{Yellow,Orange,Red}],
PlotPoints->80,Mesh->False]
\end{Verbatim}
\end{tcolorbox}
\vspace{0.7em}
\noindent


La intersección del plano y el cilindro parabólico está dada por la ecuación \linebreak
{\small $z=1+x+\sqrt{x}$}, de donde {\small $x=(-\frac{1}{2}+\sqrt{z-3/4})^2$}. Además, la otra posible sustitución nos lleva a la expresión $z=1+y+y^2$, lo cual equivale a
{\small $y=-\frac{1}{2}+\frac{1}{2}\sqrt{4z-3}$}.
\vskip0.2cm

%\vspace*{-1.5\baselineskip} % Sube el bloque sin empalmar el texto superior

\noindent
\begin{minipage}{\textwidth}
\captionsetup{font=footnotesize}
	\centering
	\captionof{figure}{Visualización del sólido y sus proyecciones para el cálculo de volumen.}
	\label{fig:volumen_solido_proyecciones}
	\begin{minipage}{0.25\linewidth}
		\centering
		\vspace{-2mm}
		\includegraphics[width=\linewidth,height=1\linewidth]{Fig/222.pdf}
	\end{minipage}
	\hspace{1em}
	\begin{minipage}{0.25\linewidth}
		\centering
		\vspace{-2mm}
		\includegraphics[width=\linewidth,height=1\linewidth]{Fig/223.pdf}
	\end{minipage}
	\hspace{1em}
	\begin{minipage}{0.25\linewidth}
		\centering
		\vspace{-2mm}
		\includegraphics[width=\linewidth,height=1\linewidth]{Fig/224.pdf}
	\end{minipage}
	\vspace{0.3cm}
	\fuentepropia
\end{minipage}
	
\vskip0.2cm
La proyección en el plano $xy$ es una región acotada por una porción de parábola y dos rectas; además $z$ recorre la región desde el plano $z=0$ hasta el plano $z=1+x+y$.
Con lo anterior, la integral triple está determinada por:
\[\int_{0}^{1}\int_{0}^{\sqrt{x}}\int_{0}^{1+x+y}6xydzdydx=
\int_{0}^{1}\int_{y^2}^{1}\int_{0}^{1+x+y}6xydzdxdy= \frac{65}{28}.\]
\vskip0.2cm
En el caso de plantear las integrales en el orden $dxdydz$ y $dydxdz$ puede ser necesario presentar cuatro integrales triples, los gráficos ilustran el porqué se da esta situación.

{\small
\begin{multline*} \int_{0}^{1}\int_{0}^{1+x}\int_{0}^{\sqrt{x}}6xydydzdx+\int_{0}^{1}\int_{1+x}^{1+x+\sqrt[2]{x}}\int_{z-1-x}^{\sqrt{x}}6xydydzdx \\
\bigintsss_{0}^{1}\bigintsss_{0}^{1}\bigintsss_{0}^{\sqrt{x}}6xydydxdz
+\bigintsss_{1}^2\bigintsss_{\left( -\frac{1}{2}+\sqrt{z-\frac{3}{4}}\right) ^2}^{z-1}\bigintsss_{z-1-x}^{\sqrt{x}}6xydydxdz\\
+\bigintsss_{1}^2\bigintsss_{z-1}^{1}\bigintsss_{0}^{\sqrt{x}}6xydydxdz +\bigintsss_{2}^{3}\bigintsss_{\left( -\frac{1}{2}+\sqrt{z-\frac{3}{4}}\right)
^2}^{1}\bigintsss_{z-1-x}^{\sqrt{x}}6xydydxdz= \frac{65}{28}.
\end{multline*}
}

Integrando con respecto a $x$, las integrales que se deben evaluar son
{\footnotesize
\begin{multline*}
\bigintsss_{0}^{1}\bigintsss_{0}^{1}\bigintsss_{y^2}^{1}6xydxdydz+\bigintsss_{1}^{3}\bigintsss_{-\frac{1}{2}+\frac{1}{2}\sqrt{4z-3}}^{1}\bigintsss_{y^2}^{1}6xydxdydz+ \\ \bigintsss_{1}^{2}\bigintsss_{0}^{-\frac{1}{2}+\frac{1}{2}\sqrt{4z-3}}\bigintsss_{z-1-y}^{1}6xydxdydz+\bigintsss_{2}^{3}\bigintsss_{z-2}^{-\frac{1}{2}+\frac{1}{2}\sqrt{4z-3}}\bigintsss_{z-1-y}^{1}6xydxdydz \\
=\bigintsss_{0}^{1}\bigintsss_{0}^{1+y+y^2}\bigintsss_{y^2}^{1}6xydxdzdy+\bigintsss_{0}^{1}\bigintsss_{1+y+y^2}^{2+y}\bigintsss_{z-1-y}^{1}6xydxdzdy= \frac{65}{28}.
\end{multline*}
}

\vskip0.3cm
Las últimas integrales se evalúan en \verb"Mathematica" mediante la ejecución de la siguiente instrucción:

\vspace{0.6em}
\noindent
\begin{tcolorbox}[breakable,enhanced,title=\bf Instrucción en \verb"Mathematica", top=2pt,bottom=2pt,boxsep=0pt,before skip=2pt,after skip=0pt]
\begin{Verbatim}[formatcom=\letraverbatim]
Integrate[6*x*y,{y,0,1},{z,0,1+y+y^2},{x,y^2,1}]+
Integrate[6*x*y,{y,0,1},{z,1+y+y^2,2+y},{x,z-1-y,1}]
Integrate[6*x*y,{z,0,1},{y,0,1},{x,y^2,1}]+
Integrate[6*x*y,{z,1,3},{y,-1/2+Sqrt[4*z-3]/2,1},
{x,y^2,1}]+Integrate[6*x*y,{z,1,2},{y,0,-1/2+
Sqrt[4*z-3]/2},{x,z-1-y,1}]+Integrate[6*x*y,
{z,2,3},{y,z-2,-1/2+Sqrt[4*z-3]/2},{x,z-1-y,1}]
\end{Verbatim}
\end{tcolorbox}
\vspace{0.7em}
\noindent


Ahora se muestra otra manera de evaluar la integral triple.
\vskip0.2cm
%\begin{tcolorbox}[title=\bfseries Parametrizando el sólido]
\noindent
\begin{minipage}{\textwidth}
	\centering
	\captionof{figure}{Sólido determinado por $0 \leq z \leq 1 + x + y$, \quad $y^2 \leq x \leq 1$, \quad $0 \leq y \leq 1$.}
	\label{fig:solido_parabolico_inclinado_225_227}
	\begin{minipage}{0.29\linewidth}
		\centering
		\vspace{-6mm}
		\includegraphics[width=\linewidth]{Fig/225.pdf}
	\end{minipage}
	\hspace{1em}
	\begin{minipage}{0.29\linewidth}
		\centering
		\vspace{-6mm}
		\includegraphics[width=\linewidth]{Fig/226.pdf}
	\end{minipage}
	\hspace{1em}
	\begin{minipage}{0.29\linewidth}
		\centering
		\vspace{-6mm}
		\includegraphics[width=\linewidth]{Fig/227.pdf}
	\end{minipage}
	
	\vspace{1ex}
	\fuentepropia
\end{minipage}
	
\vskip0.3cm
En el plano $xy,$ con dos puntos $A=(t^2,t)$ y $B=(t^2,0),$ se forma la combinación $(1-u)B+uA=\langle t^2,ut \rangle,$
que parametriza la base del sólido con $t\in [0,1]$ y $u \in [0,1].$
El sólido se parametriza con una combinación convexa de dos puntos: uno sobre el plano $xy$ y otro sobre el plano $z=1+x+y$, esto es,
$\left( 1-w\right) \langle t^{2},tu,0\rangle +w\langle t^{2},tu,1+t^{2}+tu\rangle,$
${\bf r} (t,u,w) =\left\langle t^{2},tu,w\left(1+t^{2}+tu\right) \right\rangle ,$
con $t \in [0,1],$ $u \in [0,1],$ $w\in [0,1].$
Derivando parcialmente se obtienen los siguientes vectores
\begin{eqnarray*}
{\bf r}_t &=& \left\langle 2t,u,w\left( 2t+u\right) \right\rangle \\
{\bf r}_u &=& \left\langle 0,t,wt\right\rangle \\
{\bf r}_w &=& \left\langle 0,0,1+t^{2}+tu\right\rangle
\end{eqnarray*}
El triple producto vectorial es $|{\bf r}_t\cdot({\bf r}_u\times {\bf r}_w)| = | 2t^{2}\left(1+t^{2}+tu\right) |.$
%\begin{eqnarray*}
%|{\bf r}_t \cdot ({\bf r}_t \times {\bf r}_w )| &=& | 2t^{2}\left(1+t^{2}+tu\right) |
%\end{eqnarray*}
Integrando $6xy=6t^{3}u$, con los límites establecidos  se obtiene la integral
\[\ds \iiint_{E}6xydV = \int_{0}^{1}\int_{0}^{1}\int_{0}^{1}\left( 6t^{3}u\right) \left|2t^{2}\left( 1+t^{2}+tu\right) \right| dtdudw= \frac{65}{28}.\]
%\end{tcolorbox}

\subsection*{Tres planos y un cilindro parabólico IV}

Hallar el volumen del sólido acotado por
$z=1-y, \,\, y=\sqrt{x},\,\, x=0,\,\, z=0$.
La figura de este objeto se puede generar con la siguiente instrucción:

\vspace{0.6em}
\noindent
\begin{tcolorbox}[breakable,enhanced,title=\bf Instrucción en \verb"Mathematica", top=2pt,bottom=2pt,boxsep=0pt,before skip=2pt,after skip=0pt]
\begin{Verbatim}[formatcom=\letraverbatim]
RegionPlot3D[0<z<1-y&&0<x<y^2&&0<y<1,{x,0,1},{y,0,1},
{z,0,1},PlotPoints->90,Mesh->False,PlotStyle->Orange]
\end{Verbatim}
\end{tcolorbox}
\vspace{0.7em}
\noindent

% %La figura de dichas expresiones y el sólido se esbozan a continuación. 1
% \begin{center}
% \begin{pspicture}(-2,-0.9)(3,2.2) % \psgrid
% \psset{algebraic,Alpha=60,Beta=15,xunit=1.5,yunit=1.4}
% \pstThreeDSquare[linecolor=black!60,fillcolor=blue!25,fillstyle=solid](-1,1,0)(0,-1,1)(3,0,0)
% %\pstThreeDCoor[xMin=-1.4,xMax=3,yMin=-0.8,yMax=2,zMin=0,zMax=1.7,linecolor=black,IIIDticks,IIIDlabels,spotX=90,spotY=90,spotZ=180]
% \parametricplotThreeD[yPlotpoints=30,drawStyle=xyLines,linewidth=0.3pt,linecolor=gray](-0.5,1.3)(-0,1.2){t^2|t|u}
% \parametricplotThreeD[yPlotpoints=30,drawStyle=xyLines,linewidth=0.3pt,linecolor=gray](-0,1.2)(-0.5,1.3){u^2|u|t}
% \parametricplotThreeD[linewidth=1.2pt,linecolor=black](0,1){t^2|t|1-t}
% \parametricplotThreeD[linewidth=1.2pt,linecolor=black](0,1){t^2|t|0}
% % \pstThreeDLine[linewidth=1.2pt,linecolor=black](0,1,0)(0,0,0)(0,0,1)(0,1,0)(1,1,0)
% \pstThreeDLine[linewidth=0.3pt,linecolor=blue,linestyle=dashed](2,0,1)(2,0,0)(2,1,0)
% \uput[0](0.2,1.35){$z=1-y$}
% \uput[0](-1.2,1.2){$x=y^2$}
% \parametricplotThreeD[linewidth=0.6pt,linecolor=black,linestyle=dashed](0,1){(t-1)^2|0|t}
% \pstThreeDLine[linewidth=0.6pt,linecolor=black,linestyle=dashed](1,0,0)(1,1,0)
% \end{pspicture}

% \begin{pspicture}(-1.5,-0.8)(2,2.2) % \psgrid
% \psset{algebraic,Alpha=40,Beta=15,yunit=1.75,xunit=1.5}
% \pstThreeDCoor[spotX=-90,xMin=-0.8,xMax=1.3,yMin=-0.4,yMax=2,zMin=0,zMax=1.3,spotZ=180,IIIDticks,%
% IIIDlabels]
% \pscustom[fillstyle=crosshatch,hatchwidth=0.2pt,hatchcolor=blue,opacity=0.4]{
% \parametricplotThreeD(0,1){t^2|t|1-t}
% \pstThreeDLine(1,1,0)(0,1,0)(0,0,1)
% }
% \pscustom[fillstyle=solid,fillcolor=lightgray!70,opacity=0.5]{
% \parametricplotThreeD(0,1){t^2|t|1-t}
% \parametricplotThreeD(1,0){t^2|t|0}
% \pstThreeDLine(0,0,0)(0,0,1)
% }
% \psset{linewidth=0.7pt}
% \pstThreeDLine(0,0,1)(0,1,0)(1,1,0)
% \pstThreeDLine[linestyle=dashed](1,1,0)(1,0,0)
% \parametricplotThreeD(0,1){t^2|t|1-t}
% \parametricplotThreeD[linestyle=dashed](0,1){t^2|t|0}
% \pstThreeDPut(0,-0.2,1){$1$}
% \pstThreeDPut(0.8,-0.3,0){$1$}
% \end{pspicture}
% \includegraphics[height=4.5cm]{Img/62a.pdf}
% \captionof{figure}{\small Sólido determinado por $0\leq z \leq 1-y,$ $0\leq x\leq y^2,$ $0\leq y\leq 1$}
% \end{center}

% \begin{center}
% \begin{pspicture}(-1.5,-0.8)(2,2.2) % \psgrid
% \psset{algebraic,Alpha=40,Beta=15,yunit=1.75,xunit=1.5}
% \pstThreeDCoor[spotX=-90,xMin=-0.8,xMax=1.3,yMin=-0.4,yMax=2,zMin=0,zMax=1.3,spotZ=180,IIIDticks,%
% IIIDlabels]
% \pscustom[fillstyle=crosshatch,hatchwidth=0.2pt,hatchcolor=blue,opacity=0.4]{
% \parametricplotThreeD(0,1){t^2|t|1-t}
% \pstThreeDLine(1,1,0)(0,1,0)(0,0,1)
% }
% \pscustom[fillstyle=solid,fillcolor=lightgray!70,opacity=0.5]{
% \parametricplotThreeD(0,1){t^2|t|1-t}
% \parametricplotThreeD(1,0){t^2|t|0}
% \pstThreeDLine(0,0,0)(0,0,1)
% }
% \psset{linewidth=0.7pt}
% \pstThreeDLine(0,0,1)(0,1,0)(1,1,0)
% \pstThreeDLine[linestyle=dashed](1,1,0)(1,0,0)
% \parametricplotThreeD(0,1){t^2|t|1-t}
% \parametricplotThreeD[linestyle=dashed](0,1){t^2|t|0}
% \pstThreeDPut(0,-0.2,1){$1$}
% \pstThreeDPut(0.8,-0.3,0){$1$}
% \end{pspicture}
% \includegraphics[height=4.5cm]{Img/62a.pdf}
% \captionof{figure}{Sólido determinado por $0\leq z \leq 1-y,$ $0\leq x\leq y^2,$ $0\leq y\leq 1.$}
% \end{center}
%\vspace{0.2cm}

\noindent
\begin{minipage}{\textwidth}
	\centering
	\captionof{figure}{Sólido determinado por $0\leq z \leq 1-y,$ $0\leq x\leq y^2,$ $0\leq y\leq 1.$}
	\begin{minipage}{0.30\linewidth}
		\centering
		\vspace{-15mm}
		\includegraphics[width=\linewidth]{Fig/303.pdf}
	\end{minipage}
	\hspace{0.2em}
	\begin{minipage}{0.30\linewidth}
		\centering
		\vspace{-15mm}
		\includegraphics[width=\linewidth]{Fig/302.pdf}
	\end{minipage}
	\hspace{0.2em}
	\begin{minipage}{0.20\linewidth}
		\centering
		\vspace{-15mm}
		\includegraphics[width=\linewidth]{Imgn/62a.pdf}
	\end{minipage}
	
	\vspace{-8mm}
	
	%\vspace{0.3em}
	\fuentepropia
\end{minipage}
\vskip0.2cm

En el plano $xz$, reemplazando $y=\sqrt{x}$ en la ecuación del plano $y=1-z,$
se obtiene $z=1-\sqrt{x},$ lo cual equivale a $x=(1-z)^2.$
En el plano $xy$ los límites se determinan con las expresiones $x=y^2,$ $y=1$ y $x=0$.
En el plano $yz$ la región proyectada es un triángulo acotado por los ejes coordenados y la recta $z=1-y$.
\vskip0.2cm

\noindent
\begin{minipage}{\textwidth}
	\centering
	\captionof{figure}{Proyecciones en los planos coordenados.}
	\begin{minipage}{0.23\linewidth}
		\centering
		\vspace{-10mm}
		\includegraphics[width=\linewidth]{Fig/304.pdf}
	\end{minipage}
	\hspace{0.2em}
	\begin{minipage}{0.23\linewidth}
		\centering
		\vspace{-10mm}
		\includegraphics[width=\linewidth]{Fig/305.pdf}
	\end{minipage}
	\hspace{0.2em}
	\begin{minipage}{0.23\linewidth}
		\centering
		\vspace{-10mm}
		\includegraphics[width=\linewidth]{Fig/306.pdf}
	\end{minipage}
	
	%\vspace{0.3em}
	
	\vspace{-5mm}
	\fuentepropia
\end{minipage}


Estos elementos ayudan a plantear las integrales que determinan el volumen del sólido
\begin{multline*}
\int_{0}^{1}\int_{\sqrt{x}}^{1}\int_{0}^{1-y}dzdydx=
\int_{0}^{1}\int_{0}^{y^2}\int_{0}^{1-y}dzdxdy \\
= \int_{0}^{1}\int_{0}^{1-y}\int_{0}^{y^2}dxdzdy= \int_{0}^{1}\int_{0}^{1-z}\int_{0}^{y^2}dxdydz \\
= \int_{0}^{1}\int_{0}^{\left( 1-z\right)^2}\int_{\sqrt{x}}^{1-z}dydxdz=
\int_{0}^{1}\int_{0}^{1-\sqrt{x}}\int_{\sqrt{x}}^{1-z}dydzdx= \frac{1}{12}.
\end{multline*}
Enseguida se presentan las indicaciones para evaluar estas integrales.

\vspace{0.6em}
\noindent
\begin{tcolorbox}[breakable,enhanced,title=\bf Instrucción en \verb"Mathematica", top=2pt,bottom=2pt,boxsep=0pt,before skip=2pt,after skip=0pt]
\begin{Verbatim}[formatcom=\letraverbatim]
Integrate[1,{z,0,1},{x,0,(1-z)^2},{y,Sqrt[x],1-z}]
Integrate[1,{x,0,1},{y,Sqrt[x],1},{z,0,1-y}]
Integrate[1,{y,0,1},{z,0,1-y},{x,0,y^2}]
\end{Verbatim}
\end{tcolorbox}
\vspace{0.7em}
\noindent

El mismo valor se puede conseguir siguiendo el siguiente proceso.


%\begin{tcolorbox}[enhanced,colback=white,colframe=black,title=\bfseries Parametrizando el sólido,fonttitle=\bfseries,boxrule=0.8pt,arc=2pt,left=4pt, right=4pt,top=4pt, bottom=4pt]
		
%\vspace*{-0.2em}
Sobre la región de integración del plano $xy$ con dos puntos $P(t^2,t)$ y 
$Q(t^2,1)$ se forma una combinación convexa $(1-u)P+uQ$, que describe la región sombreada. 
Simplificando, se obtiene $\langle t^2,\,u+(1-u)t\rangle$, con 
$t\in[0,1]$, $u\in[0,1].$
	
\vskip0.2cm
\noindent
\begin{minipage}{\textwidth}
	\centering
	\captionof{figure}{\small  $0 \le x \le y^2,\ 0 \le x \le 1.$}
	\begin{minipage}{0.36\linewidth}
		\centering
		\vspace{-18mm}
		\includegraphics[width=\linewidth]{Fig/307.pdf}
	\end{minipage}
	
	\vspace{-12mm}
	\fuentepropia
\end{minipage}
\vskip0.2cm
	
\vspace{0.5em}
\noindent
Sobre el plano $z=1-y$ se satisface que la coordenada $z$ tiene la forma 
$z=(1-t)(1-u).$ El sólido se parametriza mediante una combinación convexa de dos puntos de la forma 
$\langle t^2,\,u+(1-u)t,\,0\rangle$ y $\langle t^2,\,u+(1-u)t,\,(1-t)(1-u)\rangle$, 
lo cual permite definir la función
\[
\mathbf{r}(t,u,w)=\langle t^2,\,t+(1-u)t,\,(1-t)(u-1)w\rangle,
\]
con $t,u,w\in[0,1].$ Con esta función se obtienen los siguientes vectores:

\vspace{-8mm}
\begin{align*}
\mathbf{r}_t &= \langle 2t,\,1-u,\,-(1-u)w\rangle, &
\mathbf{r}_w &= \langle 0,\,0,\,(1-t)(1-u)\rangle,\\
\mathbf{r}_u &= \langle 0,\,1-t,\,-(1-t)w\rangle, &
\mathbf{r}_u\times\mathbf{r}_t &= \langle 0,\,-2tw(1-t),\,-2t(1-t)\rangle.
\end{align*}
	
\noindent
Con lo cual 
\[
|\mathbf{r}_w\cdot(\mathbf{r}_u\times\mathbf{r}_t)|
=|2t(1-t)^2(1-u)|,
\]
y su integración en los límites dados proporciona el volumen del sólido:
{\small
\[
\int_0^1\!\!\int_0^1\!\!\int_0^1
|2t(1-t)^2(1-u)|\,dt\,du\,dw
=\tfrac{1}{12}
\]
}	
%\end{tcolorbox}


Este procedimiento se realiza con la ayuda de las siguientes instrucciones:

\vspace{0.6em}
\noindent
\begin{tcolorbox}[breakable,enhanced,title=\bf Instrucción en \verb"Mathematica", top=2pt,bottom=2pt,boxsep=0pt,before skip=2pt,after skip=0pt]
\begin{Verbatim}[formatcom=\letraverbatim]
{x,y}=(1-u)*{t^2,t}+u*{t^2,1};
z=1-y;
R=Factor[(1-w)*{x,y,z}+w*{x,y,0}];
Integrate[Dot[D[R,w],Cross[D[R,u],D[R,t]]],
{w,0,1},{u,0,1},{t,0,1}]
\end{Verbatim}
\end{tcolorbox}
\vspace{0.7em}
\noindent

\subsection*{Un cilindro y dos planos III}
La ecuación $x^{2}+z^{2}=1$ determina un cilindro de radio $1$ con eje de simetría el eje $y$, esta superficie es cortada por los planos con ecuaciones $y=x$ y $y=2x.$

\vskip0.2cm
\noindent
\begin{minipage}{\textwidth}
	\centering
	\captionof{figure}{Superficies $x^{2} + z^{2} = 1$, $y = x$ y $y = 2x$, con $y \geq 0$.}
	\label{fig:cilindro_planos_ypositivo}
	\begin{minipage}{0.17\linewidth}
		\centering
		\includegraphics[width=\linewidth]{Img/27b.pdf}
	\end{minipage}
	\hspace{1em}
	\begin{minipage}{0.17\linewidth}
		\centering
		\includegraphics[width=\linewidth]{Img/27a.pdf}
	\end{minipage}
	\hspace{1em}
	\begin{minipage}{0.17\linewidth}
		\centering
		\includegraphics[width=\linewidth]{Img/27c.pdf}
	\end{minipage}
	\hspace{1em}
	\begin{minipage}{0.17\linewidth}
		\centering
		\includegraphics[width=\linewidth]{Img/27d.pdf}
	\end{minipage}
	
	\vspace{0.5em}
	\fuentepropia
\end{minipage}
\vskip0.2cm

Con las siguientes instrucciones se obtienen las figuras anteriores:

\vspace{0.6em}
\noindent
\begin{tcolorbox}[breakable,enhanced,title=\bf Instrucción en \verb"Mathematica", top=2pt,bottom=2pt,boxsep=0pt,before skip=2pt,after skip=0pt]
\begin{Verbatim}[formatcom=\letraverbatim]
ContourPlot3D[{1==x^2+z^2,y==x,y==2*x},{x,0,1},
{y,0,2},{z,-1,1},PlotPoints->50,Mesh->False,
ContourStyle->{Orange,Cyan,LightRed},
Ticks->{{0,1},{0,1,2},{-1,0}}]
RegionPlot3D[-Sqrt[1-x^2]<z<Sqrt[1-x^2]&&x<y
<2*x&&0<x<1,{x,0,1},{y,0,2},{z,-1,1},PlotPoints->90,
Mesh->False,Ticks->{{0,1},{0,1,2},{-1,0,1}}]
\end{Verbatim}
\end{tcolorbox}
\vspace{0.7em}
\noindent

Con el fin de plantear las integrales triples que permiten obtener el volumen del sólido presentamos las siguientes figuras.

\vskip0.2cm
\noindent
\begin{minipage}{\textwidth}
	\centering
	\captionof{figure}{Proyección del sólido en los planos coordenados.}
	\begin{minipage}{0.30\linewidth}
		\centering
		\vspace{-9mm}		\includegraphics[width=\linewidth]{Fig/308.pdf}
	\end{minipage}
	\hspace{1em}
	\begin{minipage}{0.30\linewidth}
		\centering
		\vspace{-9mm}
		\includegraphics[width=\linewidth]{Fig/309.pdf}
	\end{minipage}
	\hspace{1em}
	\begin{minipage}{0.30\linewidth}
		\centering
		\vspace{-9mm}
		\includegraphics[width=\linewidth]{Fig/310.pdf}
	\end{minipage}
	
	\vspace{-8mm}
	\parbox{0.95\linewidth}{\centering\fontsize{9pt}{10pt}\selectfont\textbf{Fuente:} elaboración propia}
\end{minipage}
\vskip0.2cm

% \begin{figure}[H]......2
% \centering

% \vspace{0.4cm}

% % Primer gráfico: plano xz
% \begin{pspicture}(-0.6,-1.4)(2.5,2.1)
% \psset{algebraic}
% \pscustom[fillstyle=vlines,hatchcolor=gray!40,hatchsep=1.4pt,hatchwidth=0.7pt,linestyle=none]{%
%   \parametricplot{-1.57}{1.57}{cos(t),sin(t)}%
% }
% \psaxes[linecolor=black]{->}(0,0)(-0.6,-1.5)(2,2.1)[$x$,-90][$z$,0]
% \parametricplot[linecolor=red!80!black,arrows=->]{-1.6}{1.9}{cos(t),sin(t)}
% \psline[linecolor=red!80!black,arrows=-](0,1)(0,-1)
% \rput(1.6,1.3){$x^2+z^2=1$}
% \end{pspicture}

% \vspace{0.5cm}

% % Segundo gráfico: plano yz
% \begin{pspicture}(-1.2,-1.4)(3.2,2.1)
% \psset{algebraic}
% \pscustom[fillstyle=vlines,hatchcolor=gray!40,hatchsep=1.4pt,hatchwidth=0.7pt,linestyle=none]{%
%   \parametricplot{-1.57}{1.57}{cos(t),sin(t)}%
% }
% \pscustom[fillstyle=solid,fillcolor=gray!40,linestyle=none]{%
%   \parametricplot{-1.57}{1.57}{cos(t),sin(t)}%
%   \parametricplot{1.57}{1.05}{2*cos(t),sin(t)}%
%   \psline(1,0)(1,0.87)%
%   \parametricplot{-1.05}{-1.57}{2*cos(t),sin(t)}%
% }
% \pscustom[fillstyle=hlines,hatchcolor=gray!40,linestyle=none,hatchsep=1.4pt,hatchwidth=0.7pt]{%
%   \parametricplot{-1.05}{1.05}{2*cos(t),sin(t)}%
% }
% \psaxes[linecolor=black]{->}(0,0)(-0.6,-1.2)(3,1.75)[$y$,-90][$z$,0]
% \parametricplot[linecolor=red!80!black,arrows=-]{-1.57}{1.57}{cos(t),sin(t)}
% \parametricplot[linecolor=red!80!black,arrows=->]{-1.57}{1.7}{2*cos(t),sin(t)}
% \psline[linecolor=red!80!black,arrows=-](0,1)(0,-1)
% \psline[linecolor=red!80!black,arrows=-,linewidth=0.6pt,linestyle=dashed](1,-0.87)(1,0.87)
% \psline[linecolor=black,arrows=->](0.65,-0.8)(1,-1.2)
% \rput(2,-1.2){$y^2+z^2=1$}
% \rput(1.6,1.3){$y^2+4z^2=4$}
% \end{pspicture}

% \vspace{0.5cm}

% % Tercer gráfico: plano xy
% \begin{pspicture}(-1,-0.6)(2,2.8)
% \psset{algebraic,xunit=1.2}
% \psaxes[linecolor=black,linewidth=0.6pt]{->}(0,0)(-0.5,-0.5)(2,2.5)[$x$,90][$y$,0]
% \psline[fillstyle=vlines,hatchcolor=gray!50,hatchsep=1pt](0,0)(1,1)(0.5,1)(0,0)
% \psline[fillstyle=hlines,hatchcolor=gray!50,hatchsep=1pt](1,1)(1,2)(0.5,1)
% \psline[linecolor=red!80!black,arrows=-,linewidth=0.6pt,linestyle=dashed](1,1)(0.5,1)
% \parametricplot[linecolor=red!80!black,arrows=->]{-0.25}{1.2}{t,t}
% \parametricplot[linecolor=red!80!black,arrows=->]{-0.25}{1.2}{t,2*t}
% \parametricplot[linecolor=red!80!black,arrows=-]{1}{2}{1,t}
% \rput(1.8,1.2){$y=x$}
% \rput(1.8,2.2){$y=2x$}
% \end{pspicture}

% \caption{Proyección del sólido en los planos coordenados.}
% \end{figure}

Con esta representación determinamos las integrales que nos aportan el volumen del sólido. Al integrar primero con respecto a $y$, esta recorre puntos desde el plano $y=x$ hasta el plano $y=2z$. Los límites para $x$ y $z$ se encuentran en la semicircunferencia unitaria en el plano $xz$.
\vskip0.2cm
\[\int_{-1}^{1}\int_{0}^{\sqrt{1-z^{2}}}\int_{x}^{2x}dydxdz= \int_{0}^{1}\int_{-\sqrt{1-x^{2}}}^{\sqrt{1-x^{2}}}\int_{x}^{2x}dydzdx= \frac{2}{3}\]
\vskip0.2cm
Al integrar con respecto a $z$, los planos $y=x$ y $y=2x$ son paralelos al eje $z$ por lo que los límites para $z$ van desde la parte negativa del cilindro a la parte positiva de este. Por tanto, la integral en este caso es:
\begin{multline*} \int_{0}^{1}\int_{x}^{2x}\int_{-\sqrt{1-x^{2}}}^{\sqrt{1-x^{2}}}dzdydx= \int_{0}^{1}\int_{\frac{y}{2}}^{y}\int_{-\sqrt{1-x^{2}}}^{\sqrt{1-x^{2}}}dzdxdy \\ +\int_{1}^2\int_{\frac{y}{2}}^{1}\int_{-\sqrt{1-x^{2}}}^{\sqrt{1-x^{2}}}dzdxdy= \frac{2}{3}. 
\end{multline*}

Integrando con respecto a $x$, la región se divide en varias partes, según como se observe. La proyección en el plano $yz$ ayuda a visualizar estos elementos:
{\small
\[\bigintsss_{-1}^{1}\bigintsss_{0}^{\sqrt{1-z^{2}}}\bigintsss_{\frac{y}{2}}^{y}dxdydz+\bigintsss_{-1}^{1}\bigintsss_{\sqrt{1-z^{2}}}^{\sqrt{4-4z^{2}}}\bigintsss_{\frac{y}{2}}^{\sqrt{1-z^{2}}}dxdydz= \frac{2}{3}.\]
}

En el orden $dxdzdy$ la evaluación respectiva es:

\vspace{-8mm}
{\small
\begin{multline*}
\bigintsss_{0}^{1}\bigintsss_{-\sqrt{1-y^{2}}}^{\sqrt{1-y^{2}}}\bigintsss_{\frac{y}{2}}^{y}dxdzdy+\bigintsss_{0}^{1}\bigintsss_{\sqrt{1-y^{2}}}^{\sqrt{1-\frac{y^{2}}{4}}}\bigintsss_{\frac{y}{2}}^{\sqrt{1-z^{2}}}dxdzdy +\\
\bigintsss_{0}^{1}\bigintsss_{-\sqrt{1-\frac{y^{2}}{4}}}^{-\sqrt{1-y^{2}}}\bigintsss_{\frac{y}{2}}^{\sqrt{1-z^{2}}}dxdzdy+\bigintsss_{1}^2\bigintsss_{-\sqrt{1-\frac{y^{2}}{4}}}^{\sqrt{1-\frac{y^{2}}{4}}}\bigintsss_{\frac{y}{2}}^{\sqrt{1-z^{2}}}dxdzdy= \frac{2}{3}.
\end{multline*}
}

Las integrales anteriores en los órdenes $dydxdz,$ $dzdxdy$ y $dxdzdy$, se evalúan en \verb"Mathematica"
ejecutando las siguientes instrucciones:

\vspace{0.6em}
\noindent
\begin{tcolorbox}[breakable,enhanced,title=\bf Instrucción en \verb"Mathematica", top=2pt,bottom=2pt,boxsep=0pt,before skip=2pt,after skip=0pt]
\begin{Verbatim}[formatcom=\letraverbatim]
Integrate[1,{z,-1,1},{x,0,Sqrt[1-z^2]},{y,x,2x}]
Integrate[1,{y,0,1},{x,y/2,y},{z,-Sqrt[1-x^2],
Sqrt[1-x^2]}]+Integrate[1,{y,1,2},{x,y/2,1},
{z,-Sqrt[1-x^2],Sqrt[1-x^2]}]
Integrate[1,{y,0,1},{z,-Sqrt[1-y^2],Sqrt[1-y^2]},
{x,y/2,y}]+Integrate[1,{y,0,1},{z,Sqrt[1-y^2],
Sqrt[1-y^2/4]},{x,y/2,Sqrt[1-z^2]}]+Integrate[1,
{y,0,1},{z,-Sqrt[1-y^2/4],-Sqrt[1-y^2]},{x,y/2,
Sqrt[1-z^2]}]+Integrate[1,{y,1,2},{z,-Sqrt[1-y^2/4],
Sqrt[1-y^2/4]},{x,y/2,Sqrt[1-z^2]}]
\end{Verbatim}
\end{tcolorbox}
\vspace{0.7em}
\noindent

También se puede hacer un cambio a coordenadas polares, $x=r\cos \theta $ y $z=r\sin \theta $,
de aquí se sigue lo siguiente:
{\small
\[\int_{-\pi /2}^{\pi /2}\int_{0}^{1}\int_{r\cos \theta }^{2r\cos \theta}rdydrd\theta = \frac{2}{3}.\]
}

\vspace{-2mm}
Ahora se ilustra otra posibilidad de obtener el mismo cálculo.

\vskip0.2cm
%\begin{tcolorbox}[enhanced,colback=white,colframe=black,title=\bfseries Parametrizando el sólido,fonttitle=\bfseries,boxrule=0.8pt,arc=2pt,left=4pt, right=4pt,top=4pt,bottom=4pt]


En el plano $xz$, la proyección del sólido es un semicírculo que se parametriza como:
{\small
\[ \la u\sin t, 0, u\cos t \ra, \quad t\in [-\pi /2, \pi/2], \quad u\in [0,1]. \]
}
Con un punto $P(u\sin t, u\sin t, u\cos t)$ sobre el plano $y=x$, y otro punto $Q(u\sin t, 2u\sin t, u\cos t)$ sobre el plano $y=2x.$

\vskip0.2cm
\noindent
\begin{minipage}{\textwidth}
	\centering
	\captionof{figure}{\small ${\bf r}(u,t)=\left\langle u\sin t,u\cos t\right\rangle$}
	\begin{minipage}{0.34\linewidth}
		\centering
		\vspace{-16mm}
		\includegraphics[width=\linewidth]{Fig/311.pdf}
	\end{minipage}
	
	\vspace{-10mm}
	\fuentepropia
\end{minipage}
\vskip0.2cm


\noindent
Se define una parametrización del sólido, que se obtiene mediante una combinación convexa
$\left( 1-v\right) \left\langle u\sin t,2u\sin t,u\cos t\right\rangle +v\left\langle u\sin t,u\sin t,u\cos t\right\rangle,$
entre dos puntos $P$ y $Q.$
$ {\bf r}(t,u,v) = \left\langle u\sin t,u\left( 2-v\right) \sin t,u\cos t\right\rangle ,$
en donde $t\in [ -\pi /2,\pi/2]$, $u\in [ 0,1]$ y $v\in [0,1].$

\vskip0.2cm
Esta función nos permite hallar los siguientes vectores:
%\vspace{-2mm}
{\small
\begin{eqnarray*}
{\bf r}_{t} &=& \left\langle u\cos t,u\left(2-v\right) \cos t,-u\sin t\right\rangle \\
{\bf r}_{u} &=& \left\langle \sin t,\left( 2-v\right) \sin t,\cos t\right\rangle \\
{\bf r}_{v} &=& \left\langle 0,-u\sin t,0\right\rangle.
\end{eqnarray*}
}

%\vspace{-2mm}
El producto vectorial triple esta dado por $\left|{\bf r}_{t}\cdot \left({\bf r}_{u}\times {\bf r}_{v}\right) \right| =\left| u^{2}\sin t\right|,$
que se integra con los límites establecidos y nos proporciona el volumen del sólido en mención
{\small
\[\int_{0}^{1}\int_{0}^{1}\int_{-\pi /2}^{\pi /2}\left| u^{2}\sin t\right|dtdudv= \frac{2}{3}.\]
}
%\end{tcolorbox}
%\vskip0.2cm
%\vspace{-2mm}
El cálculo anterior se obtiene con las siguientes instrucciones:

\vspace{0.6em}
\noindent
\begin{tcolorbox}[breakable,enhanced,title=\bf Instrucción en \verb"Mathematica", top=2pt,bottom=2pt,boxsep=0pt,before skip=2pt,after skip=0pt]
\begin{Verbatim}[formatcom=\letraverbatim]
{x,y,z}={u*Sin[t],u*Sin[t],u*Cos[t]};
r=(1-v)*{x,y,z}+v*{x,2*y,z};
A=Abs[Dot[D[r,v],Cross[D[r,u],D[r,t]]]];
Integrate[A,{u,0,1},{t,-Pi/2,Pi/2},{v,0,1}]
\end{Verbatim}
\end{tcolorbox}
\vspace{0.7em}
\noindent


\subsection*{Un paraboloide y cuatro planos} \index{Paraboloide}
El paraboloide con ecuación  $z = 4 - x^2 - y^2$ y el plano $x = 2-2y$ se intersecan en una región del espacio que a continuación se muestra.

\vskip0.2cm
\noindent
\begin{minipage}{\textwidth}
	\centering
	\captionof{figure}{Sólido definido por $0 \leq z \leq 4 - x^2 - y^2$, \quad $0 \leq x \leq 2 - 2y$, \quad $0 \leq y \leq 1$.}
	\label{fig:solido_parabolico_inclinado_44}
	\begin{minipage}{0.17\linewidth}
		\centering
		\includegraphics[width=\linewidth]{Img/44d.pdf}
	\end{minipage}
	\hspace{1em}
	\begin{minipage}{0.17\linewidth}
		\centering
		\includegraphics[width=\linewidth]{Img/44a.pdf}
	\end{minipage}
	\hspace{1em}
	\begin{minipage}{0.17\linewidth}
		\centering
		\includegraphics[width=\linewidth]{Img/44c.pdf}
	\end{minipage}
	
	\vspace{0.5em}
	\fuentepropia
\end{minipage}
\vskip0.2cm

Estas superficies determinan un sólido en el primer octante.
Las anteriores figuras se obtienen ejecutando las siguientes instrucciones:

%\vskip0.2cm
\vspace{0.6em}
\noindent
\begin{tcolorbox}[breakable,enhanced,title=\bf Instrucción en \verb"Mathematica", top=2pt,bottom=2pt,boxsep=0pt,before skip=2pt,after skip=0pt]
\begin{Verbatim}[formatcom=\letraverbatim]
ContourPlot3D[{z==4-x^2-y^2,x==2-2*y,y==0},{x,0,2},
{y,0,2},{z,0,4},Ticks->{{0,2},{2,4},{0,2,4}},
MeshStyle->Gray]
RegionPlot3D[0<z<4-x^2-y^2&&0<x<2-2*y&&0<y<1,{x,0,2},
{y,0,1},{z,0,4},PlotPoints->90,Mesh->False,
PlotStyle->RGBColor[0.7,0.2,0.1]]
ContourPlot3D[{z==4-x^2-y^2,x==2-2*y},{x,-2,2},{y,-2,
2},{z,0,4},ContourStyle->{Directive[Opacity[0.8]],
Directive[Cyan,Opacity[0.7]]},Mesh->False]
\end{Verbatim}
\end{tcolorbox}
\vspace{0.7em}
\noindent

Plantearemos ahora las integrales triples que se requieren evaluar para determinar el volumen de esta región. Para determinar la proyección del sólido en los planos coordenados se debe eliminar una de las variables adecuadamente escogidas. En el plano $xy,$ (que equivale a $z=0$),
la proyección está dada por las expresiones $x = 2 -2y\,$ y $\,x^2 + y^2=4$.
\vskip0.3cm

En el plano $yz,$ la expresión que determina la región de integración se obtiene eliminando $x$ de la ecuación del paraboloide, esto es 
$z =4-(2-2y)^2-y^2$ por lo tanto
$z = 8y-5y^2,$ entonces $y = \frac{4}{5}\pm \frac{1}{5}\sqrt{16-5z}.$


\vskip0.3cm
\noindent
\begin{minipage}{\textwidth}
	\centering
	\captionof{figure}{Proyección del sólido en dos planos coordenados}
	\begin{minipage}{0.37\linewidth}
		\centering
		\vspace{-19mm}
		\includegraphics[width=\linewidth]{Fig/312.pdf}
	\end{minipage}
	\hspace{1em} 
	\begin{minipage}{0.37\linewidth}
		\centering
		\vspace{-19mm}
		\includegraphics[width=\linewidth]{Fig/313.pdf}
	\end{minipage}
	
	\vspace{-14mm}
	\fuentepropia
\end{minipage}
\vskip0.3cm



Al integrar primera con respecto a $z,$ esta toma valores desde el plano $xy$ hasta el paraboloide y los límites de integración con respecto a $x$ e $y$ se pueden obtener del primero de los gráficos anteriores. Mientras que al integrar primero con respecto,
a $x,$ esta toma valores desde el plano $yz$ hasta el paraboloide, en un caso, y desde el plano $yz$ hasta el plano $x=2-2y$ en el otro caso. Por lo tanto, el volumen está dado por las siguientes integrales.

\vspace{-7mm}
{\small
\begin{multline*}
\bigintsss_{0}^{1}\bigintsss_{0}^{2-2y}\bigintsss_{0}^{4-x^2-y^2}dzdxdy= \bigintsss_{0}^2\bigintsss_{0}^{1-x/2}\bigintsss_{0}^{4-x^2-y^2}dzdydx \\
=\bigintsss_{0}^{1}\bigintsss_{8y-5y^2}^{4-y^2}\bigintsss_{0}^{\sqrt{4-z-y^2}}dxdzdy+\bigintsss_{0}^{1}\bigintsss_{0}^{8y-5y^2}\bigintsss_{0}^{2-2y}dxdzdy = \frac{19}{6}.
\end{multline*}
}

La integral en el orden $dxdydz$ requiere un poco de observación adicional.

La expresión $s = 8y - 5y^2$ determina una parábola en el plano $yz$, su valor máximo se puede hallar derivando, esto es
$z^\prime =8-10y=0,$ entonces $y=\frac{4}{5},$ por tanto el valor máximo de $z$ es.
$8 \left( \frac{4}{5} \right)- 5\left( \frac{4}{5} \right)^2 = \frac{16}{5}.$
\vskip0.3cm
Esta condición hace que la proyección en el plano $yz$ se divida en varias regiones que la siguiente figura ilustra.


\vskip0.2cm
\noindent
\begin{minipage}{\textwidth}
	\centering
	\captionof{figure}{Proyección en el plano $yz$.}
	\begin{minipage}{0.38\linewidth}
		\centering
		\vspace{-15mm}
		\includegraphics[width=\linewidth]{Fig/314.pdf}
	\end{minipage}
	
	\vspace{-10mm}
	\fuentepropia
\end{minipage}
%\vskip0.3cm


De la expresión $z=8y-5y^2$ se sigue que
$y=\frac{4 \pm \sqrt{16-5z}}{5}$ y la expresión $z= 4-y^2$
equivale $y=\sqrt{4-z},$ con los elementos anteriormente descritos, el volumen del sólido está dado por:
%\vspace{-2mm}
{\footnotesize
\begin{multline*}
\bigintsss_{0}^{\frac{16}{5}}\bigintsss_{0}^{\frac{4}{5}-\frac{1}{5}\sqrt{16-5z}}\bigintsss_{0}^{\sqrt{4-z-y^2}}dxdydz
+\bigintsss_{\frac{16}{5}}^{4}\bigintsss_{0}^{\sqrt{4-z}}\bigintsss_{0}^{\sqrt{4-z-y^2}}dxdydz \\
+\bigintsss_{3}^{\frac{16}{5}}\bigintsss_{\frac{4}{5}+\frac{1}{5}\sqrt{16-5z}}^{\sqrt{4-z}}\bigintsss_{0}^{\sqrt{4-z-y^2}}dxdydz +\bigintsss_{0}^{3}\bigintsss_{\frac{4}{5}-\frac{1}{5}\sqrt{16-5z}}^{1}\bigintsss_{0}^{2-2y}dxdydz\\
+\bigintsss_{3}^{\frac{16}{5}}\bigintsss_{\frac{4}{5}-\frac{1}{5}\sqrt{16-5z}}^{\frac{4}{5}+\frac{1}{5}\sqrt{16-5z}}\bigintsss_{0}^{2-2y}dxdydz= \frac{19}{6}
\end{multline*}
}

En \verb"Mathematica" estas integrales se evalúan con la siguiente instrucción:

\vspace{0.6em}
\noindent
\begin{tcolorbox}[breakable,enhanced,title=\bf Instrucción en \verb"Mathematica", top=2pt,bottom=2pt,boxsep=0pt,before skip=2pt,after skip=0pt]
\begin{Verbatim}[formatcom=\letraverbatim]
Integrate[1,{z,0,16/5},{y,0,(4-Sqrt[16-5*z])/5},
{x,0,Sqrt[4-z-y^2]}]+Integrate[1,{z,3,16/5},
{y,(4+Sqrt[16-5*z])/5,Sqrt[4-z]},{x,0,Sqrt[4-z-y^2]}]
+Integrate[1,{z,16/5,4},{y,0,Sqrt[4-z]},
{x,0,Sqrt[4-z-y^2]}]+Integrate[1,{z,0,3},
{y,(4-Sqrt[16-5*z])/5,1},{x,0,2-2*y}]
+Integrate[1,{z,3,16/5},{y,(4-Sqrt[16-5*z])/5,
(4+Sqrt[16-5*z])/5},{x,0,2-2*y}]
\end{Verbatim}
\end{tcolorbox}
\vspace{0.7em}
\noindent


En la expresión del paraboloide se elimina a $y,$ de modo que la curva de intersección se encuentre en el plano $xz.$ Esto es
\[z= 4- \left ( \frac{2-x}{2} \right ) ^2-x^2= 3+x- \frac{5}{4}x^2.\]
El valor máximo de $z,$ está dado por
$z^\prime= 1-5x/2$, que al igualar a cero nos lleva a que $x=2/5$, este valor está dado por
$z=3+\frac{2}{5}-\frac{5}{4}\left(\frac{2}{5}\right)^2=\frac{16}{5}.$
\vskip0.3cm
De la expresión $z=3+x- \frac{5}{4}x^2$
también se sigue que $x=\frac{2}{5} \pm \frac{2}{5} \sqrt{16-5z}.$ En los siguientes gráficos se ilustra porque en el orden $dydzdx$ se requiere evaluar dos integrales triples mientras que en el orden $dydxdz$ son necesarias cinco integrales triples.

\vskip0.3cm
\noindent
\begin{minipage}{\textwidth}
	\centering
	\captionof{figure}{Proyección en el plano $xz$}
	\vspace{-20mm}
	\begin{minipage}{0.38\linewidth}
		\centering
		\includegraphics[width=\linewidth]{Fig/315.pdf}
	\end{minipage}
	\hspace{1em}
	\begin{minipage}{0.38\linewidth}
		\centering
		\includegraphics[width=\linewidth]{Fig/316.pdf}
	\end{minipage}
	
	\vspace{-10mm}
	
	\vspace{0.3em}
	\fuentepropia
\end{minipage}
\vskip0.3cm


Apoyándonos en estas figuras presentamos las integrales triples respec\-tivas:
\vspace{-2mm}
{\small
\[\bigintsss_{0}^2\bigintsss_{0}^{3+x-\frac{5}{4}x^2}\bigintsss_{0}^{1-x/2}dydzdx+\bigintsss_{0}^2\bigintsss_{3+x-\frac{5}{4}x^2}^{4-x^2}\bigintsss_{0}^{\sqrt{4-x^2-z}}dydzdx= \frac{19}{6}\]
}

{\footnotesize
\begin{multline*}
\bigintsss_{0}^{3}\bigintsss_{0}^{\frac{2}{5}+\frac{2}{5}\sqrt{16-5z}}\bigintsss_{0}^{1-x/2}dydxdz +\bigintsss_{3}^{\frac{16}{5}}\bigintsss_{\frac{2}{5}-\frac{2}{5}\sqrt{16-5z}}^{\frac{2}{5}+\frac{2}{5}\sqrt{16-5z}}\bigintsss_{0}^{1-x/2}dydxdz \\
+\bigintsss_{3}^{\frac{16}{5}}\bigintsss_{0}^{\frac{2}{5}-\frac{2}{5}\sqrt{16-5z}}\bigintsss_{0}^{\sqrt{4-x^2-z}}dydxdz
+\bigintsss_{\frac{16}{5}}^{4}\bigintsss_{0}^{\sqrt{4-z}}\bigintsss_{0}^{\sqrt{4-x^2-z}}dydxdz\\
+\bigintsss_{0}^{\frac{16}{5}}\bigintsss_{\frac{2}{5}+\frac{2}{5}\sqrt{16-5z}}^{\sqrt{4-z}}\bigintsss_{0}^{\sqrt{4-x^2-z}}dydxdz = \frac{19}{6}.
\end{multline*}
}

Estas integrales se evalúan en \verb"Mathematica" con la siguiente instrucción:

\vspace{0.6em}
\noindent
\begin{tcolorbox}[breakable,enhanced,title=\bf Instrucción en \verb"Mathematica", top=2pt,bottom=2pt,boxsep=0pt,before skip=2pt,after skip=0pt]
\begin{Verbatim}[formatcom=\letraverbatim]
Integrate[1,{z,0,3},{x,0,(2+2*Sqrt[16-5z])/5},
{y,0,1-x/2}]+Integrate[1,{z,3,16/5},
{x,2*(1-Sqrt[16-5z])/5,2*(1+Sqrt[16-5z])/5},
{y,0,1-x/2}]+Integrate[1,{z,3,16/5},{x,0,
(2-2*Sqrt[16-5z])/5},{y,0,Sqrt[4-z-x^2]}]+
Integrate[1,{z,16/5,4},{x,0,Sqrt[4-z]},{y,0,
Sqrt[4-z-x^2]}]+Integrate[1,{z,0,16/5},{x,
(2+2*Sqrt[16-5z])/5,Sqrt[4-z]},{y,0,Sqrt[4-z-x^2]}]
\end{Verbatim}
\end{tcolorbox}
\vspace{0.7em}
\noindent


%El programa ofrece como solución $\frac{17}{6}+\frac{2\pi}{25}-\frac{2}{375}\left( 10 \pi-44 +5\arctan\left(\frac{44}{117}\right)\right)+ \frac{1}{375}\left(37-60\arctan\left(\frac{1}{2}\right)\right)$ que es equivalente a $19/6.$
Con el siguiente proceso también hallamos el volumen del sólido.
\vskip0.2cm
%\begin{tcolorbox}[title=\bfseries Parametrizando el sólido]
Sobre el plano $xy,$ el sólido se proyecta como un triángulo con vértices $(0,0),$ $(2,0)$ y $(0,1)$. El segmento que une el punto $(0,0)$ con $(2,0)$ se parametriza como $\left\langle 2t,0\right\rangle, $ con $t\in [ 0,1].$
Los segmentos que unen cualquier punto de este segmento con el punto
$(0,1),$ se parametriza como
$(1-u)(0,1)+u(2t,0)= \left(2ut,1-u\right) .$
Es decir, todo punto en este triángulo es de la forma
$\left\langle 2ut,1-u\right\rangle $, $t\in\lbrack 0,1]$, $u\in [ 0,1]$.
Un punto sobre el paraboloide posee coordenada $z,$ que está dada por
$z=4-x^{2}-y^{2}=4-(2ut)^2-(1-u)^2= 3+2u-u^{2}\left(4t^{2}+1\right) .$
\vskip0.2cm
Consideremos un punto $P$ sobre la base del sólido {\small $P\left( 2ut,1-u,0\right)$}
y un punto $Q$ sobre el paraboloide, {\small $Q\left( 2ut,1-u,3+2u-u^{2}\left( 4t^{2}+1\right) \right)$}.
Al formar el segmento que une estos dos puntos, se obtiene una parametrización del sólido:
\vspace{-2mm}
\begin{multline*}
{\bf r}\left( t,u,w\right) = (1-w)P +wQ \\
= \left\langle 2ut,1-u,w\left( 3+2u-u^{2}\left( 4t^{2}+1\right)
\right) \right\rangle, \end{multline*}
en donde $t\in [ 0,1],$ $u\in [ 0,1],$ $w\in [ 0,1].$

\vskip0.3cm
Con la función ${\bf r}$ hallamos los siguientes vectores:
\begin{eqnarray*}
{\bf r}_{t} &=& \left\langle 2u,0,-8wu^{2}t\right\rangle \\
{\bf r}_{u} &=& \left\langle 2t,-1,w\left( 2-2u\left( 4t^{2}+1\right) \right) \right\rangle \\
{\bf r}_{w} &=& \left\langle 0,0,3+2u-u^{2}\left( 4t^{2}+1\right) \right\rangle.
\end{eqnarray*}
El elemento diferencial de volumen o producto vectorial triple está dado por
$|{\bf r}_{t} \cdot ({\bf r}_{u} \times {\bf r}_{w})| =\left| -2u\left( 3+2u-\left(4t^{2}+1\right) u^{2}\right) \right| ,$
integrando con los límites establecidos, se halla el volumen deseado:
\[\int_{0}^{1}\int_{0}^{1}\int_{0}^{1}\left| -2u\left( 3+2u-\left(4t^{2}+1\right) u^{2}\right)  \frac{19}{6}\right|\]
%\end{tcolorbox}

Estos cálculos se hacen en \verb"Mathematica" ejecutando estas instrucciones:

\vspace{0.6em}
\noindent
\begin{tcolorbox}[breakable,enhanced,title=\bf Instrucción en \verb"Mathematica", top=2pt,bottom=2pt,boxsep=0pt,before skip=2pt,after skip=0pt]
\begin{Verbatim}[formatcom=\letraverbatim]
{x,y}=(1-u)*{2t,0}+u*{0,1};
r=(1-w)*{x,y,0}+w*{x,y,4-x^2-y^2};
A=Dot[D[r,t],Cross[D[r,u],D[r,w]]];
Integrate[A,{t,0,1},{u,0,1},{w,0,1}]
\end{Verbatim}
\end{tcolorbox}
\vspace{0.7em}
\noindent


\subsection*{Un cilindro y dos planos IV}
Determinar el volumen del sólido acotado por las figuras de las expresiones
$z=1-|y|$ y $z=x^2.$ La región del espacio se muestra a continuación:

\vskip0.2cm
\noindent
\begin{minipage}{\textwidth}
	\centering
	\captionof{figure}{Sólido determinado por $x^2 \leq z \leq 1 - |y|$.}
	\label{fig:solido_parabolico_absoluto_2}
	\begin{minipage}{0.23\linewidth}
		\centering
		\includegraphics[width=\linewidth]{Img/28a.pdf}
	\end{minipage}
	\hspace{1em} 
	\begin{minipage}{0.23\linewidth}
		\centering
		\includegraphics[width=\linewidth]{Img/28c.pdf}
	\end{minipage}
		\hspace{1em} 
	\begin{minipage}{0.23\linewidth}
		\centering
		\includegraphics[width=\linewidth]{Img/28d.pdf}
	\end{minipage}
	
	\vspace{0.3em}
	\fuentepropia
	\end{minipage}
\vskip0.3cm

Las figuras anteriores se consiguen usando las siguientes instrucciones:

\vspace{0.6em}
\noindent
\begin{tcolorbox}[breakable,enhanced,title=\bf Instrucción en \verb"Mathematica", top=2pt,bottom=2pt,boxsep=0pt,before skip=2pt,after skip=0pt]
\begin{Verbatim}[formatcom=\letraverbatim]
ContourPlot3D[{z==x^2,z==1-y,z==1+y},{x,-1,1},
{y,-1,1},{z,0,1},Ticks->{{-1,0,1},{-1,0,1},{0,1}},
PlotPoints->50,Mesh->False]
RegionPlot3D[z-1<y<1-z&&x^2<z<1,{x,-1,1},{y,-1,1},
{z,0,1},PlotPoints->90,Mesh->False,PlotStyle->Orange]
\end{Verbatim}
\end{tcolorbox}
\vspace{0.7em}
\noindent
\vskip0.2cm

De la ecuación $z=x^2$ al reemplazar a $y=\left| 1-z^2\right| $
se sigue que $z=1\pm y$ y también $z=x^2-1$ o $z=1-x^2$. Estas dos curvas determinan la región de integración sobre el plano $xz.$

\vskip0.2cm
\noindent
\begin{minipage}{\textwidth}
	\centering
	\captionof{figure}{Proyección del sólido en los planos coordenados.}
	\label{fig:proyeccion_solido_parabolico_absoluto_2}
	\begin{minipage}{0.30\linewidth}
		\centering
		\vspace{-8mm}
		\includegraphics[width=\linewidth]{Fig/317.pdf}
	\end{minipage}
	\hspace{1em} 
	\begin{minipage}{0.30\linewidth}
		\centering
		\vspace{-8mm}
		\includegraphics[width=\linewidth]{Fig/318.pdf}
	\end{minipage}
	\hspace{1em} 
	\begin{minipage}{0.30\linewidth}
		\centering
		\vspace{-8mm}
		\includegraphics[width=\linewidth]{Fig/319.pdf}
	\end{minipage}
	
	\vspace{-7mm}
	\fuentepropia
\end{minipage}
\vskip0.3cm

Con estos gráficos se plantean las integrales que determinan el volumen del sólido. Integrando primero con respecto a $z,$ esta recorre puntos desde el cilindro parabólico $z=x^2$ hasta los planos $z=1-y$ y $z=1+y.$
\begin{multline*}
\int_{-1}^{1}\int_{0}^{1-x^2}\int_{x^2}^{1-y}dzdydx+\int_{-1}^{1}\int_{x^2-1}^{0}\int_{x^2}^{1+y}dzdydx= \frac{16}{15}\\
\int_0^1\int_{-\sqrt{1-y}}^{\sqrt{1-y}}\int_{x^2}^{1-y}dzdxdy+\int_{-1}^{0}\int_{-\sqrt{1+y}}^{\sqrt{1+y}}\int_{x^2}^{1+y}dzdxdy=\frac{16}{15}.
\end{multline*}
Mientras que si se integra primero con respecto a $y$,
esta recorre los puntos que se hallan entre los planos $y=z-1$ y $y=1-z.$
$\int_{-1}^{1}\int_{x^2}^{1}\int_{z-1}^{1-z}dydzdx=\int_{0}^{1}\int_{-\sqrt{z}}^{\sqrt{z}}\int_{z-1}^{1-z}dydxdz=\frac{16}{15}.$
La variable $x$ recorre los puntos que van desde la parte negativa, hasta la parte positiva del cilindro $z=x^2$,
estos se describen como $x=\pm \sqrt{z}.$
\begin{multline*} = \int_{-1}^{0}\int_{0}^{1+y}\int_{-\sqrt{z}}^{\sqrt{z}}dxdzdy+ \int_{0}^{1}\int_{0}^{1-y}\int_{-\sqrt{z}}^{\sqrt{z}}dxdzdy \\
=\int_{0}^{1}\int_{z-1}^{1-z}\int_{-\sqrt{z}}^{\sqrt{z}}dxdydz= \frac{16}{15}.
\end{multline*}
En \verb"Mathematica", las integrales anteriores en los órdenes $dzdydx,$ $dydzdy$ y $dxdydz$,
se evalúan ejecutando las siguientes instrucciones:
\vskip0.2cm

\vspace{0.6em}
\noindent
\begin{tcolorbox}[breakable,enhanced,title=\bf Instrucción en \verb"Mathematica", top=2pt,bottom=2pt,boxsep=0pt,
before skip=2pt,after skip=0pt]
\begin{Verbatim}[formatcom=\letraverbatim]
Integrate[1,{x,-1,1},{y,0,1-x^2},{z,x^2,1-y}]
+Integrate[1,{x,-1,1},{y,x^2-1,0},{z,x^2,1+y}]
Integrate[1,{x,-1,1},{z,x^2,1},{y,z-1,1-z}]
Integrate[1,{z,0,1},{y,z-1,1-z},{x,-Sqrt[z],Sqrt[z]}]
\end{Verbatim}
\end{tcolorbox}
\vspace{0.7em}
\noindent

Este valor también se puede determinar de la siguiente manera.


\vskip0.2cm
%\begin{tcolorbox}[title=\bfseries Parametrizando el sólido]
	
En el plano $xz,$ la proyección del sólido se parametriza con
$(1-u)(t,t^2)+u(t,1),$ $t\in [-1,1],$ $u\in [0,1].$
Con un punto $P(t,t^2(1-u)+u-1,t^2(1-u)+u) $ sobre el plano $y=1+z$, y otro punto $Q(t,1-t^2(1-u)-u,t^2(1-u)+u) $ sobre el plano $z=1-y$.


\vskip0.2cm
\noindent
\begin{minipage}{\textwidth}
	\centering
	\captionof{figure} {$x^2 \leq y \leq 1.$}
	\begin{minipage}{0.32\linewidth}
		\centering
		\vspace{-10mm}
		\includegraphics[width=\linewidth]{Fig/319.pdf}
	\end{minipage}
	
	\vspace{-4mm}
	\fuentepropia
\end{minipage}
\vskip0.2cm


Con estos dos puntos, definimos la parametrización del sólido, que es una combinación convexa de la forma $(1-w)P+wQ$, esto es,
{\small
\begin{equation*}
{\bf r}\left( t,u,w\right) = \left\langle t,\left( t^2-1\right) \left( 1-u\right) \left( 1-2w\right),t^2\left( 1-u\right) +u\right\rangle,
\end{equation*}
}

en donde $t\in [ -1,1]$, $u\in [ 0,1]$, $w\in [ 0,1]$. Con la función
${\bf r}$ hallamos los siguientes vectores:
\vspace{-2mm}
\begin{eqnarray*}
{\bf r}_{t}&=&\left\langle 1,2t\left( 1-u\right) \left(1-2w\right) ,2t\left( 1-u\right) \right\rangle \\
{\bf r}_{u} &=& \left\langle 0,\left( 1-t^2\right) \left( 1-2w\right),1-t^2\right\rangle \\
{\bf r}_{w} &=& \left\langle 0,2\left( 1-t^2\right) \left( 1-u\right),0\right\rangle.
\end{eqnarray*}
El producto vectorial triple es
$\left| {\bf r}_{t}\cdot \left({\bf r}_{u}\times {\bf r}_{w}\right) \right| =
%\left| \left( 1,2t\left( 1-u\right) \left( 1-2w\right) ,2t\left( 1-u\right) \right) \cdot \left( \left( 0,\left( 1-t^2\right) \left( 1-2w\right) ,1-t^2\right)\times \left( 0,2\left( 1-t^2\right) \left( 1-u\right) ,0\right) \right)\right| =
2\left| \left( t^2-1\right) ^2\left( u-1\right) \right| .$
Este término se integra para encontrar el volumen del sólido.
\[\int_{0}^{1}\int_{0}^{1}\int_{-1}^{1}2\left| \left( t^2-1\right)^2\left( u-1\right) \right| dtdudw= \frac{16}{15}.\]
%\end{tcolorbox}

Estos cálculos se hacen en \verb"Mathematica" usando las siguientes instrucciones:

\vspace{0.6em}
\noindent
\begin{tcolorbox}[breakable,enhanced,title=\bf Instrucción en \verb"Mathematica", top=2pt,bottom=2pt,boxsep=0pt,before skip=2pt,after skip=0pt]
\begin{Verbatim}[formatcom=\letraverbatim]
{x,z}=(1-u)*{t,t^2}+u*{t,1};
s=(1-w)*{x,1-z,z}+w*{x,z-1,z};
A=Dot[D[s,t],Cross[D[s,u],D[s,w]]];
Integrate[Abs[A],{t,-1,1},{u,0,1},{w,0,1}]
\end{Verbatim}
\end{tcolorbox}
\vspace{0.7em}
\noindent

\subsection*{Dos cilindros y dos planos III}
Los planos con ecuaciones $z=2-y$ y $z=0$, junto con los cilindros con ecuaciones
$x^2+y^2=4$ y $x^2+y^2=2y,\,$ que denominaremos cilindro mayor y menor respectivamente, acotan el siguiente sólido.

\vskip0.2cm
\noindent
\begin{minipage}{\textwidth}
	\centering
	\captionof{figure}{Sólido determinado por $0 \leq z \leq 2 - y$, \quad $2y \leq x^2 + y^2 \leq 4$.}
	\label{fig:solido_corona_inclinada}
	\begin{minipage}{0.23\linewidth}
		\centering
		\vspace{-2mm}
		\includegraphics[width=\linewidth]{Img/31a.pdf}
	\end{minipage}
	\hspace{1em} 
	\begin{minipage}{0.22\linewidth}
		\centering
		\vspace{-2mm}
		\includegraphics[width=\linewidth]{Img/31b.pdf}
	\end{minipage}	
	
	\vspace{0.3em}
	\fuentepropia
\end{minipage}
\vskip0.2cm


Estas figuras se pueden generar con los siguientes comandos.

\vspace{0.6em}
\noindent
\begin{tcolorbox}[breakable,enhanced,title=\bf Instrucción en \verb"Mathematica", top=2pt,bottom=2pt,boxsep=0pt,before skip=2pt,after skip=0pt]
\begin{Verbatim}[formatcom=\letraverbatim]
RegionPlot3D[x^2+y^2<4&&x^2+y^2>2*y&&z>0&&z<2-y,
{x,-2,2},{y,-2,2},{z,0,4},PlotPoints->90,Mesh->False]
\end{Verbatim}
\end{tcolorbox}
\vspace{0.7em}
\noindent

Presentaremos ahora el cálculo del volumen del sólido de otras formas.

\subsubsection*{Integración en el orden $dxdydz$ o $dxdzdy$}

La proyección en el plano $yz$ del sólido es un triángulo. Además, $x$ recorre diferentes regiones, en una de ellas va desde la parte negativa del cilindro mayor a la parte positiva del mismo cilindro.

\vskip0.3cm
En otra región, el recorrido va desde la parte negativa del cilindro mayor a la parte negativa del cilindro menor y luego desde la parte positiva del cilindro menor a la parte positiva del cilindro mayor.


\vskip0.2cm
\noindent
\begin{minipage}{\textwidth}
	\centering
	\captionof{figure}{Proyección en el plano $yz$}
	\begin{minipage}{0.29\linewidth}
		\centering
		\vspace{-8mm}
		\includegraphics[width=\linewidth]{Fig/320.pdf}
	\end{minipage}
	
	\vspace{-4mm}
	\fuentepropia
\end{minipage}
\vskip0.2cm

Por esta razón, el volumen se determina evaluando las siguientes integrales:
\vspace{-8mm}

\begin{multline*}
V=\int_{-2}^{0}\int_{0}^{2-y}\int_{-\sqrt{4-y^2}}^{\sqrt{4-y^2}}dxdzdy +
\int_{0}^2\int_{0}^{2-y}\int_{-\sqrt{4-y^2}}^{-\sqrt{2y-y^2}}dxdzdy\\
+\int_{0}^2\int_{0}^{2-y}\int_{\sqrt{2y-y^2}}^{\sqrt{4-y^2}}dxdzdy \\
=\int_{0}^2\int_{-2}^{0}\int_{-\sqrt{4-y^2}}^{\sqrt{4-y^2}}dxdydz
+\int_{2}^{4}\int_{-2}^{2-z}\int_{-\sqrt{4-y^2}}^{\sqrt{4-y^2}}dxdydz+ \\
\int_{0}^2\int_{0}^{2-z}\int_{-\sqrt{4-y^2}}^{-\sqrt{2y-y^2}}dxdydz
+\int_{0}^2\int_{0}^{2-z}\int_{\sqrt{2y-y^2}}^{\sqrt{4-y^2}}dxdydz=7\pi.
\end{multline*}


Estas integrales se evaluaron usando las siguientes instrucciones:

\vspace{0.6em}
\noindent
\begin{tcolorbox}[breakable,enhanced,title=\bf Instrucción en \verb"Mathematica", top=2pt,bottom=2pt,boxsep=0pt,before skip=2pt,after skip=0pt]
\begin{Verbatim}[formatcom=\letraverbatim]
Integrate[1,{y,-2,0},{z,0,2-y},{x,-Sqrt[4-y^2],Sqrt[4
y^2]}]+Integrate[1,{y,0,2},{z,0,2-y},{x,-Sqrt[4-y^2],
Sqrt[2*y-y^2]}]+Integrate[1,{y,0,2},{z,0,2-y},{x,
Sqrt[2*y-y^2],Sqrt[4-y^2]}]
Integrate[1,{z,0,2},{y,-2,0},{x,-Sqrt[4-y^2],Sqrt[4
y^2]}]+Integrate[1,{z,2,4},{y,-2,2-z},{x,-Sqrt[4-
y^2],Sqrt[4-y^2]}]+Integrate[1,{z,0,2},{y,0,2-z},
{x,-Sqrt[4-y^2],-Sqrt[2*y-y^2]}]+Integrate[1,{z,0,2}
,{y,0,2-z},{x,Sqrt[2*y-y^2],Sqrt[4-y^2]}]
\end{Verbatim}
\end{tcolorbox}
\vspace{0.7em}
\noindent


\subsubsection*{Integración en el orden $dzdxdy$ y $dzdydx.$}

\vskip0.2cm
\noindent
\begin{minipage}{\textwidth}
	\centering
	\captionof{figure}{Integración del sólido en el plano $xy$}
	\begin{minipage}{0.45\linewidth}
		\centering
		\vspace{-18mm}
		\includegraphics[width=\linewidth]{Fig/321.pdf}
	\end{minipage}
	\hspace{1em}
	\begin{minipage}{0.35\linewidth}
		\centering
		\vspace{-18mm}
		\includegraphics[width=\linewidth]{Fig/322.pdf}
	\end{minipage}
	
	
	\vspace{-15mm}
	\fuentepropia
\end{minipage}
\vskip0.2cm


Este gráfico es útil para calcular el volumen mediante el planteamiento de las integrales triples necesarias, integrando primero con respecto a $z.$

\vspace{-4mm}
\begin{multline*}
\int_{0}^2\int_{-\sqrt{4-y^2}}^{-\sqrt{2y-y^2}}\int_{0}^{2-y}dzdxdy+\int_{0}^2\int_{\sqrt{2y-y^2}}^{\sqrt{4-y^2}}\int_{0}^{2-y}dzdxdy \\ +\int_{-2}^{0}\int_{-\sqrt{4-y^2}}^{\sqrt{4-y^2}}\int_{0}^{2-y}dzdxdy
=\int_{-2}^{-1}\int_{-\sqrt{4-x^2}}^{\sqrt{4-x^2}}\int_{0}^{2-y}dzdydx\\ + \int_{1}^2\int_{-\sqrt{4-x^2}}^{\sqrt{4-x^2}}\int_{0}^{2-y}dzdydx +\int_{-1}^{1}\int_{-\sqrt{4-x^2}}^{1-\sqrt{1-x^2}}\int_{0}^{2-y}dzdydx \\ +\int_{-1}^{1}\int_{1+\sqrt{1-x^2}}^{\sqrt{4-x^2}}\int_{0}^{2-y}dzdydx= 7\pi
\end{multline*}

\subsubsection*{Integración en el orden $dydxdz$ y $dydzdx$}

En estos casos, se trazan unas líneas auxiliares (en azul) sobre el sólido, con el propósito de evidenciar que el sólido se divide en diferentes regiones, lo cual dará origen a varias integrales triples.

\vskip0.2cm
\noindent
\begin{minipage}{\textwidth}
	\centering
	\captionof{figure}{Cono truncado oblicuo determinado por ${\bf r}(u,t)$}
	\begin{minipage}{0.40\linewidth}
		\centering
		\vspace{-18mm}
		\includegraphics[width=\linewidth]{Fig/323.pdf}
	\end{minipage}
	\hspace{1em} 
	\begin{minipage}{0.40\linewidth}
		\centering
		\vspace{-18mm}
		\includegraphics[width=\linewidth]{Fig/324.pdf}
	\end{minipage}
	
	\vspace{-8mm}
	\fuentepropia
\end{minipage}
\vskip0.2cm

Al reemplazar $y=2-z,$ de la ecuación del plano en las ecuaciones de los cilindros se obtiene
$x^2+\left( 2-z\right) ^2=4,$ por tanto $x^2+z^2=4z$ y $x^2+\left( 2-z\right) ^2=2\left( 2-z\right),$ entonces $x^2+z^2=2z.$
El sólido proyectado en el plano $xz$, junto con las líneas mencionadas antes, se observan en los siguientes gráficos.

\vskip0.3cm
\noindent
\begin{minipage}{\textwidth}
	\centering
	\captionof{figure}{Cono truncado oblicuo determinado por ${\bf r}(u,t)$}
	\begin{minipage}{0.35\linewidth}
		\centering
		\vspace{-10mm}
		\includegraphics[width=\linewidth]{Fig/325.pdf}
	\end{minipage}
	\hspace{1em}
	\begin{minipage}{0.35\linewidth}
		\centering
		\vspace{-10mm}
		\includegraphics[width=\linewidth]{Fig/326.pdf}
	\end{minipage}
	
	\vspace{-8mm}
	\fuentepropia
\end{minipage}
\vskip0.3cm

Las diferentes tonalidades y texturas sugieren que en cada una de esas regiones se debe plantear una integral triple.
\vskip0.2cm
Integrando primero con respecto a $y$, se debe observar que
$y$ recorre los puntos que van desde la parte negativa del cilindro de mayor radio al plano.
Por tal motivo, las integrales necesarias son las siguientes: 
\begin{eqnarray}
\int_{-2}^{-1}\int_{2-\sqrt{4-x^2}}^{2+\sqrt{4-x^2}}\int_{-\sqrt{4-x^2}}^{2-z}dydzdx+
\int_{-1}^{1}\int_{1+\sqrt{1-x^2}}^{2+\sqrt{4-x^2}}\int_{-\sqrt{4-x^2}}^{2-z}dydzdx \nonumber \\
+ \int_{1}^2\int_{2-\sqrt{4-x^2}}^{2+\sqrt{4-x^2}}\int_{-\sqrt{4-x^2}}^{2-z}dydzdx \approx 11.2156 \label{31a}
\end{eqnarray}
Cuando $y$ va desde el cilindro de mayor radio al mismo cilindro, son dos regiones que en el plano $xz,$ que se proyectan de color anaranjado 
\begin{eqnarray}
\int_{-2}^{-1}\int_{0}^{2-\sqrt{4-x^2}}\int_{-\sqrt{4-x^2}}^{\sqrt{4-x^2}}dydzdx
+\int_{1}^2\int_{0}^{2-\sqrt{4-x^2}}\int_{-\sqrt{4-x^2}}^{\sqrt{4-x^2}}dydzdx \nonumber \\
=\frac{16\pi}{3} -4\sqrt{3} -\frac{20}{3} \label{31b}
\end{eqnarray}
Si $y$ recorre la región desde la parte positiva del cilindro de menor radio a la parte positiva del cilindro grande se obtiene
\begin{equation}\label{31c}
\int_{-1}^{1}\int_{0}^{2-\sqrt{4-x^2}}\int_{1+\sqrt{1-x^2}}^{\sqrt{4-x^2}}dydzdx\approx 0.0445
\end{equation}
La base de este conjunto posee una región en la \emph{parte trasera} desde el cilindro de radio mayor al menor
\begin{equation}\label{31d}
\int_{-1}^{1}\int_{0}^{1-\sqrt{1-x^2}}\int_{-\sqrt{4-x^2}}^{1-\sqrt{1-x^2}}dydzdx\approx 0.9781
\end{equation}

Desde la parte positiva del cilindro peque\~{n}o al plano, cuya proyección en el plano $xz,$ corresponde a la parte en azul oscuro, tendremos la siguiente integral 
\begin{equation}\label{31e}
\int_{-1}^{1}\int_{2-\sqrt{4-x^2}}^{1-\sqrt{1-x^2}}\int_{1+\sqrt{1-x^2}}^{2-z}dydzdx\approx 0.0376
\end{equation}
\vskip0.2cm
Desde la parte negativa del cilindro grande a la parte negativa del cilindro peque\~{n}o, cuya proyección en el plano $xz$ es la región cuadriculada: 
\begin{equation}\label{31f}
\int_{-1}^{1}\int_{1-\sqrt{1-x^2}}^{1+\sqrt{1-x^2}}\int_{-\sqrt{4-x^2}}^{1-\sqrt{1-x^2}}dydzdx\approx 6.5551
\end{equation}
Sumando las expresiones (\ref{31a}), (\ref{31b}), (\ref{31c}), (\ref{31d}), (\ref{31e}) y (\ref{31f})
se obtiene un valor a aproximadamente $7\pi$.

\vskip0.2cm
\noindent
\begin{minipage}{\textwidth}
	\centering
	\captionof{figure}{Proyección en el plano $xz$.}
	\begin{minipage}{0.38\linewidth}
		\centering
		\vspace{-15mm}
		\includegraphics[width=\linewidth]{Fig/327.pdf}
	\end{minipage}
	
	\vspace{-8mm}
	\fuentepropia
\end{minipage}
\vskip0.2cm

\begin{comment}
Estos cálculos se hacen en \verb"Mathematica" como se muestra enseguida.
% \fontsize{9.5}{12}\selectfont
\begin{verbatim}
a=NIntegrate[1,{x,-2,-1},{z,2-Sqrt[4-x^2],2+
Sqrt[4-x^2]},{y,-Sqrt[4-x^2],2-z}]+Integrate[1,
{x,-1,1},{z,1+Sqrt[1-x^2],2+Sqrt[4-x^2]},
{y,-Sqrt[4-x^2],2-z}]+NIntegrate[1,{x,1,2},
{z,2-Sqrt[4-x^2],2+Sqrt[4-x^2]},{y,-Sqrt[4-x^2],2-z}]
b=Integrate[1,{x,-2,-1},{z,0,2-Sqrt[4-x^2]},
{y,-Sqrt[4-x^2],Sqrt[4-x^2]}]+Integrate[1,{x,1,2},
{z,0,2-Sqrt[4-x^2]},{y,-Sqrt[4-x^2],Sqrt[4-x^2]}]
c=NIntegrate[1,{x,-1,1},{z,0,2-Sqrt[4-x^2]},
{y,1+Sqrt[1-x^2],Sqrt[4-x^2]}]
d=NIntegrate[1,{x,-1,1},{z,0,1-Sqrt[1-x^2]},
{y,-Sqrt[4-x^2],1-Sqrt[1-x^2]}]
e=NIntegrate[1,{x,-1,1},{z,2-Sqrt[4-x^2],
1-Sqrt[1-x^2]},{y,1+Sqrt[1-x^2],2-z}]
f=NIntegrate[1,{x,-1,1},{z,1-Sqrt[1-x^2],
1+Sqrt[1-x^2]},{y,-Sqrt[4-x^2],1-Sqrt[1-x^2]}]
a+b+c+d+e+f
\end{verbatim}

\end{comment}


Planteando las integrales en el orden $dydxdz$, es necesario considerar varias regiones.
Una de ellas es la que recorre los puntos desde la parte negativa del cilindro de mayor radio al plano.
El volumen es el resultado de considerar varias integrales triples, las cuales presentamos a continuación:

\vspace{-7mm}
\begin{multline} \label{31k}
\int_{1}^2\int_{-1}^{-\sqrt{2z-z^2}}\int_{-\sqrt{4-x^2}}^{2-z}dydxdz \\
+\int_{1}^2\int_{\sqrt{2z-z^2}}^{1}\int_{-\sqrt{4-x^2}}^{2-z}dydxdz \approx 0. 8822
\end{multline}
\begin{eqnarray}
\int_{2-\sqrt{3}}^{2+\sqrt{3}}\int_{-\sqrt{4z-z^2}}^{-1}\int_{-\sqrt{4-x^2}}^{2-z}dydxdz&\approx& 3. 3333 \label{31g}
\end{eqnarray}
\begin{eqnarray}
\int_{2-\sqrt{3}}^{2+\sqrt{3}}\int_{1}^{\sqrt{4z-z^2}}\int_{-\sqrt{4-x^2}}^{2-z}dydxdz &\approx& 3. 3333 \label{31h} \\
\int_{2}^{2+\sqrt{3}}\int_{-1}^{1}\int_{-\sqrt{4-x^2}}^{2-z}dydxdz
&=& \frac{2}{3}\pi \sqrt{3} \label{31i} \\
\int_{2+\sqrt{3}}^{4}\int_{-\sqrt{4z-z^2}}^{\sqrt{4z-z^2}}\int_{-\sqrt{4-x^2}}^{2-z}dydxdz&\approx& 0.0391 \label{31j}
\end{eqnarray}
\vskip0.3cm
Desde el cilindro de mayor radio al mismo cilindro, hay dos regiones que en el plano $xz,$ se proyectan de color anaranjado.


%\vspace{-2mm}
\begin{eqnarray}
\int_{0}^{2-\sqrt{3}}\int_{-2}^{-1}\int_{-\sqrt{4-x^2}}^{\sqrt{4-x^2}}dydxdz
+\int_{2-\sqrt{3}}^2\int_{-2}^{-\sqrt{4z-z^2}}\int_{-\sqrt{4-x^2}}^{\sqrt{4-x^2}}dydxdz \nonumber \\
+\int_{0}^{2-\sqrt{3}}\int_{1}^2\int_{-\sqrt{4-x^2}}^{\sqrt{4-x^2}}dydxdz
+\int_{2-\sqrt{3}}^2\int_{\sqrt{4z-z^2}}^2\int_{-\sqrt{4-x^2}}^{\sqrt{4-x^2}}dydxdz \nonumber \\
= \frac{16\pi}{3}-\frac{20}{3}-4\sqrt{3} \label{31l}
\end{eqnarray}
\vskip0.3cm
De la parte positiva del cilindro peque\~{n}o al plano, cuya proyección en el plano $xz,$ corresponde a la parte en azul oscuro.

\begin{multline}
\int_{0}^{2-\sqrt{3}}\int_{-\sqrt{4z-z^2}}^{-\sqrt{2z-z^2}}\int_{1+\sqrt{1-x^2}}^{2-z}dydxdz + \\ \int_{2-\sqrt{3}}^{1}\int_{-1}^{-\sqrt{2z-z^2}}\int_{1+\sqrt{1-x^2}}^{2-z}dydxdz+ \int_{0}^{2-\sqrt{3}}\int_{\sqrt{2z-z^2}}^{\sqrt{4z-z^2}}\int_{1+\sqrt{1-x^2}}^{2-z}dydxdz \\
+\int_{2-\sqrt{3}}^{1}\int_{\sqrt{2z-z^2}}^{1}\int_{1+\sqrt{1-x^2}}^{2-z}dydxdz \approx 0.0376 \label{31m}
\end{multline}

\vskip0.4cm
Cuando $y$ recorre los puntos desde la parte positiva del cilindro pequeño a la parte positiva del mayor cilindro, la integral es:
\vspace{-8mm}
\begin{multline}
\int_{0}^{2-\sqrt{3}}\int_{-1}^{-\sqrt{4z-z^2}}\int_{1+\sqrt{1-x^2}}^{\sqrt{4-x^2}}dydxdz+ \\
\int_{0}^{2-\sqrt{3}}\int_{\sqrt{4z-z^2}}^{1}\int_{1+\sqrt{1-x^2}}^{\sqrt{4-x^2}}dydxdz \approx 0.0445 \label{31n}
\end{multline}

\vskip0.3cm
Existe una parte (la región en azul es parte de ella) que posee una \emph{región atras}, donde $y$ se recorre desde la parte negativa del cilindro mayor hasta la parte negativa el cilindro menor.
\vskip0.2cm
\begin{multline}
\int_{0}^{1}\int_{-1}^{-\sqrt{2z-z^2}}\int_{-\sqrt{4-x^2}}^{1-\sqrt{1-x^2}}dydxdz \\
+\int_{0}^{1}\int_{\sqrt{2z-z^2}}^{1}\int_{-\sqrt{4-x^2}}^{1-\sqrt{1-x^2}}dydxdz \approx 0.9781 \label{31o}
\end{multline}

\vskip0.4cm
Desde la parte negativa del cilindro grande hasta la parte negativa del cilindro pequeño, cuya proyección en el plano $xz$ es la región cuadriculada.
\vskip0.4cm
\begin{equation} \label{31p}
\int_{0}^2\int_{-\sqrt{2z-z^2}}^{\sqrt{2z-z^2}}\int_{-\sqrt{%
4-x^2}}^{1-\sqrt{1-x^2}}dydxdz \approx 6. 5551
\end{equation}
\vskip0.4cm
Sumando las expresiones (\ref{31k}), (\ref{31g}), (\ref{31h}), (\ref{31i}), (\ref{31j}), (\ref{31l}), (\ref{31m}), (\ref{31n}),
(\ref{31o}) y (\ref{31p}), se obtiene un valor aproximado a $7\pi.$

\vskip0.3cm
Con ayuda de \verb"Mathematica" se puede obtener estos valores, para tal fin se definen cada una de las integrales como $a$, $b$, $c$, $d$, $e$, $f$, $g$, $h$, $i$ y $j$.
La suma es aproximadamente $7\pi$.
\vskip0.3cm
\vspace{0.6em}
\noindent
\begin{tcolorbox}[breakable,enhanced,title=\bf Instrucción en \verb"Mathematica", top=2pt,bottom=2pt,boxsep=0pt,before skip=2pt,after skip=0pt]
\begin{Verbatim}[formatcom=\letraverbatim]
a=NIntegrate[1,{z,2-Sqrt[3],2+Sqrt[3]},{x,Sqrt[4*z-z^2],-1},
{y,-Sqrt[4-x^2],2-z}]
b=NIntegrate[1,{z,2-Sqrt[3],2+Sqrt[3]},{x,1,Sqrt[4*z-z^2]},
{y,-Sqrt[4-x^2],2-z}]
c=NIntegrate[1,{z,2,2+Sqrt[3]},{x,-1,1},{y,-Sqrt[4-x^2],2-z}]
d=NIntegrate[1,{z,2+Sqrt[3],4},{x,-Sqrt[4*z-z^2],
Sqrt[4*z-z^2]},{y,-Sqrt[4-x^2],2-z}]
e=NIntegrate[1,{z,1,2},{x,-1,-Sqrt[2*z-z^2]},
{y,-Sqrt[4-x^2],2-z}]+NIntegrate[1,{z,1,2},
{x,Sqrt[2*z-z^2],1},{y,-Sqrt[4-x^2],2-z}]
f=NIntegrate[1,{z,0,2-Sqrt[3]},{x,-2,-1},{y,
Sqrt[4-x^2],Sqrt[4-x^2]}]+NIntegrate[1,{z,
2-Sqrt[3],2},{x,-2,-Sqrt[4*z-z^2]},{y,
Sqrt[4-x^2],Sqrt[4-x^2]}]+NIntegrate[1,
{z,0,2-Sqrt[3]},{x,1,2},{y,-Sqrt[4-x^2],
Sqrt[4-x^2]}]+NIntegrate[1,{z,2-Sqrt[3],2},
{x,Sqrt[4*z-z^2],2},{y,-Sqrt[4-x^2],Sqrt[4-x^2]}]
g=Integrate[1,{z,0,2-Sqrt[3]},{x,-Sqrt[4*z-z^2],
Sqrt[2*z-z^2]},{y,1+Sqrt[1-x^2],2-z}]+Integrate[1,{z,2-
Sqrt[3],1},{x,-1,-Sqrt[2*z-z^2]},
{y,1+Sqrt[1-x^2],2-z}]+
Integrate[1,{z,0,2-Sqrt[3]},{x,Sqrt[2*z-z^2],
Sqrt[4*z-z^2]},
{y,1+Sqrt[1-x^2],2-z}]+Integrate[1,{z,2-Sqrt[3],1},
{x,Sqrt[2*z-z^2],1},{y,1+Sqrt[1-x^2],2-z}]
h=NIntegrate[1,{z,0,2-Sqrt[3]},{x,-1,-Sqrt[4*z-z^2]},
{y,1+Sqrt[1-x^2],Sqrt[4-x^2]}]+
NIntegrate[1,{z,0,2-Sqrt[3]},
{x,Sqrt[4*z-z^2],1},{y,1+Sqrt[1-x^2],Sqrt[4-x^2]}]
i=NIntegrate[1,{z,0,1},{x,-1,-Sqrt[2*z-z^2]},
{y,-Sqrt[4-x^2],
1-Sqrt[1-x^2]}]+NIntegrate[1,{z,0,1},
{x,Sqrt[2*z-z^2],1},
{y,-Sqrt[4-x^2],1-Sqrt[1-x^2]}]
j=NIntegrate[1,{z,0,2},{x,-Sqrt[2*z-z^2],Sqrt[2*z-z^2]},
{y,-Sqrt[4-x^2],1-Sqrt[1-x^2]}]

a+b+c+d+e+f+g+h+i+j
\end{Verbatim}
\end{tcolorbox}
\vspace{0.7em}
\noindent

\vskip0.2cm
Presentamos dos maneras de calcular el volumen del sólido, la primera de ellas es usando coordenadas cilíndricas, luego se parametriza el sólido.
\begin{align*}
\int_0^{\pi} \int_0^2 \int_{2 \sin \theta}^{2 - r \sin \theta} r\,dz\,dr\,d\theta 
+
\int_{\pi}^{2\pi} \int_0^2 \int_0^{2 - r \sin \theta} r\,dz\,dr\,d\theta &= 7\pi
\end{align*}


En coordenadas polares, los cilindros se representan como
{\small $r=2\sin \theta$} y {\small $r=2.$} Las parametrizaciones elegidas son
{\small $\langle 2\sin \theta \cos \theta , 2\sin \theta\sin \theta \rangle$} y {\small $\langle 2 \cos \theta , 2\sin \theta \rangle.$}
Para recorrer adecuadamente cada curva, iniciando en el punto $(0,2),$ se sustituye $\theta$ por $\theta + \frac{\pi}{2}.$

	
\vskip0.2cm
\noindent
\begin{minipage}{\textwidth}
	\centering
	\captionof{figure}{\textbf{$2y \leq x^2 + y^2 \leq 4.$}}
	\begin{minipage}{0.38\linewidth}
		\centering
		\vspace{-15mm}
		\includegraphics[width=\linewidth]{Fig/328.pdf}
	\end{minipage}
	
	\vspace{-10mm}
	\fuentepropia
\end{minipage}
\vskip0.2cm
	

Los puntos $P$ y $Q$ del gráfico, se determinan por medio de las funciones
$\mathbf{r}_1(\theta)= \langle -\sin(2\theta), 1+\cos(2\theta) \rangle,$
y $\mathbf{r}_2(\theta) = \langle -2 \sin(2\theta), 2\cos(2\theta)\rangle,$
res\-pectivamente.
Una combinación convexa de estos parametriza la región sombreada. Es decir,

\vspace{-8mm}
\[\mathbf{r}_3(\theta, u) = \langle (u-2) \sin(2\theta), u - (u-2) \cos(2\theta) \rangle, \theta \in [0, \pi], u \in [0, 1].\]

Con un punto $A((u-2)\sin(2\theta),u-(u-2)\cos(2\theta),0)$ el plano $z=0$ y otro punto
$B((u-2)\sin(2\theta),u-(u-2)\cos(2\theta),2-(u-(u-2)\cos(2\theta))),$
sobre el plano $z=2-y$ se forma la combinación convexa de la forma $(1-w)A+ wB,$ que describe el sólido, esto define la función:

\vspace{-8mm}
\[\mathbf{r}(\theta, u, w) = \left\langle (u-2) \sin(2\theta), u - (u-2) \cos(2\theta), 2w(2-u) \sin^2(\theta) \right\rangle,\]

en donde $\theta \in [0,\pi],\,u\in[0,1]$ y $w\in[0,1].$ 
Derivando parcialmente se obtienen los siguientes vectores

\vspace{-8mm}
{\footnotesize
\begin{align*}
{\bf r}_\theta &= \langle 2(u-1)\cos(2\theta),2(u-2)\sin(2\theta),2w(2-u)\sin(2\theta) \rangle \\
{\bf r}_u &= \langle \sin(2\theta),2\sin^2(\theta),-2w\sin^2(\theta) \rangle \\
{\bf r}_w &= \langle 0,0,2(2-u)\sin^2(\theta) \rangle \\
{\bf r}_{\theta} \times {\bf r}_u &= \langle 0, 4w(2 - u) \sin^2(t), 4 (2 - u) \sin^2(t) \rangle
\end{align*}
}

\vspace{-3mm}
Con lo cual, $|{\bf r}_w \cdot ({\bf r}_{\theta} \times {\bf r}_u)|= 8(u-2)^2\sin^4(\theta).$
Entonces, el volumen del sólido está dado por:

{\footnotesize
\[\int_0^1 \int_0^1 \int_0^\pi 8(u-2)^2\sin^4(\theta) d\theta du dw =7\pi\]
}
	
%\end{tcolorbox}

Este valor se obtiene en \verb"Mathematica" usando las siguientes instrucciones.


\vspace{0.6em}
\noindent
\begin{tcolorbox}[breakable,enhanced,title=\bf Instrucción en \verb"Mathematica", top=2pt,bottom=2pt,boxsep=0pt,before skip=2pt,after skip=0pt]
\begin{Verbatim}[formatcom=\letraverbatim]
{x,y,z}={(u-2)*Sin[2*t],u-(u-2)*Cos[2*t],2-y};
r=Simplify[(1-w)*{x,y,0}+w*{x,y,z}]
A=Simplify[Dot[D[r,w],Cross[D[r,t],D[r,u]]]]
Integrate[Abs[A],{w,0,1},{u,0,1},{t,0,Pi}]
\end{Verbatim}
\end{tcolorbox}
\vspace{0.7em}
\noindent

\section{Sólidos acotados por paraboloides} \index{Sólidos!con paraboloides}
\vskip0.5cm

\subsection*{Un paraboloide elíptico} \index{Paraboloide!elíptico}
Consideremos el paraboloide con vértice en el punto $(4,3,2)$ que abre hacia la parte negativa del eje $x$, determinado por la ecuación
$\frac{(z-2)^2}{4}+\frac{(y-3)^2}{9}=1-\frac{x}{4},$
esta superficie junto con el plano $x=0,$ acotan un sólido. La figura es la siguiente:

\vskip0.2cm
\noindent
\begin{minipage}{\textwidth}
	\centering
	\captionof{figure}{Sólido definido por $0 \leq x \leq 4 - (z - 2)^2 - \frac{4}{9}(y - 3)^2$.}
	\label{fig:solido_eliptico_desplazado_2}
	\begin{minipage}{0.20\linewidth}
		\centering
		\includegraphics[width=\linewidth]{Img/32a.pdf}
	\end{minipage}
	\hspace{1em} 
	\begin{minipage}{0.20\linewidth}
		\centering
		\includegraphics[width=\linewidth]{Img/32.pdf}
	\end{minipage}
	
	\vspace{0.3em}
	
	\vspace{0.3em}
	\fuentepropia
\end{minipage}

Las figuras anteriores se consiguen con los siguientes comandos.

\vspace{0.6em}
\noindent
\begin{tcolorbox}[breakable,enhanced,title=\bf Instrucción en \verb"Mathematica", top=2pt,bottom=2pt,boxsep=0pt,before skip=2pt,after skip=0pt]
\begin{Verbatim}[formatcom=\letraverbatim]
ContourPlot3D[(z-2)^2/4+(y-3)^2/9==1-x/4,{x,0,4},
{y,0,6},{z,0,4},PlotPoints->80,MeshStyle->Gray,
Mesh->5,ContourStyle->Directive[RGBColor[1,0.3,0.3],
Specularity[RGBColor[1,1,0],10]]]
\end{Verbatim}
\end{tcolorbox}
\vspace{0.7em}
\noindent

En el plano $yz$ $(x=0),$ se satisface que $\frac{\left( z-2\right) ^2}{4}+%
\frac{\left( y-3\right) ^2}{9}=1,$ esta ecuación representa una elipse con centro en el punto $(2,3),$ alargada en el sentido de las $y$, y puede parametrizarse en la forma
$y=3\sin t+3,\,\, z=2+2\cos t,\,\,t\in [0,2\pi ].$
\vskip0.3cm
La proyección en el plano $xy,$ se consigue cuando $z=2,$ y es representada por la ecuación
$\frac{\left( y-3\right) ^2}{9}=1-\frac{x}{4},$ que se puede reescrbir como $x=-\frac{4}{9}y^2+\frac{8}{3}y.$ Esta ecuación representa una parábola que abre hacia la parte negativa del eje $x$ y posee
vértice en el punto $\left( 4,3\right) .$ Una parametrización de esta curva es
$y=3t,\,\, x=-\frac{4}{9}\left( 3t\right) ^2+\frac{8}{3}(3t) = 8t-4t^2,$ para $t\in [0,2].$
\vskip0.3cm
Cuando $y=3,$ obtendremos la proyección del sólido en el plano $xz,$
esto es $\frac{\left( z-2\right) ^2}{4}=1-\frac{x}{4},$ de donde
$x=4(1-\frac{(z-2)^2}{4}) = 4z-z^2.$
Esta curva se puede parametrizar como $x=4t-t^2,\,\,z=t,$ para $t\in [0,4].$
\vskip0.3cm
Con los resultados anteriores se presentan ahora las figuras.

\vskip0.2cm
\noindent
\begin{minipage}{\textwidth}
	\centering
	\captionof{figure}{Proyecciones del sólido en los planos coordenados.}
	\label{fig:proyeccion_solido_eliptico_desplazado_2}
	\begin{minipage}{0.30\linewidth}
		\centering
		\vspace{-10mm}
		\includegraphics[width=\linewidth]{Fig/329.pdf}
	\end{minipage}
	\hspace{1em} 
	\begin{minipage}{0.30\linewidth}
		\centering
		\vspace{-10mm}
		\includegraphics[width=\linewidth]{Fig/330.pdf}
	\end{minipage}
	\hspace{1em} 
	\begin{minipage}{0.30\linewidth}
		\centering
		\vspace{-10mm}
		\includegraphics[width=\linewidth]{Fig/331.pdf}
	\end{minipage}
	
	\vspace{-8mm}
	\fuentepropia
\end{minipage}
\vskip0.2cm

% \begin{figure}[H]
% \centering
% % Primer gráfico
% \begin{pspicture}(-0.5,-0.4)(3.5,3) %\psgrid
% \psset{algebraic,unit=0.5}
% \psaxes[Dx=2,Dy=2,linecolor=gray]{->}(0,0)(-0.5,-1.2)(5,6.2)[$x$,-90][$z$,0]
% \pscustom[fillstyle=solid,fillcolor=gray!30,opacity=0.7,linestyle=none]{
%   \psparametricplot{0}{4}{4*t-t^2,t}
%   \psline(0,4)(0,0)
% }
% \psparametricplot{0}{4}{4*t-t^2,t}
% \rput(2.5,4.5){$\frac{(z-2)^2}{4}=1-\frac{x}{4}$}
% \end{pspicture}

% \vspace{0.5cm}

% % Segundo gráfico
% \begin{pspicture}(-0.5,-0.4)(3,3) % \psgrid
% \psset{algebraic,yunit=0.35,xunit=0.45}
% \psaxes[Dx=2,Dy=3,linecolor=gray]{->}(0,0)(-1.5,-1.2)(6.2,8)[$x$,-90][$y$,0]
% \pscustom[fillstyle=solid,fillcolor=gray!30,opacity=0.7,linestyle=none]{
%   \psparametricplot{0}{2}{8*t-4*t^2,3*t}
%   \psline(0,6)(0,0)
% }
% \psparametricplot{0}{2}{8*t-4*t^2,3*t}
% \rput(4.5,6){$\frac{(y-3)^2}{9}=1-\frac{x}{4}$}
% \end{pspicture}

% \vspace{0.5cm}

% % Tercer gráfico
% \begin{pspicture}(-1,-0.4)(3.5,3) % \psgrid
% \psset{algebraic,unit=0.5}
% \psaxes[Dx=2,Dy=2,linecolor=gray]{->}(0,0)(-0.5,-1.2)(7,6)[$y$,-90][$z$,0]
% \pscustom[fillstyle=solid,fillcolor=gray!30,opacity=0.7,linestyle=none]{
%   \psparametricplot{0}{\psPiTwo}{3+3*sin(t),2+2*cos(t)}
% }
% \psparametricplot{0}{\psPiTwo}{3+3*sin(t),2+2*cos(t)}
% \psline[linecolor=blue](5,3.5)(5,0.5)
% \rput(4,4.7){$\frac{(z-2)^2}{4}+\frac{(y-3)^2}{9}=1$}
% \rput(5,3.5){$\bullet$}
% \rput(5,0.5){$\bullet$}
% \rput(5,2.5){$\bullet$}
% \rput(4.5,0.8){$A$}
% \rput(4.5,3.2){$B$}
% \rput(5.5,2.5){$P$}
% \end{pspicture}
% \caption{Proyecciones del sólido en los planos coordenados.}
% \end{figure}

Estas figuras resultan útiles al momento de plantear las integrales triples que se requiere evaluar para calcular el volumen del sólido, veámoslo
\vskip0.3cm
A partir de la ecuación del paraboloide %$\frac{(z-2)^2}{4}+\frac{(y-3)^2}{9}=1-\frac{x}{4}$,
se obtiene {\footnotesize $x=4\left( 1-\frac{(z-2)^2}{4}-\frac{(y-3)^2}{9}\right),$} por tanto

\[
\frac{(z-2)^2}{4} + \frac{(y-3)^2}{9} = 1 - \frac{x}{4},
\]
se obtiene
{\small
\[
x = 4\left( 1 - \frac{(z-2)^2}{4} - \frac{(y-3)^2}{9} \right),
\]
}

por tanto
\[
\displaystyle
\int_{0}^{6}
\int_{2 - \frac{2}{3}\sqrt{6y - y^2}}^{2 + \frac{2}{3}\sqrt{6y - y^2}}
\int_{0}^{4 - (z - 2)^2 - \frac{4}{9}(y - 3)^2}
dx\,dz\,dy = 12\pi,
\]
también
\[
\displaystyle
\int_{0}^{4}
\int_{3 - \frac{3}{2}\sqrt{-z^2 + 4z}}^{3 + \frac{3}{2}\sqrt{-z^2 + 4z}}
\int_{0}^{4 - (z - 2)^2 - \frac{4}{9}(y - 3)^2}
dx\,dy\,dz = 12\pi.
\]

De la ecuación original, se sigue que
\[
(z - 2)^2 = 4\left( 1 - \frac{x}{4} - \frac{(y - 3)^2}{9} \right)
= \frac{1}{9}\left( 24y - 9x - 4y^2 \right),
\]
por tanto
\[
z = 2 \pm \frac{1}{3}\sqrt{-9x - 4y^2 + 24y}.
\]

Entonces, las integrales triples, integrando primero con respecto a \( z \), son:
\[
\displaystyle
\int_{0}^{4}
\int_{3 - \frac{3}{2}\sqrt{4 - x}}^{3 + \frac{3}{2}\sqrt{4 - x}}
\int_{2 - \frac{1}{3}\sqrt{-9x - 4y^2 + 24y}}^{2 + \frac{1}{3}\sqrt{-9x - 4y^2 + 24y}}
dz\,dy\,dx = 12\pi,
\]
\[
\displaystyle
\int_{0}^{6}
\int_{0}^{\frac{8}{3}y - \frac{4}{9}y^2}
\int_{2 - \frac{1}{3}\sqrt{-9x - 4y^2 + 24y}}^{2 + \frac{1}{3}\sqrt{-9x - 4y^2 + 24y}}
dz\,dx\,dy = 12\pi.
\]

De la misma forma,
\[
(y - 3)^2 = 9\left( 1 - \frac{x}{4} - \frac{(z - 2)^2}{4} \right)
= \frac{9}{4}(4z - x - z^2),
\]
entonces
\[
y = 3 \pm \frac{3}{2}\sqrt{4z - x - z^2}.
\]
Con lo cual, el volumen está dado por
{\small
\[\bigintss_{0}^{4}\bigintss_{2-\sqrt{4-x}}^{2+\sqrt{4-x}}\bigintss_{3-\frac{3}{2}\sqrt{4z-x-z^2} }^{3+\frac{3}{2}\sqrt{4z-x-z^2}}dydzdx= 12\pi\]
\[\bigintss_{0}^{4}\bigintss_{0}^{4z-z^2}\bigintss_{3-\frac{3}{2}\sqrt{4z-x-z^2}}^{3+\frac{3}{2}\sqrt{4z-x-z^2} }dydxdz=12\pi.\]
}

Algunas de estas integrales se evalúan con los siguientes comandos:

\vspace{0.6em}
\noindent
\begin{tcolorbox}[breakable,enhanced,title=\bf Instrucción en \verb"Mathematica", top=2pt,bottom=2pt,boxsep=0pt,before skip=2pt,after skip=0pt]
\begin{Verbatim}[formatcom=\letraverbatim]
Integrate[1,{x,0,4},{y,3-3*Sqrt[4-x]/2,3+
3*Sqrt[4-x]/2},{z,2-Sqrt[24*y-4y^2-9*x]/3,
2+Sqrt[24*y-4y^2-9*x]/3}]
Integrate[1,{y,0,6},{z,2-2Sqrt[6*y-y^2]/3,2
+2Sqrt[6*y-y^2]/3},{x,0,4-(z-2)^2-4*(y-3)^2/9}]
Integrate[1,{z,0,4},{x,0,4*z-z^2},{y,3-
3*Sqrt[4*z-x-z^2]/2,3+3*Sqrt[4*z-x-z^2]/2}]
Integrate[1,{z,0,4},{y,3-3/2*Sqrt[4z-z^2],3+
3/2*Sqrt[4z-z^2]},{x,0,4-(z-2)^2-4/9*(y-3)^2}]
\end{Verbatim}
\end{tcolorbox}
\vspace{0.7em}
\noindent


Se presenta ahora, otra manera de hallar el volumen.
\vskip0.2cm

%\begin{tcolorbox}[title=\bfseries Parametrizando el sólido]

La proyección del sólido en el plano $yz$ corresponde a una elipse cuyo interior se puede parametrizar como:
$y=3+3u\sin t,\ z=2+2u\cos t,$ en donde $u \in [0,1]$, $t \in [0,2\pi].$
Con dos puntos $A$ y $B$ sobre la elipse (con igual coordenada $y$) se escoge un punto $P$;
este punto se proyecta sobre el paraboloide. Con una combinación convexa entre $P$ y $Q$
se parametriza el sólido.


\vskip0.2cm
\noindent
\begin{minipage}{\textwidth}
	\centering
	\captionof{figure}{$(2uv, 2(1 - v + uv)).$}
	\label{fig:curva_parametrica_uv_3}
	\vspace{-18mm}
	\begin{minipage}{0.40\linewidth}
		\centering
		\includegraphics[width=\linewidth]{Fig/332.pdf}
	\end{minipage}
	
	\vspace{-16mm}
	\fuentepropia
\end{minipage}
\vskip0.2cm

Las coordenadas $x$ del punto $Q$ que está sobre el paraboloide satisfacen que
$x=4-4u^2.$ Por tanto, los puntos mencionados son
\[P\left(0,3+3u\sin t,2+2u\cos t\right)\] y \[Q\left( 4-4u^2,3+3u\sin t+3,2+2u\cos t\right).\]
\vskip0.2cm
El sólido se parametriza simplificando la expresión $(1-w)Q+wP$,
esto permite definir la función

\[{\bf r}(u,t,w)=\left\langle 4(1-w)(1-u^2),3+3u\sin t,2+2u\cos t\right\rangle,\]
\vskip0.2cm
con $\, u\in [0,1],\,t\in [0,2\pi ],\, w\in [0,1] $.
Con esta función, hallamos los siguientes vectores
\begin{eqnarray*}
{\bf r}_{t}&=& \langle 0,3u\cos t,-2u\sin t \rangle \\
{\bf r}_{u}&=& \langle 8\left( w-1\right) u,3\sin t,2\cos t\rangle \\
{\bf r}_{w}&=& \langle 4u^2-4,0,0\rangle \\
{\bf r}_{u}\times {\bf r}_{w} &=& \langle 0,8(1-u^2)\cos(t),12(u^2-1)\sin (t)\rangle.
\end{eqnarray*}
\vskip0.2cm
Entonces
$ \left|{\bf r}_{t}\cdot \left({\bf r}_{u}\times {\bf r}_{w}\right) \right| = |24u^{3}-24u|.$

Con lo cual, el volumen del sólido esta dado por

\[\int_{0}^{1}\int_{0}^{1}\int_{0}^{2\pi }\left| 24u^{3}-24u\right|dtdudw\\
= 12\pi.\]
%\end{tcolorbox}

\subsection*{Una esfera y un paraboloide} \index{Esfera}
La esfera con ecuación $x^2+y^2+z^2=2$ y el paraboloide
$z=x^2+y^2$ se intersecan en $z=1.$
Las siguientes figuras muestran esta situación y el sólido acotado por estas superficies.

\vskip0.3cm
\noindent
\begin{minipage}{\textwidth}
	\centering
	\captionof{figure}{Sólido definido por $x^2 + y^2 \leq z \leq \sqrt{2 - x^2 - y^2}$.}
	\label{fig:solido_paraboloide_esfera}
	\begin{minipage}{0.26\linewidth}
		\centering
		\includegraphics[width=\linewidth]{Img/81a.pdf}
	\end{minipage}
	\hspace{1em}
	\begin{minipage}{0.26\linewidth}
		\centering
		\includegraphics[width=\linewidth]{Img/81b.pdf}
	\end{minipage}
	
	\vspace{0.3em}
	\fuentepropia
\end{minipage}
\vskip0.3cm

Estas figuras se consiguen con las siguientes indicaciones:


\vspace{0.6em}
\noindent
\begin{tcolorbox}[breakable,enhanced,title=\bf Instrucción en \verb"Mathematica", top=2pt,bottom=2pt,boxsep=0pt,before skip=2pt,after skip=0pt]
\begin{Verbatim}[formatcom=\letraverbatim]
ContourPlot3D[{z==x^2+y^2,x^2+y^2+z^2==2},
{x,-1.5,1.5},{y,-1.5,1.5},{z,-1.5,1.5},Mesh->False,
ContourStyle->{Directive[Cyan,Opacity[0.9]],
Directive[Orange,Opacity[0.65]]},PlotPoints->60]
RegionPlot3D[x^2+y^2<z<Sqrt[2-x^2-y^2],{x,-1,1},
{y,-1,1},{z,0,1.5},PlotStyle->Orange,Mesh->False,
PlotPoints->80,Ticks->{{-1,0,1},{-1,0,1},{0,1}}]
\end{Verbatim}
\end{tcolorbox}
\vspace{0.7em}
\noindent

La proyección en el plano $xy$ está dada por $x^2+y^2=1,$ que corresponde a una circunferencia de radio $1$, mientras que en el plano $xz$ la proyección es una porción de parábola y una circunferencia de radio $\sqrt{2}$.

\vskip0.2cm
\noindent
\begin{minipage}{\textwidth}
	\centering
	\captionof{figure}{Regiones $x^2\leq z \leq \sqrt{2-x^2}$, $x^2+y^2\leq 1.$}
	\begin{minipage}{0.35\linewidth}
		\centering
		\vspace{-10mm}
		\includegraphics[width=\linewidth]{Fig/333.pdf}
	\end{minipage}
	\hspace{1em}
	\begin{minipage}{0.35\linewidth}
		\centering
		\vspace{-10mm}
		\includegraphics[width=\linewidth]{Fig/334.pdf}
	\end{minipage}
	
	\vspace{-8mm}
	\fuentepropia
\end{minipage}
\vskip0.2cm

El volumen del sólido acotado por estas superficies está dado por
\begin{multline*}
\int_{-1}^{1}\int_{-\sqrt{1-x^2}}^{\sqrt{1-x^2}}\int_{x^2+y^2}^{\sqrt{2-x^2-y^2}}dzdydx = \\ \int_{-1}^{1}\int_{-\sqrt{1-y^2}}^{\sqrt{1-y^2}}\int_{x^2+y^2}^{\sqrt{2-x^2-y^2}}dzdxdy= \frac{\pi}{6} ( 8\sqrt{2}-7)
\end{multline*}

Al integrar primero con respecto a $y$, el sólido queda dividido en dos porciones, en una de ellas $y$ toma valores desde la parte negativa de la esfera hasta la parte positiva de la misma; en la otra porción $y$ recorre los puntos desde la parte negativa del paraboloide hasta su parte positiva.
Por lo tanto, el volumen en este caso está dado por
%\vskip0.2cm
\[
\int_{-1}^{1} \int_{1}^{\sqrt{2 - x^2}} \int_{-\sqrt{2 - x^2 - z^2}}^{\sqrt{2 - x^2 - z^2}} dy\, dz\, dx
+ \int_{-1}^{1} \int_{x^2}^{1} \int_{-\sqrt{z - x^2}}^{\sqrt{z - x^2}} dy\, dz\, dx
\approx 2.259
\]
\vspace{-2mm}
{\small
\[
\int_{1}^{\sqrt{2}} \int_{-\sqrt{2 - z^2}}^{\sqrt{2 - z^2}} \int_{-\sqrt{2 - x^2 - z^2}}^{\sqrt{2 - x^2 - z^2}} dy\, dx\, dz
+ \int_{0}^{1} \int_{-\sqrt{z}}^{\sqrt{z}} \int_{-\sqrt{z - x^2}}^{\sqrt{z - x^2}} dy\, dx\, dz
= \frac{\pi}{6}(8\sqrt{2} - 7)
\]
}

\vskip0.2cm
Las anteriores integrales se evaluaron en \verb"Mathematica" usando las si\-guientes instrucciones:

\vspace{0.6em}
\noindent
\begin{tcolorbox}[breakable,enhanced,title=\bf Instrucción en \verb"Mathematica", top=2pt,bottom=2pt,boxsep=0pt,before skip=2pt,after skip=0pt]
\begin{Verbatim}[formatcom=\letraverbatim]
Integrate[1,{x,-1,1},{y,-Sqrt[1-x^2],Sqrt[1-x^2]},
{z,x^2+y^2,Sqrt[2-x^2-y^2]}]
Integrate[1,{y,-1,1},{x,-Sqrt[1-y^2],Sqrt[1-y^2]},
{z,x^2+y^2,Sqrt[2-x^2-y^2]}]
Integrate[1,{x,-1,1},{z,1,Sqrt[2-x^2]},
{y,-Sqrt[2-x^2-z^2],Sqrt[2-x^2-z^2]}]+Integrate[1,
{x,-1,1},{z,x^2,1},{y,-Sqrt[z-x^2],Sqrt[z-x^2]}]
Integrate[1,{z,1,Sqrt[2]},{x,-Sqrt[2-z^2],
Sqrt[2-z^2]},{y,-Sqrt[2-x^2-z^2],Sqrt[2-x^2-z^2]}]
+Integrate[1,{z,0,1},{x,-Sqrt[z],Sqrt[z]},
{y,-Sqrt[z-x^2],Sqrt[z-x^2]}]
\end{Verbatim}
\end{tcolorbox}
\vspace{0.7em}
\noindent


En el orden $dxdydz$ y $dxdzdy,$ la situación es muy similar al caso $dydzdx$,
por la simetría de la figura con respecto al plano
$xz$ o $yz.$ Con el siguiente procedimiento también se halla el volumen del sólido.

\vskip0.2cm
%\begin{tcolorbox}[title=\bfseries Parametrizando el sólido]
La proyección del sólido en el plano $xy$ en la forma $\left( u\cos v,u\sin v\right) ,$ con $u\in [0,1],$ $v\in [0,2\pi ],$ la coordenada $z$ sobre la esfera y el paraboloide están dadas, respectivamente, por $z =\sqrt{2-x^2-y^2}= \sqrt{2-u^2}$ y $z=x^2+y^2= u^2.$
El sólido se parametriza simplificando una combinación convexa de la forma $(1-w)P+wQ,$
en donde $P\left( u\cos v,u\sin v,\sqrt{2-u^2}\right)$ es un punto sobre la esfera y
$Q\left( u\cos v,u\sin v,u^2\right)$ es un punto sobre el paraboloide.
\vspace{-2mm}
\begin{eqnarray*}
{\bf r}\left( u,v,w\right) &=& \left\langle u\cos (v) ,u\sin(v),(1-w) \sqrt{2-u^2}+wu^2\right \rangle ,
\end{eqnarray*}
%\vskip0.2cm
con $u\in [0,1],$ $v\in [0,2\pi ],$ $u\in [0,1].$

\vskip0.2cm
Con esta función hallamos los vectores:
\begin{eqnarray*}
{\bf r}_{u} &=& \langle \cos (v) ,\sin (v) ,-\frac{\left( 1-w\right) u}{\sqrt{2-u^2}}+2wu\rangle \\
{\bf r}_{v}&=& \langle -u\sin (v) ,u\cos (v) ,0\rangle \\
{\bf r}_{w} &=& \langle 0,0,-\sqrt{2-u^2}+u^2\rangle.
\end{eqnarray*}
%\vskip0.2cm
El elemento diferencial de volumen es
$\left| {\bf r}_{w}\cdot \left( {\bf r}_{u}\times {\bf r}_{v}\right) \right| = \left| u^{3}-u\sqrt{2-u^2}\right|$
cuya integral nos proporciona el volumen del sólido:
\vskip0.2cm
\[\int_{0}^{1}\int_{0}^{1}\int_{0}^{2\pi }\left| u^{3}-u\sqrt{2-u^2}\right|
dvdudw= \frac{4}{3}\pi \sqrt{2}-\frac{7}{6}\pi.\]
%\end{tcolorbox}
\vskip0.2cm
En coordenadas cilíndricas este valor se obtiene evaluando la siguiente integral
\[\int_{0}^{2\pi }\int_{0}^{1}\int_{r^2}^{\sqrt{2-r^2}}rdzdrd\theta
= \frac{1}{6}\pi \left( 8\sqrt{2}-7\right).\]
Para usar coordenadas esféricas, la esfera posee ecuación $\rho =%
\sqrt{2}$ y el paraboloide $\rho \cos \phi =\left( \rho \sin \phi \cos
\theta \right) ^2+\left( \rho \sin \phi \sin \theta \right) ^2$
que se simplifica como
$\rho =\frac{\cos \phi }{\sin ^2\left( \phi \right)}= \cot \phi \csc \phi ,$
como la intersección se da en $z=1,$ con una circunferencia de radio $1,$ entonces
$\phi =\pi /4.$ Es por lo anterior, que el volumen está dado por:
\begin{multline*}
\int_{0}^{\frac{\pi}{4}}\int_{0}^{2\pi }\int_{0}^{\sqrt{2}}\rho ^2\sin \phi
d\rho d\theta d\phi \\
+\int_{\frac{\pi}{4}}^{\frac{\pi}{2}}\int_{0}^{2\pi }\int_{0}^{\cot \phi \csc \phi }\rho ^2\sin \phi d\rho d\theta d\phi = \frac{\pi}{6} \left( 8\sqrt{2}-7\right).
\end{multline*}

\subsection*{Un plano y un paraboloide}
El plano y el paraboloide con ecuaciones $z=6-y$ y $z=x^2+y^2$,
determinan el sólido que se presenta a continuación:

\vskip0.2cm
\begin{figure}[H]
\centering
\caption{Sólido determinado por $x^2 + y^2 \leq z \leq 6 - y$.}
\label{fig:solido_paraboloide_plano}
\includegraphics[height=4.5cm]{Img/14d.pdf}
\includegraphics[height=4.5cm]{Img/14a.pdf}
\includegraphics[height=4.5cm]{Img/14b.pdf}
\fuentepropia
\end{figure}


En \verb"Mathematica", este gráfico se consigue usando las siguientes instrucciones:

\vspace{0.6em}
\noindent
\begin{tcolorbox}[breakable,enhanced,title=\bf Instrucción en \verb"Mathematica", top=2pt,bottom=2pt,boxsep=0pt,before skip=2pt,after skip=0pt]
\begin{Verbatim}[formatcom=\letraverbatim]
ContourPlot3D[{z==6-y,z==x^2+y^2},{x,-4,4},{y,-4,4},
{z,0,10},Mesh->False,ContourStyle->{Orange,Cyan}]
RegionPlot3D[x^2+y^2<z<6-y&&-3<y<2,{x,-5/2,5/2},
{y,-3,2},{z,0,9},PlotPoints->90,Mesh->None]
\end{Verbatim}
\end{tcolorbox}
\vspace{0.7em}
\noindent

\vskip0.2cm
\noindent
\begin{minipage}{\textwidth}
	\centering
	\captionof{figure}{Proyecciones del sólido en los planos coordenados.}
	\begin{minipage}{0.30\linewidth}
		\centering
		\vspace{-11mm}
		\includegraphics[width=\linewidth]{Fig/335.pdf}
	\end{minipage}
	\hspace{1em} 
	\begin{minipage}{0.30\linewidth}
		\centering
		\vspace{-11mm}
		\includegraphics[width=\linewidth]{Fig/336.pdf}
	\end{minipage}
	\hspace{1em} 
	\begin{minipage}{0.30\linewidth}
		\centering
		\vspace{-11mm}
		\includegraphics[width=\linewidth]{Fig/337.pdf}
	\end{minipage}
	
	\vspace{-8mm}
	\fuentepropia
\end{minipage}
\vskip0.2cm

Al sustituir $z$ de la ecuación del plano en la ecuación del paraboloide, se sigue que $6-y=x^2 +y^2,$
con lo cual $x=\pm \sqrt{6-y^2-y}$ o también $y=-\frac{1}{2}\pm \frac{1}{2}\sqrt{25-4x^2}.$
Al reemplazar $y$ en la ecuación del paraboloide, se sigue que
$z=x^2 +y^2=x^2 +\left( 6-z\right)^2,$ esta expresión se simplifica como
$x^2+z^2-13z=-36$ y se puede reescribir como
$x^2+\left(z-\frac{13}{2}\right)^2=\frac{25}{4}.$
\vskip0.2cm
Las expresiones anteriores permiten graficar la proyección del sólido en los planos coordenados, con el propósito de plantear adecuadamente las integrales necesarias para determinar el volumen ($V$) u otros elementos relativos al sólido. En este caso,

\vspace{-6mm}
\begin{eqnarray*}
V&=& \int {-3}^2\int _{-\sqrt{6-y-y^2}}^{\sqrt{6-y-y^2}} \int _{x^2+y^2}^{6-y}dzdxdy = \int _{-3}^2\int _{y^2}^{6-y} \int _{-\sqrt{z-y^2}}^{\sqrt{z-y^2}}dxdzdy \\
&=&\int 4^9\int _{-\sqrt{z}}^{6-z} \int _{-\sqrt{z-y^2}}^{\sqrt{z-y^2}}dxdydz+\int _0^4\int_{-\sqrt{z}}^{\sqrt{z}}\int _{-\sqrt{z-y^2}}^{\sqrt{z-y^2}}dxdydz\\
&=& \int_{-\frac{5}{2}}^{\frac{5}{2}} \int_{-\frac{1}{2}-\sqrt{\frac{25}{4}-x^2}}^{-\frac{1}{2}+\sqrt{\frac{25}{4}-x^2}}\int_{x^2+y^2}^{6-y}dzdydx = \frac{625\pi}{32} \end{eqnarray*}
\vskip0.4cm
La proyección en  el plano $xz$ evidencia que en el orden de integración $dydxdz$ o $dydzdx$ puede ser necesario evaluar tres o cuatro integrales triples, según el caso.
En seguida se muestra una forma simplificada de obtener el volumen.
\begin{multline*}
\bi _{-\frac{5}{2}}^{\frac{5}{2}}\bi_{\frac{13}{2}-\sqrt{\frac{25}{4}-x^2}}^{\frac{13}{2}+\sqrt{\frac{25}{4}-x^2}}\bi_{-\sqrt{z-x^2}}^{6-z}dydzdx \\Z +\bi_{-\sqrt{6}}^{\sqrt{6}}\bi_{x^2}^{\frac{13}{2}-\sqrt{\frac{25}{4}-x^2}}\bi_{-\sqrt{z-x^2}}^{\sqrt{z-x^2}}dydzdx = \frac{625\pi}{32}
\end{multline*}

En \verb"Mathematica", algunas de estas integrales pueden evaluarse mediante las siguientes instrucciones:

\vspace{0.6em}
\noindent
\begin{tcolorbox}[breakable,enhanced,title=\bf Instrucción en \verb"Mathematica",top=2pt,bottom=2pt,boxsep=0pt,before skip=2pt,after skip=0pt]
\begin{Verbatim}[formatcom=\letraverbatim]
Integrate[1,{y,-3,2},{x,-Sqrt[6-y-y^2],Sqrt[6-y-y^2]},
{z,x^2+y^2,6-y}]
Integrate[1,{x,-5/2,5/2},{y,-1/2-Sqrt[25/4-x^2],
1/2+Sqrt[25/4-x^2]},{z,x^2+y^2,6-y}]
Integrate[1,{y,-3,2},{z,y^2,6-y},{x,-Sqrt[z-y^2],
Sqrt[z-y^2]}]
\end{Verbatim}
\end{tcolorbox}
\vspace{0.7em}
\noindent

Ahora se parametriza la región mencionada antes.

%\begin{tcolorbox}[title=\bfseries Parametrizando el sólido]
La base del sólido en el plano $xy$ está dada por
$x^2+y^2=6-y$ que equivale a $x^2+\left( y+\frac{1}{2}\right)^2=\frac{25}{4}.$
Esta región se parametriza al hacer $x=\frac{5}{2}u\cos(t)$, $y=-\frac{1}{2}+\frac{5}{2}u\sin(t)$
con $u\in [0,1],$ $t\in [0,2\pi].$ Un punto sobre el plano $z=6-y,$ tiene la forma
\vskip0.2cm
\[P\left(\frac{5}{2}\cos(t),-\frac{1}{2}+\frac{5}{2}\sin(t),6-(-\frac{1}{2}+\frac{5}{2}\sin(t))\right)\]
\vskip0.2cm
y un punto arbitrario sobre el paraboloide es de  forma
\vskip0.2cm
\[Q\left(\frac{5}{2}\cos(t),-\frac{1}{2}+\frac{5}{2}\sin(t),(\frac{5}{2}\cos(t))^2+(-\frac{1}{2}+\frac{5}{2}\sin(t))^2\right).\]
Formando segmentos rectilíneos que unan un punto del plano con un punto sobre el paraboloide se parametriza el sólido. Esto es, simplificando la expresión $(1-w)P+wQ,$ se define la función:
%\vskip0.2cm
\[{\bf r}\left(t,u,w\right)=\frac{5}{2} \left\langle u \cos (t),u\sin(t)-\frac{1}{5},\frac{5w}{2}(u^2-1)-u\sin(t)+\frac{13}{5} \right\rangle.\]
%\vskip0.2cm
con $u\in [0,1],$ $t\in [0,2\pi],$ $w\in [0,1],$
la cual parametriza el sólido. También se hallan los vectores:
{\small
\begin{eqnarray*}
{\bf r}_{t} &=& \frac{5}{2}\left\langle - u \sin (t), u \cos (t),- u \cos (t) \right\rangle \\
{\bf r}_{u} &=& \frac{5}{2}\left\langle \cos (t), \sin (t),5 u w-\sin (t)\right\rangle \\
{\bf r}_{w} &=& \frac{25}{4}\left\langle 0,0, \left(u^2-1\right) \right\rangle
\end{eqnarray*}
}
%\vskip0.2cm
El elemento diferencial de volumen es
$\left|{\bf r}_{t}\cdot\left({\bf r}_{u}\times {\bf r}_{w}\right)\right| =\frac{625}{16} u \left(1-u^2\right).$
Por último, evaluamos la integral que determina el volumen del sólido:
\[ \int_{0}^{1}\int_{0}^{2\pi}\int_{0}^{1} \left|{\bf r}_{t}\cdot\left({\bf r}_{u}\times {\bf r}_{w}\right)\right|dudtdw=\frac{625\pi}{32}\]
%\end{tcolorbox}
%\vskip0.2cm
El procedimiento descrito antes, se puede efectuar mediante la ejecución de las siguientes instrucciones: % La parametrización de la base se denomina \verb"A" y también define la función \verb"r" que parametriza el sólido. El producto triple se denomina \verb"B" que se integra en los límites establecidos.

\vspace{0.6em}
\noindent
\begin{tcolorbox}[breakable,enhanced,title=\bf Instrucción en \verb"Mathematica", top=2pt,bottom=2pt,boxsep=0pt,before skip=2pt,after skip=0pt]
\begin{Verbatim}[formatcom=\letraverbatim]
{x,y}={5/2*u*Cos[t],-1/2+5/2*u*Sin[t]};
r=Simplify[w*{x,y,x^2+y^2}+(1-w)*{x,y,6-y}];
B=Abs[Dot[D[r,t],Cross[D[r,u],D[r,w]]]]
Integrate[B,{t,0,2*Pi},{u,0,1},{w,0,1}]
\end{Verbatim}
\end{tcolorbox}
\vspace{0.7em}
\noindent

\subsection*{Un paraboloide, un cilindro y un plano}

Consideremos el sólido que se halla acotado bajo el paraboloide con ecuación $z=4-x^2-y^2$,
sobre el plano $z=0$ y fuera del cilindro con ecuación $x^2+y^2=2x$.
\vskip0.4cm
\vspace{-0.8mm}
\noindent
\begin{minipage}{\textwidth}
	\centering
	\captionof{figure}{Sólido determinado por $0\leq z\leq 4-x^2-y^2$, $2x \leq x^2+y^2$.}
	\label{fig:solido}
	\begin{minipage}{0.21\linewidth}
		\centering
		\vspace{-12mm}
		\includegraphics[width=\linewidth]{Img/3a.pdf}
	\end{minipage}
	\hspace{1em}
	\begin{minipage}{0.27\linewidth}
		\centering
		\vspace{-12mm}
		\includegraphics[width=\linewidth]{Fig/338.pdf}
	\end{minipage}
	\hspace{1em}
	\begin{minipage}{0.21\linewidth}
		\centering
		\vspace{-12mm}
		\includegraphics[width=\linewidth]{Img/3c.pdf}
	\end{minipage}
	
	\vspace{-8mm}
	\fuentepropia
\end{minipage}


\vskip0.2cm

Las instrucciones para obtener la figura anterior son las siguientes:

\vspace{0.6em}
\noindent
\begin{tcolorbox}[breakable,enhanced,title=\bf Instrucción en \verb"Mathematica", top=2pt,bottom=2pt,boxsep=0pt,before skip=2pt,after skip=0pt]
\begin{Verbatim}[formatcom=\letraverbatim]
RegionPlot3D[0<z<4-x^2-y^2&&x^2+y^2>2*x-0.05,
{x,-2,2},{y,-2,2},{z,0,4},Mesh->False,
PlotStyle->Orange,PlotPoints->90]
\end{Verbatim}
\end{tcolorbox}
\vspace{0.7em}
\noindent

La proyección de este sólido en el plano $xy$ es la siguiente:
\vskip0.2cm

\noindent
\begin{minipage}{\textwidth}
	\centering
	\captionof{figure}{Región $2x \leq x^2+y^2 \leq 4$.}
	\begin{minipage}{0.31\linewidth}
		\centering
		\vspace{-14mm}
		\includegraphics[width=\linewidth]{Fig/339.pdf}
	\end{minipage}
	\hspace{1em}
	\begin{minipage}{0.31\linewidth}
		\centering
		\vspace{-14mm}
		\includegraphics[width=\linewidth]{Fig/340.pdf}
	\end{minipage}
	
	\vspace{-8mm}
	\fuentepropia
\end{minipage}

\vspace{0.2cm} 

El relleno de esta región sugiere varios elementos: los límites de integración, en el orden $dzdydx$ se evalúa mediante tres integrales triples y en el orden $dzdxdy$ es necesario evaluar cuatro integrales triples. El volumen está dado por:

\vspace{-4mm}
{\small
\begin{multline*}
\int_{-2}^{0}\int_{-\sqrt{4-x^2}}^{\sqrt{4-x^2}}\int_{0}^{4-x^2-y^2}dzdydx+\int_{0}^2\int_{\sqrt{2x-x^2}}^{\sqrt{4-x^2}}\int_{0}^{4-x^2-y^2}dzdydx \\
+\int_{0}^2\int_{-\sqrt{4-x^2}}^{-\sqrt{2x-x^2}}\int_{0}^{4-x^2-y^2}dzdydx= \frac{11}{2}\pi
\end{multline*}
}
\vspace{-4mm}
{\small
\begin{multline*}
\int_{1}^2\int_{-\sqrt{4-y^2}}^{\sqrt{4-y^2}}\int_{0}^{4-x^2-y^2}dzdxdy
+\int_{-2}^{-1}\int_{-\sqrt{4-y^2}}^{\sqrt{4-y^2}}\int_{0}^{4-x^2-y^2}dzdxdy+ \\
\int_{-1}^{1}\int_{-\sqrt{4-y^2}}^{1-\sqrt{1-y^2}}\int_{0}^{4-x^2-y^2}dzdxdy
+\int_{-1}^{1}\int_{1+\sqrt{1-y^2}}^{\sqrt{4-y^2}}\int_{0}^{4-x^2-y^2}dzdxdy= \frac{11}{2}\pi
\end{multline*}
}

Estas integrales se evalúan en \verb"Mathematica" con las siguientes instrucciones

\vspace{0.6em}
\noindent
\begin{tcolorbox}[breakable,enhanced,title=\bf Instrucción en \verb"Mathematica", top=2pt,bottom=2pt,boxsep=0pt,before skip=2pt,after skip=0pt]
\begin{Verbatim}[formatcom=\letraverbatim]
Integrate[1,{x,-2,0},{y,-Sqrt[4-x^2],Sqrt[4-x^2]},
{z,0,4-x^2-y^2}]+Integrate[1,{x,0,2},{y,Sqrt[2*x-x^2],
Sqrt[4-x^2]},{z,0,4-x^2-y^2}]+Integrate[1,{x,0,2},
{y,-Sqrt[4-x^2],-Sqrt[2*x-x^2]},{z,0,4-x^2-y^2}]
Integrate[1,{y,1,2},{x,-Sqrt[4-y^2],Sqrt[4-y^2]},
{z,0,4-x^2-y^2}]+Integrate[1,{y,-2,-1},{x,-Sqrt[4-y^2],
Sqrt[4-y^2]},{z,0,4-x^2-y^2}]+Integrate[1,{y,-1,1},
{x,-Sqrt[4-y^2],1-Sqrt[1-y^2]},{z,0,4-x^2-y^2}]+
Integrate[1,{y,-1,1},{x,1+Sqrt[1-y^2],Sqrt[4-y^2]},
{z,0,4-x^2-y^2}]
\end{Verbatim}
\end{tcolorbox}
\vspace{0.7em}
\noindent

	
Integrando primero con respecto a $y,$ se puede plantear el volumen como la integral triple en donde $y$ va desde la parte negativa del paraboloide hasta la parte positiva del mismo y luego restar la parte correspondiente a la región que ocupa el cilindro, los límites para $x$ y $z$ se obtienen de la proyección del sólido en el plano $xz$.


\vskip0.3cm
\noindent
\begin{minipage}{\textwidth}
	\centering
	\captionof{figure}{$-\sqrt{4-z}\leq x \leq 2-z/2$}
	\begin{minipage}{0.39\linewidth}
		\centering
		\vspace{-17mm}
		\includegraphics[width=\linewidth]{Fig/341.pdf}
	\end{minipage}
	
	\vspace{-12mm}
	\fuentepropia
\end{minipage}
\vskip0.2cm


Con ayuda de este gráfico, el volumen se puede plantear mediante las siguientes integrales
\[\int_{0}^{4}\int_{-\sqrt{4-z}}^{\frac{4-z}{2}}\int_{-\sqrt{4-z-x^2}}^{\sqrt{4-z-x^2}}dydxdz
-\int_{0}^{4}\int_{0}^{\frac{4-z}{2}}\int_{-\sqrt{2x-x^2}}^{\sqrt{2x-x^2}}dydxdz \approx 17.2788.\]

\begin{multline*}
\int_{-2}^{0}\int_{0}^{4-x^2}\int_{-\sqrt{4-z-x^2}}^{\sqrt{4-z-x^2}}dydzdx
+\int_{0}^2\int_{0}^{4-2x}\int_{-\sqrt{4-z-x^2}}^{\sqrt{4-z-x^2}}dydzdx \\
-\int_{0}^2\int_{0}^{4-2x}\int_{-\sqrt{2x-x^2}}^{\sqrt{2x-x^2}}dydzdx = \frac{11}{2}\pi.
\end{multline*}
También se puede hacer uso de las coordenadas cilíndricas, cuyo planteamiento es el siguiente:
\[2 \int_{0}^{\pi /2}\int_{2\cos \theta}^2\int_{0}^{4-r^2}rdzdrd\theta +\int_{\pi /2}^{3\pi
/2}\int_{0}^2\int_{0}^{4-r^2}rdzdrd\theta = \frac{11}{2}\pi. \]

Las integrales anteriores se evaluaron con las siguientes instrucciones:

\vspace{0.6em}
\noindent
\begin{tcolorbox}[breakable,enhanced,title=\bf Instrucción en \verb"Mathematica", top=2pt,bottom=2pt,boxsep=0pt,before skip=2pt,after skip=0pt]
\begin{Verbatim}[formatcom=\letraverbatim]
Integrate[1,{z,0,4},{x,-Sqrt[4-z],(4-z)/2},
{y,-Sqrt[4-z-x^2],Sqrt[4-z-x^2]}]-Integrate[1,
{z,0,4},{x,0,(4-z)/2},{y,-Sqrt[2x-x^2],Sqrt[2x-x^2]}]
		
Integrate[1,{x,-2,0},{z,0,4-x^2},{y,-Sqrt[4-z-x^2],
Sqrt[4-z-x^2]}]+Integrate[1,{x,0,2},{z,0,4-2*x},
{y,-Sqrt[4-z-x^2],Sqrt[4-z-x^2]}]-Integrate[1,{x,0,2},
{z,0,4-2x},{y,-Sqrt[2x-x^2],Sqrt[2x-x^2]}]
		
2*Integrate[r,{\[Theta],0,Pi/2},{r,2*Cos[\[Theta]],2},
{z,0,4-r^2}]+Integrate[r,{\[Theta],Pi/2,3*Pi/2},
{r,0,2},{z,0,4-r^2}]
\end{Verbatim}
\end{tcolorbox}
\vspace{0.7em}
\noindent


\begin{ejer}
Plantear las integrales que se deben evaluar para determinar el volumen del sólido determinado por las expresiones $z\leq 4-x^2-y^2$, $x^2+y^2 \geq 2x$ y $z \geq 0$;
en las órdenes $dxdydz$ y $dxdzdy.$
\end{ejer}
\vskip0.2cm
El mismo valor obtenido antes se puede hallar de la siguiente manera. 
\vskip0.2cm

%\begin{tcolorbox}[title=\bfseries Parametrizando el sólido]
%\small
%\parskip=0pt
%\setlength{\belowdisplayskip}{0pt}
%\setlength{\abovedisplayskip}{0pt}
%\noindent
%\vspace{-20pt}

La proyección del sólido en el plano $xy$ permite ver que los puntos sobre la circunferencia interna puede representarse en la forma
$(1+\cos(t),\sin(t)),$ mientras que los puntos sobre la circunferencia externa son de la forma
$(2\cos(t),2\sin(t))$ para $t\in [0,2\pi]$.


\vskip0.2cm
\noindent
\begin{minipage}{\textwidth}
	\centering
	\captionof{figure}{Región $2x \leq x^2+y^2 \leq 4$}
	\label{fig:curva_parametrica_uv_3_2}
	\begin{minipage}{0.35\linewidth}
		\centering
		\vspace{-18mm}
		\includegraphics[width=\linewidth]{Fig/342.pdf}
	\end{minipage}
	
	\vspace{-10mm}
	\fuentepropia
\end{minipage}
\vskip0.2cm

Formando segmentos rectilíneos que unen las dos circunferencias, se parametriza la región entre las dos curvas, esto se logra simplificando
$\left( 1-u\right)\left(1+\cos(t) ,\sin(t)\right) +u\left(2\cos(t),2\sin(t)\right),$ para $u\in [0,1].$
Consideremos un punto
{\small $P\left( 1-u+(u+1)\cos (t),(1+u)\sin (t),0\right)$} en la base del sólido
y {\small $Q\left( 1-u+(u+1)\cos (t),(1+u)\sin (t),2(1-u^2)(1-\cos(t))\right)$} un \\ punto sobre el paraboloide.
Con lo anterior se pueden formar segmentos rectilíneos (paralelos al eje $z$) que unan la base del sólido con puntos sobre el paraboloide. Esto es $\left(1-v\right) P+vQ$; simplificando, obtendremos una parametrización para el sólido.
\begin{multline*}
{\bf r}(t,u,v) =\langle 1+(1+u) \cos t-u,( 1+u)\sin(t),
2v(1-u^2) ( 1-\cos t) \rangle,
\end{multline*}
\vskip0.2cm
donde $u\in [0,1],\,\, t\in [0,2\pi ],\,\, v\in [0,1].$ También puede verse que:
%\vspace{-4mm}
\begin{eqnarray*}
{\bf r}_{t}&=& \langle - \left( 1+u\right)\sin (t),\left( 1+u\right) \cos t,2v\left( 1-u^2\right) \sin t \rangle \\
{\bf r}_{u}&=& \langle -1+\cos t,\sin t,-4vu\left( 1-\cos t\right)\rangle\\
{\bf r}_{v}&=& \langle 0,0,2\left( 1-u^2\right) \left( 1-\cos t\right) \rangle.
\end{eqnarray*}
\vskip0.2cm
Con lo anterior,
${\bf r}_{u}\times {\bf r}_{v} = \big \langle 2 \left(1-u^2\right) \left( 1-\cos t\right)\sin t,2(1-\cos t) ^2(1-u^2) ,0\big \rangle .$
Además
\[\left|{\bf r}_{t}\cdot \left( {\bf r}_{u}\times {\bf r}_{v}\right)\right| = \left| 2\left( 1-\cos t\right)
^2\left( 1-u\right) \left( 1+u\right) ^2\right|,\]
\vskip0.2cm
con lo cual, el volumen del sólido esta dado por
\vskip0.2cm
\[\int_{0}^{1}\int_{0}^{1}\int_{0}^{2\pi }\left| 2\left( 1-\cos t\right)
^2\left( 1-u\right) \left( 1+u\right) ^2\right| dtdudv= \frac{11}{2}\pi.\]
%\end{tcolorbox}

\subsection*{Un paraboloide hueco
} \index{Paraboloide!hueco}
Consideraremos el sólido que se presenta en la sección \secref{p1}.

\vskip0.5cm
\noindent
\begin{minipage}{\textwidth}
	\centering
	\captionof{figure}{Región determinada por $2x \leq x^2 + y^2$ y $0 \leq z \leq 9 - x^2 - y^2$.}
	\label{fig:region_paraboloide_desplazada}
	\begin{minipage}{0.25\linewidth}
		\centering
		\includegraphics[width=\linewidth]{Img/79a.pdf}
	\end{minipage}
	\hspace{1em} 
	\begin{minipage}{0.25\linewidth}
		\centering
		\includegraphics[width=\linewidth]{Img/79b.pdf}
	\end{minipage}
	\hspace{1em} 
	\begin{minipage}{0.25\linewidth}
		\centering
		\includegraphics[width=\linewidth]{Img/79c.pdf}
	\end{minipage}
	
	\vspace{0.3em}
	\fuentepropia
\end{minipage}
\vskip0.2cm

La proyección del sólido en el plano $xy$ es una región acotada por las circunferencias
${\bf r}_1(t)= \langle 1+\cos(2t),\sin(2t)\rangle $
y ${\bf r}_2(t)= \langle 3\cos(2t),3\sin(2t)\rangle,$
que se parametriza por medio de una combinación de la forma $(1-u){\bf r}_1(t)+u{\bf r}_2(t),$ simplificando para $t\in [0,\pi]$, $u\in [0,1],$ se obtiene:
\[\langle 1-u+(1+2u)\cos (2t) ,(1+2u)\sin(2t)\rangle.\]


\vskip0.2cm
\noindent
\begin{minipage}{\textwidth}
	\centering
	\captionof{figure}{${\bf r}(t,u)$}
	\label{fig:curva_parametrica_uv_4}
	\begin{minipage}{0.34\linewidth}
		\centering
		\vspace{-18mm}
		\includegraphics[width=\linewidth]{Fig/343.pdf}
	\end{minipage}
	
	\vspace{-10mm}
	\fuentepropia
\end{minipage}
\vskip0.2cm


El sólido se parametriza mediante la función
{\small
\begin{eqnarray*}
{\bf r}\left( t,u,v\right) &=& \bigg\langle 1-u+\left( 1+2u\right) \cos (2t) ,\left( 1+2u\right)
\sin (2t) , \\ &&
\hskip2cm \left( 7-2u-5u^2+\left( 4u^2-2u-2\right) \cos(2t)\right) v\bigg\rangle,
\end{eqnarray*}
}

en donde $t\in [0,\pi]$, $u\in [0,1]$ y $v \in [0,1]$.
Con la función ${\bf r}$ hallamos los siguientes vectores:

\vspace{-4mm}
{\small
\begin{eqnarray*}
{\bf r}_{t} &=& \left\langle -2\left( 2u+1\right) \sin (2t)
,2\left( 2u+1\right) \cos (2t) ,-4v(2u^2-u-1)\sin (2t) \right\rangle \\
{\bf r}_{u} &=& \left\langle -1+2\cos (2t) ,2\sin \left(2t\right) ,v\left( -2-10u+\left( -2+8u\right) \cos (2t) \right) \right\rangle \\
{\bf r}_{v} &=& \left\langle 0,0,7-2u-5u^2+4\cos (2t)u^2-2\cos (2t) u-2\cos (2t) \right\rangle,
\end{eqnarray*}
}

\vskip0.2cm
luego,
$\left| {\bf r}_{v}\cdot ({\bf r}_{t}\times {\bf r}_{u}) \right| =\bigg| 2(\cos (2t)-2)(2u+1)(u-1) ((4u+2)\cos (2t) \linebreak -5u-7) \bigg| ,$
que corresponde al elemento diferencial de volumen, su integral en los
límites establecidos nos proporciona el volumen del sólido
\[\int_{0}^{1}\int_{0}^{1}\int_{0}^{\pi }\left| {\bf r}_{v}\cdot ({\bf r}_{t}\times {\bf r}_{u}) \right| dtdudv=33\pi.\]
\vskip0.2cm
Este cálculo se hizo en \verb"Mathematica" usando las siguientes instrucciones:

\vspace{0.6em}
\noindent
\begin{tcolorbox}[breakable,enhanced,title=\bf Instrucción en \verb"Mathematica", top=2pt,bottom=2pt,boxsep=0pt,before skip=2pt,after skip=0pt]
\begin{Verbatim}[formatcom=\letraverbatim]
{x,y}=(1-u)*{3*Cos[2*t],3*Sin[2*t]}+u*{1+Cos[2*t],
Sin[2*t]};
z=Simplify[9-x^2-y^2];
r={x,y,z*v}
A=Simplify[Dot[D[r,v],Cross[D[r,t],D[r,u]]]];
Integrate[A,{t,0,Pi},{u,0,1},{v,0,1}]
\end{Verbatim}
\end{tcolorbox}
\vspace{0.7em}
\noindent


En coordenadas rectangulares, el volumen está dado por
\begin{multline*} \int_{-3}^{3}\int_{-\sqrt{9-x^2}}^{\sqrt{9-x^2}}\left(
9-x^2-y^2\right) dydx- \\ 2\int_{0}^2\int_{0}^{\sqrt{2x-x^2}}\left(9-x^2-y^2\right)dydx= 33\pi,
\end{multline*}
y usando coordenadas polares, el planteamiento es el siguiente:
\[\int_{0}^{2\pi }\int_{0}^{3}\left( 9-r^2\right) rdrd\theta -\int_{0}^{\pi}\int_{0}^{2\cos \theta }\left( 9-r^2\right) rdrd\theta = 33\pi.\]
Estas integrales se evalúan con la siguiente rutina:

\vspace{0.6em}
\noindent
\begin{tcolorbox}[breakable,enhanced,title=\bf Instrucción en \verb"Mathematica", top=2pt,bottom=2pt,boxsep=0pt,before skip=2pt,after skip=0pt]
\begin{Verbatim}[formatcom=\letraverbatim]
Integrate[9-x^2-y^2,{x,-3,3},{y,-Sqrt[9-x^2],
Sqrt[9-x^2]}]-2*Integrate[9-x^2-y^2,{x,0,2},
{y,0,Sqrt[2*x-x^2]}]
Integrate[(9-r^2)*r,{t,0,2*Pi},{r,0,3}]-
Integrate[(9-r^2)*r,{t,0,Pi},{r,0,2*Cos[t]}]
\end{Verbatim}
\end{tcolorbox}
\vspace{0.7em}
\noindent


\subsection*{Un paraboloide y un cilindro parabólico}
Consideremos el sólido acotado por las figuras de las ecuaciones $z = x^2+y^2$ y
$z = 5 - y^2,$ que representan respectivamente un paraboloide y un cilindro parabólico; estas figuras se muestran enseguida.

\begin{figure}[H]
\centering
\caption{Sólido determinado por $x^2 + y^2 \leq z \leq 5 - y^2$.}
\label{fig:solido_paraboloide_plano_curvo}
\includegraphics[height=4.5cm]{Img/26a.pdf}
\includegraphics[height=4.5cm]{Img/26b.pdf}
\includegraphics[height=4.5cm]{Img/26c.pdf}
\fuentepropia
\end{figure}

Las instrucciones necesarias para obtenerlas son las siguientes.

\vspace{0.6em}
\noindent
\begin{tcolorbox}[breakable,enhanced,title=\bf Instrucción en \verb"Mathematica", top=2pt,bottom=2pt,boxsep=0pt,before skip=2pt,after skip=0pt]
\begin{Verbatim}[formatcom=\letraverbatim]
ContourPlot3D[{z==5-y^2,z==x^2+y^2},{x,-2.5,2.5},
{y,-2.5,2.5},{z,0,5},PlotPoints->90,
ContourStyle->{RGBColor[0.95,0.6,0.9],Cyan}]
RegionPlot3D[x^2+y^2<z<5-y^2&&-Sqrt[5-2*y^2]<x<
Sqrt[5-2*y^2],{x,-2.5,2.5},{y,-2.5,2.5},{z,0,5},
PlotPoints->90,Mesh->False]
\end{Verbatim}
\end{tcolorbox}
\vspace{0.7em}
\noindent


El sólido se proyecta en los planos coordenados como se muestra en las siguientes figuras.
\vskip0.2cm

\noindent
\begin{minipage}{\textwidth}
	\centering
	\captionof{figure}{Proyecciones del sólido en los planos coordenados.}
	\begin{minipage}{0.30\linewidth}
		\centering
		\vspace{-8mm}
		\includegraphics[width=\linewidth]{Fig/344.pdf}
	\end{minipage}
	\hspace{1em}
	\begin{minipage}{0.30\linewidth}
		\centering
		\vspace{-8mm}
		\includegraphics[width=\linewidth]{Fig/345.pdf}
	\end{minipage}
	\hspace{1em}
	\begin{minipage}{0.30\linewidth}
		\centering
		\vspace{-8mm}
		\includegraphics[width=\linewidth]{Fig/346.pdf}
	\end{minipage}
	
	\vspace{0.3em}
	
	\vspace{-8mm}
	\fuentepropia
\end{minipage}

\vskip0.2cm
Eliminando $z$ de las ecuaciones $z=x^2+y^2$ y $z=5-y^2$, se obtiene $x^2+2y^2=5;$
esta expresión en el plano $xy$ representa una elipse.
Para integrar primero con respecto a $z$ se debe tener en cuenta que $z$ recorre los puntos que se hallan entre el paraboloide y el cilindro, los límites en $y$ y $x$ se pueden obtener de la proyección en el plano $xy.$ Por tanto, el volumen del sólido está dado por las siguientes integrales:
\vspace{-2mm}
\begin{multline*} \int_{-\sqrt{5/2}}^{\sqrt{5/2}}\int_{-\sqrt{5-2y^2}}^{\sqrt{5-2y^2}}\int_{x^2+y^2}^{5-y^2}dzdxdy = \\ \int_{-\sqrt{5}}^{\sqrt{5}}\int_{-\frac{1}{2}\sqrt{10-2x^2}}^{\frac{1}{2}\sqrt{10-2x^2}}\int_{x^2+y^2}^{5-y^2}dzdydx= \frac{25\sqrt{2}\pi}{4}.
\end{multline*}
En el plano $xz,$ las curvas que describen la región de integración se determinan por las ecuaciones $2z=5+x^2,$ $z=x^2$ y $z=5.$
De la expresión $2z=5+x^2$ se sigue que $ x=\pm \sqrt{2z-5}.$ Por lo tanto, el volumen es
\[\int_{-\sqrt{5}}^{\sqrt{5}}\int_{x^2}^{\frac{5+x^2}{2}}\int_{-\sqrt{z-x^2}}^{\sqrt{z-x^2}}dydzdx+\int_{-\sqrt{5}}^{\sqrt{5}}\int_{\frac{5+x^2}{2}}^{5}\int_{-\sqrt{5-z}}^{\sqrt{5-z}}dydzdx= \frac{25}{4}\sqrt{2}\pi.\]
También se puede calcular el volumen de la siguiente manera
\begin{multline*}
\int_{5/2}^{5}\int_{-\sqrt{2z-5}}^{\sqrt{2z-5}}\int_{-\sqrt{5-z}}^{\sqrt{5-z}}dydxdz+\int_{5/2}^{5}\int_{-\sqrt{z}}^{-\sqrt{2z-5}}\int_{-\sqrt{z-x^2}}^{\sqrt{z-x^2}}dydxdz\\
+\int_{5/2}^{5}\int_{\sqrt{2z-5}}^{\sqrt{z}}\int_{-\sqrt{z-x^2}}^{\sqrt{z-x^2}}dydxdz+\int_{0}^{5/2}\int_{-\sqrt{z}}^{\sqrt{z}}\int_{-\sqrt{z-x^2}}^{\sqrt{z-x^2}}dydxdz\approx 27.768.
\end{multline*}

Al integrar primero con respecto a $x$ hay que tener presente que esta recorre los puntos que se encuentran entre la parte negativa del paraboloide (despejando $x$)
hasta la parte positiva de la misma superficie.
\vspace{-2mm}
\[
\displaystyle
\int_{0}^{5/2}
\int_{-\sqrt{z}}^{\sqrt{z}}
\int_{-\sqrt{z - y^2}}^{\sqrt{z - y^2}}
dx\,dy\,dz
+
\int_{5/2}^{5}
\int_{-\sqrt{5 - z}}^{\sqrt{5 - z}}
\int_{-\sqrt{z - y^2}}^{\sqrt{z - y^2}}
dx\,dy\,dz
\approx 27.768
\]

\[
\displaystyle
\int_{-\sqrt{5/2}}^{\sqrt{5/2}}
\int_{y^2}^{5 - y^2}
\int_{-\sqrt{z - y^2}}^{\sqrt{z - y^2}}
dx\,dz\,dy
= \frac{25}{4}\sqrt{2}\pi.
\]

\vskip0.2cm
Utilizando coordenadas cilíndricas las ecuaciones del paraboloide y el cilindro se reescriben como sigue
$z=r^2$ y $z=5-r^2\sin^2(\theta)$. La ecuación de la elipse de intersección de estas superficies está dada por
$ r^2 (3- \cos(2 \theta))=10.$
Es por esto que la integral que determina el volumen del sólido es
\[\int_0^{2\pi}\int_0^{\frac{\sqrt{10}}{\sqrt{3-\cos(2\theta)}}}\int_{r^2}^{5-r^2\sin ^2(\theta)}rdzdrd\theta=\frac{25\pi}{2\sqrt{2}}.\]
Algunas de estas integrales se evaluaron usando \verb"Mathematica" con las siguientes instrucciones:

\vspace{0.6em}
\noindent
\begin{tcolorbox}[breakable,enhanced,title=\bf Instrucción en \verb"Mathematica", top=2pt,bottom=2pt,boxsep=0pt,before skip=2pt,after skip=0pt]
\begin{Verbatim}[formatcom=\letraverbatim]
Integrate[1,{x,-Sqrt[5],Sqrt[5]},{y,-Sqrt[10-2*x^2]/2,
Sqrt[10-2*x^2]/2},{z,x^2+y^2,5-y^2}]
Integrate[1,{y,-Sqrt[5/2],Sqrt[5/2]},{x,-Sqrt[5-2y^2],
Sqrt[5-2y^2]},{z,x^2+y^2,5-y^2}]
Integrate[1,{y,-Sqrt[5/2],Sqrt[5/2]},{z,y^2,5-y^2},
{x,-Sqrt[z-y^2],Sqrt[z-y^2]}]
\end{Verbatim}
\end{tcolorbox}
\vspace{0.7em}
\noindent



\begin{comment}

\begin{verbatim}
Integrate[1,{x,-Sqrt[5],Sqrt[5]},{y,-Sqrt[10-2*x^2]/2,
Sqrt[10-2*x^2]/2},{z,x^2+y^2,5-y^2}]

Integrate[1,{y,-Sqrt[5/2],Sqrt[5/2]},{x,-Sqrt[5-2y^2],
Sqrt[5-2y^2]},{z,x^2+y^2,5-y^2}]

Integrate[1,{x,-Sqrt[5],Sqrt[5]},{z,x^2,(5+x^2)/2},
{y,-Sqrt[z-x^2],Sqrt[z-x^2]}]+Integrate[1,{x,-Sqrt[5],
Sqrt[5]},{z,(5+x^2)/2,5},{y,-Sqrt[5-z],Sqrt[5-z]}]

Integrate[1,{y,-Sqrt[5/2],Sqrt[5/2]},{z,y^2,5-y^2},
{x,-Sqrt[z-y^2],Sqrt[z-y^2]}]

Integrate[1,{z,5/2,5},{x,-Sqrt[2*z-5],Sqrt[2*z-5]},
{y,-Sqrt[5-z],Sqrt[5-z]}]+Integrate[1,{z,5/2,5},
{x,-Sqrt[z],-Sqrt[2z-5]},{y,-Sqrt[z-x^2],Sqrt[z-x^2]}]
+ Integrate[1,{z,5/2,5},{x,Sqrt[2z-5],Sqrt[z]},
{y,-Sqrt[z-x^2],Sqrt[z-x^2]}]+Integrate[1,{z,0,5/2},
{x,-Sqrt[z],Sqrt[z]},{y,-Sqrt[z-x^2],Sqrt[z-x^2]}]

Integrate[1,{z,0,5/2},{y,-Sqrt[z],Sqrt[z]},
{x,-Sqrt[z-y^2],Sqrt[z-y^2]}]+Integrate[1,{z,5/2,5},
{y,-Sqrt[5-z],Sqrt[5-z]},{x,-Sqrt[z-y^2],Sqrt[z-y^2]}]
\end{verbatim}

\verb"Integrate[r,{\[Theta],0,2*Pi},{r,0,Sqrt[10]/Sqrt[3-"
\verb"Cos[2\[Theta]]]},{z,r^2,5-r^2*Sin[\[Theta]]^2}]"
\end{comment}

\vskip0.2cm
Otra manera de hallar el volumen se presenta enseguida.
\vskip0.2cm

%\begin{tcolorbox}[title=\bfseries Parametrizando el sólido]
En el plano $xy$, la proyección del sólido se determina por
$x^2+2y^2=5,$ esta ecuación corresponde a elipse. Esta curva junto con su interior se parametrizan como {\small $x=\sqrt{5}u\cos t,$}
{\small $y=\sqrt{\frac{5}{2}}u\sin t),$} para {\small $t\in [0,2\pi ],$} {\small $u\in [0,1].$} 

\vskip0.2cm
La coordenada $z$ está dada por {\small $z = \frac{5}{2}u^2\cos ^2t+\frac{5}{2}u^2$} para el paraboloide , {\small $z=5-\frac{5}{2}u^2\sin ^2t$} para el cilindro.
Con un punto $P$ sobre el paraboloide y un punto $Q$ sobre el cilindro, se forma una combinación convexa de la forma $(1-w)P+wQ$ que parametriza el sólido, esto es,

\vspace{-5mm}
\[{\bf r}(t,u,w)= \langle \sqrt{5}u\cos t,\sqrt{\frac{5}{2}}u\sin t,\frac{5}{2}u^2( \cos ^2t+1) -5w(u^2-1)\rangle.\]

En donde $t \in [0,1],$ $u \in [0,1]$ y $w \in [0,1].$ Con esta función se hallan los siguientes vectores:
\vspace{-4mm}
\begin{eqnarray*}
{\bf r}_{t} & =& \langle -\sqrt{5}u\sin t,\frac{1}{2}\sqrt{10}u\cos t,-5u^2\cos t\sin t\rangle \\
{\bf r}_{u} &=& \langle \sqrt{5}\cos t,\frac{1}{2}\sqrt{10}\sin t,5u(\cos^2t+1)-10wu\rangle \\
{\bf r}_{w} &=& \langle0,0,-5u^2+5\rangle
\end{eqnarray*}


Con lo anterior, se calcula el producto vectorial triple que corresponde al elemento diferencial de volumen.
$| {\bf r}_{t}\cdot ( {\bf r}_{u}\times {\bf r}_{w}) | =\frac{25}{2}\sqrt{2} \left| u( u^2-1)\right|,$
que al integrarlo, nos produce el volumen del sólido mencionado
\[\frac{25\sqrt{2}}{2} \int_{0}^{1}\int_{0}^{1}\int_{0}^{2\pi } | u(u^2-1) | dtdudw= \frac{25}{4}\sqrt{2}\pi.\]
%\end{tcolorbox}

De manera más breve, \verb"Mathematica" permite hacer estos cálculos con la ejecución de las siguientes instrucciones:

\vspace{0.6em}
\noindent
\begin{tcolorbox}[breakable,enhanced,title=\bf Instrucción en \verb"Mathematica", top=2pt,bottom=2pt,boxsep=0pt,before skip=2pt,after skip=0pt]
\begin{Verbatim}[formatcom=\letraverbatim]
{x,y}={Sqrt[5]*u*Cos[t],Sqrt[5/2]*u*Sin[t]};
r=(1-w)*{x,y,5-y^2}+w*{x,y,x^2+y^2};
A=Dot[D[r,t],Cross[D[r,u],D[r,w]]];
Integrate[A,{t,0,2*Pi},{u,0,1},{w,0,1}]
\end{Verbatim}
\end{tcolorbox}
\vspace{0.7em}
\noindent


\subsection*{Dos paraboloides} \index{Paraboloide}
El siguiente ejercicio es tomado del libro \textcite{thomas_calculo_2015}.
Los paraboloides con ecuaciones $z=8-x^2 -y^2 $ y $z=x^2 +3y^2$ acotan un sólido en el espacio.


\noindent
\begin{minipage}{\textwidth}
	\centering
	\captionof{figure}{Sólido determinado por $x^2 + 3y^2 \leq z \leq 8 - x^2 - y^2$.}
	\label{fig:solido_paraboloides_elipticos}
	\begin{minipage}{0.23\linewidth}
		\centering
		\includegraphics[width=\linewidth]{Img/96a.pdf}
	\end{minipage}
	\hspace{1em}
	\begin{minipage}{0.23\linewidth}
		\centering
		\includegraphics[width=\linewidth]{Img/96b.pdf}
	\end{minipage}
	\hspace{1em}
	\begin{minipage}{0.22\linewidth}
		\centering
		\includegraphics[width=\linewidth]{Img/96c.pdf}
	\end{minipage}
	
	\vspace{0.3em}
	\fuentepropia
\end{minipage}
\vskip0.2cm

Estos gráficos se consiguen con la ayuda de las siguientes instrucciones:

\vspace{0.6em}
\noindent
\begin{tcolorbox}[breakable,enhanced,title=\bf Instrucción en \verb"Mathematica", top=2pt,bottom=2pt,boxsep=0pt,before skip=2pt,after skip=0pt]
\begin{Verbatim}[formatcom=\letraverbatim]
ContourPlot3D[{z==8-x^2-y^2,z==x^2+3*y^2},{x,-3,3},
{y,-3,3},{z,0,8},ContourStyle->{Directive[Cyan,
Opacity[0.8]],Directive[Orange,Opacity[0.9]]},
Mesh->False]
RegionPlot3D[x^2+3*y^2<z<8-x^2-y^2&&x^2+2y^2<4,
{x,-2,2},{y,-Sqrt[2],Sqrt[2]},{z,0,8},Mesh->False,
PlotPoints->90]
\end{Verbatim}
\end{tcolorbox}
\vspace{0.7em}
\noindent


La proyección en el plano $xz$ se obtiene eliminando $y$ de las ecuaciones de los paraboloides, esto es, al sumar las expresiones
$3z=24-3x^2 -3y^2$ y $z=x^2 +3y^2$ se obtiene $2z=12-x^2 $. Sumando las ecuaciones de los paraboloides, además\ de eliminar $x,$ se obtiene la ecuación que termina la proyección del sólido en el plano $yz,$ esto es $2z=8+2y^2 $ que equivale a $z=4+y^2 .$
\vskip0.2cm
Las siguientes figuras muestran la proyección del sólido en los planos $yz$ y $xz$; ellos ilustran el planteamiento de las integrales triples que determinan el volumen del sólido.

\vskip0.3cm
\noindent
\begin{minipage}{\textwidth}
	\centering
	\captionof{figure}{Cono truncado oblicuo determinado por ${\bf r}(u,t)$}
	\begin{minipage}{0.35\linewidth}
		\centering
		\vspace{-10mm}
		\includegraphics[width=\linewidth]{Fig/347.pdf}
	\end{minipage}
	\hspace{1em}
	\begin{minipage}{0.35\linewidth}
		\centering
		\vspace{-10mm}
		\includegraphics[width=\linewidth]{Fig/348.pdf}
	\end{minipage}

	\vspace{-9mm}
	\fuentepropia
\end{minipage}

\vskip0.3cm

Integrando primero con respecto a $y$, las integrales triples correspon\-dientes son:
\[\bi_{-2}^2 \bi_{6-\frac{1}{2}x^2 }^{8-x^2}\bi_{-\sqrt{8-z-x^2}}^{\sqrt{8-z-x^2}}dydzdx+\bi_{-2}^2 \bi_{x^2 }^{6-\frac{1}{2}x^2 }\bi_{-\sqrt{\frac{z-x^2 }{3}}}^{\sqrt{\frac{z-x^2 }{3}} }dydzdx= 8\sqrt{2}\pi\]

De la ecuación $2z=12-x^2 $, se sigue que $x=\pm \sqrt{12-2z}$
entonces
{\footnotesize
\begin{multline}
\bi_{4}^{6}\bi_{-\sqrt{8-z}}^{-\sqrt{12-2z}}\bi_{-\sqrt{8-z-x^2}}^{\sqrt{8-z-x^2}}dydxdz+\bi_{4}^{6}\int_{\sqrt{12-2z}}^{\sqrt{8-z}}\bi_{-\sqrt{8-z-x^2}}^{\sqrt{8-z-x^2}}dydxdz \\
+\bi_{6}^{8}\bi_{-\sqrt{8-z}}^{\sqrt{8-z}}\bi_{-\sqrt{8-z-x^2}}^{\sqrt{8-z-x^2}}dydxdz+\bi_{4}^{6}\int_{-\sqrt{12-2z}}^{\sqrt{12-2z}}\bi_{-\sqrt{\frac{z-x^2 }{3}} }^{\sqrt{\frac{z-x^2 }{3}}}dydxdz \\
+\bi_{0}^{4}\bi_{-\sqrt{z}}^{\sqrt{z}}\bi_{-\sqrt{\frac{z-x^2 }{3}}}^{\sqrt{\frac{z-x^2 }{3}} }dydxdz = 8\sqrt{2}\pi
\end{multline}
}
Estas integrales se evalúan en \verb"Mathematica" con las siguientes instrucciones:

\vspace{0.6em}
\noindent
\begin{tcolorbox}[breakable,enhanced,title=\bf Instrucción en \verb"Mathematica", top=2pt,bottom=2pt,boxsep=0pt,before skip=2pt,after skip=0pt]
\begin{Verbatim}[formatcom=\letraverbatim]
Integrate[1,{x,-2,2},{z,6-x^2/2,8-x^2},
{y,-Sqrt[8-z-x^2],Sqrt[8-z-x^2]}]+Integrate[1,
{x,-2,2},{z,x^2,6-x^2/2},{y,-Sqrt[(z-x^2)/3],
Sqrt[(z-x^2)/3]}]
Integrate[1,{z,4,6},{x,-Sqrt[8-z],-Sqrt[12-2z]},
{y,-Sqrt[8-z-x^2],Sqrt[8-z-x^2]}]+Integrate[1,{z,4,6},
{x,Sqrt[12-2z],Sqrt[8-z]},{y,-Sqrt[8-z-x^2],
Sqrt[8-z-x^2]}]+Integrate[1,{z,6,8},{x,-Sqrt[8-z],
Sqrt[8-z]},{y,-Sqrt[8-z-x^2],Sqrt[8-z-x^2]}]+
Integrate[1,{z,4,6},{x,-Sqrt[12-2z],Sqrt[12-2z]},
{y,-Sqrt[(z-x^2)/3],Sqrt[(z-x^2)/3]}]+Integrate[1,
{z,0,4},{x,-Sqrt[z],Sqrt[z]},{y,-Sqrt[(z-x^2)/3],
Sqrt[(z-x^2)/3]}]
\end{Verbatim}
\end{tcolorbox}
\vspace{0.7em}
\noindent


Integrando primero con respecto a $x$, las integrales triples que determinan el volumen del sólido son las siguientes:
{\footnotesize
\begin{multline} \bi_{-\sqrt{2}}^{\sqrt{2}}\bi_{4+y^2 }^{8-y^2}\bi_{-\sqrt{8-z-y^2}}^{\sqrt{8-z-y^2}}dxdzdy
+\bi_{-\sqrt{2}}^{\sqrt{2}}\bi_{3y^2 }^{4+y^2}\bi_{-\sqrt{z-3y^2 } }^{\sqrt{z-3y^2}}dxdzdy\\
= \bi_{6}^{8}\bi_{-\sqrt{8-z}}^{\sqrt{8-z}}\bi_{-\sqrt{8-z-y^2 }}^{\sqrt{8-z-y^2}}dxdydz+\bi_{4}^{6}\bi_{-\sqrt{z-4}}^{\sqrt{z-4}}\bi_{-\sqrt{8-z-y^2}}^{\sqrt{8-z-y^2 }}dxdydz \\
+\bi_{4}^{6}\int_{-\sqrt{z/3}}^{-\sqrt{z-4}}\bi_{-\sqrt{z-3y^2}}^{\sqrt{z-3y^2}}dxdydz+\bi_{4}^{6}\bi_{\sqrt{z-4}}^{\sqrt{z/3}}\bi_{-\sqrt{z-3y^2}}^{\sqrt{z-3y^2 }}dxdydz \\
+\bi_{0}^{4}\bi_{-\sqrt{z/3}}^{\sqrt{z/3}}\bi_{-\sqrt{z-3y^2 }}^{\sqrt{z-3y^2 } }dxdydz = 8\sqrt{2}\pi
\end{multline}
}
Con las siguientes instrucciones se pueden evaluar estas integrales.
\vskip0.2cm
\vspace{0.6em}
\noindent
\begin{tcolorbox}[breakable,enhanced,title=\bf Instrucción en \verb"Mathematica", top=2pt,bottom=2pt,boxsep=0pt,before skip=2pt,after skip=0pt]
\begin{Verbatim}[formatcom=\letraverbatim]
Integrate[1,{y,-Sqrt[2],Sqrt[2]},{z,4+y^2,8-y^2},
{x,-Sqrt[8-z-y^2],Sqrt[8-z-y^2]}]+Integrate[1,{y,
Sqrt[2],Sqrt[2]},{z,3*y^2,4+y^2},{x,-Sqrt[z-3y^2],
Sqrt[z-3y^2]}]
Integrate[1,{z,6,8},{y,-Sqrt[8-z],Sqrt[8-z]},{x,
Sqrt[8-z-y^2],Sqrt[8-z-y^2]}]+Integrate[1,{z,4,6},
{y,-Sqrt[z-4],Sqrt[z-4]},{x,-Sqrt[8-z-y^2],
Sqrt[8-z-y^2]}]+Integrate[1,{z,4,6},{y,-Sqrt[z/3],
Sqrt[z-4]},{x,-Sqrt[z-3y^2],Sqrt[z-3y^2]}]+
Integrate[1,{z,4,6},{y,Sqrt[z-4],Sqrt[z/3]},{x,
Sqrt[z-3y^2],Sqrt[z-3y^2]}]+Integrate[1,{z,0,4},{y,
Sqrt[z/3],Sqrt[z/3]},{x,-Sqrt[z-3y^2],Sqrt[z-3y^2]}]
\end{Verbatim}
\end{tcolorbox}
\vspace{0.7em}
\noindent



La proyección en el plano $xy$ está dada por
$ 8-x^2 -y^2 =x^2 +3y^2,$ que se simplifica en la forma $ 4=x^2 +2y^2;$
esta ecuación corresponde a una elipse cuya figura se muestra enseguida.

\vskip0.2cm
\noindent
\begin{minipage}{\textwidth}
	\centering
	\captionof{figure}{$x^2 + 2y^2 \leq 4$}
	\label{fig:curva_parametrica_uv_5}
	\begin{minipage}{0.36\linewidth}
		\centering
		\vspace{-18mm}
		\includegraphics[width=\linewidth]{Fig/349.pdf}
	\end{minipage}
	
	\vspace{-10mm}
	\fuentepropia
\end{minipage}
\vskip0.3cm


Con ayuda de esta figura se puede plantear las integrales que permiten hallar el volumen del sólido.
En los órdenes $dzdydx$ y $dzdxdy$, las integrales son las siguientes.

\vskip0.4cm
En los órdenes $dzdydx$ y $dzdxdy$, las integrales son las siguientes.

\begin{multline*}
\int_{-2}^2 \int_{-\frac{1}{2}\sqrt{8-2x^2 }}^{\frac{1}{2}\sqrt{8-2x^2 }}\int_{x^2 +3y^2 }^{8-x^2 -y^2 }dzdydx= \\
\int_{-\sqrt{2}}^{\sqrt{2}}\int_{-\sqrt{4-2y^2 }}^{\sqrt{4-2y^2 }}\int_{x^2 +3y^2 }^{8-x^2 -y^2 }dzdxdy= 8\sqrt{2}\pi . \end{multline*}
\vskip0.2cm
Estas integrales se evalúan con la ejecución de las siguientes instrucciones:

\vspace{0.6em}
\noindent
\begin{tcolorbox}[breakable,enhanced,title=\bf Instrucción en \verb"Mathematica", top=2pt,bottom=2pt,boxsep=0pt,before skip=2pt,after skip=0pt]
\begin{Verbatim}[formatcom=\letraverbatim]
Integrate[1,{x,-2,2},{y,-Sqrt[8-2x^2]/2,
Sqrt[8-2x^2]/2},{z,x^2+3*y^2,8-x^2-y^2}]
Integrate[1,{y,-Sqrt[2],Sqrt[2]},{x,-Sqrt[4-2*y^2],
Sqrt[4-2*y^2]},{z,x^2+3*y^2,8-x^2-y^2}]
\end{Verbatim}
\end{tcolorbox}
\vspace{0.7em}
\noindent

Otra alternativa para obtener el volumen del sólido es la siguiente.
\vskip0.3cm

% \begin{figure}[H]
% \centering
% \begin{pspicture}(-2.2,-1.2)(2.5,1.4) %\psgrid
% \psset{algebraic,unit=0.7}
% \pscustom[fillstyle=vlines,hatchcolor=gray,hatchsep=1.5pt,linestyle=none]{
%   \psparametricplot{0}{\psPiTwo}{2*cos(t)|sqrt(2)*sin(t)}
% }
% \psaxes[Dx=2,deltax=2,linecolor=black]{->}(0,0)(-2.5,-1.5)(3,1.8)[$x$,90][$y$,0]
% \psparametricplot[linecolor=black!60!red]{0}{\psPiTwo}{2*cos(t)|sqrt(2)*sin(t)}
% % \rput(2.5,1.3){$x^2+2y^2=4$}
% \end{pspicture}
% \caption{$x^2+2y^2 \leq 4$.}
% \end{figure}
% \par
% \justifying

%\begin{tcolorbox}[title=\bfseries Parametrizando el sólido]

La intersección de los paraboloides está dada por \( x^2+2y^2=4 \), que representa una elipse en el plano \( xy \). Al ser $x = 2u \sin(t)$, $y = \sqrt{2}u \cos(t)$, con $u \in [0,1]$, $t \in [0,2\pi],$ se describe esta región.


\vskip0.2cm
\noindent
\begin{minipage}{\textwidth}
	\centering
	\captionof{figure}{$x^2 + 2y^2 \leq 4$}
	\label{fig:curva_parametrica_uv_6}
	\begin{minipage}{0.38\linewidth}
		\centering
		\vspace{-18mm}
		\includegraphics[width=\linewidth]{Fig/349.pdf}
	\end{minipage}
	
	\vspace{-11mm}
	\fuentepropia
\end{minipage}
\vskip0.3cm


Con dos puntos:
\begin{align*}
P &= \left( 2u\sin(t),\ \sqrt{2}u\cos(t),\ 8 - u^2(3 - \cos(2t)) \right), \\
Q &= \left( 2u\sin(t),\ \sqrt{2}u\cos(t),\ u^2(5 + \cos(2t)) \right).
\end{align*}

Cada uno sobre un paraboloide. Una combinación convexa de la forma \( (1-w)P + wQ \) define la parametrización del sólido:
\[
\mathbf{r}(u,t,w) = \left\langle 2u\sin t,\ \sqrt{2}u\cos t,\ 8(1-w) + u^2(8w - 3 + \cos(2t)) \right\rangle.
\]

De esta función obtenemos los siguientes vectores:
\begin{align*}
\mathbf{r}_t &= \left\langle 2u\cos(t),\ -\sqrt{2}u\sin(t),\ -2u^2\sin(2t) \right\rangle \\
\mathbf{r}_u &= \left\langle 2\sin(t),\ \sqrt{2}\cos(t),\ 2u(8w - 3 + \cos(2t)) \right\rangle \\
\mathbf{r}_w &= \left\langle 0,\ 0,\ 8(u^2 - 1) \right\rangle.
\end{align*}

Por tanto,
\[
\left| \mathbf{r}_w \cdot (\mathbf{r}_u \times \mathbf{r}_t) \right| = 16\sqrt{2}u(1 - u^2).
\]

Entonces el volumen del sólido está dado por:
\vskip0.2cm
\[
\int_0^1 \int_0^1 \int_0^{2\pi} 16\sqrt{2}u(1 - u^2)\, dt\, du\, dw = 8\sqrt{2}\pi.
\]
%\end{tcolorbox}
\vskip0.2cm
El cálculo anterior resulta sencillo de evaluar en \verb"Mathematica" con ayuda de las siguientes instrucciones:


\vspace{0.6em}
\noindent
\begin{tcolorbox}[breakable,enhanced,title=\bf Instrucción en \verb"Mathematica", top=2pt,bottom=2pt,boxsep=0pt,before skip=2pt,after skip=0pt]
\begin{Verbatim}[formatcom=\letraverbatim]
{x,y}={2*u*Sin[t],Sqrt[2]*u*Cos[t]};
r=Simplify[(1-w)*{x,y,8-x^2-y^2}+w*{x,y,x^2+3*y^2}];
Integrate[Dot[D[r,w],Cross[D[r,u],D[r,t]]],{w,0,1},
{u,0,1},{t,0,2*Pi}]
\end{Verbatim}
\end{tcolorbox}
\vspace{0.7em}
\noindent


\subsection*{Dos paraboloides II}
Las ecuaciones $63x^2+252x-8+28z^2+56z+36y=0$y $9x^2+36x+16+4z^2+8z-12y=0$
representan dos paraboloides. Estas superficies acotan la región en el espacio que se muestra a continuación.

\begin{figure}[H]
\centering
\caption{Región acotada por los paraboloides.}
\label{fig:region_paraboloides_89}
\includegraphics[height=4.5cm]{Img/89b.pdf}
\includegraphics[height=4.5cm]{Img/89c.pdf}
\includegraphics[height=4.5cm]{Img/89d.pdf}
\fuentepropia
\end{figure}

Con las siguientes instrucciones se obtuvieron estos gráficos.

\vspace{0.6em}
\noindent
\begin{tcolorbox}[breakable,enhanced,title=\bf Instrucción en \verb"Mathematica", top=2pt,bottom=2pt,boxsep=0pt,before skip=2pt,after skip=0pt]
\begin{Verbatim}[formatcom=\letraverbatim]
ContourPlot3D[{(x+2)^2/4+(z+1)^2/9==(8/7-y/7),
36*(y/3+2/3)==9*(x+2)^2+4*(z+1)^2},{x,-6,2},
{y,-2,8},{z,-7,5},Mesh->False,PlotPoints->60,
ContourStyle->{Directive[Orange,Opacity[0.95]],
Directive[Cyan,Opacity[0.8]]}]
RegionPlot3D[3*(x+2)^2/4+(z+1)^2/3-2<=y<8-7*(x+2)^2/4-
7*(z+1)^2/9&&-Sqrt[9-9*(x+2)^2/4]-1<z<Sqrt[9-
9*(x+2)^2/4]-1&&-4<x<0,{x,-4,0},{y,-2,8},{z,-4,2},
PlotPoints->60,MeshStyle->RGBColor[0.4,0.5,0.4],
PlotStyle->Orange,BoxRatios->{2,2,1.5}]
\end{Verbatim}
\end{tcolorbox}
\vspace{0.7em}
\noindent


Las ecuaciones inicialmente consideradas, al factorizarlas, se reescriben como
$63(x+2)^2+28(z+1)^2=36(8-y)$ y $9\left( x+2\right) ^2+4\left( z+1\right) ^2=12(y+2).$
Eliminado $y$ se obtiene $9x^2+36x+4z^2+8z=-4,$ que corresponde a la proyección del sólido en el plano $xz$, esta ecuación equivale a
$\frac{\left( x+2\right) ^2}{4}+\frac{\left( z+1\right) ^2}{9}=1$
y representa una elipse con centro en el punto $(-2,-1)$,
alargada hacia el eje $z.$ De esta ecuación se sigue que
$x=-2\pm \sqrt{4-\frac{4\left( z+1\right) ^2}{9}} $ y $z=-1\pm
\sqrt{9-\frac{9\left( x+2\right) ^2}{4}}.$ 
\vskip 0.3cm
Si $z=-1,$ las ecuaciones se reducen a
$\frac{\left( x+2\right) ^2}{4}=\frac{8-y}{7}$ y $9\left( x+2\right) ^2=36\left( \frac{y+2}{3}\right) $
y representan dos parábolas en el plano $xy.$
Si $x=-2,$, estas ecuaciones se reducen a $\frac{\left( z+1\right) ^2}{9}=\frac{8-y}{7}$ y
$4\left( z+1\right) ^2=36\left( \frac{y+2}{3}\right) $
y representan dos parábolas en el plano $yz$.
\vskip 0.3cm
Por lo descrito anteriormente, el sólido en mención se proyecta sobre los planos coordenados como se muestra a continuación.
\vskip0.4cm

\noindent
\begin{minipage}{\textwidth}
	\centering
	\captionof{figure}{Región acotada por los paraboloides proyectada en los planos coordenados.}
	\begin{minipage}{0.30\linewidth}
		\centering
		\vspace{-10mm}	\includegraphics[width=\linewidth]{Fig/350.pdf}
	\end{minipage}
	\hspace{1em}
	\begin{minipage}{0.30\linewidth}
		\centering
		\vspace{-10mm}
		\includegraphics[width=\linewidth]{Fig/351.pdf}
	\end{minipage}
	\hspace{1em}
	\begin{minipage}{0.30\linewidth}
		\centering
		\vspace{-10mm}
		\includegraphics[width=\linewidth]{Fig/352.pdf}
	\end{minipage}
	\fuentepropia
\end{minipage}
\vskip0.3cm

Utilizando esta representación se plantean las integrales triples que determinan el volumen del sólido:
\vskip0.3cm
\begin{equation*}
\bigintsss_{-4}^{0}\bigintsss_{-1-\sqrt{9-\frac{9\left( x+2\right) ^2}{4}}}^{-1+\sqrt{9-\frac{9\left( x+2\right) ^2}{4}}}
\bigintsss_{\frac{3}{4}x^2+3x+\frac{4}{3}+\frac{1}{3}z^2+\frac{2}{3}z}^{\frac{2}{9}-\frac{7}{4}x^2-7x-\frac{7}{9}z^2-\frac{14}{9}z}
dydzdx = 30\pi
\end{equation*}

\begin{equation*}
\bigintsss_{-4}^2\bigintsss_{-2-\sqrt{4-\frac{4\left( z+1\right) ^2}{9}}}^{-2+\sqrt{4-\frac{4\left( z+1\right) ^2}{9}}}
\bigintsss_{\frac{3}{4}x^2+3x+\frac{4}{3}+\frac{1}{3}z^2+\frac{2}{3}z}^{\frac{2}{9}-\frac{7}{4}x^2-7x-\frac{7}{9}z^2-\frac{14}{9}z}
dydxdz = 30\pi.
\end{equation*}
\vskip0.4cm
Para integrar primero con respecto a $z$, se debe obtener $z$ de las expresiones originales y tener en cuenta que $z$ recorre los puntos desde la parte negativa de cada paraboloide hasta su parte positiva.

\begin{multline*}
\bigintsss_{-4}^{0}\bigintsss_{1}^{8-\frac{7}{4}\left( x+2\right) ^2}\bigintsss_{-1-3\sqrt{\frac{8-y}{7}-\frac{\left( x+2\right) ^2}{4}}}^{-1+3\sqrt{\frac{8-y}{7}-\frac{\left( x+2\right) ^2}{4}}}dzdydx \\
+\bigintsss_{-4}^{0}\bigintsss_{\frac{3}{4}\left( x+2\right) ^2-2}^{1}\bigintsss_{-1-3\sqrt{\frac{y+2}{3}-\frac{1}{4}\left(x+2\right) ^2}}^{-1+3\sqrt{\frac{y+2}{3}-\frac{1}{4}\left( x+2\right) ^2}}dzdydx= 30\pi.
\end{multline*}

Al integrar primero con respecto a $x,$ se debe presente que $x$ igual que en el caso anterior, recorre los puntos desde la parte negativa de cada paraboloide hasta la parte positiva del mismo, por tal razón las integrales en este caso son
\vspace{-2mm}
\begin{multline*}
\bigintsss_{-4}^2\bigintsss_{1}^{8-\frac{7\left( z+1\right) ^2}{9}}\bigintsss_{-2-2\sqrt{\frac{8-y}{7}-\frac{\left( z+1\right) ^2}{9}}}^{-2+2\sqrt{\frac{8-y}{7}-\frac{\left( z+1\right) ^2}{9}}}dxdydz \\
+\bigintsss_{-4}^2\bigintsss_{\frac{\left(z+1\right) ^2}{3}-2}^{1}\bigintsss_{-2-2\sqrt{\frac{1}{3}\left( y+2\right) -\frac{1}{9}\left( z+1\right) ^2}}^{-2+2\sqrt{\frac{1}{3}\left( y+2\right) -\frac{1}{9}\left( z+1\right) ^2}}dxdydz= 30\pi
\end{multline*}
Con las siguientes instrucciones se evaluó la integral anterior.

\vspace{0.6em}
\noindent
\begin{tcolorbox}[breakable,enhanced,title=\bf Instrucción en \verb"Mathematica", top=2pt,bottom=2pt,boxsep=0pt,before skip=2pt,after skip=0pt]
\begin{Verbatim}[formatcom=\letraverbatim]
Integrate[1,{z,-4,2},{y,1,8-7*(z+1)^2/9},{x,-2-
2*Sqrt[(8-y)/7-(z+1)^2/9],-2+2*Sqrt[(8-y)/7-
(z+1)^2/9]}]+Integrate[1,{z,-4,2},{y,-2+(z+1)^2/3,1},
{x,-2-Sqrt[(4/3)*(y+2)-4(z+1)^2/9],-2+Sqrt[(4/3)*(y+2)
4(z+1)^2/9]}]
\end{Verbatim}
\end{tcolorbox}
\vspace{0.7em}
\noindent


También puede usarse la siguiente opción para calcular el volumen.

\vskip0.3cm

%\begin{tcolorbox}
En el plano $xz$ la proyección del sólido es una elipse, cuyo interior se parametriza como $x=2u\cos(t)-2$, $z=3u\sin(t)-1.$ Sobre los parabo\-loides, se obtiene $y=3u^2-2$ y $y=8-7u^2;$
con dos puntos sobre estas superficies $P(-2+2u\cos (t) ,-1+3u\sin (t),3u^2-2)$ y
$Q(2u\cos (t)-2,3u\sin (t)-1,8-7u^2)$, se forma una combinación convexa de la forma $(1-w)P+wQ,$ que define la función que parametriza el sólido.
\vspace{-3mm}
\begin{equation*}
{\bf r}\left( t,u,w\right) =\langle 2u\cos t-2,3u\sin t-1,10w\left( 1-u^2\right) +3u^2-2 \rangle,
\end{equation*}
con $t\in [0,2\pi ],$ $u\in [0,1],$ $w\in [0,1];$
y obtenemos los siguientes vectores
\begin{eqnarray*}
{\bf r}_{t}&=& \langle -2u\sin t,3u\cos t,0\rangle \\
{\bf r}_{u}&=& \langle 2\cos t,3\sin t,6u-20uw\rangle\\
{\bf r}_{w}&=& \langle 0,0,10-10u^2\rangle.
\end{eqnarray*}
El elemento diferencial de volumen es
$\left\vert {\bf r}_{t}\cdot \left({\bf r}_{u}\times {\bf r}_{w}\right) \right\vert
=\left\vert 60u\left( u^2-1\right) \right\vert ,$
cuya integral proporciona el volumen del sólido mencionado
\vskip0.2cm
\[\int_{0}^{1}\int_{0}^{1}\int_{0}^{2\pi }\left\vert 60u\left( u^2-1\right)\right\vert dtdudw= 30\pi.\]
%\end{tcolorbox}

\section{Un hiperboloide y una esfera} \index{Esfera} \index{Hiperboloide}

Las ecuaciones $x^2+y^2+z^2=16$ y $x^2+z^2=9+y^2,$ representan una esfera de radio $4$
y un hiperboloide de un manto con eje de simetría el eje $y.$
La región interior a ambas superficies determinó el sólido que se muestra en seguida.

\vskip0.2cm
\noindent
\begin{minipage}{\textwidth}
	\centering
	\captionof{figure}{Región determinada por $x^2 + y^2 + z^2 \leq 16$ y $x^2 + z^2 \leq 9 + y^2$.}
	\label{fig:region_esferica_cilindrica}
	\begin{minipage}{0.20\linewidth}
		\centering
		\includegraphics[width=\linewidth]{Img/61a.pdf}
	\end{minipage}
	\hspace{1em} 
	\begin{minipage}{0.20\linewidth}
		\centering
		\includegraphics[width=\linewidth]{Img/61c.pdf}
	\end{minipage}
	\hspace{1em} 
	\begin{minipage}{0.20\linewidth}
		\centering
		\includegraphics[width=\linewidth]{Img/61b.pdf}
	\end{minipage}
	
	\vspace{0.3em}
	\fuentepropia
\end{minipage}

\vskip0.2cm
La intersección de la esfera y el hiperboloide se da en $y=\pm \frac{\sqrt{14}}{2}.$
En seguida se muestra la proyección del sólido en los planos $xz$ y $xy$.
Para la región que contiene al origen en el plano $xz$, $y$ recorre, por un lado, los puntos del sólido que van desde la parte negativa de la esfera hasta su parte positiva, y por otro lado, para la región entre las dos circunferencias, $y$
recorre los puntos entre el hiperboloide y la esfera; esto ocurre en dos regiones simétricas con respecto al plano $xz.$

\vskip0.2cm
\noindent
\begin{minipage}{\textwidth}
	\centering
	\captionof{figure}{Proyección en los planos coordenados.}
	\begin{minipage}{0.30\linewidth}
		\centering
		\vspace{-17mm}
		\includegraphics[width=\linewidth]{Fig/353.pdf}
	\end{minipage}
	\hspace{1em}
	\begin{minipage}{0.38\linewidth}
		\centering
		\vspace{-17mm}
		\includegraphics[width=\linewidth]{Fig/354.pdf}
	\end{minipage}
	
	\vspace{-14mm}
	\fuentepropia
\end{minipage}

% \begin{figure}[H]
% \centering
% % Primer gráfico
% \begin{pspicture}(-1.75,-2)(3,2) %\psgrid
% \psset{algebraic,unit=0.45}
% \pscustom[fillstyle=hlines,hatchcolor=gray!60,hatchsep=1.2pt,hatchwidth=0.5pt,opacity=0.7,linestyle=none]{
%   \psparametricplot{0}{\psPiTwo}{3*cos(t)|3*sin(t)}
% }
% \pscustom[fillstyle=solid,fillcolor=gray!60,opacity=0.7,linestyle=none]{
%   \psparametricplot{0}{\psPiTwo}{3*cos(t)|3*sin(t)}
%   \psline(3,0)(4,0)
%   \psparametricplot{\psPiTwo}{0}{4*cos(t)|4*sin(t)}
% }
% \psaxes[Dy=3,Dx=3,linecolor=black,labelFontSize=\scriptstyle]{->}(0,0)(-5,-5)(5.5,5.6)[$x$,90][$z$,0]
% \psparametricplot[linecolor=black!60!blue]{0}{\psPiTwo}{3*cos(t)|3*sin(t)}
% \psparametricplot[linecolor=black!60!blue]{0}{\psPiTwo}{4*cos(t)|4*sin(t)}
% \psline[linecolor=black!60!blue,linewidth=0.6pt,arrows=<->](-3,-3.5)(-3,3.75)
% \psline[linecolor=black!60!blue,linewidth=0.6pt,arrows=<->](3,-3.5)(3,3.75)
% \rput(0.3,0.8){$x^2+z^2=9$}
% \rput(4,4.5){$x^2+z^2=\frac{25}{2}$}
% \end{pspicture}

% \vspace{0.5cm}

% % Segundo gráfico
% \begin{pspicture}(-2.8,-2)(3,2) % \psgrid
% \psset{algebraic,unit=0.425}
% \psaxes[Dy=4,Dx=3,linecolor=black,labelFontSize=\scriptstyle]{->}(0,0)(-5,-5)(8,5.6)[$x$,90][$y$,0]
% \pscustom[fillstyle=hlines,hatchcolor=gray!60,hatchsep=1.2pt,hatchwidth=0.5pt,opacity=0.7,linestyle=none]{
%   \psparametricplot{-0.59}{0.59}{3*COSH(t)|3*SINH(t)}
%   \psparametricplot{0.49}{2.65}{4*cos(t)|4*sin(t)}
%   \psparametricplot{0.59}{-0.59}{-3*COSH(t)|3*SINH(t)}
%   \psparametricplot{3.63}{5.79}{4*cos(t)|4*sin(t)}
% }
% \psparametricplot[linecolor=black!60!blue]{0}{\psPiTwo}{4*cos(t)|4*sin(t)}
% \psparametricplot[linecolor=black!60!blue,arrows=->]{-1.1}{1.2}{3*COSH(t)|3*SINH(t)}
% \psparametricplot[linecolor=black!60!blue,arrows=->]{-1.1}{1.2}{-3*COSH(t)|3*SINH(t)}
% \psline[linecolor=black!60!blue,linewidth=0.6pt,arrows=->](-5,-1.87)(5.5,-1.87)
% \psline[linecolor=black!60!blue,linewidth=0.6pt,arrows=->](3.53,3.5)(3.53,-4)
% \rput(8.5,-2.1){$y=-\sqrt{\frac{7}{2}}$}
% \rput(3.75,-5){$x=\frac{5}{\sqrt{2}}$}
% \rput(7,5){$x^2-y^2=9$}
% % \rput(-0.4,2){$x^2+y^2=16$}
% \end{pspicture}
% \caption{Proyección en los planos coordenados.}
% \end{figure}
% \par
% \justifying

\vskip0.2cm
Sobre el plano $xz,$ se puede utilizar la simetría de la región, integrando la región para $z \geq 0,$ y se multiplica por $4$ a estas integrales. En este caso, el volumen del sólido está dado por:

\vspace{-4mm}
{\small
\begin{multline*}
\hskip-0.4cm \int_{-3}^{3}\int_{-\sqrt{9-x^2}}^{\sqrt{9-x^2}}\int_{-\sqrt{16-x^2-z^2}}^{\sqrt{16-x^2-z^2}}dydzdx + 4\bigg(\int_{-3}^{3}\int_{\sqrt{9-x^2}}^{\sqrt{\frac{25}{2}-x^2}}\int_{\sqrt{x^2+z^2-9}}^{\sqrt{16-x^2-z^2}}dydzdx+ \\
2\int_{3}^{\frac{5}{\sqrt{2}}}\int_{0}^{\sqrt{\frac{25}{2}-x^2}}\int_{\sqrt{x^2+z^2-9}}^{\sqrt{16-x^2-z^2}}dydzdx\bigg)= \frac{2\pi}{3}(128-7 \sqrt{14}).
\end{multline*}
}

%\vskip0.2cm
Si se integra primero con respecto a $z$, hay tres regiones en el espacio, en dos de ellas $z$ recorre los puntos desde la parte negativa de la esfera hasta su parte positiva; en la tercer región $z$ recorre los puntos que van desde la parte negativa del hiperboloide hasta la parte positiva del mismo, razón por la cual las integrales a evaluar en este caso son:
\vspace{-2mm}
\begin{multline*}
2\bi_{\sqrt{\frac{7}{2}}}^{4}\bi_{-\sqrt{16-y^2}}^{\sqrt{16-y^2}}\bi_{-\sqrt{16-x^2-y^2}}^{\sqrt{16-x^2-y^2}}dzdxdy+ \\
\bi_{-\sqrt{\frac{7}{2}}}^{\sqrt{\frac{7}{2}}}\bi_{-\sqrt{9+y^2}}^{\sqrt{9+y^2}}\bi_{-\sqrt{9+y^2-x^2}}^{\sqrt{9+y^2-x^2}}dzdxdy =\frac{2}{3}\pi\left( 128-7\sqrt{14}\right).
\end{multline*}
En coordenadas cilíndricas, el cálculo es el siguiente:
\begin{multline*}
\int_{0}^{2\pi }\int_{0}^{3}\int_{-\sqrt{16-r^2}}^{\sqrt{16-r^2}}rdydrd\theta \\ +2\int_{0}^{2\pi}\int_{3}^{\frac{5}{\sqrt{2}}}\int_{\sqrt{r^2-9}}^{\sqrt{16-r^2}}rdydrd\theta = \frac{2\pi}{3}(128-7 \sqrt{14}),
\end{multline*}
y su obtención en \verb"Mahematica" se consigue con la siguiente sentencia:

\vspace{0.6em}
\noindent
\begin{tcolorbox}[breakable,enhanced,title=\bf Instrucción en \verb"Mathematica", top=2pt,bottom=2pt,boxsep=0pt,before skip=2pt,after skip=0pt]
\begin{Verbatim}[formatcom=\letraverbatim]
Integrate[r,{\[Theta],0,2*Pi},{r,0,3},{y,-Sqrt[16-
r^2],Sqrt[16-r^2]}]+2*Integrate[r,{\[Theta],0,2*Pi},
{r,3,5/Sqrt[2]},{y,Sqrt[r^2-9],Sqrt[16-r^2]}]
\end{Verbatim}
\end{tcolorbox}
\vspace{0.7em}
\noindent


\begin{comment}
Estos gráficos se obtienen en \verb"Mathematica" con las siguientes instrucciones

\begin{verbatim}
ContourPlot3D[{x^2+z^2==9+y^2,x^2+y^2+z^2==16},{x,-5,5},
{y,-4,4},{z,-5,5},Mesh->False,Ticks->{{-3,0,3},{-3,0,3},
{-4,-2,2,4}},AxesLabel->{x,y,z},
ContourStyle->{Blend[{Yellow,Cyan},1/9],Orange}]
\end{verbatim}

\begin{verbatim}
RegionPlot3D[x^2+z^2<9+y^2&&x^2+y^2+z^2<16,{x,-4,4},{y,-4,4},
{z,-4,4},PlotPoints->90,Mesh->False,Ticks->{{-4,-2,0,2,4},
{-3,-1,1,3},{-3,-1,1,3}},PlotStyle->RGBColor[0.3,0.9,0.5],
AxesLabel->{x,y,z}]
\end{verbatim}


La intersección de la esfera y el hiperboloide se obtiene al sustituir $x^2+z^2=9+y^2$ en la ecuación $x^2+y^2+z^2=16,$
lo que genera la expresión $2y^2=7,$, es decir $y=\pm \frac{\sqrt{14}}{2}.$

% \begin{figure}
% \centering
% % Primer gráfico: proyección en plano xz
% \begin{pspicture}(-2,-1.4)(3,2) % \psgrid
% \psset{algebraic,unit=0.36}
% \pscustom[fillstyle=hlines,hatchcolor=gray!60,hatchsep=1.2pt,hatchwidth=0.5pt,opacity=0.7,linestyle=none]{
%   \psparametricplot{0}{\psPiTwo}{3*cos(t)|3*sin(t)}
% }
% \pscustom[fillstyle=solid,fillcolor=gray!60,opacity=0.7,linestyle=none]{
%   \psparametricplot{0}{\psPiTwo}{3*cos(t)|3*sin(t)}
%   \psline(3,0)(4,0)
%   \psparametricplot{\psPiTwo}{0}{4*cos(t)|4*sin(t)}
% }
% \psaxes[Dy=3,Dx=3,linecolor=black,labelFontSize=\scriptstyle]{->}(0,0)(-5,-5)(5.5,5.6)[$x$,90][$z$,0]
% \psparametricplot[linecolor=black!60!blue]{0}{\psPiTwo}{3*cos(t)|3*sin(t)}
% \psparametricplot[linecolor=black!60!blue]{0}{\psPiTwo}{4*cos(t)|4*sin(t)}
% \psline[linecolor=black!60!blue,linewidth=0.6pt,arrows=<->](-3,-3.5)(-3,3.75)
% \psline[linecolor=black!60!blue,linewidth=0.6pt,arrows=<->](3,-3.5)(3,3.75)
% \rput(0.3,0.8){$x^2+z^2=9$}
% \rput(4,4.5){$x^2+z^2=\frac{25}{2}$}
% \end{pspicture}

% \vspace{0.5cm}

% % Segundo gráfico: proyección en plano xy
% \begin{pspicture}(-2.5,-1.4)(3,2) % \psgrid
% \psset{algebraic,unit=0.35}
% \psaxes[Dy=4,Dx=3,linecolor=black,labelFontSize=\scriptstyle]{->}(0,0)(-5,-5)(8,5.6)[$x$,90][$y$,0]
% \pscustom[fillstyle=hlines,hatchcolor=gray!60,hatchsep=1.2pt,hatchwidth=0.5pt,opacity=0.7,linestyle=none]{
%   \psparametricplot{-0.59}{0.59}{3*COSH(t)|3*SINH(t)}
%   \psparametricplot{0.49}{2.65}{4*cos(t)|4*sin(t)}
%   \psparametricplot{0.59}{-0.59}{-3*COSH(t)|3*SINH(t)}
%   \psparametricplot{3.63}{5.79}{4*cos(t)|4*sin(t)}
% }
% \psparametricplot[linecolor=black!60!blue]{0}{\psPiTwo}{4*cos(t)|4*sin(t)}
% \psparametricplot[linecolor=black!60!blue,arrows=->]{-1.1}{1.2}{3*COSH(t)|3*SINH(t)}
% \psparametricplot[linecolor=black!60!blue,arrows=->]{-1.1}{1.2}{-3*COSH(t)|3*SINH(t)}
% \psline[linecolor=black!60!blue,linewidth=0.6pt,arrows=->](-5,-1.87)(5.5,-1.87)
% \psline[linecolor=black!60!blue,linewidth=0.6pt,arrows=->](3.53,3.5)(3.53,-4)
% \rput(7.75,-1.9){$y=-\sqrt{\frac{7}{2}}$}
% \rput(3.75,-5){$x=\frac{5}{\sqrt{2}}$}
% \rput(4.5,5){$x^2-y^2=9$}
% \rput(-0.4,2){$x^2+y^2=16$}
% \end{pspicture}
% \caption{Proyección en los planos coordenados.}
% \end{figure}



% \begin{figure}
% \centering
% \begin{pspicture}(-2.5,-1.8)(3,2.3) % \psgrid
% \psset{algebraic,unit=0.45}
% \pscustom[fillstyle=hlines,hatchcolor=gray!90,hatchsep=1.2pt,hatchwidth=0.5pt,linestyle=none]{
%   \psparametricplot{0}{\psPiTwo}{3*cos(t)|3*sin(t)}
% }
% \pscustom[fillstyle=solid,fillcolor=gray!60,opacity=0.7,linestyle=none]{
%   \psparametricplot{0}{\psPiTwo}{3*cos(t)|3*sin(t)}
%   \psline(3,0)(4,0)
%   \psparametricplot{\psPiTwo}{0}{5/sqrt(2)*cos(t)|5/sqrt(2)*sin(t)}
% }
% \psaxes[Dy=3,Dx=3,linecolor=black,labelFontSize=\scriptstyle]{->}(0,0)(-4.5,-4.5)(5.5,5.2)[$x$,90][$z$,0]
% \psparametricplot[linecolor=black!60!blue]{0}{\psPiTwo}{3*cos(t)|3*sin(t)}
% \psparametricplot[linecolor=black!60!blue]{0}{\psPiTwo}{5/sqrt(2)*cos(t)|5/sqrt(2)*sin(t)}
% %\rput(4,4.75){$x^2+z^2=\frac{25}{2}$}
% %\rput(5,3.5){$x^2+z^2=9$}
% \psline[linewidth=1.2pt,linestyle=dashed](-1.96,2.94)(0,0)
% % \psline[linewidth=0.6pt,arrows=->](2.2,2)(4,3)
% \rput(-1.96,2.94){$\bullet$}
% \rput(-1.66,2.496){$\bullet$}
% \rput(-1.8,1.65){$P$}
% \rput(-2.2,3.55){$Q$}
% \end{pspicture}
% \caption{$9 \leq x^2+z^2 \leq \frac{25}{2}$.}
% \label{figura168}
% \end{figure}

El eje $y$ es el eje de simetría del hiperboloide, por este motivo usaremos la proyección en el plano $xz,$ para construir la parametrización del sólido. El gráfico muestra la figura de curva de intersección $x^2+z^2=\frac{25}{2}$
y la traza del hiperboloide cuando $y=0.$
La región acotada por $x^2+z^2=9,$ se describe usando las expresiones $x=3u\sin v,$ $z=3u\cos v.$

Sobre el sólido, esta región se proyecta y determina los puntos que estan entre la parte negativa de la esfera hasta la parte positiva de la misma. Es decir, la expresión para $y,$ es
$y=\pm \sqrt{16-x^2-z^2}= \pm \sqrt{16-9u^2}$.
Es por lo anterior que la siguiente combinación convexa
$(1-w) \langle 3u\sin v,-\sqrt{16-9u^2},3u\cos v\rangle + w\langle 3u\sin v,\sqrt{16-9u^2},3u\cos v\rangle$
permite definir una función que parametriza el sólido
$ {\bf r} (u,v,w)= \left\langle 3u\sin v,\left(2w-1\right) \sqrt{16-9u^2},3u\cos v\right\rangle $
{\bf El volumen de la parte del sólido interior a la esfera y al cilindro} $x^2+z^2=9,$
se determina con el siguiente procedimiento. Con la función ${\bf r}$ hallamos los siguientes vectores
\begin{eqnarray*}
{\bf r}_{u}&=& \langle 3\sin v,\frac{9u(1-2w)}{\sqrt{16-9u^2}},3\cos v\rangle \\
{%\bf r}_{v}&=& \langle 3u\cos v,0,-3u\sin v \rangle \\
{\bf r}_{w}&=& \langle 0,2\sqrt{16-9u^2},0\rangle
\end{eqnarray*}
Entonces, $\left\vert {\bf r}_{w}\cdot \left({\bf r}_{u}\times {\bf r}_{v}\right) \right\vert =18u\sqrt{16-9u^2},$
Por tanto, el volumen de la mencionada parte del sólido es
\begin{equation} \label{hiperesfera1}
\int_{0}^{1}\int_{0}^{1}\int_{0}^{2\pi }2u\sqrt{16-u^2} dvdudw= \frac{4}{3}(64 - 7\sqrt{7})\pi
\end{equation}


\begin{tcolorbox}
Por otra parte, la región entre las dos circunferencias de la figura \ref{figura168}, se describe formando segmentos rectilíneos que unan dos puntos $P=\left(\frac{5}{\sqrt{2}}\sin v,\frac{5}{\sqrt{2}}\cos v\right)$
y $Q=\left(3 \sin v,3\cos v\right)$. Esto se consigue simplificando una combinación convexa de la forma
$(1-u)P +uQ$, entonces
\begin{eqnarray*}
x &=& \frac{1}{2} \left(-5\sqrt{2}u+6u+5\sqrt{2}\right)\sin(v) \\
z &=& \frac{1}{2}\left(-5\sqrt{2}u+6u+5\sqrt{2}\right)\cos (v)
\end{eqnarray*}
Las siguientes expresiones determinan, respectivamente, la coordenada $y$ sobre la esfera y el hiperboloide.
\begin{eqnarray*}
y_1&=&\sqrt{16-x^2-z^2}=\frac{1}{\sqrt{2}}\sqrt{\left(30 \sqrt{2}-43\right) u^2+\left(50-30 \sqrt{2}\right) u+7} \\
y_2 &=& \sqrt{x^2+z^2-9}=\frac{1}{\sqrt{2}}\sqrt{(1-u) \left(\left(30 \sqrt{2}-43\right) u+7\right)}
\end{eqnarray*}
La expresión $(1-w)y_1+wy_2$ determina la componente en $y$ de la función que parametriza el sólido.
Es por lo anterior que la siguiente función se utilizará como dicha parametrización.
\begin{multline*}
{\bf r}(u,v,w) = \bigg\langle \frac{1}{2} \left(-5\sqrt{2}u+6u+5\sqrt{2}\right)\sin(v),(1-w)y_1+wy_2,\\
\frac{1}{2}\left(-5\sqrt{2}u+6u+5\sqrt{2}\right)\cos (v) \bigg\rangle
\end{multline*}
Puede probarse (usando \verb"Mathematica", por ejemplo) que el elemento diferencial de volumen es
$|{\bf r_w}\cdot({\bf r_u}\times{\bf r_v})|=\bigg\vert\frac{1}{2}\left(30\sqrt{2}u-43u-15\sqrt{2}+25\right)(y_1-y_2)\bigg\vert.$
El volumen de la segunda región es
\begin{eqnarray}
\label{hiperesfera2}
2\int_0^1\int_0^1\int_0^{2\pi}|{\bf r_w}\cdot({\bf r_v}\times{\bf r_u})|dudvdw&=& \frac{14\sqrt{14}}{3}(\sqrt{2}-1)\pi \\
&\approx & 22.722 \nonumber
\end{eqnarray}
Nótese que la integral en (\ref{hiperesfera2}) se multiplicó por $2$ por la simetría del sólido con respecto al plano $xz.$
\end{tcolorbox}
El volumen del sólido es la suma de los resultados obtenidos por las expresiones (\ref{hiperesfera1}) y (\ref{hiperesfera2}).
\vskip0.2cm
Los cálculos anteriores se efectuaron en \verb"Mathematica" utilizando las siguientes instrucciones.

\begin{verbatim}
{x,z}={3*u*Sin[v],3*u*Cos[v]};
r=(1-w)*{x,-Sqrt[16-x^2-z^2],z}+w*{x,Sqrt[16-x^2-z^2],z};
Integrate[Dot[D[r,w],Cross[D[r,u],D[r,v]]],{w,0,1},
{u,0,1},{v,0,2*Pi}]
\end{verbatim}
%
\begin{verbatim}
{x1,z1}=5/Sqrt[2]*{Sin[v],Cos[v]};
{x2,z2}=3*{Sin[v],Cos[v]};
{x,z}=(1-u)*{x1,z1}+u*{x2,z2};
y1=Simplify[Sqrt[16-x^2-z^2]]
y2=Simplify[Sqrt[x^2+z^2-9]]
r={x,(1-w)*y1+w*y2,z};
A=Simplify[Dot[D[r,u],Cross[D[r,v],D[r,w]]]]
2*NIntegrate[A,{w,0,1},{u,0,1},{v,0,2*Pi}]
\end{verbatim}
\end{comment}

\begin{ejer}
Plantee y evalúe las integrales que determinan el volumen del sólido determinado por
\begin{enumerate}[font=\normalfont]
\item \textit{Las expresiones $63x^2+252x+28z^2+56z+36y=8$ y $9x^2+36x+16+4z^2+8z-12y=0$,
En las órdenes $dxdydz$ y $dzdxdy.$}
\item \textit{$x^2 + y^2 + z^2 \leq 16$ y $x^2 + z^2 \leq 9 + y^2$; en los órdenes $dydxdz$, $dzdydx$ $dxdydz$ y $dxdzdy.$}
Parametrice el sólido y obtenga el mismo valor por este método.
\end{enumerate}
\end{ejer}

\section{Un semicono y una semiesfera} \index{Semiesfera}
El semicono con ecuación $z=1-\sqrt{x^2+y^2}$ y la semiesfera determinada por la expresión
$z=-\sqrt{1-x^2-y^2}$, se intersecan en una circunferencia de radio $1,$ centrada en el origen.
La figura de este sólido se obtiene en \verb"Mathematica" con la siguiente instrucción:


\vspace{0.6em}
\noindent
\begin{tcolorbox}[breakable,enhanced,title=\bf Instrucción en \verb"Mathematica", top=2pt,bottom=2pt,boxsep=0pt,before skip=2pt,after skip=0pt]
\begin{Verbatim}[formatcom=\letraverbatim]
RegionPlot3D[x^2+y^2+z^2<=1&&z<1-Sqrt[x^2+y^2]&&x^2
+y^2<=1,{x,-1,1},{y,-1,1},{z,-1,1},PlotPoints->90,
Mesh->None]
ContourPlot3D[{x^2+y^2+z^2==1,(z-1)^2==x^2+y^2},
{x,-2,2},{y,-2,2},{z,-1,1},PlotPoints->90,Mesh->None,
BoxRatios->{1,1,1/2},ContourStyle->Opacity[0.85]]
\end{Verbatim}
\end{tcolorbox}
\vspace{0.7em}
\noindent
\vskip0.2cm

El sólido mencionado es el siguiente
\vskip0.2cm

\noindent
\begin{minipage}{\textwidth}
	\centering
	\captionof{figure}{Sólido determinado por $-\sqrt{1-x^2-y^2} \leq z\leq 1-\sqrt{x^2+y^2}.$}
	 \label{fig:solido_eps}
	\begin{minipage}{0.28\linewidth}
		\centering
		\includegraphics[width=\linewidth]{Img/13c.pdf}
	\end{minipage}
	\hspace{1em}
	\begin{minipage}{0.20\linewidth}
		\centering
		\includegraphics[width=\linewidth]{Img/13b.pdf}
	\end{minipage}
	\hspace{1em}
	\begin{minipage}{0.20\linewidth}
		\centering
		\includegraphics[width=\linewidth]{Img/13a.pdf}
	\end{minipage}
	
	\vspace{0.3em}
	\fuentepropia
\end{minipage}

\vskip0.2cm
\noindent
\begin{minipage}{\textwidth}
	\centering
	\captionof{figure}{Proyecciones del sólido en los planos coordenados.}
	 \label{fig:proyecciones}
	\begin{minipage}{0.32\linewidth}
		\centering
		\vspace{-10mm}
		\includegraphics[width=\linewidth]{Fig/355.pdf}
	\end{minipage}
	\hspace{1em}
	\begin{minipage}{0.32\linewidth}
		\centering
		\vspace{-10mm}
		\includegraphics[width=\linewidth]{Fig/356.pdf}
	\end{minipage}
	\begin{minipage}{0.28\linewidth}
		\centering
		\vspace{-10mm}
		\includegraphics[width=\linewidth]{Fig/357.pdf}
	\end{minipage}
	
	\vspace{-10mm}
	\fuentepropia
\end{minipage}

\vskip0.2cm
Con estas figuras se establecen las integrales triples, que determinan el volumen del sólido:
{\small
\[
\int_{-1}^{1} \int_{-\sqrt{1 - x^2}}^{\sqrt{1 - x^2}} \int_{-\sqrt{1 - x^2 - y^2}}^{1 - \sqrt{x^2 + y^2}} dz\,dy\,dx
=
\int_{-1}^{1} \int_{-\sqrt{1 - y^2}}^{\sqrt{1 - y^2}} \int_{-\sqrt{1 - x^2 - y^2}}^{1 - \sqrt{x^2 + y^2}} dz\,dx\,dy
= \pi.
\]
}

\vskip0.3cm
De la ecuación del cono se sigue que
{\small $y=\pm \sqrt{\left( 1-z\right) ^2-x^2}$} y con la ecuación de la semiesfera se obtiene {\small $y=\pm \sqrt{1-x^2-z^2}.$} Lo anterior permite obtener el volumen del sólido planteando la integral triple en el orden $dydxdz,$

$\ds \int_{0}^{1}\int_{z-1}^{1-z}\int_{-\sqrt{\left( 1-z\right)^2-x^2}}^{\sqrt{\left( 1-z\right) ^2-x^2}}dydxdz+\int_{-1}^{0}\int_{-\sqrt{1-z^2}}^{\sqrt{1-z^2}}\int_{-\sqrt{1-x^2-z^2}}^{\sqrt{1-x^2-z^2}}dydxdz=\pi. $
\vskip0.2cm
Esta integral se evalúa con \verb"Mathematica", ejecutando la siguiente instrucción.

\vspace{0.6em}
\noindent
\begin{tcolorbox}[breakable,enhanced,title=\bf Instrucción en \verb"Mathematica", top=2pt,bottom=2pt,boxsep=0pt,before skip=2pt,after skip=0pt]
\begin{Verbatim}[formatcom=\letraverbatim]
Integrate[1,{z,0,1},{x,z-1,1-z},{y,-Sqrt[(1-z)^2-x^2],
Sqrt[(1-z)^2-x^2]}]+Integrate[1,{z,-1,0},
{x,-Sqrt[1-z^2],Sqrt[1-z^2]},{y,-Sqrt[1-x^2-z^2],
Sqrt[1-z^2-x^2]}]
\end{Verbatim}
\end{tcolorbox}
\vspace{0.7em}
\noindent


En el caso del orden de integración $dydzdx,$, la primera de las integrales anteriores debe expresarse como la suma de dos integrales triples, es decir,
\vspace{-4mm}
{\small
\begin{multline*}
\int_{-1}^{0}\int_{0}^{1+x}\int_{-\sqrt{(1-z)^2-x^2}}^{\sqrt{(1-z)^2-x^2}}dydzdx+\int_{0}^{1}\int_{0}^{1-x}\int_{-\sqrt{\left( 1-z\right) ^2-x^2}}^{\sqrt{\left( 1-z\right) ^2-x^2}}dydzdx\\
+\int_{-1}^{1}\int_{-\sqrt{1-z^2}}^{\sqrt{1-z^2}}\int_{-\sqrt{1-x^2-z^2}}^{\sqrt{1-x^2-z^2}}dydzdx = \pi.
\end{multline*}
}
El sólido es simétrico con respecto al plano $xz$, las integrales en el orden $dydxdz$ y $dxdydz$ son idénticas.
En coordenadas cilíndricas este valor está dado por
{\small$\ds \int_{0}^{2\pi }\int_{0}^{1}\int_{-\sqrt{1-r^2}}^{1-r}rdzdrd\theta = \pi.$}
\vskip0.2cm

%\begin{tcolorbox}[title=\bfseries Parametrizando el sólido]
La intersección de las superficies está dada por
$x^2+y^2=1,$ que se parametriza con
$x= u\cos v,$ $y=u\sin v.$ Para la esfera $z= -\sqrt{1-u^2}$
y con el semicono se hace $z=1-u.$ Una combinación convexa de la forma
$\left( 1-w\right) \left( u\cos v,u\sin v,-\sqrt{1-u^2}\right) +w\left(u\cos v,u\sin v,1-u\right).$
Se define una parametrización del sólido mediante:
\vspace{-2mm}
\[
\mathbf{r}(u, v, w) =
\left\langle
u\cos v,\ 
u\sin v,\ 
\sqrt{1 - u^2}(w - 1) - w(u - 1)
\right\rangle.
\]

Con los parámetros en los siguientes rangos:

\[
u \in [0, 1], \quad v \in [0, 2\pi], \quad w \in [0, 1].
\]

Derivando parcialmente, se obtiene el elemento diferencial de volumen.
%\begin{eqnarray*} {\bf r}_{u}&=& \langle \cos(v),\sin(v),\frac{(1-w)u}{\sqrt{1-u^2}}-w\rangle \\
%{\bf r}_{v}&=&\langle -u\sin (v) ,u\cos (v) ,0\rangle \\ {\bf r}_{w}&=&\langle 0,0,\sqrt{1-u^2}-u+1\rangle . \end{eqnarray*}
$\left| {\bf r}_{w}\cdot \left( {\bf r}_{u}\times {\bf r}_{v}\right) \right| = u\left( \sqrt{1-u^2}-u+1\right),$
cuya integración nos proporciona el volumen del sólido,
%\vskip0.2cm
{\small
\[\int_{0}^{1}\int_{0}^{2\pi}\int_{0}^{1}u\left( \sqrt{1-u^2}-u+1\right) dudvdw= \pi.\]}
%\end{tcolorbox}

\section{Un elipsoide y un hiperboloide
} \index{Elipsoide} \index{Hiperboloide}

Las ecuaciones $x^2 + z^2 = 1 + y^2$ y $x^2 + y^2 + 4z^2 = 16$ representan un elipsoide y un hiperboloide de un manto; presentamos el sólido acotado por estas.


\noindent
\begin{minipage}{\textwidth}
	\centering
	\captionof{figure}{Sólido acotado por el elipsoide y el hiperboloide.}
	 \label{fig:elipsoide_hiperboloide}
	\begin{minipage}{0.20\linewidth}
		\centering
		\vspace{-2mm}
		\includegraphics[width=\linewidth]{Img/60c.pdf}
	\end{minipage}
	\hspace{1em}
	\begin{minipage}{0.20\linewidth}
		\centering
		\vspace{-2mm}
		\includegraphics[width=\linewidth]{Img/60a.pdf}
	\end{minipage}
	\hspace{1em}
	\begin{minipage}{0.20\linewidth}
		\centering
		\vspace{-2mm}
		\includegraphics[width=\linewidth]{Img/60aa.pdf}
	\end{minipage}
	
	\vspace{0.3em}
	\fuentepropia
\end{minipage}

\vskip0.2cm
El sólido se proyecta en los planos coordenados con el propósito de plantear las integrales que determinan su volumen.
\vskip0.2cm
\noindent
\begin{minipage}{\textwidth}
	\centering
	\captionof{figure}{El sólido proyectado en los planos $xy$ y $yz$.}
	\vspace{-15mm}
	\begin{minipage}{0.37\linewidth}
		\centering
		\vspace{-4mm}
		\includegraphics[width=\linewidth]{Fig/358.pdf}
	\end{minipage}
	\hspace{1em}
	\begin{minipage}{0.37\linewidth}
		\centering
		\vspace{-4mm}
		\includegraphics[width=\linewidth]{Fig/359.pdf}
	\end{minipage}
	
	%\vspace{0.3em}
	
	\vspace{-12mm}
	\fuentepropia
\end{minipage}
\vskip0.2cm


\begin{center}
	\begin{minipage}{0.40\linewidth} % Le doy un ancho adecuado para la figura
		\centering
		\captionof{figure}{Figura 6.X: Título descriptivo.} 
		\vspace{-19mm}
		\includegraphics[width=\linewidth]{Fig/360.pdf}
		\vspace{-1.2\baselineskip}
		\vspace{-2cm}
		\parbox{\linewidth}{\centering\fontsize{8pt}{10pt}\selectfont\textbf{Fuente:} elaboración propia.}
	\end{minipage}
\end{center}

\vskip0.2cm
\begin{ejer}
Plantee y evalúe las integrales que determinan el volumen del sólido determinado por las ecuaciones $x^2 + z^2 = 1 + y^2$ y $x^2 + y^2 + 4z^2 = 16.$
\end{ejer}


\section{Cilindros parabólicos} \index{Sólidos!cilindros parabólicos}
\vskip0.5cm
\subsection*{Dos cilindros parabólicos y un plano II}

Las ecuaciones $z=4-y^2$ y $y=2-x^2$ determinan dos cilindros parabólicos.
Estas superficies, junto con el plano $z=0$ acotan el sólido que se muestra a continuación, y del cual plantearemos las integrales que se requiere evaluar para determinar su volumen.

\vskip0.2cm
\noindent
\begin{minipage}{\textwidth}
	\centering
	\captionof{figure}{$0 \leq z \leq 4 - y^2$, \quad $y \leq 2 - x^2$.}
	\label{fig:proyeccion_xyz}
	\begin{minipage}{0.23\linewidth}
		\centering
		\includegraphics[width=\linewidth]{Img/34d.pdf}
	\end{minipage}
	\begin{minipage}{0.23\linewidth}
		\centering
		\includegraphics[width=\linewidth]{Img/34a.pdf}
	\end{minipage}
	\begin{minipage}{0.23\linewidth}
		\centering
		\includegraphics[width=\linewidth]{Img/34c.pdf}
	\end{minipage}
	
	\vspace{0.3em}
	\fuentepropia
\end{minipage}
\vskip0.2cm


Las figuras mostradas se obtienen con la ayuda de los siguientes comandos.

\vspace{0.6em}
\noindent
\begin{tcolorbox}[breakable,enhanced,title=\bf Instrucción en \verb"Mathematica", top=2pt,bottom=2pt,boxsep=0pt,before skip=2pt,after skip=0pt]
\begin{Verbatim}[formatcom=\letraverbatim]
ContourPlot3D[{z==4-y^2,2-x^2==y},{x,-2,2},{y,-2,2},
{z,0,4},ContourStyle->{Blend[{Yellow,Red}],
RGBColor[0.3,0.6,0.4]},Mesh->None,PlotPoints->60]
RegionPlot3D[0<z<4-y^2&&-Sqrt[2-y]<x<Sqrt[2-y]&&-2
<y<2,{x,-2,2},{y,-2,2},{z,0,4},PlotPoints->90,
Mesh->None,PlotStyle->Orange]
\end{Verbatim}
\end{tcolorbox}
\vspace{0.7em}
\noindent


Esta figura la proyectaremos sobre planos paralelos a los planos coordenados, con el fin de hacer sencilla la determinación de los límites de integración.
\vskip0.2cm
Integrando primero con respecto a $z,$ los puntos recorridos van desde $z=0$ hasta el cilindro $z=4-y^2$ y los límites para $x$ e $y$ se pueden obtener del segundo de los siguientes gráficos. Además, integrando primero con respecto a $x,$
se recorren puntos desde la parte negativa del cilindro $y = 2 - x^2$
hasta la parte positiva del mismo, pues el cilindro $z=4-y^2$ es paralelo al eje $x.$
\vskip0.2cm

\noindent
\begin{minipage}{\textwidth}
	\centering
	\captionof{figure}{Proyecciones del sólido en los planos $yz$ y $xy$.}
	\label{fig:proyecciones_yz_xy}
	\begin{minipage}{0.35\linewidth}
		\centering
		\vspace{-18mm}
		\includegraphics[width=\linewidth]{Fig/361.pdf}
	\end{minipage}
	\hspace{1em}
	\begin{minipage}{0.40\linewidth}
		\centering
		\vspace{-18mm}
		\includegraphics[width=\linewidth]{Fig/362.pdf}
	\end{minipage}
	
	
	\vspace{-15mm}
	\fuentepropia
\end{minipage}
\vskip0.2cm

Con la información observada en estos gráficos, las siguientes integrales triples determinan el volumen del sólido.

\begin{eqnarray*}
\int_{-2}^2\int_{-2}^{2-x^2}\int_{0}^{4-y^2}dzdydx=
\int_{-2}^2\int_{-\sqrt{2-y}}^{\sqrt{2-y}}\int_{0}^{4-y^2}dzdxdy &=& \frac{1024}{35} \\
\int_{0}^{4}\int_{-\sqrt{4-z}}^{\sqrt{4-z}}\int_{-\sqrt{2-y}}^{\sqrt{2-y}}dxdydz=
\int_{-2}^2\int_{0}^{4-y^2}\int_{-\sqrt{2-y}}^{\sqrt{2-y}}dxdzdy&=& \frac{1024}{35}.
\end{eqnarray*}
\vskip0.3cm


La proyección en el plano $xz$ merece varias consideraciones y la región queda dividida en dos subregiones; en una de ellas y recorre los puntos desde el cilindro
$z =4-y^2$ hasta el cilindro $y=2-x^2$ y en la otra subregion $y$ recorre puntos desde la parte negativa del cilindro $z=4-y^2$ hasta la parte positiva del mismo cilindro.


\vskip0.2cm
\noindent
\begin{minipage}{\textwidth}
	\centering
	\captionof{figure}{Proyección en el plano $xz$}
	\begin{minipage}{0.38\linewidth}
		\centering
		\vspace{-18mm}
		\includegraphics[width=\linewidth]{Fig/363.pdf}
	\end{minipage}
	
	\vspace{-14mm}
	\fuentepropia
\end{minipage}
\vskip0.2cm



Además, utilizando  la figura anterior, es pertinente considerar que, si en la ecuación, se sustituye en la otra ecuación, se obtiene lo siguiente:
\[z=4-y^2=4-(2-x^2)^2=4x^2-x^4.\]
De esta expresión se sigue que $\,\,2-x^2=\pm\sqrt{4-z},$ por lo tanto $2\pm \sqrt{4-z}=x^2$ y de ahí se obtiene las cuatro expresiones que determinan a $x,$ estas ecuaciones 
$ x = \pm \sqrt{2 \pm \sqrt{4-z}}.$
En la siguiente figura se ilustra esta situación.

\vskip0.2cm
\noindent
\begin{minipage}{\textwidth}
	\centering
	\captionof{figure}{Proyección del sólido en el plano $xz$.}
	\label{fig:proyeccion_xz}
	\begin{minipage}{0.60\linewidth}
		\centering
		\vspace{-43mm}
		\includegraphics[width=\linewidth]{Fig/364.pdf}
	\end{minipage}

	\vspace{-38mm}
	\fuentepropia
\end{minipage}
\vskip0.2cm

Ahora se presenta la integral triple faltante. Su valor fue necesario hallarlo de manera numérica utilizando un asistente computacional.

\[\int_{-\sqrt{2}}^{\sqrt{2}}\int_{4x^2-x^{4}}^{4}\int_{-\sqrt{4-z}}^{\sqrt{4-z}}dydzdx+\int_{-2}^2\int_{0}^{4x^2-x^{4}}\int_{-\sqrt{4-z}}^{2-x^2}dydzdx \approx 29.2571\]
\begin{multline*}
\int_{0}^{4}\int_{-\sqrt{2-\sqrt{4-z}}}^{\sqrt{2-\sqrt{4-z}}}\int_{-\sqrt{4-z}}^{\sqrt{4-z}}dydxdz+\int_{0}^{4}\int_{-\sqrt{2+\sqrt{4-z}}}^{-\sqrt{2-\sqrt{4-z}}}\int_{-\sqrt{4-z}}^{2-x^2}dydxdz\\
+\int_{0}^{4}\int_{\sqrt{2-\sqrt{4-z}}}^{\sqrt{2+\sqrt{4-z}}}\int_{-\sqrt{4-z}}^{2-x^2}dydxdz \approx 29.2571.
\end{multline*}

La evaluación numérica de las integrales anteriores se consigue ejecutando las siguientes instrucciones:

\vspace{0.6em}
\noindent
\begin{tcolorbox}[breakable,enhanced,title=\bf Instrucción en \verb"Mathematica", top=2pt,bottom=2pt,boxsep=0pt,before skip=2pt,after skip=0pt]
\begin{Verbatim}[formatcom=\letraverbatim]
NIntegrate[1,{x,-Sqrt[2],Sqrt[2]},{z,4*x^2-x^4,4},
{y,-Sqrt[4-z],Sqrt[4-z]}]+NIntegrate[1,{x,-2,2},
{z,0,4*x^2-x^4},{y,-Sqrt[4-z],2-x^2}]
NIntegrate[1,{z,0,4},{x,-Sqrt[2-Sqrt[4-z]],Sqrt[2-
Sqrt[4-z]]},{y,-Sqrt[4-z],Sqrt[4-z]}]+NIntegrate[1,
{z,0,4},{x,-Sqrt[2+Sqrt[4-z]],-Sqrt[2-Sqrt[4-z]]},
{y,-Sqrt[4-z],2-x^2}]+NIntegrate[1,{z,0,4},{x,Sqrt[2-
Sqrt[4-z]],Sqrt[2+Sqrt[4-z]]},{y,-Sqrt[4-z],2-x^2}]
\end{Verbatim}
\end{tcolorbox}
\vspace{0.7em}
\noindent


Otra manera de hallar el volumen del sólido es la siguiente.
%\vskip0.2cm

%\begin{tcolorbox}[title=\bfseries Parametrizando el sólido]
En el plano $xy$, la proyección del sólido se parametriza mediante con una combinación convexa de dos puntos
$A(t,-2)$ y $B(t,2-t^2)$; esto es $(1-u)A+uB=(t,(4-t^2)u-2).$
Con un punto $P\left(t,2-t^2-4u+ut^2,0\right) $ en la base del sólido y un punto $Q\left( t,(4-t^2)u-2,u(4-t^2)((t^2-4)u+4)\right)$
sobre el cilindro se forma la combinación $(1-w)P+wQ,$ que define la función
{\small\[{\bf r}(u,v,w) = \langle t,(4-t^2)u-2,-(t^2-4)uw((t^2-4)u+4)\rangle,\]}
que parametriza el sólido con $t\in [ -2,2]$, $u\in [0,1]$, $w\in [0,1].$

\vskip0.3cm
De esta función se obtienen los siguientes vectores
{\small
\begin{eqnarray*}
{\bf r}_{t}&=& \langle 1,2t(u-1),4wt(1-u)((u-1)t^2-4u+2) \rangle\\
{\bf r}_{u} &=& \langle 0,t^2-4,w(t^2-4) ((1-u)t^2+4u) +w(1-u)(t^2-4)^2 \rangle \\
{\bf r}_{w} &=& \langle 0,0,(u-1)( t^2-4) \left(t^2-ut^2+4u\right) \rangle.
\end{eqnarray*}
}
Con lo anterior, el elemento diferencial de volumen está dado por {\small \[\left| {\bf r}_{t}\cdot \left({\bf r}_{u}\times {\bf r}_{w}\right) \right|
=\left| \left( t^2-4\right) ^2\left( u-1\right) \left(t^2-ut^2+4u\right) \right|,\]}

entonces, el volumen del sólido esta dado por

\[\int_{0}^{1}\int_{0}^{1}\int_{-2}^2\left| \left( t^2-4\right)
^2\left( u-1\right) \left( t^2-ut^2+4u\right) \right|
dtdudw= \frac{1024}{35}.\]
%\end{tcolorbox}
\vskip0.2cm
Con las siguientes instrucciones, se realiza el procedimiento anterior:

\begin{tcolorbox}
\begin{Verbatim}[formatcom=\letraverbatim]
{x,y}=u*{t,-2}+(1-u)*{t,2-t^2};
r=Simplify[(1-w)*{x,y,0}+w*{x,y,4-y^2}]
A=Abs[Dot[D[r,t],Cross[D[r,u],D[r,w]]]];
Integrate[A,{t,-2,2},{u,0,1},{w,0,1}]
\end{Verbatim}
\end{tcolorbox}

\subsection*{Un cilindro parabólico y uno trigonométrico}

Hallar el volumen del sólido acotado por las figuras de las expresiones $z=1-x^2$ y $z=\cos(y)$, con
$y\in [0,2\pi]$. Estas superficies y el sólido que determinan se muestran enseguida.
\vskip0.2cm

\noindent
\begin{minipage}{\textwidth}
	\centering
	\captionof{figure}{$\cos y \leq z \leq 1 - x^2$, \quad $y \in [0, 2\pi]$.}
	\label{fig:region_cos_1mx2}
	\begin{minipage}{0.21\linewidth}
		\centering
		\includegraphics[width=\linewidth]{Img/84b.pdf}
	\end{minipage}
	\hspace{1em}
	\begin{minipage}{0.21\linewidth}
		\centering
		\includegraphics[width=\linewidth]{Img/84c.pdf}
	\end{minipage}
	\hspace{1em}
	\begin{minipage}{0.21\linewidth}
		\centering
		\includegraphics[width=\linewidth]{Img/84d.pdf}
	\end{minipage}
	
	\vspace{0.3em}
	\fuentepropia
\end{minipage}


\vskip0.2cm
Con las siguientes instrucciones se consiguen estas figuras:

\vspace{0.6em}
\noindent
\begin{tcolorbox}[breakable,enhanced,title=\bf Instrucción en \verb"Mathematica", top=2pt,bottom=2pt,boxsep=0pt,before skip=2pt,after skip=0pt]
\begin{Verbatim}[formatcom=\letraverbatim]
ContourPlot3D[{z==1-x^2,z==Cos[y]},{x,-Sqrt[2],
Sqrt[2]},{y,0,2*Pi},{z,-1,1},Mesh->False,
ContourStyle->{Orange,Green}]
RegionPlot3D[Cos[y]<z<1-x^2&&-2<x<2,{x,-Sqrt[2],
Sqrt[2]},{y,0,2*Pi},{z,-1,1},PlotPoints->30,
Mesh->False]
\end{Verbatim}
\end{tcolorbox}
\vspace{0.7em}
\noindent

\noindent
\begin{minipage}{\textwidth}
	\centering
	\captionof{figure}{Proyecciones del sólido en los planos coordenados.}
	\label{fig:proyecciones_coordenadas_2}
	\begin{minipage}{0.26\linewidth}
		\centering
		\vspace{-14mm}
		\includegraphics[width=\linewidth]{Fig/365.pdf}
	\end{minipage}
	\hspace{1em}
	\begin{minipage}{0.26\linewidth}
		\centering
		\vspace{-14mm}
		\includegraphics[width=\linewidth]{Fig/366.pdf}
	\end{minipage}
	\hspace{1em}
	\begin{minipage}{0.26\linewidth}
		\centering
		\vspace{-14mm}
		\includegraphics[width=\linewidth]{Fig/367.pdf}
	\end{minipage}
	
	\vspace{-8mm}
	\fuentepropia
\end{minipage}


Con estas figuras se plantean las integrales que determinan el volumen del sólido:

\vspace{-8mm}
\begin{multline*}
\int_{0}^{2\pi}\int_{-\sqrt{1-\cos y}}^{\sqrt{1-\cos y}}\int_{\cos y}^{1-x^2}dzdxdy= \int_{-\sqrt{2}}^{\sqrt{2}}\int_{\arccos \left( 1-x^2\right) }^{2\pi -\arccos \left( 1-x^2\right) }\int_{\cos y}^{1-x^2}dzdydx \\
\int_{-1}^{1}\int_{-\sqrt{1-z}}^{\sqrt{1-z}}\int_{\arccos z}^{2\pi -\arccos z}dydxdz= \int_{-\sqrt{2}}^{\sqrt{2}}\int_{-1}^{1-x^2}\int_{\arccos z}^{2\pi -\arccos z}dydzdx \\
\int_{0}^{2\pi }\int_{\cos y}^{1}\int_{-\sqrt{1-z}}^{\sqrt{1-z}}dxdzdy=\int_{-1}^{1}\int_{\arccos z}^{2\pi -\arccos z}\int_{-\sqrt{1-z}}^{\sqrt{1-z}}dxdydz = \frac{64\sqrt{2}}{9}.
\end{multline*}

Las anteriores integrales se evaluaron usando \verb"Mathematica";
el código requerido en tres de ellas, se presenta a continuación.
%\vskip0.2cm

\vspace{0.6em}
\noindent
\begin{tcolorbox}[breakable,enhanced,title=\bf Instrucción en \verb"Mathematica", top=2pt,bottom=2pt,boxsep=0pt,before skip=2pt,after skip=0pt]
\begin{Verbatim}[formatcom=\letraverbatim]
Integrate[1,{y,0,2*Pi},{x,-Sqrt[1-Cos[y]],
Sqrt[1-Cos[y]]},{z,Cos[y],1-x^2}]
Integrate[1,{z,-1,1},{x,-Sqrt[1-z],Sqrt[1-z]},
{y,ArcCos[z],2*Pi-ArcCos[z]}]
Integrate[1,{y,0,2*Pi},{z,Cos[y],1},{x,-Sqrt[1-z],
Sqrt[1-z]}]
\end{Verbatim}
\end{tcolorbox}
\vspace{0.7em}
\noindent

El volumen también se encuentra con el siguiente procedimiento.
\vskip0.2cm
% \begin{figure}[H]
% \centering
% \begin{pspicture}(-2,-0.6)(2,3.7) % \psgrid
% \psset{algebraic,xunit=0.9,yunit=0.4}
% \pscustom[fillstyle=solid,fillcolor=gray!40]{
%   \parametricplot{0}{\psPiTwo}{sqrt(1-cos(t)),t}
%   \parametricplot{\psPiTwo}{0}{-sqrt(1-cos(t)),t}
% }
% \psaxes[arrows=->,showorigin=false,dx=1.4,xLabels={-\sqrt{2},,,\sqrt{2}},yunit=\psPi,ytrigLabels,trigLabelBase=1](0,0)(-2,-0.3)(2.35,2.6)[$x$,90][$y$,180]
% \parametricplot[arrows=->,linewidth=1.2pt]{-0.3}{7}{sqrt(1-cos(t)),t}
% \parametricplot[arrows=->,linewidth=1.2pt]{-0.3}{7}{-sqrt(1-cos(t)),t}
% \rput{0}(-1.31,2.36){$\bullet$}
% \rput{0}(1.31,2.36){$\bullet$}
% \psline[linewidth=1pt,linecolor=black](-1.31,2.4)(1.31,2.4)
% \rput{0}(-1.65,2.5){$A$}
% \rput{0}(1.65,2.5){$B$}
% \end{pspicture}
% \caption{$\cos(y) = 1 - x^2$.}
% \label{fig:curva_cosyx2}
% \end{figure}
% \par
% \justifying
%\begin{tcolorbox}[title=\bfseries Parametrizando el sólido]
\noindent 
Con dos puntos $A\left( -\sqrt{1-\cos t},t\right) $
y $B\left( \sqrt{1-\cos t},t\right)$ en la proyección del sólido en el plano $xy$,
se forma la expresión
$(1-u)A+uB,$ que parametriza la región mencionada. % esto es $ \left((1-2u)\sqrt{1-\cos t},t\right).$
Sobre las dos superficies que determinan el sólido se satisface que
$z= 4u\left( 1-u\right) +\left( 1-2u\right)^2\cos t$ y $z=\cos (t).$


\vskip0.2cm
\noindent
\begin{minipage}{\textwidth}
	\centering
	\captionof{figure}{$\cos(y) = 1 - x^2$}
	\label{fig:curva_cosyx2}
	\begin{minipage}{0.25\linewidth}
		\centering
		\vspace{-8mm}
		\includegraphics[width=\linewidth]{Fig/368.pdf}
	\end{minipage}
	
	\vspace{-4mm}
	\fuentepropia
\end{minipage}
\vskip0.2cm


El sólido se parametriza mediante una función $\mathbf{r}$, que es combinación convexa con dos puntos de la forma $P\left((2u - 1)\sqrt{1 - \cos t}, t, \cos t\right)$ y $Q\left((1 - 2u)\sqrt{1 - \cos t}, 4u(1 - u) + (1 - 2u)^2 \cos t\right)$. Esto es
%\vskip0.2cm
{\small
\[
\begin{aligned}
	\mathbf{r}(t, u, w) = \left\langle (1 - 2w)(2u - 1)\sqrt{1 - \cos t}, t, 4uw(1 - u) + \right. \\
	\left. (1 - 4wu + 4wu^2) \cos t \right\rangle.
\end{aligned}
\]
}

%\vskip0.2cm
Con esta función hallamos los siguientes vectores
{\small
\begin{eqnarray*}
{\bf r}_{t} &=& \left\langle \frac{1}{2}\frac{\left(
1-2w\right) \left( 2u-1\right) }{\sqrt{1-\cos t}}\sin t,1,\left(
4wu-4wu^2-1\right) \sin t\right\rangle \\
{\bf r}_{u}&=& \left\langle 2\left( 1-2w\right) \sqrt{1-\cos
t},0,4w\left( 1-2u\right) \left( 1-\cos t\right) \right\rangle \\
{\bf r}_{w}&=& \left\langle 2\left( 1-2u\right) \sqrt{1-\cos t},0,4u\left( u-1\right) \left( \cos t-1\right) \right\rangle
\end{eqnarray*}
}

Además, $\left| {\bf r}_{w}\cdot ( {\bf r}_{t}\times {\bf r}_{u}) \right| =8(w+u(2w+1)(u-1))\sqrt{\left( 1-\cos t\right) ^{3}};$
el volu\-men del sólido esta dado por
\vspace{-2mm}
{\small
\[\int_{0}^{1}\int_{0}^{1}\int_{0}^{2\pi }8(w+u(2w+1)(u-1)) \sqrt{\left( 1-\cos t\right) ^{3}}%
dtdudw= \frac{64}{9}\sqrt{2}\]
}
%\end{tcolorbox}
%\vskip0.2cm
La rutina descrita antes se puede hacer en \verb"Mathematica" con ayuda de las siguientes instrucciones.

\vspace{0.6em}
\noindent
\begin{tcolorbox}[breakable,enhanced,title=\bf Instrucción en \verb"Mathematica", top=2pt,bottom=2pt,boxsep=0pt,before skip=2pt,after skip=0pt]
\begin{Verbatim}[formatcom=\letraverbatim]
{x,y,z}={(2*u-1)*Sqrt[1-Cos[t]],t,1-x^2};
r=(1-w)*{x,y,z}+w*{x,y,Cos[t]};
A=Dot[D[r,t],Cross[D[r,u],D[r,w]]];
Integrate[A,{t,0,2*Pi},{u,0,1},{w,0,1}]
\end{Verbatim}
\end{tcolorbox}
\vspace{0.7em}
\noindent


\subsection*{Dos cilindros parabólicos y un plano}
Los cilindros parabólicos con ecuaciones
$z= 4x-x^2,$ $z= 4y-y^2$ y el plano $z=0,$ acotan un sólido en el primer octante, su gráfico se muestra a continuación.

\vskip0.2cm
\noindent
\begin{minipage}{\textwidth}
	\centering
	\captionof{figure}{$0 \leq z \leq 4x - x^2$, \quad $0 \leq z \leq 4y - y^2$.}
	\label{fig:parabolas_z_2}
	\begin{minipage}{0.22\linewidth}
		\centering
		\includegraphics[width=\linewidth]{Img/10b.pdf}
	\end{minipage}
	\hspace{1em}
	\begin{minipage}{0.22\linewidth}
		\centering
		\includegraphics[width=\linewidth]{Img/10a.pdf}
	\end{minipage}
	
	\vspace{0.5mm}
	\fuentepropia
\end{minipage}
\vskip0.2cm

Las instrucciones para obtener los gráficos del sólido que se consideran en este ejemplo, son las siguientes:

\vspace{0.6em}
\noindent
\begin{tcolorbox}[breakable,enhanced,title=\bf Instrucción en \verb"Mathematica", top=2pt,bottom=2pt,boxsep=0pt,before skip=2pt,after skip=0pt]
\begin{Verbatim}[formatcom=\letraverbatim]
ContourPlot3D[{z==4*x-x^2,z==4*y-y^2},{x,0,4},{y,0,4},
{z,0,4},PlotPoints->80,ContourStyle->{Yellow,Green}]
RegionPlot3D[0<z<4*x-x^2&&0<z<4*y-y^2,{x,0,4},{y,0,4},
{z,0,4},PlotStyle->Blend[{Red,Orange,Gray}],
PlotPoints->90,Mesh->False]
\end{Verbatim}
\end{tcolorbox}
\vspace{0.7em}
\noindent


En la base de la región, es decir en el plano $xy$, se satisface que
$4x-x^2=4y-y^2$ de donde $(x-2)^2=(y-2)^2$, con lo cual $y-2= \pm(x-2)$, esto nos conduce a que $y=x$ o $y=4-x$. Esto divide la base en cuatro partes, por tal razón al integrarse en el orden $dzdydx$
se puede plantear una integral triple y multiplicarla por cuatro.

\vskip0.2cm
\noindent
\begin{minipage}{\textwidth}
	\centering
	\captionof{figure}{Proyección en el plano $xy$ y un cuarto del sólido.}
	\label{fig:proyeccion_xy_cuarto}
	\begin{minipage}{0.35\linewidth}
		\centering
		\vspace{-10mm}
		\includegraphics[width=\linewidth]{Fig/369.pdf}
	\end{minipage}
	\hspace{1em}
	\begin{minipage}{0.25\linewidth}
		\centering
		\vspace{-10mm}
		\includegraphics[width=\linewidth]{Img/10c.pdf}
	\end{minipage}
	
	\vspace{-4mm}
	\fuentepropia
\end{minipage}


\vskip0.2cm
La siguiente instrucción permite obtener la porción señalada:

\vspace{0.6em}
\noindent
\begin{tcolorbox}[breakable,enhanced,title=\bf Instrucción en \verb"Mathematica", top=2pt,bottom=2pt,boxsep=0pt,before skip=2pt,after skip=0pt]
\begin{Verbatim}[formatcom=\letraverbatim]
RegionPlot3D[0<z<2*x-x^2&&0<z<2*y-y^2&&2-x<y<x&&1<x<2,
{x,0,2},{y,0,2},{z,0,1},PlotPoints->90,Mesh->False,
PlotStyle->Orange]
\end{Verbatim}
\end{tcolorbox}
\vspace{0.7em}
\noindent


Cualquiera de las siguientes posibilidades permite calcular el volumen de la figura en mención.

\begin{itemize}
\item $4\ds \int_{2}^{4}\int_{4-x}^{x}\int_{0}^{4x-x^2}dzdydx=32$
\item $2\left( \ds \int_{0}^2\int_{x}^{4-x}\int_{0}^{4x-x^2}dzdydx+\int_{2}^{4}\int_{4-x}^{x}\int_{0}^{4x-x^2}dzdydx\right) =32$
\item $2\left( \ds \int_{0}^2\int_{y}^{4-y}\int_{0}^{4y-y^2}dzdxdy+\int_{2}^{4}\int_{4-y}^{y}\int_{0}^{4y-y^2}dzdxdy\right) =32$
\end{itemize}
\vskip0.2cm

\noindent
\begin{minipage}{\textwidth}
	\centering
	\captionof{figure}{Proyecciones en los planos $xz$ y $yz$.}
	\label{fig:proyecciones_xz_yz}
	\begin{minipage}{0.30\linewidth}
		\centering
		\vspace{-10mm}
		\includegraphics[width=\linewidth]{Fig/370.pdf}
	\end{minipage}
	\hspace{1em}
	\begin{minipage}{0.30\linewidth}
		\centering
		\vspace{-10mm}
		\includegraphics[width=\linewidth]{Fig/371.pdf}
	\end{minipage}
	
	\vspace{-8mm}
	\fuentepropia
\end{minipage}
\vskip0.2cm


Estas figuras permiten el planteamiento de las integrales triples en los órdenes restantes.
\[\int_{0}^{4}\int_{2-\sqrt{4-z}}^{2+\sqrt{4-z}}\int_{2-\sqrt{4-z}}^{2+\sqrt{4-z}}dxdydz= \int_{0}^{4}\int_{0}^{4y-y^2}\int_{2-\sqrt{4-z}}^{2+\sqrt{4-z}}dxdzdy= 32\]
\[\int_{0}^{4}\int_{2-\sqrt{4-z}}^{2+\sqrt{4-z}}\int_{2-\sqrt{4-z}}^{2+\sqrt{4-z}}dydxdz=\int_{0}^{4}\int_{0}^{4x-x^2}\int_{2-\sqrt{4-z}}^{2+\sqrt{4-z}}dydzdx=32\]
\vskip0.2cm
Algunos de estos cálculos pueden hacerse con las siguientes instrucciones:

\vspace{0.6em}
\noindent
\begin{tcolorbox}[breakable,enhanced,title=\bf Instrucción en \verb"Mathematica", top=2pt,bottom=2pt,boxsep=0pt,before skip=2pt,after skip=0pt]
\begin{Verbatim}[formatcom=\letraverbatim]
4*Integrate[1,{x,2,4},{y,4-x,x},{z,0,4*x-x^2}]
Integrate[1,{z,0,4},{x,2-Sqrt[4-z],2+Sqrt[4-z]},
{y,2-Sqrt[4-z],2+Sqrt[4-z]}]
\end{Verbatim}
\end{tcolorbox}
\vspace{0.7em}
\noindent


Este valor también puede obtenerse con el siguiente proceso.


\vskip0.2cm
%\begin{tcolorbox}[title=\bfseries Parametrizando el sólido]
\noindent

En el plano $xz$ se forma una combinación convexa entre dos puntos de la forma
$P(2-\sqrt{4-t}, t)$ y $Q(2 + \sqrt{4-t}, t)$, obteniendo $(2+(2u-1)\sqrt{4-t}, t)$ que parametriza la porción interior a la parábola sobre el eje $z$, con $t \in[0,4],$ $u \in[0,1].$


\vskip0.2cm
\noindent
\begin{minipage}{\textwidth}
	\centering
	\captionof{figure}{$0 \leq y \leq 4x - x^2$}
	\label{fig:region_y_vs_x}
	\begin{minipage}{0.32\linewidth}
		\centering
		\vspace{-15mm}
		\includegraphics[width=\linewidth]{Fig/372.pdf}
	\end{minipage}
	
	\vspace{-10mm}
	\fuentepropia
\end{minipage}
\vskip0.2cm



El sólido se parametriza simplificando una combinación convexa entre dos puntos de la forma
$(2+(2u-1)\sqrt{4-t},2-\sqrt{4-t},t)$ y $(2+(2u-1)\sqrt{4-t},2+\sqrt{4-t},t).$
Por tanto, la parametrización escogida para el sólido es
${\bf r}(u,t,w)= \langle 2+(2u-1)\sqrt{4-t},2+(2w-1)\sqrt{4-t},t \rangle, $
en donde $t\in [0,4],$ $u\in [0,1]$ y $w\in [0,1].$
De esta función se obtiene los siguientes vectores:
{\small
\begin{align*}
	\mathbf{r}_u &= \left\langle 2\sqrt{4-t},0,0 \right\rangle, \\
	\mathbf{r}_w &= \left\langle 0,2\sqrt{4-t},0 \right\rangle, \\
	\mathbf{r}_t &= \left\langle \frac{1-2u}{2\sqrt{4-t}}, \frac{1-2w}{2\sqrt{4-t}},1 \right\rangle.
\end{align*}
}
Entonces, ${\bf r}_w \cdot ({\bf r}_t \times {\bf r}_u)= 4(4-t)$, el volumen del sólido está dado por
%\vskip0.2cm
{\small
\[\int_0^1 \int_0^1 \int_0^4 4(4-t)dtdudw = 32.\]}
%\end{tcolorbox}
\vskip0.2cm
La rutina anteriormente descrita se puede obtener en \verb"Mathematica", ejecutando las siguientes instrucciones:

\vspace{0.6em}
\noindent
\begin{tcolorbox}[breakable,enhanced,title=\bf Instrucción en \verb"Mathematica", top=2pt,bottom=2pt,boxsep=0pt,before skip=2pt,after skip=0pt]
\begin{Verbatim}[formatcom=\letraverbatim]
r=Factor[(1-w){2+(2*u-1)*Sqrt[4-t],2-Sqrt[4-t],t}+
w{2+(2*u-1)*Sqrt[4-t],2+Sqrt[4-t],t}]
Integrate[Dot[D[r,w],Cross[D[r,t],D[r,u]]],{u,0,1},
{t,0,4},{w,0,1}]
\end{Verbatim}
\end{tcolorbox}
\vspace{0.7em}
\noindent


\subsection*{Un cilindro parabólico y un plano inclinado}
Consideremos el sólido acotado por el gráfico de las ecuaciones $y=4-x^2,$ $z=2-y/2,$ $y=0,$ $z=0$,
hallaremos su volumen de varias maneras.
\vskip0.2cm
\noindent
\begin{minipage}{\textwidth}
	\centering
	\captionof{figure}{$0 \leq z \leq 4x - x^2$, \quad $0 \leq z \leq 4y - y^2$.}
	\label{fig:parabolas_z}
	\begin{minipage}{0.25\linewidth}
		\centering
		\vspace{-15mm}
		\includegraphics[width=\linewidth]{Img/T_0a.pdf}
	\end{minipage}
	\hspace{1em}
	\begin{minipage}{0.35\linewidth}
		\centering
		\vspace{-15mm}
		\includegraphics[width=\linewidth]{Fig/373.pdf}
	\end{minipage}
	
	\vspace{-10mm}
	\fuentepropia
\end{minipage}
\vskip0.2cm


Esta región nos lleva a presentar las siguientes figuras.
\vskip0.2cm
\noindent
\begin{minipage}{\textwidth}
	\centering
	\captionof{figure}{Proyecciones en los planos coordenados.}
	\label{fig:proyecciones_coordenadas}
	\begin{minipage}{0.29\linewidth}
		\centering
		\vspace{-10mm}
		\includegraphics[width=\linewidth]{Fig/374.pdf}
	\end{minipage}
	\hspace{1em}
	\begin{minipage}{0.29\linewidth}
		\centering
		\vspace{-10mm}
		\includegraphics[width=\linewidth]{Fig/375.pdf}
	\end{minipage}
	\hspace{1em}
	\begin{minipage}{0.29\linewidth}
		\centering
		\vspace{-10mm}
		\includegraphics[width=\linewidth]{Fig/376.pdf}
	\end{minipage}
	
	\vspace{-5mm}
	\fuentepropia
\end{minipage}
\vskip0.2cm

Las integrales que se evalúan para determinar el volumen del sólido son:
\[\int_{-2}^2\int_{0}^{4-x^2}\int_{0}^{2-y/2}dzdydx = \int_{0}^{4}\int_{-\sqrt{4-y}}^{\sqrt{4-y}}\int_{0}^{2-y/2}dzdxdy= \frac{64}{5}\]
\[\int_{0}^{4}\int_{0}^{2-y/2}\int_{-\sqrt{4-y}}^{\sqrt{4-y}}dxdzdy= \int_{0}^2\int_{0}^{4-2z}\int_{-\sqrt{4-y}}^{\sqrt{4-y}}dxdydz= \frac{64}{5}.\]
\vskip0.2cm
La proyección del sólido en el plano $xz$ es un rectángulo dividido por la curva
$z=x^2/2$. En el orden $dydxdz$ es necesario plantear tres integrales triples
\begin{multline*} \int_{0}^2\int_{-\sqrt{2z}}^{\sqrt{2z}}\int_{0}^{4-2z}dydxdz
+\int_{0}^2\int_{-2}^{-\sqrt{2z}}\int_{0}^{4-x^2}dydxdz \\ +\int_{0}^2\int_{\sqrt{2z}}^2\int_{0}^{4-x^2}dydxdz= \frac{64}{5} 
\end{multline*}
\[\int_{-2}^2\int_{x^2/2}^2\int_{0}^{4-2z}dydzdx +\int_{-2}^2\int_0^{x^2/2}\int_{0}^{4-x^2}dydzdx=\frac{64}{5}.\]
Estas integrales se evalúan en \verb"Mathematica" con las siguientes indicaciones:

\vspace{0.6em}
\noindent
\begin{tcolorbox}[breakable,enhanced,title=\bf Instrucción en \verb"Mathematica", top=2pt,bottom=2pt,boxsep=0pt,before skip=2pt,after skip=0pt]
\begin{Verbatim}[formatcom=\letraverbatim]
Integrate[1,{z,0,2},{y,0,4-2*z},{x,-Sqrt[4-y],
Sqrt[4-y]}]
Integrate[1,{x,-2,2},{z,x^2/2,2},{y,0,4-2*z}]+
Integrate[1,{x,-2,2},{z,0,x^2/2},{y,0,4-x^2}]
\end{Verbatim}
\end{tcolorbox}
\vspace{0.7em}
\noindent


El valor obtenido puede ser calculado con el siguiente método.
\vskip0.2cm

%\begin{tcolorbox}[title=\bfseries Parametrizando el sólido]
La región sobre el plano $xy$ se parametriza uniendo dos puntos, uno sobre el eje $x$ que tiene la forma $(t,0) $ y el otro sobre la curva $y=4-x^2$ que posee la forma $(t,4-t^2),$ con $t\in [-2,2].$
Simplificando $( 1-u) ( t,0) +u( t,4-t^2)$ se sigue, por un lado, que $z=2+y/2= 2-2u+\frac{1}{2}ut^2.$

Por tanto, la función ${\bf r}$ es una combinación convexa entre un punto sobre el plano $xy$
y otros sobre el plano $z=2-y/2;$ simplificando
$(1-w) \langle t,u( 4-t^2),0\rangle +w\langle t,u( 4-t^2) ,2-2u+\frac{1}{2}ut^2\rangle$
se define
$ {\bf r}( t,u,w) = \langle t,4u-ut^2,\frac{1}{2}w( 4-4u+ut^2)\rangle ,$
con $t \in [-2,2], \, u \in [0,1],\, w \in [0,1].$ Además,
{\footnotesize
\begin{eqnarray*} {\bf r}_{t} &=& \langle 1,-2tu,wtu\rangle \\ {\bf r}_{u}&=& \langle 0,4-t^2,w(t^2-4)/2 \rangle \\
{\bf r}_{w}&=& \langle 0,0,2-2u+ut^2/2 \rangle. \end{eqnarray*}
}
%\vskip0.2cm

El triple producto vectorial
$|{\bf r}_t\cdot ( {\bf r}_u \times {\bf r}_w)|=\frac{1}{2}| ( 4-t^2) ( 4-4u+ut^2) | $
se integra para obtener el volumen del sólido
%\vskip0.2cm
{\footnotesize
\[\int_{0}^{1}\int_{0}^{1}\int_{-2}^2\frac{1}{2}\left|(4-t^2)(4-4u+ut^2)\right|dtdudw= \frac{64}{5}\]
}
%\end{tcolorbox}


\subsection*{Un cilindro parabólico y un plano inclinado II}


Hallar el volumen del sólido acotado por
las figuras de las ecuaciones $y=0,$ $y=2-z,$ $z=2-x^2,$ $z=0$. Esta región del espacio se muestra a continuación.
\vskip0.2cm

\noindent
\begin{minipage}{\textwidth}
	\centering
	\captionof{figure}{Visualización tridimensional del sólido limitado por $z=2-x^2$ y $z=2-y$.}
	\label{fig:solido_3D_parabolico}
	\begin{minipage}{0.35\linewidth}
		\centering
		\vspace{-15mm}
		\includegraphics[width=\linewidth]{Fig/378.pdf}
	\end{minipage}
	\hspace{1em}
	\begin{minipage}{0.35\linewidth}
		\centering
		\vspace{-15mm}
		\includegraphics[width=\linewidth]{Fig/379.pdf}
	\end{minipage}
	
	\vspace{-8mm}
	\fuentepropia
\end{minipage}
\vskip0.2cm


Ahora se proyecta el sólido sobre los planos coordenados.
\vskip0.2cm

\noindent
\begin{minipage}{\textwidth}
	\centering
	\captionof{figure}{Proyecciones del sólido en los planos coordenados.}
	\label{fig:proyecciones_sólido}
	\begin{minipage}{0.25\linewidth}
		\centering
		\vspace{-12mm}
		\includegraphics[width=\linewidth]{Fig/380.pdf}
	\end{minipage}
	\hspace{1em}
	\begin{minipage}{0.25\linewidth}
		\centering
		\vspace{-12mm}
		\includegraphics[width=\linewidth]{Fig/381.pdf}
	\end{minipage}
	\hspace{1em}
	\begin{minipage}{0.25\linewidth}
		\centering
		\vspace{-12mm}
		\includegraphics[width=\linewidth]{Fig/382.pdf}
	\end{minipage}
	
	\vspace{-8mm}
	\fuentepropia
\end{minipage}
\vskip0.2cm


Con estos elementos, el volumen en mención está dado por las siguientes integrales triples en los seis órdenes posibles y utilizando coordenadas cartesianas.
\vspace{-1mm}
\begin{eqnarray*}
\int_{-\sqrt{2}}^{\sqrt{2}}\int_{0}^{2-x^2}\int_{0}^{2-z}dydzdx=
\int_{0}^2\int_{-\sqrt{2-z}}^{\sqrt{2-z}}\int_{0}^{2-z}dydxdz&=& \frac{16}{5}\sqrt{2}\\
\int_{0}^2\int_{0}^{2-z}\int_{-\sqrt{2-z}}^{\sqrt{2-z}}dxdydz=
\int_{0}^2\int_{0}^{2-y}\int_{-\sqrt{2-z}}^{\sqrt{2-z}}dxdzdy&=& \frac{16}{5}\sqrt{2}
\end{eqnarray*}

Al integrar primero con respecto a $z$ se debe tener en cuenta que en el plano $xy$ la región de integración está
dividida en dos partes y dependiendo de la forma como se recorra tendremos dos o tres integrales triples.
\[\int_{-\sqrt{2}}^{\sqrt{2}}\int_{0}^{x^2}\int_{0}^{2-x^2}dzdydx + \int_{-\sqrt{2}}^{\sqrt{2}}\int_{x^2}^2\int_{0}^{2-y}dzdydx= \frac{16}{5}\sqrt{2}\]
\begin{multline*} \int_{0}^2\int_{-\sqrt{2}}^{-\sqrt{y}}\int_{0}^{2-x^2}dzdxdy+\int_{0}^2\int_{\sqrt{y}}^{\sqrt{2}}\int_{0}^{2-x^2}dzdxdy \\
+\int_{0}^2\int_{-\sqrt{y}}^{\sqrt{y}}\int_{0}^{2-y}dzdxdy= \frac{16}{5}%
\sqrt{2}.
\end{multline*}
Las siguientes son las indicaciones a usar en \verb"Mathematica" para evaluar estas últimas integrales.

\vspace{0.6em}
\noindent
\begin{tcolorbox}[breakable,enhanced,title=\bf Instrucción en \verb"Mathematica", top=2pt,bottom=2pt,boxsep=0pt,before skip=2pt,after skip=0pt]
\begin{Verbatim}[formatcom=\letraverbatim]
Integrate[1,{x,-Sqrt[2],Sqrt[2]},{y,0,x^2},{z,0,2-x^2}]
+Integrate[1,{x,-Sqrt[2],Sqrt[2]},{y,x^2,2},{z,0,2-y}]
Integrate[1,{y,0,2},{x,-Sqrt[2],-Sqrt[y]},{z,0,2-x^2}]
+Integrate[1,{y,0,2},{x,Sqrt[y],Sqrt[2]},{z,0,2-x^2}]+
Integrate[1,{y,0,2},{x,-Sqrt[y],Sqrt[y]},{z,0,2-y}]
\end{Verbatim}
\end{tcolorbox}
\vspace{0.7em}
\noindent


Con el siguiente procedimiento se muestra este cálculo de otra manera.

%\begin{tcolorbox}[title=\bfseries Parametrizando el sólido]
La proyección en el plano $xz$ se parametriza con $x=t,$
$z=( 2-t^2) u,$ por tanto un punto del sólido en el plano
$xz$ es de la forma $( t,0,2-t^2) $ y un punto sobre el plano $y=2-z$ es de la forma $( t,2-(2-t^2) u,(2-t^2) u).$
Simplificando la expresión $(1-w)\langle t,0,(2-t^2) u\rangle +w\langle t,2-( 2-t^2)u,( 2-t^2) u\rangle$
se parametriza el sólido, esto es
%\vskip0.2cm
\[{\bf r}( t,u,w) =\langle t,w(2-2u+ut^2) ,2u-ut^2\rangle,\]
\vskip0.2cm
con $t\in [ -\sqrt{2},\sqrt{2}],$ $u\in [0,1],$ $w\in
\lbrack 0,1].$ Derivando parcialmente se obtiene
\[{\bf r}_{t}= \langle 1,2wtu,-2tu\rangle,\] \[{\bf r}_{u}= \langle 0,w( -2+t^2),2-t^2\rangle,\] 
\[{\bf r}_{w}= \langle 0,2-2u+ut^2,0\rangle,\]
\[|{\bf r}_{t}\cdot({\bf r}_{u}\times {\bf r}_{w})|=|(2-t^2)(2u-ut^2-2)|.\]
Este último valor se integra y proporciona el volumen del sólido
\[\int_{0}^{1}\int_{0}^{1}\int_{-\sqrt{2}}^{\sqrt{2}}| ( 2-t^2) ( 2u-ut^2-2) |dtdudw= \frac{16}{5}\sqrt{2}.\]
%\end{tcolorbox}

El cálculo anterior se hace en \verb"Mathematica" utilizando las siguientes sentencias:

\vspace{0.6em}
\noindent
\begin{tcolorbox}[breakable,enhanced,title=\bf Instrucción en \verb"Mathematica", top=2pt,bottom=2pt,boxsep=0pt,before skip=2pt,after skip=0pt]
\begin{Verbatim}[formatcom=\letraverbatim]
r=(1-w)*{t,0,u*(2-t^2)}+w*{t,2-u*(2-t^2),u*(2-t^2)};
A=Simplify[Dot[D[r,t],Cross[D[r,u],D[r,w]]]]
Integrate[Abs[A],{t,-Sqrt[2],Sqrt[2]},{u,0,1},
{w,0,1}]
\end{Verbatim}
\end{tcolorbox}
\vspace{0.7em}
\noindent


\subsection*{Tres cilindros parabólicos}

Las expresiones $y=x^2$, $y=8-x^2$ y $z=4-x^2$ representan en el espacio tres cilindros parabólicos. Estas superficies y el plano $xy$ determinan un sólido que a continuación se muestra.



\begin{figure}[H]
\centering
\caption{$0 \leq z \leq 4 - x^2$, \quad $x^2 \leq y \leq 8 - x^2$.}
\label{fig:solido_parabolico}
\vspace{-3mm}
\includegraphics[height=4.5cm]{Img/25a.pdf} \hspace{0.2cm}
\includegraphics[height=4.5cm]{Img/25b.pdf} \hspace{0.2cm}
\includegraphics[height=4.5cm]{Img/25c.pdf}
\fuentepropia
\end{figure}
%\vskip0.2cm

Estas figuras se pueden generar usando las siguientes instrucciones:

\vspace{0.6em}
\noindent
\begin{tcolorbox}[breakable,enhanced,title=\bf Instrucción en \verb"Mathematica", top=2pt,bottom=2pt,boxsep=0pt,before skip=2pt,after skip=0pt]
\begin{Verbatim}[formatcom=\letraverbatim]
ContourPlot3D[{y==x^2,y==8-x^2,z==4-x^2},{x,-2,2},
{y,0,8},{z,0,4},PlotPoints->80]
RegionPlot3D[0<z<4-x^2&&x^2<y<8-x^2,{x,-2,2},{y,0,8},
{z,0,4},PlotPoints->70,Mesh->False]
\end{Verbatim}
\end{tcolorbox}
\vspace{0.7em}
\noindent


Las integrales sobre el sólido se plantean a partir de las siguientes figuras:
\vskip0.2cm

\noindent
\begin{minipage}{\textwidth}
	\centering
	\captionof{figure}{Proyección del sólido en los planos $xy$ y $xz$.}
	\label{fig:proyeccion_xy_xz_2}
	\begin{minipage}{0.36\linewidth}
		\centering
		\vspace{-15mm}
		\includegraphics[width=\linewidth]{Fig/383.pdf}
	\end{minipage}
	\hspace{1em} 
	\begin{minipage}{0.36\linewidth}
		\centering
		\vspace{-15mm}
		\includegraphics[width=\linewidth]{Fig/384.pdf}
	\end{minipage}
	
	\vspace{-10mm}
	\fuentepropia
\end{minipage}
\vskip0.2cm

Esta representación nos ayuda a determinar las integrales triples que determinan su volumen.
Al integrar primera con respecto a $z$, esta recorre los puntos desde el plano $z=0$
hasta el cilindro $z=4-x^2$ y los límites de integración para $x$ e $y$ se obtienen del primero de los gráficos anteriores. Entonces, el volumen del sólido está dado por:
{\small
\[\int_{-2}^2\int_{x^2}^{8-x^2}\int_{0}^{4-x^2}dzdydx= \frac{1024}{15}\]
\[\int_{0}^{4}\int_{-\sqrt{y}}^{\sqrt{y}}\int_{0}^{4-x^2}dzdxdy+
\int_{4}^{8}\int_{-\sqrt{8-y}}^{\sqrt{8-y}}\int_{0}^{4-x^2}dzdxdy= \frac{1024}{15}.\]
}
Si se integra con respecto a $y$ primero, se debe tener en cuenta que esta recorre los puntos desde el cilindro $y=x^2$ hasta el cilindro $y=8-x^2$, que equivalen respectivamente a
$x= \pm \sqrt{y}$ y $x= \pm \sqrt{8-z}$. Los límites de integración para $x$ y $z$ se obtienen del segundo de los gráficos anteriores.
\vspace{-1mm}
{\small
\[\int_{-2}^2\int_{0}^{4-x^2}\int_{x^2}^{8-x^2}dydzdx=\int_{0}^{4}\int_{-\sqrt{4-z}}^{\sqrt{4-z}}\int_{x^2}^{8-x^2}dydxdz= \frac{1024}{15}.\]
}


En el plano $yz$ el sólido se proyecta como un rectángulo. Al sustituir en $z=4-x^2$
las dos expresiones que involucran $x^2$, se sigue $z=4-x^2=4-y,$ por una parte, y $z=4-(8-y)=y-4,$
por la parte restante.


\vskip0.2cm
\noindent
\begin{minipage}{\textwidth}
	\centering
	\captionof{figure}{Proyección en el plano $yz$}
	\label{fig:proyeccion_yz_2}
	\begin{minipage}{0.40\linewidth}
		\centering
		\vspace{-23mm}
		\includegraphics[width=\linewidth]{Fig/385.pdf}
	\end{minipage}
	
	\vspace{-17mm}
	\fuentepropia
\end{minipage}
\vskip0.2cm



En este caso, el volumen del sólido está determinado por:
{\small
\begin{multline*}
\int_{0}^{4}\int_{0}^{4-z}\int_{-\sqrt{y}}^{\sqrt{y}}dxdydz+
\int_{0}^{4}\int_{4-z}^{z+4}\int_{-\sqrt{4-z}}^{\sqrt{4-z}}dxdydz\\
+\int_{0}^{4}\int_{z+4}^{8}\int_{-\sqrt{8-y}}^{\sqrt{8-y}}dxdydz= \frac{1024}{15}
\end{multline*}
}
y en el orden $dxdzdy,$ el planteamiento es el siguiente:
\vspace{-2mm}
{\small
\begin{multline*}
\int_{0}^{4}\int_{0}^{4-y}\int_{-\sqrt{y}}^{\sqrt{y}}dxdzdy+
\int_{0}^{4}\int_{4-y}^{4}\int_{-\sqrt{4-z}}^{\sqrt{4-z}}dxdzdy+ \\
\int_{4}^{8}\int_{y-4}^{4}\int_{-\sqrt{4-z}}^{\sqrt{4-z}}dxdzdy+\int_{4}^{8}%
\int_{0}^{y-4}\int_{-\sqrt{8-y}}^{\sqrt{8-y}}dxdzdy=\frac{1024}{15}.
\end{multline*}
}
\vskip0.2cm
Algunas de estas integrales se evalúan en \verb"Mathematica" mediante la ejecución de las siguientes instrucciones:


\vspace{0.6em}
\noindent
\begin{tcolorbox}[breakable,enhanced,title=\bf Instrucción en \verb"Mathematica", top=2pt,bottom=2pt,boxsep=0pt,before skip=2pt,after skip=0pt]
\begin{Verbatim}[formatcom=\letraverbatim]
Integrate[1,{x,-2,2},{z,0,4-x^2},{y,x^2,8-x^2}]
Integrate[1,{z,0,4},{x,-Sqrt[4-z],Sqrt[4-z]},
{y,x^2,8-x^2}]
Integrate[1,{y,0,4},{x,-Sqrt[y],Sqrt[y]},
{z,0,4-x^2}]+Integrate[1,{y,4,8},{x,-Sqrt[8-y],
Sqrt[8-y]},{z,0,4-x^2}]
Integrate[1,{z,0,4},{y,0,4-z},{x,-Sqrt[y],Sqrt[y]}]+
Integrate[1,{z,0,4},{y,4-z,z+4},{x,-Sqrt[4-z],
Sqrt[4-z]}]+Integrate[1,{z,0,4},{y,z+4,8},
{x,-Sqrt[8-y],Sqrt[8-y]}]
\end{Verbatim}
\end{tcolorbox}
\vspace{0.7em}
\noindent


%Integrate[1,{y,0,4},{z,0,4-y},{x,-Sqrt[y],Sqrt[y]}]+
%Integrate[1,{y,0,4},{z,4-y,4},{x,-Sqrt[4-z],Sqrt[4-z]}]+
%Integrate[1,{y,4,8},{z,y-4,4},{x,-Sqrt[4-z],Sqrt[4-z]}]+
%Integrate[1,{y,4,8},{z,0,y-4},{x,-Sqrt[8-y],Sqrt[8-y]}]

Otra manera de obtener el mismo valor es la descrita a continuación.


%\begin{tcolorbox}[title=\bfseries Parametrizando el sólido]
\noindent

En el plano $xy$, con dos puntos $P(t,t^2)$ y $Q(t,8-t^2)$
sobre las curvas $y=x^2$ y $y=8-x^2$.
La región mostrada se parametriza simplificando la expresión $(1-u)Q+uP,$ es decir,
$(t,8(1-u)+(2u-1)t^2),$ con $t \in [-2,2]$, $u \in [0,1].$


\vskip0.2cm
\noindent
\begin{minipage}{\textwidth}
	\centering
	\captionof{figure}{Región delimitada por $x^2 \leq y \leq 8 - x^2$}
	\label{fig:region_parabolica}
	\begin{minipage}{0.38\linewidth}
		\centering
		\vspace{-18mm}
		\includegraphics[width=\linewidth]{Fig/386.pdf}
	\end{minipage}
	
	\vspace{-10mm}
	\fuentepropia
\end{minipage}
\vskip0.2cm


Sobre el cilindro $z=4-x^2,$ se satisface que $z=4-t^2$
y los puntos con coordenadas
$R(t,8(1-u)+(2u-1)t^2,0)$ y $S(t,8(1-u)+(2u-1)t^2,4-t^2)$
se unen mediante un segmento rectilíneos $(1-w)R+wS,$ que permite definir la función que parametriza el sólido.
\begin{align*}
\mathbf{r}(t,u,w) &= \langle t,8(1-u)+(2u-1)t^2,w(4-t^2)\rangle,
\end{align*}
con $t\in [-2,2],\, u\in [0,1], \, w\in [0,1].$
Esta función nos permite hallar los siguientes vectores
\vspace{-4mm}
\begin{align*}
\mathbf{r}_{t}&=\langle 1,2(2u-1)t,-2wt\rangle \\
\mathbf{r}_{u}&=\langle 0,-8+2t^2,0\rangle\\
\mathbf{r}_{w}&=\langle 0,0,4-t^2\rangle.
\end{align*}
Con estos se halla el producto vectorial triple
$|\mathbf{r}_{t}\cdot(\mathbf{r}_{u}\times \mathbf{r}_{w})| =2(t^2-4)^2,$
que se integra con los límites establecidos y nos proporciona el volumen del sólido
\[\int_{0}^{1}\int_{0}^{1}\int_{-2}^22(t^2-4)^2dtdudw=\frac{1024}{15}.\]

%\end{tcolorbox}
%\vskip0.2cm
Este proceso se hace en \verb"Mathematica" ejecutando las siguientes instrucciones:


\vspace{0.6em}
\noindent
\begin{tcolorbox}[breakable,enhanced,title=\bf Instrucción en \verb"Mathematica", top=2pt,bottom=2pt,boxsep=0pt,before skip=2pt,after skip=0pt]
\begin{Verbatim}[formatcom=\letraverbatim]
R={x,y}=Factor[u*{t,t^2}+(1-u)*{t,8-t^2}]
r=Simplify[(1-w)*{x,y,0}+w*{x,y,4-x^2}]
Integrate[Abs[Dot[D[r,t],Cross[D[r,u],D[r,w]]]],
{t,-2,2},{u,0,1},{w,0,1}]
\end{Verbatim}
\end{tcolorbox}
\vspace{0.7em}
\noindent


\subsection*{Dos cilindros parabólicos}

Consideraremos el sólido acotado por las figuras de las expresiones $z=4x-x^2$, $y=4x-x^2$, $y=0$ y $z=0$.
Un bosquejo de dicha región del espacio se muestra en la siguiente figura.
Las ecuaciones cuadráticas corresponden a dos cilindros parabólicos.
%\vskip0.4cm
\vskip0.2cm
\noindent
\begin{minipage}{\textwidth}
	\centering
	\captionof{figure}{Visualización del sólido definido por $0 \leq z \leq 4x - x^2$, \quad $0 \leq y \leq 4x - x^2$.}
	\label{fig:solido_4x_x2}
	\vspace{-10pt}
	\begin{minipage}{0.33\linewidth}
		\centering
		\vspace{-13mm}
		\includegraphics[width=\linewidth]{Fig/387.pdf}
	\end{minipage}
	\hspace{1em}
	\begin{minipage}{0.24\linewidth}
		\centering
		\vspace{-13mm}
		\includegraphics[width=\linewidth]{Fig/388.pdf}
	\end{minipage}
	\hspace{1em}
	\begin{minipage}{0.24\linewidth}
		\centering
		\vspace{-13mm}
		\includegraphics[width=\linewidth]{Img/7b.pdf}
	\end{minipage}
	
	\vspace{-10mm}
	\fuentepropia
\end{minipage}
\vskip0.4cm

La línea roja corresponde a la curva de intersección de los dos cilindros parabólicos; en esta curva se usó la parametrización
$\overrightarrow{r}(t)= \langle t,4t-t^2,4t-t^2\rangle,$ $t\in [0,4].$
Esta línea se proyecta en el plano $yz$ por medio de la recta $z=y.$
Estas figuras se consiguen con las siguientes instrucciones:


\vspace{0.6em}
\noindent
\begin{tcolorbox}[breakable,enhanced,title=\bf Instrucción en \verb"Mathematica", top=2pt,bottom=2pt,boxsep=0pt,before skip=2pt,after skip=0pt]
\begin{Verbatim}[formatcom=\letraverbatim]
ContourPlot3D[{z==4*x-x^2,y==4*x-x^2},{x,0,4},
{y,0,4},{z,0,4}]
RegionPlot3D[0<z<4*x-x^2&&0<y<4*x-x^2,{x,0,4},
{y,0,4},
{z,0,4},PlotStyle->{Blend[{Red,Orange,Gray}]},
Mesh->False,PlotPoints->50]
\end{Verbatim}
\end{tcolorbox}
\vspace{0.7em}
\noindent



La proyección del sólido en el plano $yz$ corresponde a un cuadrado, mientras en los planos $xy$ y $xz$ coincide con sectores parabólicos.

\vskip0.2cm
\noindent
\begin{minipage}{\textwidth}
	\centering
	\captionof{figure}{Proyección del sólido.}
	\label{fig:proyeccion_solido}
	\begin{minipage}{0.30\linewidth}
		\centering
		\vspace{-10mm}
		\includegraphics[width=\linewidth]{Fig/389.pdf}
	\end{minipage}
	\hspace{1em}
	\begin{minipage}{0.30\linewidth}
		\centering
		\vspace{-10mm}
		\includegraphics[width=\linewidth]{Fig/390.pdf}
	\end{minipage}
	\hspace{1em}
	\begin{minipage}{0.25\linewidth}
		\centering
		\vspace{-10mm}
		\includegraphics[width=\linewidth]{Fig/391.pdf}
	\end{minipage}
	
	\vspace{-8mm}
	\fuentepropia
\end{minipage}
\vskip0.2cm


Con ayuda de los elementos anteriores se presentan las integrales triples necesarias para calcular el volumen del sólido en mención. \\[0.2cm]
$\ds \int_{0}^{4}\int_{0}^{z}\int_{2-\sqrt{4-z}}^{2+\sqrt{4-z}}dxdydz+\int_{0}^{4}\int_{z}^{4}\int_{2-\sqrt{4-y}}^{2+\sqrt{4-y}}dxdydz= \frac{512}{15}$\\[0.1cm]
$\ds \int_{0}^{4}\int_{y}^{4}\int_{2-\sqrt{4-z}}^{2+\sqrt{4-z}}dxdzdy+\int_{0}^{4}\int_{0}^{y}\int_{2-\sqrt{4-y}}^{2+\sqrt{4-y}}dxdzdy= \frac{512}{15}$\\[0.1cm]
$\ds \int_{0}^{4}\int_{0}^{4x-x^2}\int_{0}^{4x-x^2}dzdydx = \int_{0}^{4}\int_{2-\sqrt{4-y}}^{2+\sqrt{4-y}}\int_{0}^{4x-x^2}dzdxdy= \frac{512}{15}$\\[0.1cm]
$\ds \int_{0}^{4}\int_{0}^{4x-x^2}\int_{0}^{4x-x^2}dydzdx= \int_{0}^{4}\int_{2-\sqrt{4-z}}^{2+\sqrt{4-z}}\int_{0}^{4x-x^2}dydxdz= \frac{512}{15}.$\\[0.15cm]
\vskip0.2cm
Algunas de estas integrales se evalúan en \verb"Mathematica" con las siguientes instrucciones:

\vspace{0.6em}
\noindent
\begin{tcolorbox}[breakable,enhanced,title=\bf Instrucción en \verb"Mathematica", top=2pt,bottom=2pt,boxsep=0pt,before skip=2pt,after skip=0pt]
\begin{Verbatim}[formatcom=\letraverbatim]
Integrate[1,{x,0,4},{y,0,4*x-x^2},{z,0,4*x-x^2}]
Integrate[1,{x,0,4},{z,0,4*x-x^2},{y,0,4*x-x^2}]
Integrate[1,{z,0,4},{x,2-Sqrt[4-z],2+Sqrt[4-z]},
{y,0,4*x-x^2}]
Integrate[1,{z,0,4},{y,0,z},{x,2-Sqrt[4-z],2+
Sqrt[4-z]}]+Integrate[1,{z,0,4},{y,z,4},{x,2-
Sqrt[4-y],2+Sqrt[4-y]}]
Integrate[1,{y,0,4},{z,y,4},{x,2-Sqrt[4-z],
2+Sqrt[4-z]}]+Integrate[1,{y,0,4},{z,0,y},
{x,2-Sqrt[4-y],2+Sqrt[4-y]}]
\end{Verbatim}
\end{tcolorbox}
\vspace{0.7em}
\noindent


Otra manera de calcular el volumen del sólido es la siguiente.


%\begin{tcolorbox}[title=\bfseries Parametrizando el sólido]
\noindent

En la proyección sobre el plano $xy$,
consideremos un punto $P(t,0)$ y un punto $Q(t,4t-t^2)$,
con $t \in [0,4]$. El segmento que los une corresponde a la expresión
$\left( 1-v\right) \left( t,0\right) +v\left( t,4t-t^2\right),$
que se simplifica como $(t,vt(4-t)),$ $v \in[0,1].$


\vskip0.2cm
\noindent
\begin{minipage}{\textwidth}
	\centering
	\captionof{figure}{Región delimitada por $0 \leq y \leq 4x - x^2$}
	\label{fig:region_4x_x2}
	\begin{minipage}{0.37\linewidth}
		\centering
		\vspace{-18mm}
		\includegraphics[width=\linewidth]{Fig/392.pdf}
	\end{minipage}
	
	\vspace{-14mm}
	\fuentepropia
\end{minipage}
\vskip0.2cm


Un punto en la base del sólido tiene la forma
$\left( t,vt(4-t),0\right).$ El cilindro con ecuación $y=4x-x^2$ es paralelo al eje $z,$
sobre esta región, un punto tiene la forma $\left( t,vt(4-t),4t-t^2\right)$.
Formamos segmentos rectilíneos que unen dichos puntos, usaremos el parámetro $w.$ Por lo tanto, al simplificar
$\left( 1-w\right) \left( t,vt(4-t),0\right) +w\left( t,vt(4-t),4t-t^2\right)$
obtendremos la función
\[{\bf r}(u,v,w)=\left\langle t,vt(4-t),-wt\left(t -4\right) \right\rangle,\]
que parametriza el sólido con $t\in [0,4],\,v\in [0,1],\,w\in [0,1]. $
A partir de esta se pueden obtener 
{\small
\begin{gather*}
	\mathbf{r}_{t}= \left\langle 1,4v-2vt,4w-2wt\right\rangle \\
	\mathbf{r}_{v}= \left\langle 0,4t-t^2,0\right\rangle \\
	\mathbf{r}_{w}= \left\langle 0,0,t\left( 4-t\right) \right\rangle \\
	\mathbf{r}_{v}\times \mathbf{r}_{w}= \left\langle t^2\left( t-4\right) ^2,0,0\right\rangle.
\end{gather*}
}
%\vskip0.2cm
Con esto se obtiene
$\left| {\bf r}_{t}\cdot \left( {\bf r}_{v}\times {\bf r}_{w}\right) \right| = t^2(t-4) ^2,$
el volumen está dado por:
%\vskip0.2cm
{\small
\[\int_{0}^{4}\int_{0}^{1}\int_{0}^{1}t^2\left( t-4\right)^2dvdwdt= \frac{512}{15}.\]
}
%\end{tcolorbox}
\vskip0.2cm
El cálculo anterior se hizo en \verb"Mathematica" con las siguientes instrucciones:


\vspace{0.6em}
\noindent
\begin{tcolorbox}[breakable,enhanced,title=\bf Instrucción en \verb"Mathematica", top=2pt,bottom=2pt,boxsep=0pt,before skip=2pt,after skip=0pt]
\begin{Verbatim}[formatcom=\letraverbatim]
{x,y}=Simplify[u*{t,t^2}+(1-u)*{t,8-t^2}]
r=Simplify[(1-w)*{x,y,0}+w*{x,y,4-x^2}]
A=Abs[Dot[D[r,w],Cross[D[r,t],D[r,u]]]];
Integrate[A,{t,-2,2},{u,0,1},{w,0,1}]
\end{Verbatim}
\end{tcolorbox}
\vspace{0.7em}
\noindent


\subsection*{Un cilindro parabólico y un plano inclinado III}

Las expresiones $z=2-y$, $y=x^2$ y $z=0$,
determinan una región en el espacio, su gráfico se muestra a continuación:
\vskip0.2cm

\noindent
\begin{minipage}{\textwidth}
	\centering
	\captionof{figure}{Visualización del sólido definido por $0 \leq z \leq 2 - y$, \quad $x^2 \leq y \leq 2$.}
	\label{fig:solido_parabolico_2y}
	\begin{minipage}{0.42\linewidth}
		\centering
		\vspace{-24mm}
		\includegraphics[width=\linewidth]{Fig/393.pdf}
	\end{minipage}
	\hspace{1em}
	\begin{minipage}{0.42\linewidth}
		\centering
		\vspace{-24mm}
		\includegraphics[width=\linewidth]{Fig/394.pdf}
	\end{minipage}
	
	\vspace{-15mm}
	\fuentepropia
\end{minipage}
\vskip0.2cm


Sobre los planos coordenados, el sólido se proyecta como sigue:
\vskip0.2cm

\noindent
\begin{minipage}{\textwidth}
	\centering
	\captionof{figure}{Proyecciones sobre los planos coordenados.}
	\label{fig:proyecciones_coordenadas_final}
	\begin{minipage}{0.30\linewidth}
		\centering
		\vspace{-10mm}
		\includegraphics[width=\linewidth]{Fig/395.pdf}
	\end{minipage}
	\hspace{1em}
	\begin{minipage}{0.30\linewidth}
		\centering
		\vspace{-10mm}
		\includegraphics[width=\linewidth]{Fig/396.pdf}
	\end{minipage}
	\hspace{1em}
	\begin{minipage}{0.30\linewidth}
		\centering
		\vspace{-10mm}
		\includegraphics[width=\linewidth]{Fig/397.pdf}
	\end{minipage}
	
	\vspace{-8mm}
	\fuentepropia
\end{minipage}
\vskip0.2cm


Estos gráficos nos ayudan a plantear las integrales necesarias para calcular el volumen del sólido.

\begin{multline*} \int_{-\sqrt{2}}^{\sqrt{2}}\int_{x^2}^2\int_{0}^{2-y}dzdydx=
\int_{0}^2\int_{-\sqrt{y}}^{\sqrt{y}}\int_{0}^{2-y}dzdxdy \\
= \int_{-\sqrt{2}}^{\sqrt{2}}\int_{0}^{2-x^2}\int_{x^2}^{2-z}dydzdx = \frac{32\sqrt{2}}{15} \end{multline*}

\begin{multline*} \int_{0}^2\int_{-\sqrt{2-z}}^{\sqrt{2-z}}\int_{x^2}^{2-z}dydxdz =
\int_{0}^2\int_{0}^{2-z}\int_{-\sqrt{y}}^{\sqrt{y}}dxdydz \\
= \int_{0}^2\int_{0}^{2-y}\int_{-\sqrt{y}}^{\sqrt{y}}dxdzdy = \frac{32\sqrt{2}}{15}. \end{multline*}
Se pueden evaluar unas de estas integrales usando \verb"Mathematica" con ayuda de las siguientes sentencias.


\vspace{0.6em}
\noindent
\begin{tcolorbox}[breakable,enhanced,title=\bf Instrucción en \verb"Mathematica", top=2pt,bottom=2pt,boxsep=0pt,before skip=2pt,after skip=0pt]
\begin{Verbatim}[formatcom=\letraverbatim]
Integrate[1,{x,-Sqrt[2],Sqrt[2]},{y,x^2,2},{z,0,2-y}]
Integrate[1,{z,0,2},{y,0,2-z},{x,-Sqrt[y],Sqrt[y]}]
Integrate[1,{y,0,2},{z,0,2-y},{x,-Sqrt[y],Sqrt[y]}]
Integrate[1,{x,-Sqrt[2],Sqrt[2]},{y,x^2,2},{z,0,2-y}]
\end{Verbatim}
\end{tcolorbox}
\vspace{0.7em}
\noindent


El siguiente procedimiento es una alternativa al cálculo del volumen.

\vskip0.2cm
%\begin{tcolorbox}[title=\bfseries Parametrizando el sólido]
En el plano $xy,$, la base del sólido se parametriza formando una combinación convexa
$\left( 1-u\right) \left( t,t^2\right) +u\left(t,2\right) = \left( t,2u-t^2u+t^2\right),$
entre dos puntos de la forma $\left( t,t^2\right) $ y $\left( t,2\right).$
Sobre el plano $z=2-y$ se satisface que $z=(t^2-2)(u-1) .$ Por tanto, el sólido se parametriza con una combinación convexa de la forma
\[\left( 1-w\right) \left\langle t,2u-t^2u+t^2,0\right\rangle +w\left\langle t,2u-t^2u+t^2,\left(t^2-2\right) \left( u-1\right) \right\rangle,\]
que al simplificarse nos permiten definir la función
\[{\bf r}(t,u,w)=\left\langle t,2u-t^2u+t^2,w\left( t^2-2\right) \left( u-1\right) \right\rangle,\]
con $t\in [-\sqrt{2},\sqrt{2}],$ $u\in [0,1],$ $w\in [0,1].$ 
\vskip0.2cm
Derivando parcialmente se obtienen los siguientes vectores
\begin{gather*}
	\mathbf{r}_{t} = \left\langle 1,2t-2tu,2tw\left( u-1\right) \right\rangle, \\
	\mathbf{r}_{u}= \left\langle 0,2-t^2,w\left( t^2-2\right)\right\rangle, \\
	\mathbf{r}_{w} =\left\langle 0,0,\left( t^2-2\right) \left( u-1\right)\right\rangle.
\end{gather*}

De modo que $\left\vert {\bf r}_{w}\cdot \left({\bf r}_{t}\times {\bf r}_{u}\right) \right\vert
=\left\vert \left( t^2-2\right) ^2\left( 1-u\right)\right\vert .$
Por tanto, el volumen del sólido está dado por
%\vskip0.2cm
\[\int_{0}^{1}\int_{0}^{1}\int_{-\sqrt{2}}^{\sqrt{2}}\left\vert
\left( t^2-2\right) ^2\left( 1-u\right) \right\vert dtdudw= \frac{32}{15}\sqrt{2}.\]
%\vskip0.2cm
%\end{tcolorbox}

El procedimiento anterior se hace en \verb"Mathematica" ejecutando las siguientes instrucciones:


\vspace{0.6em}
\noindent
\begin{tcolorbox}[breakable,enhanced,title=\bf Instrucción en \verb"Mathematica", top=2pt,bottom=2pt,boxsep=0pt,before skip=2pt,after skip=0pt]
\begin{Verbatim}[formatcom=\letraverbatim]
R={x,y}=(1-u)*{t,t^2}+u{t,2};
r=Simplify[(1-w)*{x,y,0}+w*{x,y,2-y}]
Integrate[Dot[D[r,w],Cross[D[r,t],D[r,u]]],
{t,-Sqrt[2],Sqrt[2]},{u,0,1},{w,0,1}]
\end{Verbatim}
\end{tcolorbox}
\vspace{0.7em}
\noindent



\subsection*{Un cilindro parabólico y un plano oblicuo}

A continuación, se muestra el sólido acotado por las figuras de las expresiones:
$z=2x-x^2, \,\, y=2-x,\,\, x=1,\,\, z=0, \,\, y=0 $.
\vskip0.4cm


\noindent
\begin{minipage}{\textwidth}
	\centering
	\captionof{figure}{Visualización del sólido definido por $0 \leq z \leq 2x - x^2$, \quad $1 \leq x \leq 2 - y$.}
	\label{fig:solido_2x_x2}
	\begin{minipage}{0.37\linewidth}
		\centering
		\vspace{-10mm}
		\includegraphics[width=\linewidth]{Fig/398.pdf}
	\end{minipage}
	\hspace{1em}
	\begin{minipage}{0.37\linewidth}
		\centering
		\vspace{-10mm}
		\includegraphics[width=\linewidth]{Fig/399.pdf}
	\end{minipage}
	
	\vspace{-10mm}
	\fuentepropia
\end{minipage}
\vskip0.2cm

Este sólido proyectado determina las siguientes regiones.

\vskip0.2cm
\noindent
\begin{minipage}{\textwidth}
	\centering
	\captionof{figure}{Proyecciones en cada uno de los planos coordenados.}
	\label{fig:proyecciones_planes}
	\begin{minipage}{0.28\linewidth}
		\centering
		\vspace{-9mm}
		\includegraphics[width=\linewidth]{Fig/400.pdf}
	\end{minipage}
	\hspace{1em}
	\begin{minipage}{0.28\linewidth}
		\centering
		\vspace{-9mm}
		\includegraphics[width=\linewidth]{Fig/401.pdf}
	\end{minipage}
	\hspace{1em}
	\begin{minipage}{0.28\linewidth}
		\centering
		\vspace{-9mm}
		\includegraphics[width=\linewidth]{Fig/402.pdf}
	\end{minipage}
	
	\vspace{-4mm}
	\fuentepropia
\end{minipage}
\vskip0.2cm

El volumen del sólido está dado por
\[\int_{1}^2\int_{0}^{2-x}\int_{0}^{2x-x^2}dzdydx = \int_{0}^{1}\int_{1}^{2-y}\int_{0}^{2x-x^2}dzdxdy= \frac{5}{12}\]
\[\int_{1}^2\int_{0}^{2x-x^2}\int_{0}^{2-x}dydzdx=
\int_{0}^{1}\int_{1}^{1+\sqrt{1-z}}\int_{0}^{2-x}dydxdz= \frac{5}{12}.\]

La curva de intersección proyectada en el plano $yz,$ satisface que
$z=2y-y^2,$ con lo cual $y=1\pm \sqrt{1-z}.$ La figura de esta ecuación divide en dos partes  la proyección del sólido que es un cuadrado, esto nos lleva a plantear dos integrales en el orden $dxdzdy$ o $dxdydz$.
\begin{eqnarray*}
\int_{0}^{1}\int_{0}^{2y-y^2}\int_{1}^{2-y}dxdzdy+\int_{0}^{1}%
\int_{2y-y^2}^{1}\int_{1}^{1+\sqrt{1-z}}dxdzdy &=& \frac{5}{12} \\
\int_{0}^{1}\int_{1-\sqrt{1-z}}^{1}\int_{1}^{2-y}dxdydz+\int_{0}^{1} \int_{0}^{1-\sqrt{1-z}}\int_{1}^{1+\sqrt{1-z}}dxdydz &=& \frac{5}{12}.
\end{eqnarray*}
En \verb"Mathematica" se pueden evaluar algunas de estas integrales ejecutando las siguientes instrucciones.


\vspace{0.6em}
\noindent
\begin{tcolorbox}[breakable,enhanced,title=\bf Instrucción en \verb"Mathematica", top=2pt,bottom=2pt,boxsep=0pt,before skip=2pt,after skip=0pt]
\begin{Verbatim}[formatcom=\letraverbatim]
Integrate[1,{y,0,1},{x,1,2-y},{z,0,2x-x^2}]
Integrate[1,{z,0,1},{x,1,1+Sqrt[1-z]},{y,0,2-x}]
Integrate[1,{x,1,2},{y,0,2-x},{z,0,2*x-x^2}]
Integrate[1,{x,1,2},{z,0,2*x-x^2},{y,0,2-x}]
\end{Verbatim}
\end{tcolorbox}
\vspace{0.7em}
\noindent


La otra posibilidad de obtener el volumen en mención es la siguiente.
%\begin{tcolorbox}[title=\bfseries Parametrizando el sólido]
La base del sólido es un triángulo en el plano $xy$ que se parametriza primero formando segmentos que unen los puntos $(1,0,0) $ .
$\left(2,0,0\right),$ posteriormente se forman otros segmentos que unan con $\left( 1,1,0\right) ,$
esto es $(1-u)(1+t,0,0)+u(1,1,0)= ((1-u)(1+t)+u,u,0).$
Con lo anterior, la coordenada $z$ sobre el cilindro satisface que $z=w (t+1-ut)(tu-t+1).$
Es por lo anterior, que el sólido se parametriza como sigue
%\vskip0.2cm
\[{\bf r}(t,u,w) =\left\langle t+1-ut,u,w(t+1-tu)(tu-t+1)\right\rangle,\]
con $t\in [0,1],$ $u\in [0,1],$ $w\in [0,1].$ Con esta función se determinan los siguientes vectores:
\vspace{-2mm}
\begin{eqnarray*}
{\bf r}_{t}&=& \left\langle 1-u,0,-2tw(u-1)^2\right\rangle \\
{\bf r}_{u}&=&\left\langle -t,1,2wt^2(1-u) \right\rangle \\
{\bf r}_{w}&=& \left\langle 0,0,1-t^2(u-1)^2\right\rangle
\end{eqnarray*}

El producto vectorial triple está dado por

\[
\left| {\bf r}_{t}\cdot \left( {\bf r}_{u}\times {\bf r}_{w}\right) \right| =
\left| (u-1)\, (t (1-u)+1)\, (t (u-1)+1) \right|,
\]

cuya integración nos proporciona el volumen del sólido:

\[
\int_{0}^{1} \int_{0}^{1} \int_{0}^{1}
\left| (u-1)\, (t (1-u)+1)\, (t (u-1)+1) \right|
\, dt\, du\, dw = \frac{5}{12}.
\]
\vskip0.2cm
%\end{tcolorbox}

Con las siguientes sentencias se hace el procedimiento descrito anteriormente.


\vspace{0.6em}
\noindent
\begin{tcolorbox}[breakable,enhanced,title=\bf Instrucción en \verb"Mathematica", top=2pt,bottom=2pt,boxsep=0pt,before skip=2pt,after skip=0pt]
\begin{Verbatim}[formatcom=\letraverbatim]
{x,y}=(1-u)*{1+t,0}+u*{1,1};
r=Simplify[{x,y,w*(2*x-x^2)}]
A=Simplify[Dot[D[r,w],Cross[D[r,u],D[r,t]]]]
Integrate[Abs[A],{t,0,1},{u,0,1},{w,0,1}]
\end{Verbatim}
\end{tcolorbox}
\vspace{0.7em}
\noindent


\subsection*{Dos cilindros parabólicos II}
La región del espacio acotada por las figuras de las expresiones
$x=z^2$ y $x=4y-y^2$, determinan el sólido que a continuación se presenta.
\vskip0.3cm

\noindent
\begin{minipage}{\textwidth}
	\centering
	\captionof{figure}{$z^2 \leq x \leq 4y - y^2$, \quad $y^2 + z^2 \leq 4y$.}
	\label{fig:solido_curvo}
	\begin{minipage}{0.20\linewidth}
		\centering
		\includegraphics[width=\linewidth]{Imgn/53a.pdf}
	\end{minipage}
	\hspace{1em}
	\begin{minipage}{0.21\linewidth}
		\centering
		\includegraphics[width=\linewidth]{Img/53b.pdf}
	\end{minipage}
	\hspace{1em}
	\begin{minipage}{0.21\linewidth}
		\centering
		\includegraphics[width=\linewidth]{Img/53c.pdf}
	\end{minipage}
	\hspace{1em}
	\begin{minipage}{0.21\linewidth}
		\centering
		\includegraphics[width=\linewidth]{Img/53d.pdf}
	\end{minipage}
	
	\vspace{-2mm}
	\fuentepropia
\end{minipage}
\vskip0.3cm

Estos gráficos se obtienen con la ayuda de las siguientes expresiones.


\vspace{0.6em}
\noindent
\begin{tcolorbox}[breakable,enhanced,title=\bf Instrucción en \verb"Mathematica", top=2pt,bottom=2pt,boxsep=0pt,before skip=2pt,after skip=0pt]
\begin{Verbatim}[formatcom=\letraverbatim]
ContourPlot3D[{x==z^2,x==4*y-y^2},{x,0,4},{y,0,4},
{z,-2,2},PlotPoints->50,Mesh->False,Ticks->{{0,2,4},
{0,1,3},{0,2}},ContourStyle->{Blend[{Yellow,Pink,
Green}],Orange}]
RegionPlot3D[-Sqrt[x]<z<Sqrt[x]&&0<x<4*y-y^2&&0<y<4,
{x,0,4},{y,0,4},{z,-2,2},PlotPoints->80,Mesh->False,
PlotStyle->Directive[Orange,Specularity[White,40]]]
\end{Verbatim}
\end{tcolorbox}
\vspace{0.7em}
\noindent


Al igualar estas expresiones para determinar su proyección en el plano $yz,$ obtenemos $(y-2)^2+z^2=4,$
que corresponde a una circunferencia de radio $2$ centrada en el punto $(2,0).$
\vskip 0.2cm
En el plano $xz$ la proyección es la porción de parábola $x=z^2$ para $x\in[0,4];$ mientras que para el plano $xy$ la región de integración es la porción de parábola $x=4y-y^2$ para $x\in[0,4].$ En los siguientes gráficos se ilustran estas regiones.
\vskip0.4cm

\noindent
\begin{minipage}{\textwidth}
	\centering
	\captionof{figure}{Proyecciones en los planos coordenados.}
	\label{fig:proyecciones_curvas}
	\begin{minipage}{0.29\linewidth}
		\centering
		\vspace{-10mm}
		\includegraphics[width=\linewidth]{Fig/403.pdf}
	\end{minipage}
	\hspace{1em}
	\begin{minipage}{0.29\linewidth}
		\centering
		\vspace{-10mm}
		\includegraphics[width=\linewidth]{Fig/404.pdf}
	\end{minipage}
	\hspace{1em}
	\begin{minipage}{0.29\linewidth}
		\centering
		\vspace{-10mm}
		\includegraphics[width=\linewidth]{Fig/405.pdf}
	\end{minipage}
	
	\vspace{-8mm}
	\fuentepropia
\end{minipage}
\vskip0.2cm

Con ayuda de los gráficos anteriores se pueden plantear las integrales que determinan el volumen del sólido; en un caso conviene escribir la expresión $x=4y-y^2$ en la forma
$y=2\pm \sqrt{4-x}.$
{\small
\begin{eqnarray*}
\int_{0}^{4}\int_{-\sqrt{4y-y^2}}^{\sqrt{4y-y^2}}\int_{z^2}^{4y-y^2}dxdzdy= \int_{-2}^2\int_{2-\sqrt{4-z^2}}^{2+\sqrt{4-z^2}}\int_{z^2}^{4y-y^2}dxdydz &=& 8\pi \\
\int_{-2}^2\int_{z^2}^{4}\int_{2-\sqrt{4-x}}^{2+\sqrt{4-x}}dydxdz=\int_{0}^{4}\int_{-\sqrt{x}}^{\sqrt{x}}\int_{2-\sqrt{4-x}}^{2+\sqrt{4-x}}dydzdx &=& 8\pi \\
\int_{0}^{4}\int_{0}^{4y-y^2}\int_{-\sqrt{x}}^{\sqrt{x}}dzdxdy= \int_{0}^{4}\int_{2-\sqrt{4-x}}^{2+\sqrt{4-x}}\int_{-\sqrt{x}}^{\sqrt{x}}dzdydx &= & 8\pi
\end{eqnarray*}
}
Algunas de estas integrales se evalúan como sigue:


\vspace{0.6em}
\noindent
\begin{tcolorbox}[breakable,enhanced,title=\bf Instrucción en \verb"Mathematica", top=2pt,bottom=2pt,boxsep=0pt,before skip=2pt,after skip=0pt]
\begin{Verbatim}[formatcom=\letraverbatim]
Integrate[1,{y,0,4},{z,-Sqrt[4y-y^2],Sqrt[4y-y^2]},
{x,z^2,4y-y^2}]
Integrate[1,{z,-2,2},{y,2-Sqrt[4-z^2],2+Sqrt[4-z^2]},
{x,z^2,4y-y^2}]
Integrate[1,{y,0,4},{x,0,4y-y^2},{y,-Sqrt[x],Sqrt[x]}]
Integrate[1,{z,-2,2},{x,z^2,4},{y,2-Sqrt[4-x],
2+Sqrt[4-x]}]
\end{Verbatim}
\end{tcolorbox}
\vspace{0.7em}
\noindent

El mismo valor se puede obtener de la siguiente manera.

%\begin{tcolorbox}[title=\bfseries Parametrizando el sólido]
En el plano $xz$ dos puntos sobre la región de integración poseen coordenadas
$( t,-\sqrt{t})$ y $(t,\sqrt{t})$; con una combinación convexa entre ellos se describe completamente la región.
Es decir, simplificando
%\vskip0.2cm
\[( 1-u) \langle t,-\sqrt{t}\rangle +u\langle t,\sqrt{t}\rangle;\]
%\vskip0.2cm
para $t\in [0,4]$ y $u\in [0,1]$,
se obtiene $\langle t,\sqrt{t}(2u-1)\rangle\!.$ Por lo tanto, un punto sobre la región de integración posee coordenadas $( t,\sqrt{t}( 2u-1)),$
mientras que dos puntos $P$ y $Q$ sobre el cilindro poseen coordenadas
$( t,2-\sqrt{4-t},\sqrt{t}( 2u-1) ) $ y $( t,2+\sqrt{4-t},\sqrt{t}( 2u-1) ) .$ % 
Una combinación convexa de estos puntos parametriza el sólido, esto es
%\vskip0.2cm
{\small
\[{\bf r}(t,u,w) = \langle t,2+( 1-2w) \sqrt{4-t},( 2u-1) \sqrt{t}\rangle.\]
}
%\vskip0.2cm
Con esta función hallamos los siguientes vectores:
{\small
\begin{eqnarray*}
{\bf r}_t & =& \langle 1,-\frac{1}{2}\frac{1-2w}{\sqrt{4-t}},\frac{2u-1}{2\sqrt{t}}\rangle \\
{\bf r}_u & =& \langle 0,0,2\sqrt{t}\rangle \\
{\bf r}_w & =& \langle 0,-2\sqrt{4-t},0\rangle.
\end{eqnarray*}
}
Es sencillo ver que $|{\bf r}_t \cdot ({\bf r}_u \times {\bf r}_w )| = 4\sqrt{t}\sqrt{4-t}.$
Integrando este término obtenemos el volumen deseado.
%\vskip0.2cm
\[\int_{0}^{1}\int_{0}^{1}\int_{0}^{4}( 4\sqrt{t}\sqrt{4-t}) dtdudw =8\pi.\]
%\vskip0.2cm
%\end{tcolorbox}

Ejecutando las siguientes instrucciones se obtiene el volumen del sólido.


\vspace{0.6em}
\noindent
\begin{tcolorbox}[breakable,enhanced,title=\bf Instrucción en \verb"Mathematica", top=2pt,bottom=2pt,boxsep=0pt,before skip=2pt,after skip=0pt]
\begin{Verbatim}[formatcom=\letraverbatim]
r=Simplify[(1-w)*{t,2+Sqrt[4-t],Sqrt[t]*(2*u-1)}+
w*{t,2-Sqrt[4-t],Sqrt[t]*(2*u-1)}]
Integrate[Dot[D[r,t],Cross[D[r,u],D[r,w]]],{w,0,1},
{u,0,1},{t,0,4}]
\end{Verbatim}
\end{tcolorbox}
\vspace{0.7em}
\noindent

\subsection*{Dos cilindros parabólicos y dos planos}

Las expresiones $z=y^2$, $x=z^2/4$, $x=0$ y $y=2,$ determinan un sólido cuya figura se muestra a continuación.
\vskip0.2cm

\noindent
\begin{minipage}{\textwidth}
	\centering
	\captionof{figure}{Visualización del sólido definido por $0 \leq 4x \leq z^2$, \quad $0 \leq z \leq y^2$.}
	\label{fig:solido_xz_y2}
	\begin{minipage}{0.32\linewidth}
		\centering
		\vspace{-15mm}
		\includegraphics[width=\linewidth]{Fig/406.pdf}
	\end{minipage}
	\hspace{1em}
	\begin{minipage}{0.28\linewidth}
		\centering
		\vspace{-15mm}
		\includegraphics[width=\linewidth]{Fig/407.pdf}
	\end{minipage}
	
	\vspace{-8mm}
	\fuentepropia
\end{minipage}

\vskip0.2cm
En el plano $yz$ la proyección se determina por las ecuaciones $y=2$, $z=y^2,$ $z=0,$
y en el plano $xz$ la proyección la determinan las expresiones $4x=z^2,\,\, z=4$ y $x=0$.
Al sustituir las ecuaciones
$4x=z^2$ y $z=y^2$, en términos de $x$ e $y$ se puede representar la proyección en el plano
$xy.$ Es decir, $4x=z^2=(y^2)^2$ luego $4x=y^4,$ con lo anterior se obtiene las siguientes figuras.


\vskip0.2cm
\noindent
\begin{minipage}{\textwidth}
	\centering
	\captionof{figure}{Proyección sobre los planos coordenados.}
	\label{fig:proyecciones_planes_curvas}
	\begin{minipage}{0.25\linewidth}
		\centering
		\vspace{-10mm}
		\includegraphics[width=\linewidth]{Fig/408.pdf}
	\end{minipage}
	\hspace{1em}
	\begin{minipage}{0.25\linewidth}
		\centering
		\vspace{-10mm}
		\includegraphics[width=\linewidth]{Fig/409.pdf}
	\end{minipage}
	\hspace{1em}
	\begin{minipage}{0.25\linewidth}
		\centering
		\vspace{-10mm}
		\includegraphics[width=\linewidth]{Fig/410.pdf}
	\end{minipage}
	
	\vspace{-4mm}
	\fuentepropia
\end{minipage}
\vskip0.2cm

Con ayuda de estos gráficos se pueden plantear las seis integrales triples que permiten calcular el volumen del sólido en mención.
\vspace{-2mm}
\begin{multline*}
\int_{0}^2\int_{0}^{y^{4}/4}\int_{2\sqrt{x}}^{y^2}dzdxdy=\int_{0}^{4}\int_{\sqrt[4]{4x}}^2\int_{2\sqrt{x}}^{y^2}dzdydx \\
= \int_{0}^{4}\int_{0}^{z^2/4}\int_{\sqrt{z}}^2dydxdz=
\int_{0}^{4}\int_{2\sqrt{x}}^{4}\int_{\sqrt{z}}^2dydzdx \\= \int_{0}^{4}\int_{\sqrt{z}}^2\int_{0}^{z^2/4}dxdydz= \int_{0}^2\int_{0}^{y^2}\int_{0}^{z^2/4}dxdzdy= \frac{32}{21}.
\end{multline*}
Se presentan las instrucciones para evaluar cuatro de estas integrales.


\vspace{0.6em}
\noindent
\begin{tcolorbox}[breakable,enhanced,title=\bf Instrucción en \verb"Mathematica", top=2pt,bottom=2pt,boxsep=0pt,before skip=2pt,after skip=0pt]
\begin{Verbatim}[formatcom=\letraverbatim]
Integrate[1,{y,0,2},{z,0,y^2},{x,0,z^2/4}]
Integrate[1,{z,0,4},{y,Sqrt[z],2},{x,0,z^2/4}]
Integrate[1,{x,0,4},{z,Sqrt[4*x],4},{y,Sqrt[z],2}]
Integrate[1,{y,0,2},{x,0,y^4/4},{z,2*Sqrt[x],y^2}]
\end{Verbatim}
\end{tcolorbox}
\vspace{0.7em}
\noindent


Los siguientes gráficos ofrecen otras vistas del mismo sólido.

\begin{figure}[H]
\centering
\caption{Visualización del sólido definido por $0 \leq 4x \leq z^2$, \quad $0 \leq z \leq y^2$.}
\label{fig:solido_80_xyz}
\includegraphics[height=4.5cm]{Img/80c.pdf}
\includegraphics[height=4.5cm]{Img/80b.pdf}
\includegraphics[height=4.5cm]{Img/80a.pdf}
\fuentepropia
\end{figure}

Estos se consiguen ejecutando las siguientes instrucciones.


\vspace{0.6em}
\noindent
\begin{tcolorbox}[breakable,enhanced,title=\bf Instrucción en \verb"Mathematica", top=2pt,bottom=2pt,boxsep=0pt,before skip=2pt,after skip=0pt]
\begin{Verbatim}[formatcom=\letraverbatim]
RegionPlot3D[0<x<z^2/4&&Sqrt[z]<y<2,{x,0,4},{y,0,2},
{z,0,4},PlotPoints->90,PlotStyle->Orange,Boxed->False,
MeshStyle->Gray,AxesStyle->False,Axes->None]
\end{Verbatim}
\end{tcolorbox}
\vspace{0.7em}
\noindent


Una alternativa para obtener el volumen del sólido descrito es la siguiente.



%\begin{tcolorbox}[title=\bfseries Parametrizando el sólido]
\noindent

Consideremos dos puntos $P$ y $Q$ con coordenadas $(t,t^2)$ y $(2,t^2)$ respectivamente, en la proyección del sólido en el plano $yz,$
con $t \in [0,2]$. Simplificando una combinación convexa de estos se llena la región, es decir \linebreak
$(1-u)(t,t^2)+u(2,t^2)=( t+u(2-t),t^2) $
$t \in [0,2],\, u \in [0,1].$


\vskip0.3cm
\noindent
\begin{minipage}{\textwidth}
	\centering
	\captionof{figure}{Proyección de la parábola $z = y^2$ y puntos $P$, $Q$ sobre la recta $z = 2$}
	\label{fig:proyeccion_parabola_pq}
	\begin{minipage}{0.32\linewidth}
		\centering
		\vspace{-15mm}
		\includegraphics[width=\linewidth]{Fig/411.pdf}
	\end{minipage}
	
	\vspace{-8mm}
	\fuentepropia
\end{minipage}
\vskip0.2cm


Ahora, el punto $\left( 0,t+u\left( 2-t\right) ,t^2\right) $ está
en el plano $x=0,$ mientras que el punto
$\left( \frac{1}{4}t^{4},t+u\left( 2-t\right) ,t^2\right) $ yace en el cilindro $x=\frac{1}{4}z^2.$ Uniendo estos dos puntos con un segmento rectilíneo se parametriza el sólido, es decir simplificando la expresión
\vspace{-2mm}
\[\left( 1-w\right) \left (t^{4}/4,t+u(2-t),t^2\right) +w\left( 0,t+u(2-t) ,t^2\right),\] 
\[{\bf r}\left( t,u,w\right) = \left\langle \frac{1}{4}\left(1-w\right) t^{4},t+2u-tu,t^2\right\rangle,\]
en donde $t\in [0,2],$ $u\in [0,1],$ $w\in [0,1].$
Derivando parcialmente, se hallan los siguientes vectores

\vspace{-4mm}
{\small
\begin{align*}
	\mathbf{r}_t &= \frac{\partial}{\partial t} \left\langle \frac{1}{4}(1 - w) t^4, t + 2u - tu, t^2 \right\rangle = \left\langle (1 - w) t^3, 1 - u, 2t \right\rangle \\
	\mathbf{r}_u &= \frac{\partial}{\partial u} \left\langle \frac{1}{4}(1 - w) t^4, t + 2u - tu, t^2 \right\rangle = \left\langle 0, 2 - t, 0 \right\rangle \\
	\mathbf{r}_w &= \frac{\partial}{\partial w} \left\langle \frac{1}{4}(1 - w) t^4, t + 2u - tu, t^2 \right\rangle = \left\langle -\frac{1}{4}t^4, 0, 0 \right\rangle.
\end{align*}
}

\vskip0.2cm
Con ellos se calcula el triple producto vectorial
{\small
\[\left| {\bf r}_{t}\cdot \left( {\bf r}_{u}\times {\bf r}_{w}\right) \right| =
\frac{1}{2}\left| t^{5}\left( t-2\right) \right|,\]
}
que corresponde al elemento diferencial de volumen. Integrándolo, nos proporciona el volumen del sólido
%\vskip0.2cm
{\small
\[\int_{0}^{1}\int_{0}^{1}\int_0^2\frac{1}{2}\left|t^{5}(t-2)\right|dtdudw= \frac{32}{21}.\]
}
%\vskip0.2cm
%\end{tcolorbox}

Este valor se obtiene en \verb"Mathematica" ejecutando las siguientes instrucciones.
\vskip0.2cm
\begin{tcolorbox}[breakable,enhanced,title=\bf Instrucción en \verb"Mathematica", top=2pt,bottom=2pt,boxsep=0pt,before skip=2pt,after skip=0pt]
\begin{Verbatim}[formatcom=\letraverbatim]
r=(1-w)*{t^4/4,t+u*(2-t),t^2}+w*{0,t+u*(2-t),t^2}
Integrate[Dot[D[r,w],Cross[D[r,t],D[r,u]]],{t,0,2},
{u,0,1},{w,0,1}]
\end{Verbatim}
\end{tcolorbox}

\subsection*{Un cilindro parabólico y dos planos II}

Las ecuaciones $z=y^2$, $x=-1$ y $2z+x=3$ acotan un sólido en el espacio. Enseguida se representan estas expresiones junto con la región acotada.
\vskip0.4cm

\noindent
\begin{minipage}{\textwidth}
	\centering
	\captionof{figure}{Visualización del sólido definido por $-1 \leq x \leq 3 - 2z$, \quad $-\sqrt{z} \leq y \leq \sqrt{z}$.}
	\label{fig:solido_88_xyz}
	\begin{minipage}{0.22\linewidth}
		\centering
		\includegraphics[width=\linewidth]{Img/88a.pdf}
	\end{minipage}
	\hspace{1em} 
	\begin{minipage}{0.22\linewidth}
		\centering
		\includegraphics[width=\linewidth]{Img/88b.pdf}
	\end{minipage}
	\hspace{1em} 
	\begin{minipage}{0.22\linewidth}
		\centering
		\includegraphics[width=\linewidth]{Img/88d.pdf}
	\end{minipage}
	
	\vspace{0.3em}
	\fuentepropia
\end{minipage}
\vskip0.3cm

las figuras se obtienen con las siguientes instrucciones:

\vspace{0.6em}
\noindent
\begin{tcolorbox}[breakable,enhanced,title=\bf Instrucción en \verb"Mathematica", top=2pt,bottom=2pt,boxsep=0pt,before skip=2pt,after skip=0pt]
\begin{Verbatim}[formatcom=\letraverbatim]
ContourPlot3D[{z==y^2,x==-1,z==(3-x)/2},{x,-1,3},
{y,-1.5,1.5},{z,0,2},ContourStyle->{Orange,Cyan,
RGBColor[0.2,0.4,0.4]},Mesh->False,BoxRatios->{2,2,1}]
RegionPlot3D[y^2<=z<(3-x)/2&&-1<x<3,{x,-1,3},
{y,-1.5,1.5},{z,0,2},PlotPoints->90,Mesh->False,
BoxRatios->{2,2,1},PlotStyle->RGBColor[0.65,0.1,0.2]]
\end{Verbatim}
\end{tcolorbox}
\vspace{0.7em}
\noindent


Este sólido se proyecta sobre los planos coordenados.
\vskip0.2cm

\noindent
\begin{minipage}{\textwidth}
	\centering
	\captionof{figure}{Proyecciones del sólido definido por $-1 \leq x \leq 1 - 3z$, \quad $y^2 \leq z \leq 2$, \quad $-1 \leq x \leq 3 - 2y^2$.}
	\label{fig:proyecciones_sistema_curvo}
	\begin{minipage}{0.27\linewidth}
		\centering
		\vspace{-10mm}
		\includegraphics[width=\linewidth]{Fig/412.pdf}
	\end{minipage}
	\hspace{1em}
	\begin{minipage}{0.27\linewidth}
		\centering
		\vspace{-10mm}
		\includegraphics[width=\linewidth]{Fig/413.pdf}
	\end{minipage}
	\hspace{1em}
	\begin{minipage}{0.27\linewidth}
		\centering
		\vspace{-10mm}
		\includegraphics[width=\linewidth]{Fig/414.pdf}
	\end{minipage}
	
	\vspace{-4mm}
	\fuentepropia
\end{minipage}
\vskip0.2cm

En coordenadas rectangulares, los límites de las integrales que determinan el volumen se establecen con ayuda de estas figuras,
\begin{multline*}
\bigintsss_{0}^2\bigintsss_{-1}^{3-2z}\bigintsss_{-\sqrt{z}}^{\sqrt{z}}dydxdz = \bigintsss_{-1}^{3} \bigintsss_{0}^{\frac{3-x}{2}} \bigintsss_{-\sqrt{z}}^{\sqrt{z}}dydzdx \\
= \bigintsss_{-\sqrt{2}}^{\sqrt{2}}\bigintsss_{-1}^{3-2y^2}\bigintsss_{y^2}^{\frac{3-x}{2}}dzdxdy= \bigintsss_{-1}^{3}\bigintsss_{-\sqrt{\frac{3-x}{2}}}^{\sqrt{\frac{3-x}{2}}}\bigintsss_{y^2}^{\frac{3-x}{2}}dzdydx \\
= \bigintsss_{-\sqrt{2}}^{\sqrt{2}}\bigintsss_{y^2}^2 \bigintsss_{-1}^{3-2z}dxdzdy=\bigintsss_{0}^2\bigintsss_{-\sqrt{z}}^{\sqrt{z}}\bigintsss_{-1}^{3-2z}dxdydz= \frac{64}{15}\sqrt{2}
\end{multline*}

Cinco de estas integrales se evalúan en \verb"Mathematica" ejecutando las siguien\-tes instrucciones:


\vspace{0.6em}
\noindent
\begin{tcolorbox}[breakable,enhanced,title=\bf Instrucción en \verb"Mathematica", top=2pt,bottom=2pt,boxsep=0pt,before skip=2pt,after skip=0pt]
\begin{Verbatim}[formatcom=\letraverbatim]
Integrate[1,{y,-Sqrt[2],Sqrt[2]},{z,y^2,2},{x,-1,3-2z}]
Integrate[1,{y,-Sqrt[2],Sqrt[2]},{x,-1,3-2y^2},
{z,y^2,(3-x)/2}]
Integrate[z^2,{z,0,2},{x,-1,3-2z},{y,-Sqrt[z],Sqrt[z]}]
Integrate[1,{z,0,2},{y,-Sqrt[z],Sqrt[z]},{x,-1,3-2z}]
Integrate[1,{x,-1,3},{y,-Sqrt[(3-x)/2],Sqrt[(3-x)/2]},
{z,y^2,(3-x)/2}]
\end{Verbatim}
\end{tcolorbox}
\vspace{0.7em}
\noindent


El siguiente proceso ofrece otra posibilidad de obtener el mismo valor.

%\begin{tcolorbox}[title=\bfseries Parametrizando el sólido]
%\footnotesize 
%\parskip=0pt 
%\setlength{\belowdisplayskip}{0pt} 
%\setlength{\abovedisplayskip}{0pt} 

En el plano $xy,$ un punto $M$ sobre la curva $x=3-2y^2$ tiene la forma $(3-2t^2,t)$ y sobre la recta $x=-1$, un punto $N$
tiene la forma $(-1,t)$. Por tanto, la proyección del sólido se parametriza mediante una combinación convexa de $M$ y $N$,
esto es $(1-u)N+uM=\left( 2u\left( 2-t^2\right) -1,t\right) ,$ $t \in [-\sqrt{2},\sqrt{2}],$ $u \in [0,1].$
\vskip0.4cm

\begin{center}
	% 1. Ampliamos la minipage al ancho total para que el texto NO se comprima
	\begin{minipage}{1.0\linewidth} 
		\centering
		% La leyenda ahora tiene espacio para extenderse horizontalmente
		\captionof{figure}{Región delimitada por $-1 \leq x \leq 3 - 2y^2$}
		\label{fig:region_parabolica_x} 
		
		\vspace{-20mm} % Espacio estético entre leyenda y figura
		
		% 2. Mantenemos la figura pequeña (ajusta este valor a tu gusto)
		\includegraphics[width=0.45\linewidth]{Fig/415.pdf}
		
		\vspace{-18mm}
		% Fuente centrada respecto a la página
		\parbox{\linewidth}{\centering\fontsize{9pt}{10pt}\selectfont\textbf{Fuente:} elaboración propia.}
	\end{minipage}
\end{center}
\vskip0.4cm
Supongamos un punto $P\left( 2u\left( 2-t^2\right) -1,t,t^2\right) $ sobre el cilindro $z=y^2,$
y un punto $Q$ sobre el plano $2z+x=3,$ cuya coordenada $z$ satisface
%\vskip0.2cm
\[
z = \frac{3 - x}{2} = \frac{3}{2} - \frac{1}{2} \left(2u \left(2 - t^2\right) - 1\right) = u\left(t^2 - 2\right) + 2.
\]
%\vskip0.2cm
Por tanto, $Q$ es de la forma $Q( 2u\left( 2-t^2\right) -1,t,u\left( t^2-2\right)+2).$ Formando una combinación convexa entre $P$ y $Q$, se parametriza el sólido, esto es,
%\vskip0.2cm
\begin{equation*}
{\bf r}\left( t,u,w\right) = \langle 2u\left( 2-t^2\right)-1,t,t^2\left( 1-w+uw\right) -2w\left( u-1\right) \rangle.
\end{equation*}
%\vskip0.2cm
Con esta función hallamos los vectores
\begin{gather*}
	\mathbf{r}_{t} = \left\langle -4tu,1,2t\left( uw-w+1\right) \right\rangle, \\
	\mathbf{r}_{u} = \left\langle 4-2t^2,0,t^2w-2w\right\rangle, \\
	\mathbf{r}_{w} = \left\langle 0,0,\left( u-1\right) t^2-2u+2\right\rangle.
\end{gather*}
\vskip0.2cm
El elemento diferencial de volumen está dado por
\[\left\vert {\bf r}_{t}\cdot \left( {\bf r}_{u}\times {\bf r}_{w}\right) \right\vert
=\left\vert \left( 2t^2-4\right) \left( \left( u-1\right)t^2-2u+2\right) \right\vert,\]
cuya integral determina el volumen del sólido
%\vskip0.2cm
\[\int_{-\sqrt{2}}^{\sqrt{2}}\int_{0}^{1}\int_{0}^{1}\left\vert \left(2t^2-4\right) \left( \left( u-1\right) t^2-2u+2\right) \right\vert dwdudt= \frac{64}{15}\sqrt{2}.\]
%\vskip0.2cm
%\end{tcolorbox}

Con las siguientes instrucciones se hacen los cálculos anteriores.
%\vskip0.2cm

\vspace{0.6em}
\noindent
\begin{tcolorbox}[breakable,enhanced,title=\bf Instrucción en \verb"Mathematica", top=2pt,bottom=2pt,boxsep=0pt,before skip=2pt,after skip=0pt]
\begin{Verbatim}[formatcom=\letraverbatim]
r=Simplify[(1-w)*{2u*(2-t^2)-1,t,t^2}+w*{2u*(2-t^2)-1,
t,u*(t^2-2)+2}]
A=Dot[D[r,w],Cross[D[r,t],D[r,u]]]
Integrate[Abs[A],{w,0,1},{u,0,1},{t,-Sqrt[2],Sqrt[2]}]
\end{Verbatim}
\end{tcolorbox}
\vspace{0.7em}
\noindent


\subsection*{Un cilindro parabólico y dos planos}

Consideremos el sólido $E$, acotado por las figuras de las expresiones
$x=1-z^2$, $z=1-y$, $y=0$. Estas superficies y el sólido se muestran enseguida.

\vskip0.4cm
\noindent
\begin{minipage}{\textwidth}
	\centering
	\captionof{figure}{Visualización del sólido definido por $0 \leq x \leq 1 - z^2$, \quad $0 \leq y \leq 1 - z$.}
	\label{fig:solido_30_xyz}
	\begin{minipage}{0.24\linewidth}
		\centering
		\vspace{-15mm}
		\includegraphics[width=\linewidth]{Img/30a.pdf}
	\end{minipage}
	\hspace{1em}
	\begin{minipage}{0.40\linewidth}
		\centering
		\vspace{-15mm}
		\includegraphics[width=\linewidth]{Fig/416.pdf}
	\end{minipage}
	\hspace{1em}
	\begin{minipage}{0.24\linewidth}
		\centering
		\vspace{-15mm}
		\includegraphics[width=\linewidth]{Img/30d.pdf}
	\end{minipage}
	
	\vspace{-12mm}
	\parbox{\linewidth}{\centering\fontsize{9pt}{10pt}\selectfont\textbf{Fuente:} elaboración propia.}
\end{minipage}

\vskip0.4cm
Con las siguientes instrucciones se obtienen las figuras anteriores:


\vspace{0.6em}
\noindent
\begin{tcolorbox}[breakable,enhanced,title=\bf Instrucción en \verb"Mathematica", top=2pt,bottom=2pt,boxsep=0pt,before skip=2pt,after skip=0pt]
\begin{Verbatim}[formatcom=\letraverbatim]
ContourPlot3D[{x==1-z^2,z==1-y},{x,0,1},{y,0,2},
{z,-1,1},Mesh->False,ContourStyle->{Orange,Gray,
Opacity[0.85]},PlotPoints->80]
RegionPlot3D[0<x<1-z^2&&0<y<1-z,{x,0,1},{y,0,2},
{z,-1,1},PlotStyle->Blend[{Yellow,Orange,Gray,Red}],
Mesh->False,PlotPoints->90]
\end{Verbatim}
\end{tcolorbox}
\vspace{0.7em}
\noindent


Con el propósito de plantear las integrales triples que se requieren para calcular el volumen de $E$, se muestran ahora las proyecciones en los planos coordenados.
Para el caso del plano $xy$, se elimina la variable $z$ de las expresiones que determinan la intersección, es decir, se hace la sustitución $x=1-z^2=1-(1-y)^2=2y-y^2.$

\vskip0.3cm
\noindent
\begin{minipage}{\textwidth}
	\centering
	\captionof{figure}{Proyecciones del sólido definido por $0 \leq x \leq 1 - z^2$, \quad $0 \leq y \leq 1 - z$, \quad $-1 \leq z \leq 1$.}
	\label{fig:proyecciones_1menosz}
	\begin{minipage}{0.28\linewidth}
		\centering
		\vspace{-10mm}
		\includegraphics[width=\linewidth]{Fig/417.pdf}
	\end{minipage}
	\hspace{1em}
	\begin{minipage}{0.28\linewidth}
		\centering
		\vspace{-10mm}
		\includegraphics[width=\linewidth]{Fig/418.pdf}
	\end{minipage}
	\hspace{1em}
	\begin{minipage}{0.28\linewidth}
		\centering
		\vspace{-10mm}
		\includegraphics[width=\linewidth]{Fig/419.pdf}
	\end{minipage}
	
	\vspace{-6mm}
	\fuentepropia
\end{minipage}
\vskip0.3cm


Ahora se presentan las integrales que determinan el volumen del sólido $E$:
\begin{eqnarray*}
\ds \int_{-1}^{1}\int_{0}^{1-z}\int_{0}^{1-z^2}dxdydz=\ds\int_{0}^2\int_{-1}^{1-y}\int_{0}^{1-z^2}dxdzdy&=& \frac{4}{3}\\
\ds \int_{-1}^{1}\int_{0}^{1-z^2}\int_{0}^{1-z}dydxdz=\ds\int_{0}^{1}\int_{-\sqrt{1-x}}^{\sqrt{1-x}}\int_{0}^{1-z}dydzdx&=& \frac{4}{3}\\
\ds \int_{0}^{1}\int_{1-\sqrt{1-x}}^{1+\sqrt{1-x}}\int_{-\sqrt{1-x}}^{1-y}dzdydx+\int_{0}^{1}\int_{0}^{1-\sqrt{1-x}}\int_{-\sqrt{1-x}}^{\sqrt{1-x}}dzdydx&=& \frac{4}{3}\\
\ds \int_{0}^2\int_{0}^{2y-y^2}\int_{-\sqrt{1-x}}^{1-y}dzdxdy+\int_{0}^{1} \int_{2y-y^2}^{1}\int_{-\sqrt{1-x}}^{\sqrt{1-x}}dzdxdy&=& \frac{4}{3}
\end{eqnarray*}
Las integrales anteriores se evaluaron en \verb"Mathematica" mediante la ejecución de las siguientes instrucciones:


\vspace{0.6em}
\noindent
\begin{tcolorbox}[breakable,enhanced,title=\bf Instrucción en \verb"Mathematica", top=2pt,bottom=2pt,boxsep=0pt,before skip=2pt,after skip=0pt]
\begin{Verbatim}[formatcom=\letraverbatim]
Integrate[1,{z,-1,1},{y,0,1-z},{x,0,1-z^2}]
Integrate[1,{y,0,2},{z,-1,1-y},{x,0,1-z^2}]
Integrate[1,{z,-1,1},{x,0,1-z^2},{y,0,1-z}]
Integrate[1,{x,0,1},{z,-Sqrt[1-x],Sqrt[1-x]},{y,0,1-z}]
Integrate[1,{x,0,1},{y,1-Sqrt[1-x],1+Sqrt[1-x]},
{z,-Sqrt[1-x],1-y}]+Integrate[1,{x,0,1},
{y,0,1-Sqrt[1-x]},{z,-Sqrt[1-x],Sqrt[1-x]}]
Integrate[1,{y,0,2},{x,0,2*y-y^2},{z,-Sqrt[1-x],1-y}]+
Integrate[1,{y,0,1},{x,2*y-y^2,1},{z,-Sqrt[1-x],
Sqrt[1-x]}]
\end{Verbatim}
\end{tcolorbox}
\vspace{0.7em}
\noindent


Si agregamos la condición de que la densidad de $E$ en cualquier punto $(x,y,z)$ es
dada por $\delta(x,y,z)=x^2y$, entonces su masa estará dada por
\begin{eqnarray*}
\ds \int_{-1}^{1}\int_{0}^{1-z}\int_{0}^{1-z^2}x^2ydxdydz=\ds\int_{0}^2\int_{-1}^{1-y}\int_{0}^{1-z^2}x^2ydxdzdy&=&\frac{32}{189}
%\\ \ds \int_{-1}^{1}\int_{0}^{1-z^2}\int_{0}^{1-z}x^2ydydxdz= \ds \int_{0}^{1}\int_{-\sqrt{1-x}}^{\sqrt{1-x}}\int_{0}^{1-z}x^2ydydzdx&=& \frac{32}{189}\\ \ds \int_{0}^{1}\int_{1-\sqrt{1-x}}^{1+\sqrt{1-x}}\int_{-\sqrt{1-x}}^{1-y}x^2ydzdydx+\int_{0}^{1}\int_{0}^{1-\sqrt{1-x}}\int_{-\sqrt{1-x}}^{\sqrt{1-x}}x^2ydzdydx &=& \frac{32}{189} \\ \ds \int_{0}^2\int_{0}^{2y-y^2}\int_{-\sqrt{1-x}}^{1-y}x^2ydzdxdy+ \int_{0}^{1}\int_{2y-y^2}^{1}\int_{-\sqrt{1-x}}^{\sqrt{1-x}}x^2ydzdxdy&=& \frac{32}{189}
\end{eqnarray*}
El cálculo del volumen y la masa también se obtienen como sigue.




%\begin{tcolorbox}[title=\bfseries Parametrizando el sólido]
\noindent

En el plano $xz$, sobre la curva $x=1-z^2$ consideramos dos puntos {\small $P(u,\sqrt{1-u})$} y {\small $Q(u,-\sqrt{1-u});$}
Formando segmentos rectilíneos de la forma {\small $\left(1-v\right)P +vQ$,} con {\small $v\in [0,1];$} así, se simplifica como {\small $\left(u,\sqrt{1-u}\left(2v-1\right)\right),$} para {\small $u \in [0,1], \,\, v \in [0,1].$}


\vskip0.3cm
\noindent
\begin{minipage}{\textwidth}
	\centering
	\captionof{figure}{Región delimitada por $0 \leq x \leq 1 - z^2$}
	\label{fig:region_parabolica_xz}
	\begin{minipage}{0.38\linewidth}
		\centering
		\vspace{-18mm}
		\includegraphics[width=\linewidth]{Fig/420.pdf}
	\end{minipage}
	
	\vspace{-10mm}
	\fuentepropia
\end{minipage}
\vskip0.2cm


Formando segmentos rectilíneos desde la base del sólido al plano con ecuación $z=1-y$;
esto es, una combinación convexa de la forma
$(1-w)(u,0,\sqrt{1-u}( 2v-1))+w(u,1-\sqrt{1-u}\left( 2v-1\right) ,\sqrt{1-u}\left( 2v-1\right)),$
que al simplificarse nos permite definir la función:

%se lleva a la forma
%$ \left( u,w\left( 1+\sqrt{1-u}\right) -2wv\sqrt{1-u},\left( 2v-1\right) \sqrt{1-u}\right) , $
\vspace{-6mm}
\[{\bf r}(u,v,w)= \left\langle u,w\left( 1+\sqrt{1-u}\right) -2wv\sqrt{1-u},\left( 2v-1\right) \sqrt{1-u}\right\rangle,\]
%\vskip0.2cm
la cual parametriza el sólido, para $u\in [0,1],\,v\in [0,1],\,w\in [0,1].$
Con lo anterior se obtiene:
\begin{gather*}
	\mathbf{r}_{u}=\left\langle 1,\frac{w}{2}\frac{2v-1}{\sqrt{1-u}},\frac{1-2v}{2\sqrt{1-u}} \right\rangle, \\
	\mathbf{r}_{v} = \left\langle 0,-2w\sqrt{1-u},2\sqrt{1-u}\right\rangle, \\
	\mathbf{r}_{w} = \left\langle 0,1+\left( 1-2v\right) \sqrt{1-u},0\right\rangle.
\end{gather*}

%\vskip0.2cm
Además, $ \left| {\bf r}_{u}\cdot \left( {\bf r}_{v}\times {\bf r}_{w}\right) \right|
= \left| -2\sqrt{1-u}+2\left( 2v-1\right) \left( 1-u\right) \right| .$
\vskip0.3cm
El volumen de $E$ está dado por
\[\ds
\int_{0}^{1}\int_{0}^{1}\int_{0}^{1}\left| -2\sqrt{1-u}+2\left( 2v-1\right)
\left( 1-u\right) \right| dudvdw= \frac{4}{3}.\]
%\end{tcolorbox}
%\vskip0.2cm

Las siguientes instrucciones permiten evaluar los cálculos anteriores.

%\begin{tcolorbox}
Si la densidad del sólido está dada por {\footnotesize $\delta \left( x,y,z\right) =x^2y,$} entonces,
{\footnotesize $\delta \left( u,v,w\right) =u^2w\left( \left( 1+\sqrt{1-u}\right) -2v\sqrt{1-u}\right);$}
por lo tanto, la masa del sólido esta dada por
{\footnotesize $\ds \int_{0}^{1}\int_{0}^{1}\int_{0}^{1} 2 u^2 w\sqrt{1-u} \left((1-2v)\sqrt{1-u}+1\right)^2 dudvdw= \frac{32}{189}.$}
%\end{tcolorbox}
\vskip0.4cm
Las siguientes instrucciones permiten evaluar los cálculos anteriores.


\vspace{0.6em}
\noindent
\begin{tcolorbox}[breakable,enhanced,title=\bf Instrucción en \verb"Mathematica", top=2pt,bottom=2pt,boxsep=0pt,before skip=2pt,after skip=0pt]
\begin{Verbatim}[formatcom=\letraverbatim]
{x,z}=(1-v)*{u,Sqrt[1-u]}+v*{u,-Sqrt[1-u]};
r={x,y,z}=Simplify[(1-w)*{x,0,z}+w*{x,1-z,z}]
A=Abs[Dot[D[r,w],Cross[D[r,u],D[r,v]]]];
Integrate[A,{u,0,1},{v,0,1},{w,0,1}]
Integrate[x^2*y*A,{u,0,1},{v,0,1},{w,0,1}]
\end{Verbatim}
\end{tcolorbox}
\vspace{0.7em}
\noindent


\section{Sólidos acotados por cilindros} \index{Sólidos!con cilindros}
\vskip0.4cm
\subsection*{Dos cilindros y dos planos} \label{2cilindros2planos}
Las figuras de las ecuaciones $z=10-x-y^2/8$, $z=0$, $x^3=y$ y $x=0,$ acotan el sólido que se muestra a continuación.
\vskip0.2cm
\noindent
\begin{minipage}{\textwidth}
	\centering
	\captionof{figure}{Visualización del sólido definido por $0 \leq z \leq 10 - x - \frac{y^2}{8}$ \quad y \quad $x^3 \leq y \leq \sqrt{8(10 - x)}$.}
	\label{fig:solido_11_xyz}
	\begin{minipage}{0.20\linewidth}
		\centering
		\includegraphics[width=\linewidth]{Img/11d.pdf}
	\end{minipage}
	\hspace{1em} 
	\begin{minipage}{0.20\linewidth}
		\centering
		\includegraphics[width=\linewidth]{Img/11c.pdf}
	\end{minipage}
	\hspace{1em} 
	\begin{minipage}{0.20\linewidth}
		\centering
		\includegraphics[width=\linewidth]{Img/11b.pdf}
	\end{minipage}
	
	\vspace{2mm}
	\fuentepropia
\end{minipage}
\vskip0.2cm

Las siguientes son instrucciones que se requieren para obtener la figura:

\vspace{0.6em}
\noindent
\begin{tcolorbox}[breakable,enhanced,title=\bf Instrucción en \verb"Mathematica", top=2pt,bottom=2pt,boxsep=0pt,before skip=2pt,after skip=0pt]
\begin{Verbatim}[formatcom=\letraverbatim]
ContourPlot3D[{z==10-x-y^2/8,y==x^3,z==0,x==0},
{x,0,2},{y,0,9},{z,0,10},ContourStyle->{Orange,Cyan,
LightGray,Yellow},Mesh->False]
RegionPlot3D[0<=z<=10-x-y^2/8&&x^3<=y<=Sqrt[8*(10-x)]
&&0<=x<=2,{x,0,2},{y,0,9},{z,0,10},PlotPoints->90,
Mesh->False]
\end{Verbatim}
\end{tcolorbox}
\vspace{0.7em}
\noindent


Para hallar la curva de intersección de las superficies, se debe observar que al reemplazar
$y=x^{3},$ en la expresión $z=8-x-y^2/8,$ se obtiene $z=8-x-x^{6}/8$ o $z=8-\sqrt[3]{y}-y^2/8.$

\vskip0.2cm
\noindent
\begin{minipage}{\textwidth}
	\centering
	\captionof{figure}{Proyecciones del sólido en los planos coordenados.}
	\label{figura115}
	\begin{minipage}{0.27\linewidth}
		\centering
		\vspace{-10mm}
		\includegraphics[width=\linewidth]{Fig/421.pdf}
	\end{minipage}
	\hspace{1em}
	\begin{minipage}{0.27\linewidth}
		\centering
		\vspace{-10mm}
		\includegraphics[width=\linewidth]{Fig/422.pdf}
	\end{minipage}
	\hspace{1em}
	\begin{minipage}{0.27\linewidth}
		\centering
		\vspace{-10mm}
		\includegraphics[width=\linewidth]{Fig/423.pdf}
	\end{minipage}
	
	\vspace{-4mm}
	\fuentepropia
\end{minipage}
\vskip0.2cm

En el plano $yz$, la proyección del sólido se divide en dos regiones y en el plano $xy$ la región de integración se determina por las ecuaciones $y=x^3,$ $x+y^2/8=10$ y $x=0;$ el volumen del sólido está dado por
\vspace{-10mm}

{\small
\begin{multline*}
\int_0^2\int_{x^3}^{\sqrt{8 (10-x)}}\int_0^{10-x-\frac{y^2}{8}}dzdydx=
\bigintsss _0^8\bigintsss _0^{\sqrt[3]{y}}\bigintsss _0^{10-x-\frac{y^2}{8}}dzdxdy \\
+\bigintsss_8^{\sqrt{80}}\bigintsss_0^{10-\frac{y^2}{8}}\bigintsss_0^{10-x-\frac{y^2}{8}}dzdxdy=\frac{320 \sqrt{5}}{3}-\frac{2488}{15}
\end{multline*}
}

{\footnotesize
\[\bigintsss_0^2\bigintsss_0^{10-x-\frac{x^6}{8}}\bigintsss_{x^3}^{\sqrt{8(10-x-z)}}dydzdx=\frac{320\sqrt{5}}{3}-\frac{2488}{15}.\]
}

Al integrar primero con respecto a $x$ se debe observar que en una región esta recorre los puntos desde el plano $yz$ hasta el cilindro $x=\sqrt[3]{y},$ y en otro caso hasta la superficie {\small $x=10-z-y^2/8.$}
\vspace{-3mm}
{\footnotesize
\begin{multline*}
\bigintsss _0^8\bigintsss _0^{10-\sqrt[3]{y}-\frac{y^2}{8}}\bigintsss _0^{\sqrt[3]{y}}dxdzdy+\bigintsss _0^8\bigintsss_{10-\sqrt[3]{y}-\frac{y^2}{8}}^{10-\frac{y^2}{8}}\bigintsss _0^{10-z-\frac{y^2}{8}}dxdzdy\\ +\bigintsss_8^{\sqrt{80}}\bigintsss
_0^{10-\frac{y^2}{8}}\bigintsss _0^{10-z-\frac{y^2}{8}}dxdzdy = \frac{8}{15} \left(200 \sqrt{5}-311\right).
\end{multline*}
}

%\vskip0.2cm
\vspace{-2mm}
Estos cálculos se evalúan en \verb"Mathematica" con las siguientes instrucciones.


%\vspace{0.6em}

\noindent
\begin{tcolorbox}[breakable,enhanced,title=\bf Instrucción en \verb"Mathematica", top=2pt,bottom=2pt,boxsep=0pt,before skip=2pt,after skip=0pt]
\begin{Verbatim}[formatcom=\letraverbatim]
Integrate[1,{x,0,2},{y,x^3,Sqrt[8*(10-x)]},
{z,0,10-x-y^2/8}]
Integrate[1,{y,0,8},{x,0,y^(1/3)},
{z,0,10-x-y^2/8}]+Integrate[1,{y,8,Sqrt[80]},
{x,0,10-y^2/8},{z,0,10-x-y^2/8}]
Integrate[1,{x,0,2},{z,0,10-x-x^6/8},{y,x^3,
Sqrt[8*(10-x-z)]}]
Integrate[1,{y,0,8},{z,0,10-y^(1/3)-y^2/8},
{x,0,y^(1/3)}]+Integrate[1,{y,0,8},{z,10-y^(1/3)
y^2/8,10-y^2/8},{x,0,10-z-y^2/8}]+Integrate[1,
{y,8,Sqrt[80]},{z,0,10-y^2/8},{x,0,10-z-y^2/8}]
\end{Verbatim}
\end{tcolorbox}
\vspace{0.7em}
\noindent


\subsection*{Cálculo del centro de masa}

Si la densidad del sólido en cualquier punto {\small $(x,y,z)$} está dada por {\small $\delta(x,y,z)=xy$,} entonces la masa del sólido es inversa.

{\small
\[M=\int _0^2\int {x^3}^{\sqrt{8 (10-x)}}\int 0^{10-x-\frac{y^2}{8}}xydzdydx=\frac{13000}{63}.\]
}

Los momentos con respecto a los planos coordenados están dados por:

\vspace{-5mm}
{\footnotesize
\begin{align*}
	\mu_{yz} &= \int _0^2\int _{x^3}^{\sqrt{8 (10-x)}}\int _0^{10-x-\frac{y^2}{8}} x^2y\,dz\,dy\,dx \\
	&= \frac{9952}{45} \\[1ex]
	\mu_{xz} &= \int _0^2\int _{x^3}^{\sqrt{8 (10-x)}}\int _0^{10-x-\frac{y^2}{8}} xy^2\,dz\,dy\,dx \\
	&= \frac{5120 \left(18700 \sqrt{5}-34721\right)}{35343} \\[1ex]
	\mu_{xy} &= \int _0^2\int _{x^3}^{\sqrt{8 (10-x)}}\int _0^{10-x-\frac{y^2}{8}} xyz\,dz\,dy\,dx. \\
	&= \frac{177584}{315}
\end{align*}
}

\vskip0.2cm
Con los elementos anteriores se obtienen las coordenadas del centro de masa:
$(\overline{x},\overline{y},\overline{z})=\left(\frac{8708}{8125},\frac{512 \sqrt{5}}{39}-\frac{4444288}{182325},\frac{22198}{8125}\right)$.
Este proceso se hace en \verb"Mathematica" como sigue:


\vspace{0.6em}
\noindent
\begin{tcolorbox}[breakable,enhanced,title=\bf Instrucción en \verb"Mathematica", top=2pt,bottom=2pt,boxsep=0pt,before skip=2pt,after skip=0pt]
\begin{Verbatim}[formatcom=\letraverbatim]
M=Integrate[x*y,{x,0,2},{y,x^3,Sqrt[8*(10-x)]},
{z,0,10-x-y^2/8}];
Mxy=Integrate[x*y*z,{x,0,2},{y,x^3,Sqrt[8*(10-x)]},
{z,0,10-x-y^2/8}];
Mxz=Integrate[x*y^2,{x,0,2},{y,x^3,Sqrt[8*(10-x)]},
{z,0,10-x-y^2/8}];
Myz=Integrate[x^2*y,{x,0,2},{y,x^3,Sqrt[8*(10-x)]},
{z,0,10-x-y^2/8}];
Simplify[{Myz/M,Mxz/M,Mxy/M}]
\end{Verbatim}
\end{tcolorbox}
\vspace{0.7em}
\noindent

El volumen también se puede obtener con el siguiente procedimiento.

%\begin{tcolorbox}[title=\bfseries Parametrizando el sólido]
%\small
La proyección del sólido en el plano $xy$, ver figura (\ref{figura115}),
se parametriza considerando dos puntos
{\footnotesize $A(t,t^{3})$} y {\footnotesize $B\left( t,\sqrt{8\left( 10-t\right) }\right) $} para {\footnotesize $t\in\lbrack 0,2].$}
Simplificando la expresión {\footnotesize $(1-u)A+uB$}, se obtiene
{\footnotesize $\left\langle t,t^{3}(1-u)+u\sqrt{8\left(10-t\right) }\right\rangle.$}
\vskip0.2cm
La coordenada $z,$ sobre la superficie $z=10-x-\frac{y^{2}}{8},$
que abreviaremos como $\overline{z}$, satisface que
{\footnotesize $\overline{z}=10-t-\frac{1}{8}\left( t^{3}(1-u)+u\sqrt{8\left( 10-t\right) }\right) ^{2}.$}
\vskip0.2cm

Con dos puntos de la forma
{\small \[P\left( t,t^{3}(1-u)+u\sqrt{8\left( 10-t\right) },0\right) \]} y
{\small \[Q\left( t,t^{3}(1-u)+u\sqrt{8\left( 10-t\right) },\overline{z}\right)\]}
se forma la combinación $(1-v) P+vQ,$ que define la función
%\vskip0.2cm
{\small
\[{\bf r}(t,u,v)=\left\langle t,t^{3}(1-u)+u\sqrt{8(10-t) },v \,\overline{z} \right\rangle,\]}
%\vskip0.2cm
que parametriza el sólido, con $t\in [0,2],\,u\in [0,1],\,v\in [0,1].$
	% Ahora bien, %\begin{eqnarray*}
	%{\bf r}_{t} &=& \bigg\langle 1,3t^{2}(1-u) -\frac{\sqrt{2}u}{\sqrt{10-t}},\\ && v\left( \frac{1}{4}\left( t^{3}(1-u) +u\sqrt{8(10-t)}\right) \left( 3t^{2}(u-1) +\frac{\sqrt{2}u}{\sqrt{10-t}}\right) -1\right) \bigg\rangle \\
	%{\bf r}_{u} &=&\left\langle 0,\sqrt{8(10-t)}-t^{3},\frac{v}{4}\left( t^{3}(1-u) +u\sqrt{8(10-t)}\right) \left( t^{3}-\sqrt{8(10-t)}\right) \right\rangle \\ {\bf r}_{v} &=&\left\langle 0,0,10-\frac{1}{8}\left( t^{3}(u-1) -u\sqrt{8(10-t)}\right) ^{2}-t\right\rangle
	%\end{eqnarray*}
El elemento diferencial de volumen
$\left\vert {\bf r}_{t}\cdot \left({\bf r}_{u}\times {\bf r}_{v}\right) \right\vert$
esta dado por
\vspace{-4mm}
%\vskip0.2cm

{\footnotesize
\[\bigg\vert \dfrac{(u-1)(t^{3}-\sqrt{8(10-t)})}{8}\bigg(8(u+1)(10-t)
+t^6(u-1)-2t^{3}u\sqrt{8(10-t)}\bigg) \bigg\vert.\]
}

%\vskip0.2cm
Su integración determina el volumen del sólido
%\vskip0.2cm
{\footnotesize
\begin{equation*}
\int_{0}^{1}\int_{0}^{1}\int_{0}^{2}\left\vert {\bf r}_{t}\cdot \left({\bf r}_{u}\times {\bf r}_{v}\right) \right\vert dtdudv=\frac{320}{3}\sqrt{5}-\frac{2488}{15}.
\end{equation*}
}
%\vskip0.2cm
%\end{tcolorbox}


Con las siguientes indicaciones se pueden efectuar estos cálculos.
\vskip0.2cm

%\vspace{0.6em}
\noindent
\begin{tcolorbox}[breakable,enhanced,title=\bf Instrucción en \verb"Mathematica", top=2pt,bottom=2pt,boxsep=0pt,before skip=2pt,after skip=0pt]
\begin{Verbatim}[formatcom=\letraverbatim]
{x,y}=(1-u)*{t,t^3}+u*{t,Sqrt[8*(10-t)]};
r=Simplify[{x,y,w*(10-x-y^2/8)}];
A=Abs[Dot[D[r,w],Cross[D[r,u],D[r,t]]]];
Integrate[A,{t,0,2},{u,0,1},{w,0,1}]
\end{Verbatim}
\end{tcolorbox}
\vspace{0.7em}
\noindent



\begin{ejer}
Plantear y evaluar las integrales que determinan el volumen del sólido $E$ de la sección \secref{2cilindros2planos}
En los órdenes $dydxdz$ y $dxdydz.$
\end{ejer}

\subsection*{Dos cilindros}
Las ecuaciones $x^2+z^2=16$ y $2x^2+y^2=8x$ representan un cilindro circular y un cilindro elíptico. Estas superficies acotan en el primer octante, el sólido que se muestra enseguida.


\vskip0.4cm
\noindent
\begin{minipage}{\textwidth}
	\centering
	\captionof{figure}{Visualización del sólido definido por $x^2 + z^2 \leq 16$ \quad y \quad $2x^2 + y^2 \leq 8x$}
	\label{fig:solido_cilindro_elipse}
	\begin{minipage}{0.22\linewidth}
		\centering
		\includegraphics[width=\linewidth]{Img/29a.pdf}
	\end{minipage}
	\begin{minipage}{0.22\linewidth}
		\centering
		\includegraphics[width=\linewidth]{Img/29b.pdf}
	\end{minipage}
	\begin{minipage}{0.22\linewidth}
		\centering
		\includegraphics[width=\linewidth]{Img/29c.pdf}
	\end{minipage}
	\begin{minipage}{0.22\linewidth}
		\centering
		\includegraphics[width=\linewidth]{Img/29d.pdf}
	\end{minipage}
	
	\vspace{0.3em}
	\fuentepropia
\end{minipage}
\vskip0.4cm



En \verb"Mathematica", estas figuras se consiguen con las siguientes sentencias:

\vspace{0.6em}
\noindent
\begin{tcolorbox}[breakable,enhanced,title=\bf Instrucción en \verb"Mathematica", top=2pt,bottom=2pt,boxsep=0pt,before skip=2pt,after skip=0pt]
\begin{Verbatim}[formatcom=\letraverbatim]
ContourPlot3D[{16==x^2+z^2,2*x^2+y^2==8*x},{x,0,4},
{y,0,3},{z,0,4},PlotPoints->60,Ticks->{{0,2,4},
{0,1,2},{0,2,4}},Mesh->False,ContourStyle->{Orange,
Blend[{Yellow,Pink,Green}]}]
RegionPlot3D[0<z<Sqrt[16-x^2]&&0<y<Sqrt[8*x-2*x^2]&&0
<x<4,{x,0,4},{y,0,3},{z,0,4},PlotPoints->90,
PlotStyle->Directive[Orange,Specularity[White,40]],
Mesh->False]
\end{Verbatim}
\end{tcolorbox}
\vspace{0.7em}
\noindent


Para establecer con claridad los límites en las integrales triples respectivas, que determinan el volumen del sólido, se verifica que en el plano $xz$ la región es un cuarto del círculo centrado en el origen y radio $4.$
Entre tanto, en el plano $xy$, la expresión $2x^2+y^2=8x$ equivale a $2(x-2)^2+y^2=8$ y representa una elipse con centro en el punto $(2,0).$

\vskip0.4cm
\noindent
\begin{minipage}{\textwidth}
	\centering
	\captionof{figure}{Proyecciones en los planos coordenados.}
	\label{fig:proyecciones_cilindro_elipse_curva}
	\begin{minipage}{0.25\linewidth}
		\centering
		\vspace{-18mm}
		\includegraphics[width=\linewidth]{Fig/424.pdf}
	\end{minipage}
	\hspace{1em}
	\begin{minipage}{0.25\linewidth}
		\centering
		\vspace{-18mm}
		\includegraphics[width=\linewidth]{Fig/425.pdf}
	\end{minipage}
	\hspace{1em}
	\begin{minipage}{0.38\linewidth}
		\centering
		\vspace{-18mm}
		\includegraphics[width=\linewidth]{Fig/426.pdf}
	\end{minipage}
	
	\vspace{-12mm}
	\fuentepropia
\end{minipage}
\vskip0.2cm

Con lo anterior, presentamos las integrales triples según el orden de integración señalado.
\vskip0.2cm
Al integrar primero con respecto a $x,$ se debe tener en cuenta las dos regiones proyectadas en el plano $yz$, en un caso se parte de la parte negativa del cilindro elíptico y se llega hasta la parte positiva del mismo cilindro.
Mientras que en la otra región $x$ recorre los puntos que van desde la parte negativa del cilindro elíptico hasta el cilindro circular.
De la expresión $x^2+z^2=16$, se sigue que {\footnotesize $x=\pm \sqrt{16-z^2}.$} Del reemplazo en la otra ecuación se obtiene
{\footnotesize $16-z^2+2y^2=4\sqrt{16-z^2},$} que equivale a
{\footnotesize $\left( z^2-y^2\right) ^2=16\left( z^2+2y^2\right).$}
De esta expresión se obtiene {\footnotesize $z=\sqrt{8+\frac{y^2}{2}\pm 2\sqrt{16-2y^2}}.$}
% \begin{eqnarray*} z &=& \frac{1}{2}\sqrt{32+2y^2\pm 8\sqrt{16-2y^2}} =\sqrt{8+\frac{y^2}{2}\pm 2\sqrt{16-2y^2}} \end{eqnarray*}
Con esto, el volumen del sólido está dado por
\vspace{-4mm}

{\small
\begin{multline*}
\bigintsss_{0}^{\sqrt{8}}\bigintsss_{0}^{\sqrt{8+\frac{y^2}{2}-2\sqrt{16-2y^2}}}\bigintsss_{2-\frac{1}{2}\sqrt{16-2y^2}}^{2+\frac{1}{2}\sqrt{16-2y^2}}dxdzdy \\
+\bigintsss_{0}^{\sqrt{8}}\bigintsss_{\sqrt{8+\frac{y^2}{2}-\sqrt{64-8y^2}}}^{\sqrt{8+\frac{y^2}{2}+\sqrt{64-8y^2}}}\bigintsss_{2-\frac{1}{2}\sqrt{16-2y^2}}^{\sqrt{16-z^2}}dxdzdy = \frac{176}{3}-12\sqrt{2}\ln( 3+2\sqrt{2}).
\end{multline*}
}
\vskip0.2cm
De la expresión $\left( z^2-y^2\right) ^2=16\left( z^2+2y^2\right),$ se puede obtener \\
$y=\sqrt{2z^2-32\pm 8\sqrt{16-z^2}},$ entonces tenemos otra forma de hallar el mismo volumen.
\begin{multline*}
\bigintsss_{0}^{2\sqrt{3}}\bigintsss_{\sqrt{2z^2-32+8\sqrt{16-z^2}}}^{\sqrt{8}}\bigintsss_{2-\frac{1}{2}\sqrt{16-2y^2}}^{2+\frac{1}{2}\sqrt{16-2y^2}}dxdydz \\
+\bigintsss_{0}^{4}\bigintsss_{0}^{\sqrt{2z^2-32+8\sqrt{16-z^2}}}\bigintsss_{2-\frac{1}{2}\sqrt{16-2y^2}}^{\sqrt{16-z^2}}dxdydz \approx 28.7519.
\end{multline*}
\vskip0.2cm
Se debe observar que de la expresión $2( x-2) ^2+y^2=8$ se sigue que $x=2\pm \frac{1}{2}\sqrt{16-2y^2}$,
por tanto, el volumen del sólido esta dado por
\vspace{-1mm}

{\small
\[
\int_{\scriptscriptstyle 0}^{\scriptscriptstyle 4}\int_{\scriptscriptstyle 0}^{\scriptscriptstyle\sqrt{8x-2x^2}}\int_{\scriptscriptstyle 0}^{\scriptscriptstyle\sqrt{16-x^2}}dzdydx \! = \! \int_{\scriptscriptstyle 0}^{\scriptscriptstyle\sqrt{8}}\int_{\scriptscriptstyle 2-\frac{1}{2}\sqrt{16-2y^2}}^{\scriptscriptstyle 2+\frac{1}{2}\sqrt{16-2y^2}}\int_{\scriptscriptstyle 0}^{\scriptscriptstyle\sqrt{16-x^2}}dzdxdy \! \approx \! 28.75.
\] 
}
\\ 
Para el orden $dydxdz$, $y$ recorre los puntos desde el plano $xz$ hasta el cilindro elíptico; en este caso, las integrales correspondientes son las siguientes.

\vspace{-4mm}
{\small
\[\int_{0}^{4}\int_{0}^{\sqrt{16-z^2}}\int_{0}^{\sqrt{8x-2x^2}}dydxdz=\int_{0}^{4}\int_{0}^{\sqrt{16-x^2}}\int_{0}^{\sqrt{8x-2x^2}}dydzdx \approx 28.75.\]
}

\vskip0.2cm
Una de las integrales anteriores se evaluó numéricamente en \verb"Mathematica"
ejecutando la siguiente instrucción.

\vspace{0.6em}
\noindent
\begin{tcolorbox}[breakable,enhanced,title=\bf Instrucción en \verb"Mathematica", top=2pt,bottom=2pt,boxsep=0pt,before skip=2pt,after skip=0pt]
\begin{Verbatim}[formatcom=\letraverbatim]
NIntegrate[1,{y,0,Sqrt[8]},{z,0,Sqrt[8+y^2/2
Sqrt[64-8*y^2]]},{x,2-Sqrt[16-2*y^2]/2,2+
Sqrt[16-2*y^2]/2}]+NIntegrate[1,{y,0,Sqrt[8]},
{z,Sqrt[8+y^2/2-Sqrt[64-8*y^2]],Sqrt[8+y^2/2+
Sqrt[64-8*y^2]]},{x,2-Sqrt[16-2*y^2]/2,Sqrt[16-z^2]}]
\end{Verbatim}
\end{tcolorbox}
\vspace{0.7em}
\noindent


La otra alternativa para calcular el volumen se muestra a continuación.
\vskip0.2cm

%\begin{tcolorbox}[title=\bfseries Parametrizando el sólido]
\noindent
La proyección del sólido en el plano $xz$ se describe mediante las ecuaciones
$x=4u\cos(t),\,\, z=4u\sin(t),$ para $u\in [0,1]$
y $t\in [0,\pi/2 ].$ De estas expresiones se obtiene $y$ en la ecuación del cilindro parabólico, es decir $y = \sqrt{8x-2x^2} =4\sqrt{2}\sqrt{u(1-u\cos(t)) \cos(t)}$.

\vskip0.2cm
\noindent
\begin{minipage}{\textwidth}
	\centering
	\captionof{figure}{Visualización del sólido definido por $x^2 + z^2 = 16$, \quad $2x^2 + y^2 = 8x$}
	\label{fig:solido_cilindro_elipse_3D_2}
	\begin{minipage}{0.35\linewidth}
		\centering
		\vspace{-15mm}
		\includegraphics[width=\linewidth]{Fig/427.pdf}
	\end{minipage}
	
	\vspace{-8mm}
	\fuentepropia
\end{minipage}
\vskip0.2cm


%\vskip0.8cm

El sólido se parametriza considerando un punto
%\vskip0.2cm 
\[P(4u\cos(t),0,4u\sin(t).\]
%\vskip0.2cm 
en el plano $xz$
y otro punto sobre el cilindro elíptico
%\vskip0.2cm
\[Q(4u\cos(t),4\sqrt{2}\sqrt{u(1-u\cos(t))\cos(t)},4u\sin(t)).\]
%\vskip0.2cm
Simplificando la combinación convexa $(1-w)P+wQ$, se define la siguiente función:
%\vskip0.2cm
\[{\bf r}(t,u,w) = \langle 4u\cos(t),4w\sqrt{2}\sqrt{u(1-u\cos(t))\cos(t)},4u\sin(t)\rangle,\]
%\vskip0.2cm
con $u\in [0,1],$ $t\in [ \pi/2 ]$ y $w\in [0,1]$.
% Ahora hallamos los siguientes vectores \begin{eqnarray*}
%{\bf r}_{t} &=& \langle -4u\sin(t),2w\frac{\sqrt{2}u\sin(t)(2u\cos(t)-1)}{\sqrt{u(1-u\cos(t))\cos(t)}},4u\cos(t)\rangle \\
%{\bf r}_{u} &=&\langle 4\cos(t),2w\frac{\sqrt{2}\cos(t)(1-2u\cos(t))}{\sqrt{(u(1-u\cos(t))\cos t))}},4\sin(t)\rangle \\
%{\bf r}_{w} &=&\langle 0,4\sqrt{2}\sqrt{u(1-u\cos(t))\cos(t)},0 \rangle
%\end{eqnarray*}
% además ${\bf r}_{t}\times {\bf r}_{u}=\langle -8w\sqrt{2}\sqrt{u}\frac{2u\cos(t)-1}{\sqrt{(1-u\cos(t))\cos(t)}},16u,0\rangle $
El triple producto vectorial proporciona el elemento diferencial de volumen,
%\vskip0.2cm
\[|{\bf r}_{w}\cdot({\bf r}_{t}\times {\bf r}_{u})|=64\sqrt{2}u\sqrt{u(1-u\cos (t))\cos(t)},\]
%\vskip0.2cm
que se integra y nos proporciona el volumen del sólido
\[\int_{0}^{1}\int_{0}^{1}\int_{0}^{\pi/2}64\sqrt{2}u\sqrt{u(1-u\cos(t))\cos(t)}dtdudw \approx 28.752.\]
%\end{tcolorbox}
\vskip0.2cm
Este cálculo se obtiene de manera numérica en \verb"Mathematica" ejecutando las siguientes instrucciones:


\vspace{0.6em}
\noindent
\begin{tcolorbox}[breakable,enhanced,title=\bf Instrucción en \verb"Mathematica", top=2pt,bottom=2pt,boxsep=0pt,before skip=2pt,after skip=0pt]
\begin{Verbatim}[formatcom=\letraverbatim]
{x,z,y}={4*u*Cos[t],4*u*Sin[t],Sqrt[8*x-2*x^2]};
r=Simplify[(1-w)*{x,0,z}+w*{x,y,z}]
B=Simplify[Abs[Dot[D[r,w],Cross[D[r,t],D[r,u]]]]]
NIntegrate[B,{t,0,Pi/2},{u,0,1},{w,0,1}]
\end{Verbatim}
\end{tcolorbox}
\vspace{0.7em}
\noindent


\subsection*{Dos cilindros y un plano}
Hallar el volumen del sólido acotado por los gráficos de las ecuaciones.
$z=1-x^2,$ $x^2+y^2=1$ y $z=0$. Las siguientes figuras ilustran esta situación.


\vskip0.4cm
\noindent
\begin{minipage}{\textwidth}
	\centering
	\captionof{figure}{Visualización del sólido definido por $x^2 + y^2 \leq 1$ \quad y \quad $0 \leq z \leq 1 - x^2$.}
	\label{fig:solido_disco_parabolico}
	\begin{minipage}{0.24\linewidth}
		\centering
		\vspace{-8mm}
		\includegraphics[width=\linewidth]{Img/85c.pdf}
	\end{minipage}
	\hspace{1em}
	\begin{minipage}{0.28\linewidth}
		\centering
		\vspace{-8mm}
		\includegraphics[width=\linewidth]{Fig/428.pdf}
	\end{minipage}
	\hspace{1em}
	\begin{minipage}{0.18\linewidth}
		\centering
		\vspace{-8mm}
		\includegraphics[width=\linewidth]{Img/85b.pdf}
	\end{minipage}
	
	\vspace{-4mm}
	\fuentepropia
\end{minipage}
\vskip0.2cm

Estas figuras se obtienen con la ayuda de las siguientes instrucciones:


\vspace{0.6em}
\noindent
\begin{tcolorbox}[breakable,enhanced,title=\bf Instrucción en \verb"Mathematica", top=2pt,bottom=2pt,boxsep=0pt,before skip=2pt,after skip=0pt]
\begin{Verbatim}[formatcom=\letraverbatim]
ContourPlot3D[{z==1-x^2,x^2+y^2==1},{x,-1,1},{y,-1,1},
{z,0,1},Mesh->False,ContourStyle->{Cyan,Orange}]
A=RegionPlot3D[0<z<1-x^2&&x^2+y^2<1,{x,-1,1},{y,-1,1},
{z,0,1},PlotStyle->RGBColor[0.8,0.2,0.2],Mesh->False,
PlotPoints->90,Axes->False,Boxed->False];
B=ParametricPlot3D[{{Cos[t],Sin[t],1-Cos[t]^2},
{Cos[t],Sin[t],0}},{t,0,2Pi},PlotStyle->{Gray,Gray}];
Show[A,B]
\end{Verbatim}
\end{tcolorbox}
\vspace{0.7em}
\noindent


\noindent
\begin{minipage}{\textwidth}
	\centering
	\captionof{figure}{Proyección del sólido en los planos coordenados.}
	\label{fig:proyeccion_parabolico_disco}
	\begin{minipage}{0.29\linewidth}
		\centering
		\vspace{-14mm}
		\includegraphics[width=\linewidth]{Fig/429.pdf}
	\end{minipage}
	\hspace{1em}
	\begin{minipage}{0.29\linewidth}
		\centering
		\vspace{-14mm}
		\includegraphics[width=\linewidth]{Fig/430.pdf}
	\end{minipage}
	\begin{minipage}{0.29\linewidth}
		\centering
		\vspace{-14mm}
		\includegraphics[width=\linewidth]{Fig/431.pdf}
	\end{minipage}
	
	\vspace{-7mm}
	\fuentepropia
\end{minipage}


\vskip0.2cm
En el plano $xy$ la ecuación $x^2+y^2=1$ determina una proyección del sólido, esto ayuda a plantear las integrales que permiten obtener el volumen.
\vspace{-2mm}
\[\int_{-1}^{1}\int_{-\sqrt{1-x^2}}^{\sqrt{1-x^2}}\int_{0}^{1-x^2}dzdydx= \int_{-1}^{1}\int_{-\sqrt{1-y^2}}^{\sqrt{1-y^2}}\int_{0}^{1-x^2}dzdxdy= \frac{3}{4}\pi.\]
%\vskip0.2cm
En el plano $yz,$ se satisface que $1-z+y^2=1,$ es decir $z=y^2$ cuya figura es una parábola. Por tal motivo, el volumen está dado por:
\begin{multline*}
\int_{0}^{1}\int_{-\sqrt{z}}^{\sqrt{z}}\int_{-\sqrt{1-z}}^{\sqrt{1-z}}dxdydz +\int_{0}^{1}\int_{-1}^{-\sqrt{z}}\int_{-\sqrt{1-y^2}}^{\sqrt{1-y^2}}dxdydz \\ +\int_{0}^{1}\int_{\sqrt{z}}^{1}\int_{-\sqrt{1-y^2}}^{\sqrt{1-y^2}}dxdydz
=\int_{-1}^1\int_{y^2}^1\int_{-\sqrt{1-z}}^{\sqrt{1-z}}dxdzdy \\ +\int_{-1}^{1}\int_0^{y^2}\int_{-\sqrt{1-y^2}}^{\sqrt{1-y^2}}dxdzdy=\frac{3}{4}\pi.
\end{multline*}

% \begin{multline*} \int_{0}^{1}\int_{-\sqrt{z}}^{\sqrt{z}}\int_{-\sqrt{1-z}}^{\sqrt{1-z}}dxdydz+\int_{0}^{1}\int_{-1}^{-\sqrt{z}}\int_{-\sqrt{1-y^2}}^{\sqrt{1-y^2}}%dxdydz \\ +\int_{0}^{1}\int_{\sqrt{z}}^{1}\int_{-\sqrt{1-y^2}}^{\sqrt{1-y^2}}dxdydz \approx 2.3562 \end{multline*}
En el plano $xz$, la proyección es una región acotada por una porción de parábola y en el plano $x$,
el volumen está dado por
%\vskip0.2cm
\[\int_{-1}^{1}\int_{0}^{1-x^2}\int_{-\sqrt{1-x^2}}^{\sqrt{1-x^2}}dydzdx = \int_{0}^{1}\int_{-\sqrt{1-z}}^{\sqrt{1-z}}\int_{-\sqrt{1-x^2}}^{\sqrt{1-x^2}}dydxdz \approx 2.3562.\]
%\vskip0.2cm
En coordenadas cilíndricas, esta integral se evalúa como sigue:
%\vskip0.2cm
\[\int_0^{2\pi} \int_0^1 \int_0^{1-r^2\cos^2 \theta}, r dzdrd \theta =\frac{3\pi}{4}.\]
%\vskip0.2cm
Algunas de estas integrales se evalúan mediante el uso de las siguientes instrucciones:

\vspace{0.6em}
\noindent
\begin{tcolorbox}[breakable,enhanced,title=\bf Instrucción en \verb"Mathematica", top=2pt,bottom=2pt,boxsep=0pt,before skip=2pt,after skip=0pt]
\begin{Verbatim}[formatcom=\letraverbatim]
Integrate[1,{x,-1,1},{y,-Sqrt[1-x^2],Sqrt[1-x^2]},
{z,0,1-x^2}]
Integrate[1,{y,-1,1},{x,-Sqrt[1-y^2],Sqrt[1-y^2]},
{z,0,1-x^2}]
Integrate[1,{y,-1,1},{z,y^2,1},{x,-Sqrt[1-z],
Sqrt[1-z]}]+Integrate[1,{y,-1,1},{z,0,y^2},
{x,-Sqrt[1-y^2],Sqrt[1-y^2]}]
Integrate[1,{x,-1,1},{z,0,1-x^2},{y,-Sqrt[1-x^2],
Sqrt[1-x^2]}]
Integrate[1,{z,0,1},{x,-Sqrt[1-z],Sqrt[1-z]},
{y,-Sqrt[1-x^2],Sqrt[1-x^2]}]
\end{Verbatim}
\end{tcolorbox}
\vspace{0.7em}
\noindent


El volumen también se puede obtener de la siguiente manera.

%\begin{tcolorbox}[title=\bfseries Parametrizando el sólido]
La base del sólido es un círculo unitario que se parametriza {\small $\left(u\cos \theta ,u\sin \theta ,0\right),$} sobre el cilindro parabólico se satisface que $z=1-x^2=1-u^2\cos ^2\theta .$ El sólido se parametriza formando una combinación convexa entre un punto en la base y otro punto sobre el cilindro parabólico, esto es
$ (1-w)(u\cos \theta ,u\sin \theta ,0) +w( u\cos \theta ,u\sin \theta ,1-u^2\cos ^2\theta);$
al simplificarse, nos permite definir la función
%\vskip0.2cm
{\small
\[{\bf r}\left( u,\theta ,w\right) =\langle u\cos \theta , u\sin \theta ,w(1-u^2\cos ^2\theta )\rangle,\]
}
%\vskip0.2cm
donde $\theta \in [0,2\pi]$, $u \in [0,1]$ y $w \in [0,1]$.

\vskip0.3cm
Con esta función hallamos los siguientes vectores
\begin{eqnarray*}
{\bf r}_{u} &=& \left\langle \cos \theta,\sin \theta, -2uw\cos ^2\theta \right\rangle \\
{\bf r}_{\theta } &=& \left\langle -u\sin \theta , u\cos \theta , 2u^2w\cos \theta \sin \theta \right\rangle \\
{\bf r}_{w} &=& \left\langle 0 , 0 , 1-u^2\cos ^2\theta \right\rangle
\end{eqnarray*}
El elemento diferencial de volumen está dado por

\[\left\vert {\bf r}_{u}\cdot \left( {\bf r}_{\theta }\times {\bf r}_{w}\right) \right\vert =\left\vert u\left( 1-u^2\cos ^2\theta \right) \right\vert,\]

entonces, el volumen del sólido esta dado por
\vspace{-2mm}
{\small
\[\int_{0}^{1}\int_{0}^{1}\int_{0}^{2\pi }u\left( 1-u^2\cos ^2\theta
\right) d\theta dudw= \frac{3\pi}{4}.\]
}
%\end{tcolorbox}

Este cálculo se hace en \verb"Mathematica" ejecutando las siguientes instrucciones.



\vspace{0.6em}
\noindent
\begin{tcolorbox}[breakable,enhanced,title=\bf Instrucción en \verb"Mathematica", top=2pt,bottom=2pt,boxsep=0pt,before skip=2pt,after skip=0pt]
\begin{Verbatim}[formatcom=\letraverbatim]
{x,y,z}={u*Cos[t],u*Sin[t],1-x^2};
r=(1-w)*{x,y,0}+w*{x,y,z}
A=Simplify[Abs[Dot[D[r,t],Cross[D[r,u],D[r,w]]]]]
[A,{t,0,2*Pi},{u,0,1},{w,0,1}]
\end{Verbatim}
\end{tcolorbox}
\vspace{0.7em}
\noindent


\subsection*{Tres cilindros y un plano}
Las ecuaciones $z=4-x^2,\,\, 2x=y^2$ y $4x=y^4-2y^2$ para $y \geq 0,$ acotan un sólido.
Las figuras de las expresiones que lo determinan y el sólido se muestran a continuación.




\vskip0.2cm
\noindent
\begin{minipage}{\textwidth}
	\centering
	\captionof{figure}{Visualización del sólido definido por $0 \leq z \leq 4 - x^2$ \quad y \quad $y^4 - 2y^2 \leq 4x \leq 2y^2$}
	\label{fig:solido_curvas_parabolicas}
	\begin{minipage}{0.20\linewidth}
		\centering
		\includegraphics[width=\linewidth]{Img/17c.pdf}
	\end{minipage}
	\begin{minipage}{0.20\linewidth}
		\centering
		\includegraphics[width=\linewidth]{Img/17d.pdf}
	\end{minipage}
	\begin{minipage}{0.20\linewidth}
		\centering
		\includegraphics[width=\linewidth]{Img/17b.pdf}
	\end{minipage}
	
	
	\vspace{0.3em}
	\fuentepropia
\end{minipage}
\vskip0.2cm

También se presentan las instrucciones que se requieren para graficarlo.


\vspace{0.6em}
\noindent
\begin{tcolorbox}[breakable,enhanced,title=\bf Instrucción en \verb"Mathematica", top=2pt,bottom=2pt,boxsep=0pt,before skip=2pt,after skip=0pt]
\begin{Verbatim}[formatcom=\letraverbatim]
ContourPlot3D[{z==4-x^2,2*x==y^2,4*x==y^4-2*y^2},
{x,-1,2},{y,0,2},{z,0,4},ContourStyle->{Orange,
Blend[{Yellow,Pink,Green}],LightGray},Mesh->None,
PlotPoints->60]
RegionPlot3D[0<z<4-x^2&&(y^4-2*y^2)/4<x<y^2/2,
{x,0,2},
{y,0,2},{z,0,4},PlotStyle->Blend[{Yellow,Orange,Red}],
PlotPoints->90,Mesh->None]
\end{Verbatim}
\end{tcolorbox}
\vspace{0.7em}
\noindent


La ecuación $z = 4-x^2$ en el plano $xz$ permite ver la proyección en dicho plano de la región de integración. Lo mismo sucede con las ecuaciones $4x = y^4 -2y^2$ y $2x = y^2$, que se dibujan en el plano $xy.$ En este caso $s = 4 - x^2$ equivale a $x = \pm\sqrt{4-z},$ mientras que $2x = y^2$ equivale a $y = \pm \sqrt{2x}$. La expresión $4x =y^4- 2y^2$ nos lleva a que $y = \pm \sqrt{1\pm \sqrt{1+ 4x}},$ es decir hay cuatro funciones que determinan la curva, con $x$ como variable independiente, es por esta razón que el gráfico posee diferentes colores para enfatizar dicha situación.


\vskip0.2cm
\noindent
\begin{minipage}{\textwidth}
	\centering
	\captionof{figure}{Proyecciones del sólido definido por $y^4 - 2y^2 \leq 4x \leq 2y^2$, \quad $0 \leq z \leq 4 - x^2$.}
	\label{fig:proyecciones_curvas_parabolicas}
	\begin{minipage}{0.58\linewidth}
		\centering
		\vspace{-32mm}
		\includegraphics[width=\linewidth]{Fig/432.pdf}
	\end{minipage}
	\hspace{1em}
	\begin{minipage}{0.32\linewidth}
		\centering
		\vspace{-32mm}
		\includegraphics[width=\linewidth]{Fig/433.pdf}
	\end{minipage}
	
	\vspace{-27mm}
	\fuentepropia
\end{minipage}


%\vskip0.4cm
Con ayuda de estos gráficos, podemos plantear las integrales triples que determinan el volumen del sólido.
\vspace{-1mm}
{\small
\begin{multline*}
\hskip-0.4cm \bigintss_{-\frac{1}{4}}^{0}\bigintss_{\sqrt{1-\sqrt{1+4x}}}^{\sqrt{1+\sqrt{1+4x}}}\bigintss_{0}^{4-x^2}dzdydx+\bigintss_{0}^2\bigintss_{\sqrt{2x}}^{\sqrt{1+\sqrt{1+4x}}}\bigintss_{0}^{4-x^2}dzdydx \\
\hskip-0.4cm \bigintss_{-\frac{1}{4}}^{0}\bigintss_{0}^{4-x^2}\bigintss_{\sqrt{1-\sqrt{1+4x}}}^{\sqrt{1+\sqrt{1+4x}}}dydzdx+\bigintss_{0}^2\bigintss_{0}^{4-x^2}\bigintss_{\sqrt{2x}}^{\sqrt{1+\sqrt{1+4x}}}dydzdx \\
= \bigintss_{0}^2 \bigintss_{\frac{y^{4}-2y^2}{4}}^{\frac{y^2}{2}}\bigintss_{0}^{4-x^2}dzdxdy= \frac{169472}{45045}.
\end{multline*}
}

En el plano $yz$ existen tres regiones en las que se divide la proyección del sólido; una de ellas es muy pequeña, está sombreada en negro y está
acotada por la recta $z=4$ y la curva $z=4- (y^4 -2y^2)^2/16,$ para $y \in [0, \sqrt{2}].$
\vskip0.2cm

\noindent
\begin{minipage}{\textwidth}
	\centering
	\captionof{figure}{Proyección del sólido en el plano $yz$.}
	\label{fig:proyeccion_yz_curvas}
	\begin{minipage}{0.45\linewidth}
		\centering
		\vspace{-20mm}
		\includegraphics[width=\linewidth]{Fig/434.pdf}
	\end{minipage}
	\hspace{1em}
	\begin{minipage}{0.45\linewidth}
		\centering
		\vspace{-20mm}
		\includegraphics[width=\linewidth]{Fig/435.pdf}
	\end{minipage}
	
	\vspace{-14mm}
	\parbox{\linewidth}{\centering\fontsize{8pt}{10pt}\selectfont\textbf{Fuente:} elaboración propia.}
\end{minipage}

\vskip0.2cm
En este caso, el volumen está dado por:

\vspace{-8mm}
\begin{multline*}
\bigintss_{0}^2\bigintss_{0}^{4-\frac{y^{4}}{4}}\bigintss_{\frac{y^{4}-2y^2}{4}}^{\frac{y^2}{2}}dxdzdy+\bigintss_{0}^2\bigintss_{4-\frac{y^{4}}{4}}^{4-\frac{\left(y^{4}-2y^2\right) ^2}{16}}\bigintss_{\frac{y^{4}-2y^2}{4}}^{\sqrt{4-z}}dxdzdy \\
+\bigintss_{0}^{\sqrt{2}}\bigintss_{4-\frac{\left( y^{4}-2y^2\right) ^2}{16}}^{4}\bigintss_{-\sqrt{4-z}}^{\sqrt{4-z}}dxdzdy \approx 3.7623.
\end{multline*}

La evaluación numérica de la integral anterior se puede obtener ejecutando las siguientes instrucciones.



\vspace{0.6em}
\noindent
\begin{tcolorbox}[breakable,enhanced,title=\bf Instrucción en \verb"Mathematica", top=2pt,bottom=2pt,boxsep=0pt,before skip=2pt,after skip=0pt]
\begin{Verbatim}[formatcom=\letraverbatim]
NIntegrate[1,{y,0,2},{z,0,4-y^4/4},{x,(y^4-2*y^2)/4,
y^2/2}]+NIntegrate[1,{y,0,2},{z,4-y^4/4,4-(y^4-
2*y^2)^2/16},{x,(y^4-2*y^2)/4,Sqrt[4-z]}]+NIntegrate[1,
{y,0,Sqrt[2]},{z,4-(y^4-2*y^2)^2/16,4},{x,-Sqrt[4-z],
Sqrt[4-z]}]
\end{Verbatim}
\end{tcolorbox}
\vspace{0.7em}
\noindent


También se puede obtener dicho valor como se muestra enseguida.



%\begin{tcolorbox}[title=\bfseries Parametrizando el sólido]
\noindent
Sobre el plano $xy,$ las ecuaciones $4x=y^{4}-2y^2$ y $2x=y^2,$ se parametrizan como
$\overrightarrow{\alpha}(t)=\left\langle \frac{1}{2}t^2,t\right\rangle $
y $\overrightarrow{\beta}(t)=\left\langle \frac{1}{4}\left(t^{4}-2t^2\right),t\right\rangle, $ $t\in [0,2].$

\vskip0.2cm
\noindent
\begin{minipage}{\textwidth}
	\centering
	\captionof{figure}{Región delimitada por $y^4 - 2y^2 \leq 4x \leq 2y^2$}
	\label{fig:solido_cilindro_elipse_3D}
	\begin{minipage}{0.34\linewidth}
		\centering
		\vspace{-15mm}
		\includegraphics[width=\linewidth]{Fig/436.pdf}
	\end{minipage}
	
	\vspace{-12mm}
	\fuentepropia
\end{minipage}
\vskip0.2cm


Una combinación convexa de la forma $(1-u)\overrightarrow{\alpha}(t) +u\overrightarrow{\beta}(t)$
parametriza la base del sólido, esto es
%\vskip0.2cm
{\small
\[\left\langle \frac{1}{4}t^2\left( \left( t^2-4\right) u+2\right) ,t\right\rangle\!.\]
}

La coordenada $z$ de un punto sobre el cilindro $z=4-x^2$ está dada por $z=4-\frac{t^{4}}{16}\left( \left( t^2-4\right) u+2\right) ^2.$
Por tanto, un punto $Q$ sobre el cilindro es de la forma
{\small $Q\left( \frac{1}{4}t^2((t^2-4)u+2),t,4-\frac{t^{4}}{16}\left( \left( t^2-4\right) u+2\right) ^2\right)\!,$}
con un punto en la base $P\left( \frac{1}{4}t^2\left( \left( t^2-4\right) u+2\right) ,t,0\right)\!,$
se forma la combinación convexa $(1-w)P+wQ$ que parametriza el sólido; esto define la siguiente función
\[
\mathbf{r} (t,u,w) = \left\langle \tfrac{t^2}{4}((t^2-4)u+2) ,t,w\left(4-\tfrac{1}{16}t^4((t^2-4)u+2)^2\right)\right\rangle\!;
\] \\
con $t\in [0,2],$ $u\in [0,1],$ $w\in [0,1].$
% luego hallamos los siguientes vectores. \begin{eqnarray*}
%{\bf r}_{t} &=& \langle t\left( \left( t^2-2\right) u+1\right) ,1,\frac{1}{2}wt^{3}\left( 4u-t^2u-2\right)\left(\left(t^2-2\right) u+1\right)\rangle \\
%{\bf r}_{u} &=& \langle \frac{1}{4}t^2\left( t^2-4\right) ,0,-\frac{1}{8}wt^{4}\left( t^2-4\right) \left( \left( t^2-4\right) u+2\right)\rangle \\
%{\bf r}_{w}&=& \langle 0,0,\frac{1}{16}\left( 4-t^2\right) \left(t^2\left( \left( t^2-4\right) u+2\right) +8\right) \left(t^2u+2\right) \rangle
%\end{eqnarray*}
El elemento diferencial de volumen $|{\bf r}_{t}\cdot ({\bf r}_{u}\times {\bf r}_{w})|$
es $\frac{t^2}{64}\left(t^2-4\right) ^2\left( t^2u+2\right) \left( t^2\left(
\left( t^2-4\right) u+2\right) \right)\!,$
término que integramos para conseguir el volumen del sólido
\[\int_{0}^{1}\int_{0}^{1}\int_{0}^2 |{\bf r}_{t}\cdot ({\bf r}_{u}\times {\bf r}_{w})| dtdudw= \frac{169472}{45045}.\]
%\end{tcolorbox}

Lo anterior se puede evaluar ejecutando las siguientes instrucciones:

\vspace{0.6em}
\noindent
\begin{tcolorbox}[breakable,enhanced,title=\bf Instrucción en \verb"Mathematica", top=2pt,bottom=2pt,boxsep=0pt,before skip=2pt,after skip=0pt]
\begin{Verbatim}[formatcom=\letraverbatim]
{x,y,z}=(1-u){t^2/2,t,0}+u{(t^4-2*t^2)/4,t,0};
r=Simplify[(1-w)*{x,y,0}+w*{x,y,4-x^2}]
A=Simplify[Dot[D[r,t],Cross[D[r,u],D[r,w]]]]
Integrate[A,{t,0,2},{u,0,1},{w,0,1}]
\end{Verbatim}
\end{tcolorbox}
\vspace{0.7em}
\noindent


\begin{ejer}
Plantear y evaluar las integrales triples en los órdenes $dxdydz$ y $dydxdz$,
necesarias para hallar el volumen del sólido acotado por las figuras de las expresiones $z = 4- x^2,$
$2x=y^2\,$ y $\,\, 4x = y^4 -2y^2$ para $y \geq 0.$
\end{ejer}

\subsection*{Tres cilindros rectos}

Los cilindros con ecuaciones
$x^2+y^2=2x$, $x^2+y^2=4x$ y $x^2+z^2=16$ generan el siguiente sólido.

\begin{figure}[H]
\centering
\caption{Visualización del sólido definido por $2x \leq x^2 + y^2 \leq 4x$ \quad y \quad $x^2 + z^2 \leq 16$.}
\label{fig:solido_anillo_cilindro}
\includegraphics[height=4.5cm]{Img/83a.pdf}\hskip0.25cm
\includegraphics[height=4.5cm]{Img/83b.pdf}\hskip0.25cm
\includegraphics[height=4.5cm]{Img/83c.pdf}
\parbox{\linewidth}{\centering\fontsize{9pt}{10pt}\selectfont\textbf{Fuente:} elaboración propia.}
\end{figure}

Estas figuras se consiguen ejecutando las siguientes instrucciones, se incluye una animación que muestra el recorrido de estas curvas en el plano $xy;$ esto se usa para obtener la parametrización de esta región.


\vspace{0.6em}
\noindent
\begin{tcolorbox}[breakable,enhanced,title=\bf Instrucción en \verb"Mathematica", top=2pt,bottom=2pt,boxsep=0pt,before skip=2pt,after skip=0pt]
\begin{Verbatim}[formatcom=\letraverbatim]
Manipulate[ParametricPlot[{{Cos[t]*Cos[t],
Cos[t]*Sin[t]},{2*Cos[t]*Cos[t],2*Cos[t]*Sin[t]}},
{t,-Pi/2,k},PlotRange->{{0,2},{-2,2}}],
{k,-Pi/2+Pi/40,Pi/2}]
ContourPlot3D[{16==x^2+z^2,x^2+y^2==2*x,x^2+y^2==4*x},
{x,-4,4},{y,-4,4},{z,-4,4},PlotPoints->90,Mesh->False,
Ticks->{{-4,-2,0,2,4},{-3,0,3},{-4,-2,0,2,4}},
ContourStyle->{Orange,LightGray,Cyan}]
RegionPlot3D[x^2+z^2<16&&x^2+y^2>2*x&&x^2+y^2<4*x&&0<
x<4,{x,0,4},{y,-2,2},{z,-4,4},Mesh->False]
\end{Verbatim}
\end{tcolorbox}
\vspace{0.7em}
\noindent


Con la proyección del sólido en el plano $xy$ se hallan los límites de integración en los
órdenes $dzdydx$ y $dzdxdy$; se muestra esta región, las diferentes tonalidades de sombreado determinan las integrales necesarias para su evaluación.
\vskip0.2cm


\noindent
\begin{minipage}{\textwidth}
	\centering
	\captionof{figure}{Región delimitada por $2x \leq x^2 + y^2 \leq 4x$.}
	\label{fig:region_corona_desplazada}
	\begin{minipage}{0.30\linewidth}
		\centering
		\vspace{-10mm}
		\includegraphics[width=\linewidth]{Fig/437.pdf}
	\end{minipage}
	\hspace{1em}
	\begin{minipage}{0.30\linewidth}
		\centering
		\vspace{-10mm}
		\includegraphics[width=\linewidth]{Fig/438.pdf}
	\end{minipage}
	
	\vspace{-8mm}
	\parbox{\linewidth}{\centering\fontsize{9pt}{10pt}\selectfont\textbf{Fuente:} elaboración propia.}
\end{minipage}
\vskip0.2cm

En coordenadas rectangulares, el volumen del mencionado sólido se puede determinar por medio de las siguientes integrales.
{\small
\begin{multline*}
\int_{-1}^{1}\int_{2-\sqrt{4-y^2}}^{1-\sqrt{1-y^2}}\int_{-\sqrt{16-x^2}}^{\sqrt{16-x^2}}dzdxdy+\int_{-1}^{1}\int_{1+\sqrt{1-y^2}}^{2+\sqrt{4-y^2}}\int_{-\sqrt{16-x^2}}^{\sqrt{16-x^2}}dzdxdy+\\
\int_{-2}^{-1}\int_{2-\sqrt{4-y^2}}^{2+\sqrt{4-y^2}}\int_{-\sqrt{16-x^2}}^{\sqrt{16-x^2}}dzdxdy+\int_{1}^2\int_{2-\sqrt{4-y^2}}^{2+\sqrt{4-y^2}}\int_{-\sqrt{16-x^2}}^{\sqrt{16-x^2}}dzdxdy \approx 57.207.
\end{multline*}
}
\vspace{-5mm}
{\small
\begin{multline*}
\int_{0}^2\int_{-\sqrt{4x-x^2}}^{-\sqrt{2x-x^2}}\int_{-\sqrt{16-x^2}}^{\sqrt{16-x^2}}dzdydx+\int_{0}^2\int_{\sqrt{2x-x^2}}^{\sqrt{4x-x^2}}\int_{-\sqrt{16-x^2}}^{\sqrt{16-x^2}}dzdydx\\
+\int_{2}^{4}\int_{-\sqrt{4x-x^2}}^{\sqrt{4x-x^2}}\int_{-\sqrt{16-x^2}}^{\sqrt{16-x^2}}dzdydx \approx 57.207
\end{multline*}
}

Estas integrales se evaluaron numéricamente en \verb"Mathematica";
se presentan las instrucciones para el orden $dzdydx$.


\vspace{0.6em}
\noindent
\begin{tcolorbox}[breakable,enhanced,title=\bf Instrucción en \verb"Mathematica", top=2pt,bottom=2pt,boxsep=0pt,before skip=2pt,after skip=0pt]
\begin{Verbatim}[formatcom=\letraverbatim]
NIntegrate[1,{x,0,2},{y,-Sqrt[4*x-x^2],-Sqrt[2*x-x^2]},
{z,-Sqrt[16-x^2],Sqrt[16-x^2]}]+NIntegrate[1,{x,0,2},
{y,Sqrt[2*x-x^2],Sqrt[4*x-x^2]},{z,-Sqrt[16-x^2],
Sqrt[16-x^2]}]+NIntegrate[1,{x,2,4},{y,-Sqrt[4*x-x^2],
Sqrt[4*x-x^2]},{z,-Sqrt[16-x^2],Sqrt[16-x^2]}]
\end{Verbatim}
\end{tcolorbox}
\vspace{0.7em}
\noindent

Haciendo un cambio a coordenadas cilíndricas el volumen esta dado por:
%\vskip0.2cm 
\[
\int_{0}^{\pi} \int_{2 \cos \theta}^{4 \cos \theta} \int_{-\sqrt{16-(r \cos \theta)^{2}}}^{\sqrt{16-(r \cos \theta)^{2}}} r\,dz\,dr\,d\theta \approx 57.207.
\]
%\vskip0.2cm


El mismo valor se puede obtener con el siguiente procedimiento.


%\begin{tcolorbox}[title=\bfseries Parametrizando el sólido]
\noindent
En el plano $xy$, un punto $P$ sobre el círculo interno tiene la forma
{\small
\[\left( 2\cos (\theta) \cos (\theta) ,2\cos (\theta) \sin (\theta) \right) = (2 \cos^2 (\theta), \sin(2\theta)),\]}
y un punto $Q$ sobre el círculo mayor posee la forma
{\small
\[\left( 4\cos (\theta) \cos (\theta) ,4\cos \left(
\theta \right) \sin (\theta) \right) = (4 \cos^2 (\theta), 2\sin(2\theta)).\]}

%\vspace{1ex} 

%\noindent
%\vskip0.4cm
\vspace{-2mm}
La combinación convexa $\left( 1-u\right)P +uQ$ se simplifica como
$ (2( 1+u)\cos^2(\theta), \linebreak
( 1+u)\sin (2\theta) ),$ y parametriza la región acotada por estas dos curvas. Sobre el cilindro $x^2+z^2 = 16$ se satisface que $z = \pm 2\sqrt{4-(1+u) ^2\cos ^{4}(\theta) }.$ Sobre el cilindro circular, consideremos dos puntos $P_1$ y $Q_1.$


\vskip0.3cm
\noindent
\begin{minipage}{\textwidth}
	\centering
	\captionof{figure}{Región delimitada por $2x \leq x^2 + y^2 \leq 4x$}
	\label{fig:corona_desplazada_xy}
	\begin{minipage}{0.37\linewidth}
		\centering
		\vspace{-15mm}
		\includegraphics[width=\linewidth]{Fig/439.pdf}
	\end{minipage}
	
	\vspace{-10mm}
	\fuentepropia
\end{minipage}
\vskip0.3cm

{\small
\[P_1= \big(2(1+u) \cos^2(\theta),\left( 1+u\right) \sin \left( 2\theta \right) ,-2\sqrt{4-(1+u)^2\cos ^{4}(\theta) }\big)\]  \[Q_1=\big(2(1+u) \cos ^2(\theta) ,\left( 1+u\right) \sin \left( 2\theta \right) ,2\sqrt{4-\left( 1+u\right) ^2\cos ^{4}(\theta) }\big);\]
}
%\vspace{7mm}
simplificando la expresión
$\left(1-w\right)P_1+wQ_1 $ se obtiene la función.
%\vspace{-4mm}
{\small
\[\mathbf{r}(\theta,u,w) = \left\langle 2(1+u)\cos^2\theta, (1+u)\sin(2\theta), 2(2w-1)\sqrt{4-(1+u)^2\cos^4\theta}\right\rangle\]
}
que parametriza el sólido, con $\theta \in [0,2\pi],$ $u \in [0,1]$ y $w \in [0,1].$
%% ${\bf r}_{\theta } = \left\langle -4\cos (\theta) \left(1+u\right) \sin (\theta) ,2\left( 1+u\right) \cos \left(2\theta \right) ,\frac{4\left( 2w-1\right) \left( 1+u\right) ^2\cos^{3}(\theta) \sin (\theta) }{\sqrt{4-\left(1+u\right) ^2\cos ^{4}(\theta) }}\right\rangle $\\
%%${\bf r}_{u} = \left\langle 2\cos ^2(\theta),\sin \left( 2\theta \right) ,\frac{-2\left( 2w-1\right) \left( 1+u\right)\cos ^{4}(\theta) }{\sqrt{4-\left( 1+u\right) ^2\cos^{4}(\theta) }} \right\rangle $ \\
%%${\bf r}_{w} = \left\langle 0,0,4\sqrt{\left( 4-\left( 1+u\right)^2\cos ^{4}(\theta) \right) }\right\rangle .$ \\
Con esta función se obtiene el elemento diferencial de volumen
%\vskip0.2cm
{\small
\[\left| {\bf r}_{w}\cdot \left({\bf r}_{\theta }\times {\bf r}_{u}\right) \right| =16\left| \cos ^2(\theta) \left( 1+u\right) \sqrt{4-\left( 1+u\right) ^2\cos ^{4}(\theta) }\right|,\]
}
%\vskip0.2cm
cuya integral es el volumen del sólido que se considera.
{\small
\[\int_{0}^{\pi }\int_{0}^{1}\int_{0}^{1}\left| {\bf r}_{w}\cdot \left({\bf r}_{\theta }\times {\bf r}_{u}\right) \right| dudwd\theta \approx 57.207.\]}
%\end{tcolorbox}
%\vskip0.2cm
Enseguida se presenta una forma de hacer parametrización de la base del sólido y los cálculos que permiten evaluar numéricamente la integral.


\vspace{0.6em}
\noindent
\begin{tcolorbox}[breakable,enhanced,title=\bf Instrucción en \verb"Mathematica", top=2pt,bottom=2pt,boxsep=0pt,before skip=2pt,after skip=0pt]
\begin{Verbatim}[formatcom=\letraverbatim]
{x,y}=Simplify[(1-u)*{2*Cos[t]*Cos[t],2*Cos[t]*Sin[t]}
+u*{4*Cos[t]*Cos[t],4*Cos[t]*Sin[t]}];
z=Simplify[(1-w)*Sqrt[16-x^2]+w*(-Sqrt[16-x^2])];
A=Dot[D[{x,y,z},t],Cross[D[{x,y,z},u],D[{x,y,z},w]]]
NIntegrate[A,{t,0,Pi},{u,0,1},{w,0,1}]
\end{Verbatim}
\end{tcolorbox}
\vspace{0.7em}
\noindent


Haciendo un cambio a coordenadas cilíndricas, el volumen está dado por:
\vspace{-2mm}
\[\int_{0}^{\pi }\int_{2\cos \theta }^{4\cos \theta }\int_{-\sqrt{16-\left(r\cos \theta \right)^2}}^{\sqrt{16-(r\cos\theta)^2}}rdzdrd\theta \approx 57.207.\]


\subsection*{Tres cilindros}
Las figuras de las ecuaciones $y^2 +2z^2=4,$ $2x=y^2$ y $4x=y^4-2y^2$ se muestran a continuación:


\vskip0.2cm
\noindent
\begin{minipage}{\textwidth}
	\centering
	\captionof{figure}{Tres cilindros}
	\label{fig:tres_cilindros}
	\begin{minipage}{0.20\linewidth}
		\centering
		\includegraphics[width=\linewidth]{Img/18d.pdf}
	\end{minipage}
	\begin{minipage}{0.20\linewidth}
		\centering
		\includegraphics[width=\linewidth]{Img/18c.pdf}
	\end{minipage}
	\begin{minipage}{0.20\linewidth}
		\centering
		\includegraphics[width=\linewidth]{Img/18e.pdf}
	\end{minipage}
	
	\vspace{0.3em}
	\fuentepropia
\end{minipage}
\vskip0.2cm



Estas superficies con $y\geq 0$ determinan la región que se muestra enseguida.

\vskip0.2cm
\noindent
\begin{minipage}{\textwidth}
	\centering
	\captionof{figure}{Sólido determinado por los tres cilindros.}
	\begin{minipage}{0.18\linewidth}
		\centering
		\includegraphics[width=\linewidth]{Img/18a.pdf}
	\end{minipage}
	\hspace{1em}
	\begin{minipage}{0.18\linewidth}
		\centering
		\includegraphics[width=\linewidth]{Img/18b.pdf}
	\end{minipage}
	
	\vspace{0.3em}
	\fuentepropia
\end{minipage}

\vskip0.2cm
Estas figuras se obtienen con ayuda de \verb"Mathematica", ejecutando los siguientes comandos.



\vspace{0.6em}
\noindent
\begin{tcolorbox}[breakable,enhanced,title=\bf Instrucción en \verb"Mathematica", top=2pt,bottom=2pt,boxsep=0pt,before skip=2pt,after skip=0pt]
\begin{Verbatim}[formatcom=\letraverbatim]
ContourPlot3D[{y^2+2*z^2==4,2*x==y^2,4*x==y^4-2*y^2},
{x,-1,2},{y,0,2},{z,-2,2},,PlotPoints->60,Mesh->False,
ContourStyle->{LightPurple,Orange,Blend[{Green,
Orange,Red}]}]
RegionPlot3D[-Sqrt[4-y^2]<z<Sqrt[4-y^2]&&(y^4-2*y^2)/4
<x<y^2/2,{x,-1,2},{y,0,2},{z,-2,2},PlotPoints->90,
Mesh->None,PlotStyle->Blend[{Green,Orange,Red}]]
\end{Verbatim}
\end{tcolorbox}
\vspace{0.7em}
\noindent


Con las siguientes tres instrucciones además de la figura de la superficie se obtienen las curvas de intersección.


\vspace{0.6em}
\noindent
\begin{tcolorbox}[breakable,enhanced,title=\bf Instrucción en \verb"Mathematica", top=2pt,bottom=2pt,boxsep=0pt,before skip=2pt,after skip=0pt]
\begin{Verbatim}[formatcom=\letraverbatim]
a=ContourPlot3D[{y^2+2*z^2==4,2*x==y^2,4*x==y^4
2*y^2},{x,-0.5,2},{y,0,2},{z,-2,2},PlotPoints->80,
Mesh->False,ContourStyle->{Directive[LightPurple,
Opacity[0.9]],Directive[Orange,Opacity[0.75]],
Directive[Blend[{Green,Orange,Red}],Opacity[0.75]],
Opacity[0.75]}];
b=ParametricPlot3D[{{2-t^2,Sqrt[4-2*t^2],t},
{(t^2-1)*(t^2-2),Sqrt[4-2*t^2],t}},{t,-Sqrt[2],
Sqrt[2]},PlotPoints->80,PlotStyle->{Directive[Blue,
Thickness->0.012],Directive[Red,Thickness->0.012]}];
Show[a,b]
\end{Verbatim}
\end{tcolorbox}
\vspace{0.7em}
\noindent


La expresión $y^2 +2z^2 = 4$ equivale a $y = \pm \sqrt{4-2z^2}$ o $z = \pm \sqrt{\frac{4-y^2}{2}} $ y de la expresión $4x = y^4 -2y^2$ podemos plantear cuatro funciones con $y$ como variable dependiente de $x$, esto es $y = \pm \sqrt{1\pm \sqrt{1+4x}}.$ \vskip0.3cm
En los planos $yz$ y $xy$ el sólido se proyecta como se muestra a continuación.


\vskip0.3cm
\noindent
\begin{minipage}{\textwidth}
	\centering
	\captionof{figure}{Proyección del sólido en los planos $yz$ y $xy$.}
	\label{fig:proyeccion_yz_xy}
	\begin{minipage}{0.35\linewidth}
		\centering
		\vspace{-15mm}
		\includegraphics[width=\linewidth]{Fig/440.pdf}
	\end{minipage}
	\hspace{1em}
	\begin{minipage}{0.35\linewidth}
		\centering
		\vspace{-15mm}
		\includegraphics[width=\linewidth]{Fig/441.pdf}
	\end{minipage}	
	
	\vspace{-10mm}
	\fuentepropia
\end{minipage}
\vskip0.2cm


Con las figuras anteriores planteamos las integrales triples que determinan el volumen del sólido, integrando primero con respecto a $x$ o $z$.
\vspace{-1mm}
{\small
\begin{multline*}
\hskip-0.3cm \bigintsss_{-\frac{1}{4}}^{0}\bigintsss_{\sqrt{1-\sqrt{1+4x}}}^{\sqrt{1+\sqrt{1+4x}}}\bigintsss_{-\sqrt{\frac{4-y^2}{2}}}^{\sqrt{\frac{4-y^2}{2}}}dzdydx+\bigintsss_{0}^2\bigintsss_{\sqrt{2x}}^{\sqrt{1+\sqrt{1+4x}}}\bigintsss_{-\sqrt{\frac{4-y^2}{2}}}^{\sqrt{\frac{4-y^2}{2}}}dzdydx\\
=\bigintsss_{-\sqrt{2}}^{\sqrt{2}}\bigintsss_{0}^{\sqrt{4-2z^2}}\bigintsss_{\frac{y^{4}-2y^2}{4}}^{\frac{y^2}{2}}dxdydz=\bigintsss_{0}^2\bigintsss_{-\sqrt{\frac{4-y^2}{2}}}^{\sqrt{\frac{4-y^2}{2}}}\bigintsss_{\frac{y^{4}-2y^2}{4}}^{\frac{y^2}{2}}dxdzdy \\
= \bigintsss_{0}^2\bigintsss_{\frac{y^{4}-2y^2}{4}}^{ \frac{y^2}{2}}\bigintsss_{-\sqrt{\frac{4-y^2}{2}}}^{\sqrt{\frac{4-y^2}{2}}}dzdxdy= \frac{\sqrt{2}\pi}{2}.
\end{multline*}
}

Dos de estas integrales se obtienen con ayuda de las siguientes sentencias.


\vspace{0.6em}
\noindent
\begin{tcolorbox}[breakable,enhanced,title=\bf Instrucción en \verb"Mathematica", top=2pt,bottom=2pt,boxsep=0pt,before skip=2pt,after skip=0pt]
\begin{Verbatim}[formatcom=\letraverbatim]
Integrate[1,{z,-Sqrt[2],Sqrt[2]},{y,0,Sqrt[4-2z^2]},
{x,(y^4-2y^2)/4,y^2/2}]
Integrate[1,{y,0,2},{x,(y^4-2y^2)/4,y^2/2},
{z,-Sqrt[(4-y^2)/2],Sqrt[(4-y^2)/2]}]
\end{Verbatim}
\end{tcolorbox}
\vspace{0.7em}
\noindent


De la ecuación del cilindro se sigue que $y^2= 4-2z^2$, al reemplazar esta expresión en las otras dos ecuaciones se obtienen las expresiones que definen la proyección de la figura en el plano $xz$. Esto es, para $2x = y^2 \to x = 2- z^2$ y de la ecuación
$4x=y^{4}-2y^2$ se sigue que $4x=\left( 4-2z^2\right) ^2-2\left( 4-2z^2\right),$ entonces
$x= (z^2-1)(z^2-2).$
\vskip0.2cm
Para integrar en el orden $dydxdz$ nos apoyamos en la siguiente figura.
Debemos tener en cuenta que en una región, $y$ recorre los puntos desde la curva con ecuación
$y = a $ hasta $y =b;$
en la otra región $y$ recorre puntos desde $y=\sqrt{2x}$ hasta
$y = \sqrt{4 -2z^2}$ y en la tercer región $y$ toma valores desde $y = \sqrt{2x}$ hasta
$y =b.$
\vskip0.4cm
Por la simetría de la figura con respecto al plano $xy$, la segunda gráfica de la figura \ref{fig:proyeccion_xzx} de las siguientes figuras se ha dividido en seis regiones que corresponden, respectivamente, a las seis integrales que se enuncian; luego de plantearlas basta multiplicar por $2.$

\vskip0.3cm
\noindent
\begin{minipage}{\textwidth}
	\centering
	\vspace{3mm}
	\captionof{figure}{Proyección del sólido en el plano $xz$.}
	\label{fig:proyeccion_xzx}
	\begin{minipage}{0.45\linewidth}
		\centering
		\vspace{-15mm}
		\includegraphics[width=\linewidth]{Fig/442.pdf}
	\end{minipage}
	\hspace{1em}
	\begin{minipage}{0.40\linewidth}
		\centering
		\vspace{-15mm}
		\includegraphics[width=\linewidth]{Fig/443.pdf}
	\end{minipage}
	
	\vspace{-12mm}
	\fuentepropia
\end{minipage}
\vskip0.5cm

Las regiones se enumeran en el mismo orden en que se presentan las siguien\-tes integrales, para abreviar usaremos la notación {\small $a= \sqrt{1-\sqrt{1+4x}}$} y {\small $b=\sqrt{1+\sqrt{1+4x}}.$}


{\small
\begin{multline*}
\hskip-0.3cm \bigintsss_{1}^{\sqrt{2}}\bigintsss_{(z^2-1)(z^2-2)}^{0} \bigintsss_{ a }^{\sqrt{4-2z^2}}dydxdz
+\bigintsss_{1}^{\sqrt{\frac{3}{2}}} \bigintsss_{-\frac{1}{4}}^{(z^2-1)(z^2-2)}\bigintsss_{a}^{b}dydxdz \\
+\bigintsss_{0}^{1}\bigintsss_{-\frac{1}{4}}^{0}\bigintsss_{ a }^{b}dydxdz+\bigintsss_{1}^{\sqrt{2}}\bigintsss_{0}^{2-z^2}\bigintsss_{\sqrt{2x}}^{\sqrt{4-2z^2}}dydxdz + \\
\hskip-0.3cm \bigintsss_{0}^{1}\bigintsss_{(z^2-1)(z^2-2)}^{2-z^2} \bigintsss_{\sqrt{2x}}^{\sqrt{4-2z^2}}dydxdz
+\bigintsss_{0}^{1}\bigintsss_{0}^{(z^2-1)(z^2-2)} \bigintsss_{\sqrt{2x}}^{b}dydxdz \approx \frac{\pi}{\sqrt{2}}.
\end{multline*}
}

\vskip0.2cm
Las integrales anteriores se pueden evaluar numéricamente mediante la ejecución de las siguientes instrucciones:


\vspace{0.6em}
\noindent
\begin{tcolorbox}[breakable,enhanced,title=\bf Instrucción en \verb"Mathematica", top=2pt,bottom=2pt,boxsep=0pt,before skip=2pt,after skip=0pt]
\begin{Verbatim}[formatcom=\letraverbatim]
2*(NIntegrate[1,{z,1,Sqrt[2]},{x,(z^2-1)*(z^2-2),0},
{y,Sqrt[1-Sqrt[1+4*x]],Sqrt[4-2*z^2]}]+NIntegrate[1,
{z,1,Sqrt[3/2]},{x,-1/4,(z^2-1)*(z^2-2)},{y,Sqrt[1
Sqrt[1+4*x]],Sqrt[1+Sqrt[1+4*x]]}]+NIntegrate[1,
{z,0,1},{x,-1/4,0},{y,Sqrt[1-Sqrt[1+4*x]],Sqrt[1
+Sqrt[1+4*x]]}]+Integrate[1,{z,1,Sqrt[2]},{x,0,
2-z^2},{y,Sqrt[2*x],Sqrt[4-2*z^2]}]+NIntegrate[1,
{z,0,1},{x,(z^2-1)*(z^2-2),2-z^2},{y,Sqrt[2*x],
Sqrt[4-2*z^2]}]+NIntegrate[1,{z,0,1},{x,0,
(2-z^2)(1-z^2)},{y,Sqrt[2*x],Sqrt[1+Sqrt[1+4*x]]}])
\end{Verbatim}
\end{tcolorbox}
\vspace{0.7em}
\noindent

Para plantear la integral triple en el orden $dydzdx,$ hay que obtener $z$ de la ecuación $x = (z^2-1)(z^2-2),$ esto es
$z = \pm \frac{1}{2} \sqrt{6 \pm 2 \sqrt{1+4x}};$
el signo depende de la región del plano que se considere. Usando la simetría de la figura, las integrales triples que se requieren evaluar para determinar el volumen del sólido son:

{\footnotesize
\begin{multline} \label{trescilindros1}
2\Bigg( \bigintsss_{-\frac{1}{4}}^{0} \bigintsss_{0}^{\frac{1}{2}\sqrt{6-2\sqrt{1+4x}}}\bigintsss_{ a }^{b}dydzdx+\\
\bigintsss_{-\frac{1}{4}}^0 \bigintsss_{\frac{1}{2}\sqrt{6-2\sqrt{1+4x}}}^{\frac{1}{2}\sqrt{6+2\sqrt{1+4x}}}\bigintsss_{ a }^{\sqrt{4-2z^2}}dydzdx+ \bigintsss_0^2 \bigintsss_{\frac{1}{2}\sqrt{6-2\sqrt{1+4x}}}^{\sqrt{2-x}}\bigintsss_{\sqrt{2x}}^{\sqrt{4-2z^2}}dydzdx
\\ + \bigintsss_{0}^2\bigintsss_{0}^{\frac{1}{2}\sqrt{6-2\sqrt{1+4x}}} \bigintsss_{\sqrt{2x}}^{b}dydzdx\Bigg) \approx \frac{\pi}{\sqrt{2}} \end{multline}
}

\vskip0.2cm
El volumen también se puede determinar de la siguiente manera.

%\begin{tcolorbox}[title=\bfseries Parametrizando el sólido]
En el plano $xy$ las curvas que determinan la proyección del sólido se parametrizan como
{\small $\overrightarrow{\alpha }(t)=\left\langle \frac{1}{4}\left(
t^{4}-2t^2\right) ,t\right\rangle$} y {\small $\overrightarrow{\beta }(t)=\left\langle \frac{1}{2}t^2,t\right\rangle;$} con $t\in [0,2].$
\vskip0.2cm
Con la expresión {\small $\left( 1-u\right) \overrightarrow{\alpha }(t)+u\overrightarrow{\beta }(t),$}
se parametriza la región entre estas curvas, esto es {\small $\left\langle \frac{1}{4}t^2\left( \left(1-u\right)t^2-2+4u\right) ,t\right\rangle.$}
La expresión {\small $z=\pm \sqrt{\frac{4-y^2}{2}},$} se reescribe como {\small $z=\pm \sqrt{2-\frac{1}{2}t^2}.$}
Dos puntos sobre el cilindro $y^2+2z^2=4$ tienen la forma
\[P\left( \frac{1}{4}t^2\left( \left( 1-u\right) t^2-2+4u\right) ,t,-\sqrt{2-\frac{1}{2}t^2}\right)\]
y \[Q\left( \frac{1}{4}t^2\left( \left(1-u\right) t^2-2+4u\right) ,t,\sqrt{2-\frac{1}{2}t^2}\right).\]
El sólido se parametriza simplificando la expresión $(1-w)P+wQ,$ que determina la siguiente función
\vskip0.2cm
\begin{equation*}
{\bf r}\left( t,u,w\right) = \left\langle \frac{1}{4}t^2\left( \left( 1-u\right)
t^2-2+4u\right) ,t,\left( w-\frac{1}{2}\right) \sqrt{8-2t^2}\right\rangle,
\end{equation*}
\vskip0.2cm

en donde $t\in [0,2],$ $u\in [0,1]$ y $w\in [0,1]$. %\vskip0.2cm
% Derivando parcialmente, hallamos los siguientes vectores
%% \begin{eqnarray*} {\bf r}_{t} &=& \left\langle \frac{1}{2}t\left( \left(1-u\right) t^2-2+4u\right) +\frac{1}{2}t^{3}\left( 1-u\right) ,1,\frac{t\left( 1-2w\right) }{\sqrt{8-2t^2}} \right\rangle \\ {\bf r}_{u} &=& \left\langle \frac{1}{4}t^2\left( 4-t^2\right) ,0,0\right\rangle \\
%%{\bf r}_{w}&=& \left\langle 0,0,\sqrt{8-2t^2}\right\rangle \end{eqnarray*}
Con esta función se obtiene el elemento diferencial de volumen
%\vskip0.2cm
\[\left| {\bf r}_{t}\cdot \left({\bf r}_{u}\times {\bf r}_{w}\right) \right| =\left| -\frac{1}{4}t^2\left( 4-t^2\right) \sqrt{8-2t^2}\right|,\]
%\vskip0.2cm 
cuya integración proporciona el volumen del sólido
%\vskip0.2cm
\[\int_{0}^{1}\int_{0}^{1}\int_{0}^2\left| -\frac{1}{4}t^2\left( 4-t^2\right) \sqrt{8-2t^2}\right|dtdudw= \frac{1}{2}\sqrt{2}\pi.\]
%\vskip0.2cm
%\end{tcolorbox}

Estos cálculos se realizan en \verb"Mathematica" con las siguientes instrucciones:


\vspace{0.6em}
\noindent
\begin{tcolorbox}[breakable,enhanced,title=\bf Instrucción en \verb"Mathematica", top=2pt,bottom=2pt,boxsep=0pt,before skip=2pt,after skip=0pt]
\begin{Verbatim}[formatcom=\letraverbatim]
r={x,y}=(1-u)*{(t^4-2*t^2)/4,t}+u*{t^2/2,t};
s=(1-w)*{x,y,-Sqrt[2-y^2/2]}+w*{x,y,Sqrt[2-y^2/2]};
A=Dot[D[s,t],Cross[D[s,u],D[s,w]]];
Integrate[Abs[A],{t,0,2},{u,0,1},{w,0,1}]
\end{Verbatim}
\end{tcolorbox}
\vspace{0.7em}
\noindent


Las integrales de la expresión (\ref{trescilindros1}) se evalúan ejecutando las siguientes instrucciones:


\vspace{0.6em}
\noindent
\begin{tcolorbox}[breakable,enhanced,title=\bf Instrucción en \verb"Mathematica", top=2pt,bottom=2pt,boxsep=0pt,before skip=2pt,after skip=0pt]
\begin{Verbatim}[formatcom=\letraverbatim]
(NIntegrate[1,{x,-1/4,0},{z,0,Sqrt[6-2Sqrt[1+4x]]/2},
{y,Sqrt[1-Sqrt[1+4x]],Sqrt[1+Sqrt[1+4x]]}]+
NIntegrate[1,{x,-1/4,0},{z,Sqrt[6-2Sqrt[1+4x]]/2,
Sqrt[6+2Sqrt[1+4x]]/2},{y,Sqrt[1-Sqrt[1+4x]],
Sqrt[4-2z^2]}]+NIntegrate[1,{x,0,2},{z,
Sqrt[6-2Sqrt[1+4x]]/2,Sqrt[2-x]},{y,Sqrt[2x],
Sqrt[4-2z^2]}]+NIntegrate[1,{x,0,2},{z,0,Sqrt[6-
2Sqrt[1+4x]]/2},{y,Sqrt[2x],Sqrt[1+Sqrt[1+4x]]}])*2
\end{Verbatim}
\end{tcolorbox}
\vspace{0.7em}
\noindent


\subsection*{Un cilindro y dos planos}

El cilindro elíptico con ecuaciones $4x^2+z^2=4$ y los planos $z=3+2y$ y
$z=5-x-3y$ acotan el sólido que se muestra a continuación:

\vskip0.2cm

\noindent
\begin{minipage}{\textwidth}
	\centering
	\captionof{figure}{Sólido definido por $4x^2 + z^2 \leq 4$ \quad y \quad $\dfrac{z - 3}{2} \leq y \leq \dfrac{5 - x - z}{3}$.}
	\label{fig:solido_cilindro_planos}
	\begin{minipage}{0.20\linewidth}
		\centering
		\includegraphics[width=\linewidth]{Img/55a.pdf}
	\end{minipage}
	\hspace{1em}
	\begin{minipage}{0.20\linewidth}
		\centering
		\includegraphics[width=\linewidth]{Img/55b.pdf}
	\end{minipage}
	\hspace{1em}
	\begin{minipage}{0.20\linewidth}
		\centering
		\includegraphics[width=\linewidth]{Img/55c.pdf}
	\end{minipage}
	
	\vspace{0.3em}
	\fuentepropia
\end{minipage}
\vskip0.2cm


Con las siguientes instrucciones se consiguen las figuras presentadas.

%\vskip0.4cm
\vspace{0.6em}
\noindent
\begin{tcolorbox}[breakable,enhanced,title=\bf Instrucción en \verb"Mathematica", top=2pt,bottom=2pt,boxsep=0pt,before skip=2pt,after skip=0pt]
\begin{Verbatim}[formatcom=\letraverbatim]
ContourPlot3D[{4*x^2+z^2==4,z==3+2*y,z==5-x-3*y},
{x,-2,2},{y,-3,4},{z,-2,2.5},Mesh->False,
ContourStyle->{Orange,RGBColor[0.2,0.6,0.4],Cyan}]
RegionPlot3D[(z-3)/2<y<(5-x-z)/5&&-Sqrt[1-z^2/4]<x<
Sqrt[1-z^2/4]&&-2<x<2,{x,-1,1},{y,-3,2},{z,-2,2},
Mesh->False,PlotStyle->{Blend[{Orange,Yellow},2/3]},
PlotPoints->90,Ticks->{{-1,0},{-2,0,1},{-1,0,1,2}}]
\end{Verbatim}
\end{tcolorbox}
\vspace{0.7em}
\noindent



En el plano $xz$ el sólido se proyecta como una elipse y los límites de integración para $y,$ van desde un plano al otro.
\vskip0.2cm
La proyección en el plano $xy,$ se determina reemplazando $z$ en la ecuación del cilindro $4x^2+z^2=4,$ por el respectivo término de $z$ en las ecuaciones de los planos, de tal manera que la expresión resultante solo dependa de $x$ y $y.$

\noindent
\begin{minipage}{\textwidth}
	\centering
	\captionof{figure}{Proyecciones del sólido en los planos $xy$ y $xz$.}
	\label{fig:proyeccion_xy_xz}
	\begin{minipage}{0.40\linewidth}
		\centering
		\vspace{-17mm}
		\includegraphics[width=\linewidth]{Fig/444.pdf}
	\end{minipage}
	\hspace{1em}
	\begin{minipage}{0.35\linewidth}
		\centering
		\vspace{-17mm}
		\includegraphics[width=\linewidth]{Fig/445.pdf}
	\end{minipage}
	
	\vspace{-12mm}
	\fuentepropia
\end{minipage}
\vskip0.2cm


Con lo anterior podemos plantear las integrales que determinan el volumen del sólido en los órdenes $dydxdz$ y $dydzdx,$
{\small
\begin{multline*} \bigintss_{-1}^{1}\bigintss_{-\sqrt{4-4x^2}}^{\sqrt{4-4x^2}}\bigintss_{\frac{z-3}{2}}^{\frac{5-x-z}{3}}dydzdx \\ =\bigintss_{-2}^2\bigintss_{-\sqrt{1-z^2/4}}^{\sqrt{1-z^2/4}}\bigintss_{\frac{z-3}{2}}^{\frac{5-x-z}{3}}dydxdz= \frac{19}{3}\pi.
\end{multline*}
}
%\vskip0.2cm
Con la ecuación del plano $z =3+2y$ se reemplaza $z$ en la ecuación del cilindro y se sigue que $4x^2+(3+2y)^2=4,$ que equivale a $4x^2+4(y+\frac{3}{2})^2=4;$
esta ecuación representa una circunferencia centrada en el punto $(0,-\frac{3}{2})$ de radio $1.$
Ahora, para el plano $z = 5-x-3y,$ la sustitución nos lleva a la ecuación
$4x^2 + (5-x-3y)^2 = 4,$ que corresponde a una elipse rotada en el plano $xy.$ De la expresión $4x^2+( 5-x-3y) ^2=4$, se obtiene que
$x=1-\frac{3}{5}y\pm \frac{2}{5}\sqrt{30y-20-9y^2},$
también
$y=\frac{5}{3}-\frac{1}{3}x\pm \frac{2}{3}\sqrt{1-x^2}.$
%\vskip0.4cm
De la expresión $4x^2+4( y+\frac{3}{2}) ^2=4$,
se obtiene que $x=\pm \frac{1}{2}\sqrt{ 4-( 3+2y) ^2}$
$y=-\frac{3}{2}\pm \sqrt{1-x^2}.$
Es de anotar, que en los extremos los valores de $z$ son menores en el cilindro que en los planos.
Con la información anterior podemos plantear las integrales que determinan el volumen del sólido en mención, en el orden $dzdydx$.

\vspace{-6mm}
{\small
\begin{multline*}
\bigintss_{-1}^{1}\bigintss_{-\frac{3}{2}-\sqrt{1-x^2}}^{-\frac{3}{2}+\sqrt{1-x^2}}\bigintss_{-2\sqrt{1-x^2}}^{3+2y}dzdydx \\ +\bigintss_{-1}^{1}\bigintss_{-\frac{3}{2}+\sqrt{1-x^2}}^{\frac{5}{3}-\frac{1}{3}x-\frac{2}{3}\sqrt{1-x^2}}\bigintss_{-\sqrt{4-4x^2}}^{\sqrt{4-4x^2}}dzdydx \\
+\bigintss_{-1}^{1}\bigintss_{\frac{5}{3}-\frac{1}{3}x-\frac{2}{3}\sqrt{1-x^2}}^{\frac{5}{3}-\frac{1}{3}x+\frac{2}{3}\sqrt{%
1-x^2}}\bigintss_{-\sqrt{4-4x^2}}^{5-x-3y}dzdydx= \frac{19}{3}\pi.
\end{multline*}
}
%\vskip0.2cm
Sobre el plano $yz,$ se hacen varias consideraciones, pues la proyección del sólido en este plano incluye una elipse rotada; además, de la expresión $4(5-z-3y)^2+z^2=4$ se sigue que
{\small
\[z=-\frac{12}{5}y+4\pm \frac{2}{5}\sqrt{5-9\left( y-\frac{5}{3}\right) ^2}.\]
}
La integral en el orden $dxdzdy$ corresponde a la suma de cinco integrales, esto se debe a que $x$ recorre los puntos entre la parte negativa del cilindro y su parte positiva; sin embargo, en otra región $x$ va desde la parte negativa del cilindro al plano $z=5-x-3y;$ esta última región se proyecta en el plano $yz$ como una elipse rotada.


\vskip0.3cm
\noindent
\begin{minipage}{\textwidth}
	\centering
	\captionof{figure}{Proyección del sólido en el plano $yz$}
	\label{fig:proyeccion_yz}
	\begin{minipage}{0.75\linewidth}
		\centering
		\vspace{-55mm}
		\includegraphics[width=\linewidth]{Fig/446.pdf}
	\end{minipage}
	
	\vspace{-50mm}
	\fuentepropia
\end{minipage}
\vskip0.3cm


Con lo anterior, y considerando las abreviaturas.
$a(y)=-\frac{12}{5}y+4+\frac{2}{5}\sqrt{5-9\left( y-\frac{5}{3}\right) ^2}$ y $b(y)=-\frac{12}{5}y+4-\frac{2}{5}\sqrt{5-9\left( y-\frac{5}{3}\right) ^2},$ el planteamiento es el siguiente:

\vspace{-6mm}
{\small
\begin{multline*}
\bigintsss_{-\frac{5}{2}}^{-\frac{1}{2}}\bigintsss_{-2}^{3+2y}\bigintsss_{-\sqrt{1-z^2/4}}^{\sqrt{1-z^2/4}}dxdzdy+\bigintsss_{-\frac{1}{2}}^{\frac{5-\sqrt{5}}{3}}\bigintsss_{-2}^2\bigintsss_{-\sqrt{1-z^2/4}}^{\sqrt{1-z^2/4}}dxdzdy\\
+\bigintsss_{\frac{5-\sqrt{5}}{3}}^{\frac{7}{3}}\bigintsss_{-2}^{a\left( y\right) }\bigintsss_{-\sqrt{1-z^2/4}}^{\sqrt{1-z^2/4}}dxdzdy+\bigintsss_{\frac{5-\sqrt{5}}{3}}^{1}\bigintsss_{b\left( y\right) }^2\bigintsss_{-\sqrt{1-z^2/4}}^{\sqrt{1-z^2/4}}dxdzdy\\
+\bigintsss_{\frac{5-\sqrt{5}}{3}}^{\frac{5+\sqrt{5}}{3}}\bigintsss_{a\left(y\right) }^{b\left( y\right) }\bigintsss_{-\sqrt{1-z^2/4}}^{5-z-3y}dxdzdy \approx 19.897.
\end{multline*}
}
\vskip0.2cm
Usando coordenadas cilíndricas, $x=r\cos \theta ,$ $z=r\sin \theta ,$ la ecuación del cilindro toma la forma
$4( r\cos \theta ) ^2+( r\sin \theta ) ^2=4$, que se simplifica
%\vskip0.2cm
{\small
\[r=\frac{2}{\sqrt{4-3\sin ^2\theta }}.\]
}
%\vskip0.2cm
Además, las ecuaciones de los planos toman la forma
$y=\frac{1}{2}r\sin \theta -\frac{3}{2}$ y
$y=-\frac{2}{3}r\sin \theta +\frac{5}{3}.$
Por lo tanto, el volumen del sólido está determinado por
\[\bigintsss_{0}^{2\pi }\bigintsss_{0}^{\frac{2}{\sqrt{4-3\sin ^2\theta }}}\bigintsss_{\frac{1}{2}r\sin \theta -\frac{3}{2}}^{-\frac{2}{3}r\sin\theta +\frac{5}{3}}rdydrd\theta = \frac{19}{3}\pi.\]
%\vskip0.2cm
Algunas de estas integrales se pueden evaluar numéricamente usando las siguientes instrucciones:


\vspace{0.6em}
\noindent
\begin{tcolorbox}[breakable,enhanced,title=\bf Instrucción en \verb"Mathematica", top=2pt,bottom=2pt,boxsep=0pt,before skip=2pt,after skip=0pt]
\begin{Verbatim}[formatcom=\letraverbatim]
a=Sqrt[1-z^2/4];
NIntegrate[1,{y,-5/2,-1/2},{z,-2,3+2*y},
{x,-a,a}]+
NIntegrate[1,{y,-1/2,(5-Sqrt[5])/3},{z,-2,2},{x,-a,a}]
+NIntegrate[1,{y,(5-Sqrt[5])/3,7/3},{z,-2,-12*y/5+4
2*Sqrt[30*y-20-9y^2]/5},{x,-a,a}]+NIntegrate[1,
{y,(5-Sqrt[5])/3,1},{z,-12*y/5+4+2*Sqrt[30*y-20-
9y^2]/5,2},{x,-a,a}]+NIntegrate[1,{y,(5-Sqrt[5])/3,
(5+Sqrt[5])/3},{z,-12*y/5+4-2*Sqrt[30*y-20-9y^2]/5,
12*y/5+4+2*Sqrt[30*y-20-9y^2]/5},{x,-a,5-z-3*y}]
\end{Verbatim}
\end{tcolorbox}
\vspace{0.7em}
\noindent



Este valor también se obtiene con el siguiente procedimiento:
\vskip0.2cm

%\begin{tcolorbox}[title=\bfseries Parametrizando el sólido]
%\footnotesize 
%\parskip=0pt 
%\setlength{\belowdisplayskip}{0pt}
%\setlength{\abovedisplayskip}{0pt} 
		
En el plano $xz$ el cilindro se proyecta como una elipse que se puede parametrizar en la forma
$x=\cos(t),$ $z=2\sin (t),$ para $t\in [0,2\pi ].$
La región interior a esta curva se parametriza de la siguiente manera: $x = u \cos(t),$ $z = 2u\sin(t),$ para
$t\in[0, 2\pi]$, $u\in[0, 1].$ La coordenada $y$ se puede determinar reemplazando en la ecuación de los planos.
De esta manera, para $z=3+2y$ se sigue $y = u\sin(t)-\frac{3}{2},$, mientras que en
$z = 5-x-3y,$ se obtiene que $y = -\frac{2}{3}u\sin(t)+\frac{5}{3}-\frac{1}{3}u \cos(t);$ por lo tanto, es natural que se definan las funciones:
\vspace{-2mm}
{\small
\begin{eqnarray*}
\overrightarrow{\alpha}(t,u)&=&\langle u \cos(t), u\sin(t)-\frac{3}{2},2u\sin(t)\rangle \\
\overrightarrow{\beta}(t,u) &=& \langle u \cos(t),-\frac{2}{3}u\sin(t)+\frac{5}{3}-\frac{1}{3}u \cos(t) ,2u\sin(t)\rangle.
\end{eqnarray*}
}
%\vskip0.2cm

Como parametrizaciones de la región sobre cada uno de los planos que acotan el cilindro. Con las funciones
$\overrightarrow{\alpha}$ y $\overrightarrow{\beta}$ se forma una combinación convexa que permite parametrizar el sólido. Es decir, al simplificar:
$(1-w)\overrightarrow{\alpha}(t,u)+w\overrightarrow{\beta}(t,u)$
se obtiene
\begin{multline*}
{\bf r}( t,u,w) = \big \langle u\cos (t),u\sin(t)-\frac{3}{2}-\frac{5}{3}wu\sin (t)+\frac{19}{6}w- \\ \frac{1}{3}wu\cos (t),2u\sin (t) \big \rangle, 
\end{multline*}
%\vskip0.2cm
en donde $t\in[0, 2\pi],\, u\in[0, 1],, w\in[0, 1], w\in[0, 1].$ También se obtienen los siguientes vectores:
{\small
\begin{eqnarray*}
{\bf r}_{t} &=& \langle -u\sin (t),u\cos (t)-\frac{5}{3}wu\cos (t)+\frac{1}{3}wu\sin (t),2u\cos (t)\rangle \\
{\bf r}_{u}&=& \langle \cos (t),\sin (t)-\frac{5}{3}w\sin (t)-\frac{1}{3}w\cos (t),2\sin (t)\rangle \\
{\bf r}_{w}&=& \langle 0,-\frac{5}{3}u\sin (t)+\frac{19}{6}-\frac{1}{3}u\cos (t),0\rangle.
\end{eqnarray*}
}
%\vskip0.2cm
Formando el producto vectorial triple, se obtiene
%\vskip0.2cm
\[| {\bf r}_{t}\cdot ( {\bf r}_{u}\times {\bf r}_{w}) | = \left| -\frac{10}{3}u^2\sin (t)+\frac{19}{3}u-\frac{2}{3}u^2\cos (t)\right |,\]
%\vskip0.2cm
que se integra con los límites definidos antes y nos proporciona el volumen del sólido
%\vskip0.2cm
\[\int_{0}^{1}\int_{0}^{1}\int_{0}^{2\pi }\left | -\frac{10}{3}u^2\sin (t)
+\frac{19}{3}u-\frac{2}{3}u^2\cos (t)\right | dtdudw= \frac{19}{3}\pi.\]
%\end{tcolorbox}
\vskip0.2cm
\begin{ejer}
Para el sólido acotado por las figuras de las ecuaciones $4x^2 + z^2 = 4,$ $z = 3 + 2y$ y $z = 5 -x-3y.$ Plantee y evalúe numéricamente las integrales triples en los órdenes $dzdxdy$ y $dxdydz.$
\end{ejer}

\section{Tres superficies y un plano}

El sólido acotado por las figuras de las expresiones $y=z^2$, $z=x^2$, $z=x^3$ y $y=0$,
se muestra a continuación. También se aprecian algunas líneas que corresponden a las proyecciones en los planos coordenados de las líneas de intersección de la superficie que lo conforman.

\vskip0.2cm
\noindent
\begin{minipage}{\textwidth}
	\centering
	\captionof{figure}{Visualización tridimensional del sólido definido por $0 \leq y \leq z^2$ \quad y \quad $x^2 \leq z \leq x^2$.}
	\label{fig:solido_parametrico_xyz}
	\begin{minipage}{0.28\linewidth}
		\centering
		\vspace{-12mm}
		\includegraphics[width=\linewidth]{Fig/447.pdf}
	\end{minipage}
	\hspace{1em}
	\begin{minipage}{0.28\linewidth}
		\centering
		\vspace{-12mm}
		\includegraphics[width=\linewidth]{Fig/448.pdf}
	\end{minipage}
	
	\vspace{-8mm}
	\fuentepropia
\end{minipage}
\vskip0.2cm


Las proyecciones del sólido en los planos coordenados son útiles para establecer los límites y evaluar la integral triple sobre dicha región. Enseguida se presentan las regiones correspondientes a los planos $xy$ y $yz$.

\vskip0.2cm
\noindent
\begin{minipage}{\textwidth}
	\centering
	\captionof{figure}{Proyecciones del sólido en los planos $xy$ y $yz$, definido por $x^6 \leq y \leq x^4$ \quad y \quad $0 \leq y \leq z^2$.}
	\label{fig:proyeccion_xy_yz}
	\begin{minipage}{0.30\linewidth}
		\centering
		\vspace{-17mm}
		\includegraphics[width=\linewidth]{Fig/449.pdf}
	\end{minipage}
	\hspace{1em}
	\begin{minipage}{0.30\linewidth}
		\centering
		\vspace{-17mm}
		\includegraphics[width=\linewidth]{Fig/450.pdf}
	\end{minipage}
	
	\vspace{-13mm}
	\fuentepropia
\end{minipage}


\vskip0.2cm
Utilizando la información anterior, el volumen del sólido está dado por las siguientes integrales
\vskip0.2cm
{\small
$\ds \int_{0}^{1}\int_{0}^{x^{6}}\int_{x^{3}}^{x^{2}}dzdydx+\int_{0}^{1}\int_{x^{6}}^{x^{4}}\int_{\sqrt{y}}^{x^{2}}dzdydx= \frac{1}{70}$
}

{\small
$\ds \int_{0}^{1}\int_{\root{6}\of{y}}^{1}\int_{x^{3}}^{x^{2}}dzdxdy+\int_{0}^{1}\int_{\root{4}\of{y}}^{\root{6}\of{y}}\int_{\sqrt{y}}^{x^{2}}dzdxdy= \frac{1}{70}$
}

{\small
$\ds \int_{0}^{1}\int_{0}^{z^{2}}\int_{\root{2}\of{z}}^{\root{3}\of{z}}dxdydz= \int_{0}^{1}\int_{\sqrt{y}}^{1}\int_{\root{2}\of{z}}^{\root{3}\of{z}}dxdzdy= \frac{1}{70}.$
}
\vskip0.4cm
En el plano $xz$ el sólido se proyecta como se muestra enseguida:

\vskip0.2cm
\noindent
\begin{minipage}{\textwidth}
	\centering
	\captionof{figure}{Visualización del sólido definido por $0 \leq y \leq z^2$ \quad y \quad $x^3 \leq z \leq x^2$.}
	\label{fig:solido_xz_yz}
	\begin{minipage}{0.30\linewidth}
		\centering
		\vspace{-13mm}
		\includegraphics[width=\linewidth]{Fig/451.pdf}
	\end{minipage}
	\hspace{1em}
	\begin{minipage}{0.30\linewidth}
		\centering
		\vspace{-13mm}
		\includegraphics[width=\linewidth]{Fig/452.pdf}
	\end{minipage}
	
	\vspace{-8mm}
	\fuentepropia
\end{minipage}


\vskip0.2cm
Las integrales triples que determinan el volumen del sólido en los órdenes $dydxdz$ y $dydzdx$
se presentan a continuación:

\[\ds \int_{0}^{1}\int_{\root{2}\of{z}}^{\root{3}\of{z}}\int_{0}^{z^{2}}dydxdz =\int_{0}^{1}\int_{x^{3}}^{x^{2}}\int_{0}^{z^{2}}dydzdx= \frac{1}{70}.\]
\vskip0.2cm
Las integrales presentadas antes se pueden evaluar en \verb"Mathematica" mediante la ejecución de las siguientes sentencias.


\vspace{0.6em}
\noindent
\begin{tcolorbox}[breakable,enhanced,title=\bf Instrucción en \verb"Mathematica", top=2pt,bottom=2pt,boxsep=0pt,before skip=2pt,after skip=0pt]
\begin{Verbatim}[formatcom=\letraverbatim]
Integrate[1,{x,0,1},{y,0,x^6},{z,x^3,x^2}]+
Integrate[1,{x,0,1},{y,x^6,x^4},{z,Sqrt[y],x^2}]
Integrate[1,{y,0,1},{x,y^(1/6),1},{z,x^3,x^2}]+
Integrate[1,{z,0,1},{y,0,z^2},{z,Sqrt[z],z^(1/3)}]
Integrate[1,{y,0,1},
{z,Sqrt[y],1},{x,Sqrt[z],z^(1/3)}]
Integrate[1,{z,0,1},{x,Sqrt[z],z^(1/3)},{y,0,z^2}]
Integrate[1,{x,0,1},{z,x^3,x^2},{y,0,z^2}]
\end{Verbatim}
\end{tcolorbox}
\vspace{0.7em}
\noindent


El volumen de esta figura también se puede hallar de la siguiente manera.


%\begin{tcolorbox}[title=\bfseries Parametrizando el sólido]
%\noindent
\vskip0.3cm
Sobre el sólido, consideremos los puntos
$A=(\sqrt[2]{u},0,u),$ $B=(\sqrt[3]{u},0,u),$ $C=(\sqrt[2]{u},u^{2},u)$ y $D=(\sqrt[3]{u},u^{2},u),$
para $u \in [0,1].$
Los segmentos $\overline{AB}$ y $\overline{CD}$ se parametrizan respectivamente como:
$\langle(1-v)\sqrt{u}(1-v)+v\sqrt[3]{u},0,u\rangle$ y
$\langle(1-v)\sqrt{u}+v\sqrt[3]{u},u^{2},u\rangle,$ $v \in [0,1].$
 

\vskip0.2cm
\noindent
\begin{minipage}{\textwidth}
	\centering
	\captionof{figure}{Visualización tridimensional del sólido definido por $x^3 \leq z \leq x^2$}
	\label{fig:solido_x3_x2}
	\begin{minipage}{0.36\linewidth}
		\centering
		\includegraphics[width=\linewidth]{Fig/9.pdf}
	\end{minipage}
	
	\vspace{0.3em}
	\fuentepropia
\end{minipage}
\vskip0.2cm


El sólido se recorre completamente uniendo un punto del segmento {\small $\overline{AB}$} con un punto del segmento {\small $\overline{CD},$} esto se consigue simplificando \\
{\small $\left( 1-w\right) \langle(1-v)\sqrt{u}(1-v)+v\sqrt[3]{u},0,u\rangle +w\langle(1-v)\sqrt{u}+v\sqrt[3]{u},u^{2},u\rangle$}
con lo cual, la función que parametriza el sólido es
%\vskip0.2cm
\vspace{-2mm}
\begin{equation} \label{32a}
{\bf r}\left( u,v,w\right) =\left\langle \sqrt{u}-\sqrt{u}v+v\sqrt[3]{u},wu^{2},u \right\rangle,
\end{equation}
%\vskip0.2cm
para $u\in [0,1],$ $v\in [0,1],w\in [0,1].$
También hallamos los vectores
{\small
\begin{eqnarray*}
{\bf r}_{u} &=& \left\langle \frac{1-v}{2\sqrt{u}}+\frac{1}{3}\frac{v}{\left(\sqrt[3]{u}\right) ^{2}},2wu,1\right\rangle \\
{\bf r}_{v} &=& \left\langle-\sqrt{u}+\sqrt[3]{u},0,0\right\rangle \\
{\bf r}_{w} &=& \left\langle 0,u^{2},0\right\rangle.
\end{eqnarray*}
}
%\vskip0.2cm
Además, $\left| {\bf r}_{u}\cdot \left( {\bf r}_{v}\times {\bf r}_{w}\right) \right| =
\left| \left( -\sqrt{u}+\sqrt[3]{u}\right) u^{2}\right|,$
por lo tanto, el volumen del sólido es
\[\int_{0}^{1}\int_{0}^{1}\int_{0}^{1}\left|\left(-\sqrt{u}+\sqrt[3]{u}\right) u^{2}\right|dudvdw= \frac{1}{70}.\]
%\end{tcolorbox}
%\vskip0.2cm
Si la densidad del sólido en mención en cualquier punto está dada por
$\delta \left( x,y,z\right) =xy^{2}z^{3},$ entonces la masa $M$ es:
%\vskip0.2cm
\[M= \ds \int_{0}^{1}\int_{x^{3}}^{x^{2}}\int_{0}^{z^{2}}xy^{2}z^{3}dydzdx=\ds \int_{0}^{1}\int_{\sqrt[2]{z}}^{\sqrt[3]{z}}\int_{0}^{z^{2}}xy^{2}z^{3}dydxdz= \frac{1}{2112}.\]
%\vskip0.2cm
Usando la parametrización escogida (\ref{32a}) se obtiene
%\vskip0.2cm
\[\delta \left( u,v,w\right) =\left( \sqrt{u}-\sqrt{u}v+v\sqrt[3]{u}\right) \left( wu^{2}\right) ^{2}\left( u\right) ^{3}.\]
%\vskip0.2cm
Entonces, la masa del sólido está dada por
%\vskip0.2cm
\[
\resizebox{\linewidth}{!}{
	$M=\int_{0}^{1}\int_{0}^{1}\int_{0}^{1} \left( \sqrt{u}(1-v)+v\sqrt[3]{u} \right) \left( wu^{2}\right) ^{2}u ^3 \left|\left( -\sqrt{u}+\sqrt[3]{u}\right) u^{2}\right| \,du\,dv$
}.
\]
%\vskip0.2cm
Los momentos con respecto a los planos coordenados están determinados por

\vspace{-10mm}
{\small
\begin{eqnarray*}
\mu_{yz} &=&\int _0^1\int _{x^3}^{x^2}\int _0^{z^2}x^2 y^2 z^3dydzdx =\frac{1}{2277} \\
\mu_{xz} &=& \int _0^1\int _{x^3}^{x^2}\int _0^{z^2}x y^3 z^3dydzdx =\frac{1}{3952} \\
\mu_{xy} &=& \int _0^1\int _{x^3}^{x^2}\int _0^{z^2}x y^2 z^4dydzdx=\frac{1}{2520}.
\end{eqnarray*}
}
\vskip0.2cm
Por lo tanto, las coordenadas $(\overline{x},\overline{y},\overline{z})$ del centro de masa del sólido están dadas por
$\left(\frac{\mu{yz}}{M}, \frac{\mu{xz}}{M},\frac{\mu{xy}}{M}\right)=\left(\frac{64}{69}, \frac{132}{247},\frac{88}{105}\right).$
\vskip0.2cm
\subsection*{Intersección de cilindros perpendiculares}
Consideremos el sólido acotado por dos cilindros de igual radio con ejes de simetría perpendicular. En particular, se puede suponer que poseen ecuaciones $x^{2}+z^{2}=r^{2}$ y
$x^{2}+y^{2}=r^{2},$ con $r>0$. 

%\vskip0.2cm
\vskip0.3cm
\noindent
\begin{minipage}{\textwidth}
	\centering
	\captionof{figure}{Sólido definido por $x^2 + z^2 \leq r^2$ \quad y \quad $x^2 + y^2 \leq r^2$.}
	\label{fig:bicilindro}
	\begin{minipage}{0.20\linewidth}
		\centering
		\includegraphics[width=\linewidth]{Img/51a.pdf}
	\end{minipage}
	\hspace{1em}
	\begin{minipage}{0.23\linewidth}
		\centering
		\includegraphics[width=\linewidth]{Img/51c.pdf}
	\end{minipage}
	\hspace{1em}
	\begin{minipage}{0.23\linewidth}
		\centering
		\includegraphics[width=\linewidth]{Img/51d.pdf}
	\end{minipage}
	
	\vspace{0.3em}
	\fuentepropia
\end{minipage}

\vskip0.3cm
Estos gráficos se obtienen ejecutando las siguientes instrucciones, en ellas se consideró $r=1$.


\vspace{0.6em}
\noindent
\begin{tcolorbox}[breakable,enhanced,title=\bf Instrucción en \verb"Mathematica", top=2pt,bottom=2pt,boxsep=0pt,before skip=2pt,after skip=0pt]
\begin{Verbatim}[formatcom=\letraverbatim]
ContourPlot3D[{x^2+y^2==1,x^2+z^2==1},{x,-1,1},
{y,-1,1},{z,-1,1},Mesh->False,PlotPoints->70,
ContourStyle->{Blend[{Cyan,Yellow},1/5],Yellow}]
RegionPlot3D[-Sqrt[1-y^2]<z<Sqrt[1-y^2]&&-Sqrt[1-
x^2]<y<Sqrt[1-x^2]&&-1<x<1,{x,-1,1},{y,-1,1},
{z,-1,1},PlotStyle->RGBColor[0.8,0.3,0.3],
PlotPoints->90,Mesh->False]
\end{Verbatim}
\end{tcolorbox}
\vspace{0.7em}
\noindent



En el plano $xy$, el sólido se proyecta como una circunferencia de radio $r,$ y en el plano
$yz$ el gráfico correspondiente es un cuadrado. Eliminado $x$ de las ecuaciones
$x^2+y^2=r^2$ y $x^2+z^2=r^2$ se obtienen las ecuaciones de las rectas $z = \pm y$,
razón por la cual, el cuadrado se divide en cuatro regiones.


\noindent
\begin{minipage}{\textwidth}
	\centering
	\captionof{figure}{Proyecciones del sólido definido por $x^2 + y^2 \leq r^2$, \quad $|x| \leq 1$ \quad y \quad $|z| \leq 1$.}
	\label{fig:cilindro_recortado}
	\begin{minipage}{0.25\linewidth}
		\centering
		\vspace{-10mm}
		\includegraphics[width=\linewidth]{Fig/454.pdf}
	\end{minipage}
	\hspace{1em}
	\begin{minipage}{0.25\linewidth}
		\centering
		\vspace{-10mm}
		\includegraphics[width=\linewidth]{Fig/455.pdf}
	\end{minipage}
	
	\vspace{-8mm}
	\fuentepropia
\end{minipage}
\vskip0.2cm

Estas figuras son útiles para plantear las integrales que determinan el volumen del sólido.
Integrando primero con respecto a $z$, esta recorre los puntos desde la parte negativa del cilindro con ecuación $x^2+z^2=r^2$ hasta su parte positiva y los límites en $x$ e $y$ se determinan con la circunferencia mostrada antes.
%\vskip0.2cm
\[\int_{-r}^{r}\int_{-\sqrt{r^{2}-x^{2}}}^{\sqrt{r^{2}-x^{2}}}\int_{-\sqrt{r^{2}-x^{2}}}^{\sqrt{r^{2}-x^{2}}}dzdydx=\int_{-r}^{r}\int_{-\sqrt{r^{2}-y^{2}}}^{\sqrt{r^{2}-y^{2}}}\int_{-\sqrt{r^{2}-x^{2}}}^{\sqrt{r^{2}-x^{2}}}dzdxdy= \frac{16}{3}r^{3}.\]
\vskip0.2cm
Para el planteamiento en  orden $dxdydz$, $x$ recorre los puntos entre la parte negativa del cilindro $x^2+z^2=r^2$ hasta su parte positiva y los límites para $y$ y $z$
se determinan con el cuadrado.
%\vskip0.2cm
\begin{multline*}
\int_{0}^{r}\int_{-z}^{z}\int_{-\sqrt{r^{2}-z^{2}}}^{\sqrt{r^{2}-z^{2}}}dxdydz+\int_{-r}^{0}\int_{z}^{-z}\int_{-\sqrt{r^{2}-z^{2}}}^{\sqrt{r^{2}-z^{2}}}dxdydz \\
+\int_{0}^{r}\int_{z}^{r}\int_{-\sqrt{r^{2}-y^{2}}}^{\sqrt{r^{2}-y^{2}}}dxdydz +\int_{-r}^{0}\int_{-z}^{r}\int_{-\sqrt{r^{2}-y^{2}}}^{\sqrt{r^{2}-y^{2}}}dxdydz\\
+\int_{0}^{r}\int_{-r}^{-z}\int_{-\sqrt{r^{2}-y^{2}}}^{\sqrt{r^{2}-y^{2}}}dxdydz+\int_{-r}^{0}\int_{-r}^{z}\int_{-\sqrt{r^{2}-y^{2}}}^{\sqrt{r^{2}-y^{2}}}dxdydz=\frac{16}{3}r^{3}.
\end{multline*}
\vskip0.2cm
Por la simetría de la figura, el planteamiento en los órdenes $dxdzdy$, $dydzdx$ y $dydxdz$ es muy similar al descrito en los casos ya vistos.

Estas integrales se evaluaron en \Verb"Mathematica" usando el comando \\ \Verb"Assuming",
debido a que el programa aporta una solución que depende de la naturaleza de $r$
(positiva, negativa, real o compleja). Esto se logra con las siguientes instrucciones.


\vspace{0.6em}
\noindent
\begin{tcolorbox}[breakable,enhanced,title=\bf Instrucción en \verb"Mathematica", top=2pt,bottom=2pt,boxsep=0pt,before skip=2pt,after skip=0pt]
\begin{Verbatim}[formatcom=\letraverbatim]
Integrate[1,{x,-r,r},{y,-Sqrt[r^2-x^2],Sqrt[r^2-x^2]},
{z,-Sqrt[r^2-x^2],Sqrt[r^2-x^2]}]
Assuming[r\[Element]Reals&&r>0,Integrate[1,{y,-r,r},
{x,-Sqrt[r^2-y^2],Sqrt[r^2-y^2]},{z,-Sqrt[r^2-x^2],
Sqrt[r^2-x^2]}]]

Assuming[r\[Element]Reals&&r>0,Integrate[1,{z,0,r},
{y,-z,z},{x,-Sqrt[r^2-z^2],Sqrt[r^2-z^2]}]+
Integrate[1,{z,-r,0},{y,z,-z},{x,-Sqrt[r^2-z^2],
Sqrt[r^2-z^2]}]+Integrate[1,{z,0,r},{y,z,r},
{x,-Sqrt[r^2-y^2],Sqrt[r^2-y^2]}]+Integrate[1,
{z,-r,0},{y,-z,r},{x,-Sqrt[r^2-y^2],Sqrt[r^2-y^2]}]
+Integrate[1,{z,0,r},{y,-r,-z},{x,-Sqrt[r^2-y^2],
Sqrt[r^2-y^2]}]+Integrate[1,{z,-r,0},{y,-r,z},
{x,-Sqrt[r^2-y^2],Sqrt[r^2-y^2]}]]
\end{Verbatim}
\end{tcolorbox}
\vspace{0.7em}
\noindent


En coordenadas cilíndricas, considerando el radio $a,$, el volumen se determina por la siguiente integral
%\vskip0.2cm
\[\int 0^{2 \pi }\int _0^a\int _{-\sqrt{a^2-r^2 \cos ^2(\theta)}}^{\sqrt{a^2-r^2\cos^2(\theta)}}rdzdrd\theta=\frac{16}{3}a^{3}.\]
%\vskip0.2cm
Esta se evalúa como sigue:


\vspace{0.6em}
\noindent
\begin{tcolorbox}[breakable,enhanced,title=\bf Instrucción en \verb"Mathematica", top=2pt,bottom=2pt,boxsep=0pt,before skip=2pt,after skip=0pt]
\begin{Verbatim}[formatcom=\letraverbatim]
Integrate[r,{\[Theta],0,2*Pi},{r,0,a},{z,-Sqrt[a^2
r^2*Cos[\[Theta]]^2],Sqrt[a^2-r^2*Cos[\[Theta]]^2]}]
\end{Verbatim}
\end{tcolorbox}
\vspace{0.7em}
\noindent


También se puede obtener el volumen usando el siguiente proceso.
\vskip0.3cm
%\begin{tcolorbox}[title=\bfseries Parametrizando el sólido]
El sólido se proyecta en el plano $xy$ como un círculo unitario; esta región se describe con las ecuaciones $x=u\cos v$, $y=u\sin v$, con $u \in [0,r]$, $v\in [0,2\pi]$.
Sobre el cilindro $x^2+z^2=r^2$ se satisface que $z =\pm \sqrt{r^2-u^2 \cos ^2v}.$
Por tanto, el sólido se parametriza mediante una combinación convexa de dos puntos de la forma
%\vskip0.2cm
\[P(u \cos (v),u \sin (v),-\sqrt{r^2-u^2 \cos ^2v})\] y \[Q(u \cos (v),u \sin (v),\sqrt{r^2-u^2 \cos ^2v}),\]
esto es, $(1-w)P+wQ.$

Simplificando esta expresión, se define la función
%\vskip0.2cm
\[{\boldsymbol \alpha}(u,v,w)= \langle u \cos (v),u \sin (v),(2w-1)\sqrt{r^2-u^2 \cos ^2(v)} \rangle,\]
%\vskip0.2cm
en donde $u\in [0,r]$, $v\in [0,2\pi]$ y $w\in [0,1].$
Derivando parcialmente se obtienen los siguientes vectores
\vspace{-3mm}
{\small
\begin{eqnarray*}
{\boldsymbol \alpha}_u &=& \langle \cos(v), \sin(v), \frac{u(1-2w)\cos^2(v)}{\sqrt{r^2-u^2 \cos^2(v)}} \rangle \\
{\boldsymbol \alpha}_v &=& \langle-u \sin(v), u\cos(v), \frac{u^2(2w-1)\sin(v)\cos(v)}{\sqrt{r^2-u^2 \cos^2v}}\rangle \\
{\boldsymbol \alpha}_w &=& \langle 0, 0, 2\sqrt{r^2-u^2 \cos^2v}\rangle.
\end{eqnarray*}
}
%\vskip0.2cm
Entonces,
{\small $| {\boldsymbol \alpha}_w \cdot ({\boldsymbol \alpha}_u \times {\boldsymbol \alpha}_v)| = 2u\sqrt{r^2-u^2 \cos^2(v)},$}
con lo cual el volumen del mencionado esta dado por
%\vskip0.2cm
{\small
\[\int_0^1 \int_0^{2\pi} \int_0^r 2u\sqrt{r^2-u^2 \cos^2(v)} du dv dw= \frac{16}{3}r^3\]
}
%\end{tcolorbox}
\vskip0.2cm
En \verb"Mathematica" el cálculo anterior se realiza con las siguientes instrucciones:


\vspace{0.6em}
\noindent
\begin{tcolorbox}[breakable,enhanced,title=\bf Instrucción en \verb"Mathematica", top=2pt,bottom=2pt,boxsep=0pt,before skip=2pt,after skip=0pt]
\begin{Verbatim}[formatcom=\letraverbatim]
{x,y,z}={u*Cos[v],u*Sin[v],Sqrt[r^2-x^2]};
R=Simplify[(1-w)*{x,y,-z}+w*{x,y,z}];
A=Assuming[r\[Element]Reals&&r>0,
Simplify[Dot[D[R,w],Cross[D[R,u],D[R,v]]]]]
Integrate[A,{v,0,2Pi},{u,0,r},{w,0,1}]
\end{Verbatim}
\end{tcolorbox}
\vspace{0.7em}
\noindent


Si la intersección se hace con tres cilindros del mismo radio, cuyos ejes de simetría son perpendiculares entre sí. La situación es un poco similar al procedimiento presentado cuando se trata de dos cilindros.
Consideremos los cilindros con ecuaciones
$x^2+y^2=r^2$, $y^2+z^2=r^2,$ y $x^2+z^2=r^2.$
Con las siguientes instrucciones se obtienen los gráficos presentados, en ellos se usó $r=1$.


\vspace{0.6em}
\noindent
\begin{tcolorbox}[breakable,enhanced,title=\bf Instrucción en \verb"Mathematica", top=2pt,bottom=2pt,boxsep=0pt,before skip=2pt,after skip=0pt]
\begin{Verbatim}[formatcom=\letraverbatim]
ContourPlot3D[{x^2+y^2==1,y^2+z^2==1,x^2+z^2==1},
{x,-1,1},{y,-1,1},{z,-1,1},Mesh->False,
ContourStyle->{Blend[{Cyan,Yellow},2/5],
RGBColor[0.9,0.8,0.6],Orange}]
RegionPlot3D[-Sqrt[1-y^2]<z<Sqrt[1-y^2]&&-Sqrt[1-
x^2]<y<Sqrt[1-x^2]&&-1<x<1,{x,-1,1},{y,-1,1},
{z,-1,1},PlotStyle->RGBColor[0.8,0.3,0.3],
PlotPoints->90,Mesh->False]
\end{Verbatim}
\end{tcolorbox}
\vspace{0.7em}
\noindent


El producto es el siguiente:

\vskip0.2cm
\noindent
\begin{minipage}{\textwidth}
	\centering
	\captionof{figure}{Intersección de tres cilindros: $x^2 + y^2 \leq r^2$, \quad $y^2 + z^2 \leq r^2$, \quad $x^2 + z^2 \leq r^2$}
	\label{fig:tricilindro}
	\begin{minipage}{0.21\linewidth}
		\centering
		\includegraphics[width=\linewidth]{Img/52a.pdf}
	\end{minipage}
	\begin{minipage}{0.25\linewidth}
		\centering
		\includegraphics[width=\linewidth]{Img/52b.pdf}
	\end{minipage}
	
	
	\vspace{0.3em}
	\fuentepropia
\end{minipage}
\vskip0.2cm





\noindent

La simetría de la figura con respecto a los planos coordenados, permite hacer el planteamiento en el primer octante o en la mitad de este.
La siguiente figura muestra la proyección del sólido en el plano $yz$.


\vskip0.2cm
\noindent
\begin{minipage}{\textwidth}
	\centering
	\captionof{figure}{Región definida por $y^2 + z^2 \leq r^2$, \quad $x > 0$ \quad y \quad $y > 0$}
	\label{fig:cuarto_disco_yz}
	\begin{minipage}{0.31\linewidth}
		\centering
		\vspace{-14mm}
		\includegraphics[width=\linewidth]{Fig/456.pdf}
	\end{minipage}
	
	\vspace{-9mm}
	\fuentepropia
\end{minipage}
\vskip0.2cm


Presentamos las integrales que permiten calcular el volumen, integrando primero con respecto a $z.$
%\vskip0.2cm
\vspace{-2mm}
{\small
\begin{multline*}
8\bigg( \int_{0}^{r/\sqrt{2}}\int_{0}^{x}\int_{0}^{\sqrt{r^{2}-x^{2}}}dzdydx+\int_{r/\sqrt{2}}^{r}\int_{0}^{\sqrt{r^{2}-x^{2}}}\int_{0}^{\sqrt{r^{2}-x^{2}}}dzdydx \\
+\int_{0}^{r/\sqrt{2}}\int_{x}^{\sqrt{r^{2}-x^{2}}}\int_{0}^{\sqrt{r^{2}-y^{2}}}dzdydx\bigg) = 8\left( 2-\sqrt{2}\right) r^{3} \\
8\bigg( \int_{0}^{r/\sqrt{2}}\int_{y}^{\sqrt{r^{2}-y^{2}}}\int_{0}^{\sqrt{r^{2}-x^{2}}}dzdxdy+\int_{0}^{r/\sqrt{2}}\int_{0}^{y}\int_{0}^{\sqrt{r^{2}-y^{2}}}dzdxdy \\
+\int_{r/\sqrt{2}}^{r}\int_{0}^{\sqrt{r^{2}-y^{2}}}\int_{0}^{\sqrt{r^{2}-y^{2}}}dzdxdy\bigg) = 8\left( 2-\sqrt{2}\right) r^{3}.
\end{multline*}
}
%\vskip0.2cm
Para evaluar las integrales anteriores, se ejecutaron estas instrucciones:


\vspace{0.6em}
\noindent
\begin{tcolorbox}[breakable,enhanced,title=\bf Instrucción en \verb"Mathematica", top=2pt,bottom=2pt,boxsep=0pt,before skip=2pt,after skip=0pt]
\begin{Verbatim}[formatcom=\letraverbatim]
Assuming[r\[Element]Reals&&r>0,8*(Integrate[1,{x,0,
r/Sqrt[2]},{y,0,x},{z,0,Sqrt[r^2-x^2]}]+Integrate[1,
{x,r/Sqrt[2],r},{y,0,Sqrt[r^2-x^2]},{z,0,Sqrt[r^2-
x^2]}]+Integrate[1,{x,0,r/Sqrt[2]},{y,x,Sqrt[r^2
x^2]},{z,0,Sqrt[r^2-y^2]}])]
Assuming[r\[Element]Reals&&r>0,8*(Integrate[1,{y,0,
r/Sqrt[2]},{x,y,Sqrt[r^2-y^2]},{z,0,Sqrt[r^2-x^2]}]+
Integrate[1,{y,0,r/Sqrt[2]},{x,0,y},{z,0,Sqrt[r^2-
y^2]}]+Integrate[1,{y,r/Sqrt[2],r},{x,0,Sqrt[r^2-
y^2]},{z,0,Sqrt[r^2-y^2]}])]
\end{Verbatim}
\end{tcolorbox}
\vspace{0.7em}
\noindent



Un cambio a coordenadas cilíndricas, usando $a$ como el radio de los cilindros y observando solamente la mitad del primer octante, nos lleva al mismo valor
\[16\int_{0}^{\pi/4}\int_{0}^{a}\int_{0}^{\sqrt{a^{2}-r^{2}\cos^{2}\theta}}rdzdrd\theta=8(2-\sqrt{2})a^3.\]

\section{Dos esferas} \index{Esfera}
Las ecuaciones $x^2+y^2+(z-1)^2=5$ y $x^2+y^2+z^2=4$
representan dos esferas que se intersectan en $z=0$. Consideraremos el sólido acotado por los hemisferios negativos en $z$.

\begin{figure}[H]
\centering
\caption{Región definida por $2\rho \cos \varphi + 4 \leq \rho^2 \leq 4$.}
\label{fig:region_polar_desplazada}
\includegraphics[height=4.5cm]{Img/90b.pdf}
\hspace{0.2cm}
\includegraphics[height=4.5cm]{Img/90a.pdf}
\fuentepropia
\end{figure}

Estos gráficos se obtienen con las siguientes instrucciones:


\vspace{0.6em}
\noindent
\begin{tcolorbox}[breakable,enhanced,title=\bf Instrucción en \verb"Mathematica", top=2pt,bottom=2pt,boxsep=0pt,before skip=2pt,after skip=0pt]
\begin{Verbatim}[formatcom=\letraverbatim]
ContourPlot3D[{x^2+y^2+(z-1)^2==5,x^2+y^2+z^2==4},
{x,-3,3},{y,-3,3},{z,-3,4},Mesh->False,
ContourStyle->{Directive[Cyan,Opacity[0.8]],
Directive[Orange,Opacity[0.9]]},PlotPoints->60]
\end{Verbatim}
\end{tcolorbox}
\vspace{0.7em}
\noindent


En coordenadas rectangulares, el volumen está dado por la integral triple en el orden $dzdydx:$
%\vskip0.2cm
\[\int_{-2}^{2}\int_{-\sqrt{4-x^{2}}}^{\sqrt{4-x^{2}}}\int_{-\sqrt{4-x^{2}-y^{2}}}^{1-\sqrt{5-x^{2}-y^{2}}}dzdydx=\frac{10}{3} \left(3-\sqrt{5}\right) \pi\]
%\vskip0.2cm
La proyección del sólido en el plano $xz$ se muestra a continuación.
\vskip0.2cm

\noindent
\begin{minipage}{\textwidth}
	\centering
	\captionof{figure}{Región definida por $2z + 4 \leq x^2 + z^2 \leq 4$, \quad $x^2 + (z - 1)^2 \leq 4$.}
	\label{fig:region_corona_desplazada_2}
	\begin{minipage}{0.35\linewidth}
		\centering
		\vspace{-20mm}
		\includegraphics[width=\linewidth]{Fig/457.pdf}
	\end{minipage}
	
	
	\vspace{-18mm}
	\fuentepropia
\end{minipage}

\vskip0.3cm
Con esta figura se plantea la integral que determina el volumen del sólido en el orden $dydzdx$, se usa la abreviatura
$a=\sqrt{5-x^{2}-(z-1)^{2}}.$
{\small
\begin{multline*}
\hskip-0.2cm \int_{-2}^{2}\int_{1-\sqrt{5-x^{2}}}^{0}\int_{-\sqrt{4-x^{2}-z^{2}}}^{-a}dydzdx +\int_{-2}^{2}\int_{1-\sqrt{5-x^{2}}}^{0}\int_{a}^{\sqrt{4-x^{2}-z^{2}}}dydzdx \\
+\int_{-2}^{2}\int_{-\sqrt{4-x^{2}}}^{1-\sqrt{5-x^{2}}}\int_{-\sqrt{4-x^{2}-z^{2}}}^{\sqrt{4-x^{2}-z^{2}}}dydzdx \approx 7.999.
\end{multline*}
}
%\vskip0.2cm
Por la simetría de la figura con respecto al plano $yz$, la integral en el orden $dxdzdy$ es similar.
Las integrales anteriores se evaluaron usando \verb"Mathematica" ejecutando las siguientes instrucciones:


\vspace{0.6em}
\noindent
\begin{tcolorbox}[breakable,enhanced,title=\bf Instrucción en \verb"Mathematica", top=2pt,bottom=2pt,boxsep=0pt,before skip=2pt,after skip=0pt]
\begin{Verbatim}[formatcom=\letraverbatim]
Integrate[1,{x,-2,2},{y,-Sqrt[4-x^2],Sqrt[4-x^2]},
{z,-Sqrt[4-x^2-y^2],1-Sqrt[5-x^2-y^2]}]
Integrate[r,{\[Theta],0,2Pi},{r,0,2},
{z,-Sqrt[4-r^2],1-Sqrt[5-r^2]}]
NIntegrate[2,{x,-2,2},{z,1-Sqrt[5-x^2],0},
{y,-Sqrt[4-x^2-z^2],-Sqrt[5-x^2-(z-1)^2]}]
+NIntegrate[2,{x,-2,2},{z,1-Sqrt[5-x^2],0},
{y,Sqrt[5-x^2-(z-1)^2],Sqrt[4-x^2-z^2]}]
+NIntegrate[1,{x,-2,2},{z,-Sqrt[4-x^2],1-Sqrt[5-x^2]},
{y,-Sqrt[4-x^2-z^2],Sqrt[4-x^2-z^2]}]
\end{Verbatim}
\end{tcolorbox}
\vspace{0.7em}
\noindent


En coordenadas cilíndricas, esta integral se escribe como
%\vskip0.2cm
\[\int_{0}^{2\pi }\int_{0}^{2}\int_{-\sqrt{4-r^{2}}}^{1-\sqrt{5-r^{2}}}rdzdrd\theta=\frac{10}{3}\pi\left(3-\sqrt{5}\right).\]
%\vskip0.2cm
En coordenadas esféricas, la ecuación $x^{2}+y^{2}+z^{2}=4$ se escribe como $\rho =2$ y la ecuación
$x^{2}+y^{2}+\left( z-1\right) ^{2}=5$ se transforma en $\rho ^{2}-2\rho \cos \phi=4,$ de donde
$\rho =\cos \phi +\sqrt{\cos ^{2}\phi +4}$.
La intersección entre las esferas se determina por $\phi = \pi /2.$
Entonces, el volumen del sólido esta dado por
%\vskip0.2cm
\[\int_{\pi /2}^{\pi }\int_{0}^{2\pi }\int_{\cos \phi +\sqrt{\cos ^{2}\phi +4}}^{2}\rho ^{2}\sin \phi d\rho d\theta d\phi = \frac{10}{3}\pi \left( 3-\sqrt{5}\right).\]
%\vskip0.2cm
La evaluación de esta integral se puede hacer como sigue.


\vspace{0.6em}
\noindent
\begin{tcolorbox}[breakable,enhanced,title=\bf Instrucción en \verb"Mathematica", top=2pt,bottom=2pt,boxsep=0pt,before skip=2pt,after skip=0pt]
\begin{Verbatim}[formatcom=\letraverbatim]
Integrate[\[Rho]^2*Sin[\[Phi]],{\[Phi],Pi/2,Pi},
{\[Theta],0,2*Pi},{\[Rho],Cos[\[Phi]]+Sqrt[4+
Cos[\[Phi]]^2],2}]
\end{Verbatim}
\end{tcolorbox}
\vspace{0.7em}
\noindent


El siguiente procedimiento permite obtener el mismo valor.

%\begin{tcolorbox}[title=\bfseries Parametrizando el sólido]
La curva de intersección se determina por $z=0$ y está dada por $x^{2}+y^{2}=4,$, cuyo interior se parametriza como
$x=u\cos v,$ $y=u\sin v,$ con $u\in [0,2],$ $v\in [0,2\pi ]$.
La coordenada $z$, en las dos esferas, es
$z= -\sqrt{4-u^{2}}$, $z= 1-\sqrt{5-u^{2}}$, respectivamente.
El sólido se parametriza formando una combinación convexa entre dos puntos de las esferas con iguales coordenadas $x$ e $y.$ Esto es,
$(1-w) \left\langle u\cos v,u\sin v,-\sqrt{4-u^{2}}\right\rangle +w\left\langle u\cos v,u\sin v,1-\sqrt{5-u^{2}}\right\rangle$
que se simplifica y define la función
{\small
\[
\resizebox{\linewidth}{!}{
	$\mathbf{r}(u,v,w)=\left\langle u\cos v,u\sin v,w\left(1-\sqrt{5-u^{2}}\right)-(1-w)\sqrt{4-u^{2}}\right\rangle.$
}
\]
}
Esta función nos permite obtener los siguientes vectores y el elemento diferencial de volumen
{\small
\begin{eqnarray*}
\mathbf{r} &=& \left\langle \cos v, \sin v, \frac{1}{1+u^2} \sqrt{4-u^2} \right\rangle + \frac{u}{\sqrt{5-u^2}} \mathbf{r} \\
\mathbf{r}_u &=& \left\langle \text{(completa aquí)} \right\rangle \\
\mathbf{r}_v &=& \left\langle -\sin v, \cos v, 0 \right\rangle \\
\mathbf{r}_u &=& \left\langle 0, 0, \sqrt{4-u^2} \right\rangle \\
dV &=& \left| \mathbf{r}_u \times \mathbf{r}_v \right| = \left| \sqrt{4-u^2} \right|.
\end{eqnarray*}
}
Entonces, el volumen del sólido está dado por
\vspace{-2mm}
\[\int_{0}^{1}\int_{0}^{2\pi }\int_{0}^{2}\left\vert u\left( \sqrt{5-u^{2}}-
\sqrt{4-u^{2}}-1\right)\right\vert dudvdw=\frac{10(3-\sqrt{5})}{3}\pi.\]
%\end{tcolorbox}

\section{Sólidos con cilindros cardióidicos} \index{Cardioide}

En esta sección se consideran cilindros cuyos cortes perpendiculares al eje de simetría son cardioides.

\index{Cilindro cardioidico}
\subsection*{Cilindros cardióidicos y el plano z=2-x/2}

En el plano, las ecuaciones $r = 1-\cos(t)$ y $r = 2- 2\cos(t)$ representan dos cardioides, mientras que en el espacio las mismas ecuaciones representan dos cilindros cardióidecos.



\vskip0.3cm
\noindent
\begin{minipage}{\textwidth}
	\centering
	\captionof{figure}{Región delimitada por $r = 1 - \cos(t)$, $r = 2 - 2\cos(t)$, \quad y \quad $z = 2 - \frac{x}{2}$.}
	\label{fig:cardioide_revolucion}
	\begin{minipage}{0.28\linewidth}
		\centering
		\vspace{-10mm}
		\includegraphics[width=\linewidth]{Fig/458.pdf}
	\end{minipage}
	\hspace{1em}
	\begin{minipage}{0.25\linewidth}
		\centering
		\vspace{-10mm}
		\includegraphics[width=\linewidth]{Img/36a.pdf}
	\end{minipage}
	
	\vspace{-8mm}
	\fuentepropia
\end{minipage}


\vskip0.3cm
Consideremos el sólido acotado por las figuras de las expresiones $z = 2-x/2,$
$z=0$ y los cilindros $r = 1- \cos(t)$ y $r = 2-2\cos(t).$




\vskip0.2cm
\noindent
\begin{minipage}{\textwidth}
	\centering
	\captionof{figure}{Sólido definido por $1 - \cos(t) \leq r \leq 2 - 2\cos(t)$ \quad y \quad $0 \leq z \leq 2 - \frac{x}{2}$}
	\label{fig:cardioide_z_inclinado}
	\begin{minipage}{0.22\linewidth}
		\centering
		\includegraphics[width=\linewidth]{Img/36c.pdf}
	\end{minipage}
	\begin{minipage}{0.22\linewidth}
		\centering
		\includegraphics[width=\linewidth]{Img/36d.pdf}
	\end{minipage}
	\begin{minipage}{0.22\linewidth}
		\centering
		\includegraphics[width=\linewidth]{Img/36e.pdf}
	\end{minipage}
	
	\vspace{0.3em}
	\fuentepropia
\end{minipage}
\vskip0.2cm



Las figuras de las cardioides en una animación que muestra el recorrido de esta curva se obtienen ejecutando las siguientes instrucciones:


\vspace{0.6em}
\noindent
\begin{tcolorbox}[breakable,enhanced,title=\bf Instrucción en \verb"Mathematica", top=2pt,bottom=2pt,boxsep=0pt,before skip=2pt,after skip=0pt]
\begin{Verbatim}[formatcom=\letraverbatim]
Manipulate[ParametricPlot[{{(1-Cos[t])*Cos[t],
(1-Cos[t])*Sin[t]},{2*(1-Cos[t])*Cos[t],
2*(1-Cos[t])*Sin[t]}},{t,0,k},PlotRange->{{-4,1},
{-3,3}}],{k,Pi/20,2*Pi}]
\end{Verbatim}
\end{tcolorbox}
\vspace{0.7em}
\noindent



Las tres superficies básicas que intervienen, se consiguen con las siguientes instrucciones:



\vspace{0.6em}
\noindent
\begin{tcolorbox}[breakable,enhanced,title=\bf Instrucción en \verb"Mathematica", top=2pt,bottom=2pt,boxsep=0pt,before skip=2pt,after skip=0pt]
\begin{Verbatim}[formatcom=\letraverbatim]
a=ParametricPlot3D[{{(1-Cos[t])*Cos[t],
(1-Cos[t])*Sin[t],u},{(2-2*Cos[t])*Cos[t],
(2-2*Cos[t])*Sin[t],u}},{u,0,4},{t,0,2*Pi},
Mesh->False,PlotPoints->80,PlotStyle->{Orange,
RGBColor[0.7,0.1,0.1]}];
b=ContourPlot3D[{z==2-x/2},{x,-5,5},{y,-5,5},{z,-2,4},
Mesh->False,ContourStyle->RGBColor[0.2,0.4,0.3]];
Show[a,b]
\end{Verbatim}
\end{tcolorbox}
\vspace{0.7em}
\noindent


Además, la parte del plano que se halla entre los dos cilindros cardioidicos se consigue con estas sentencias:


\vspace{0.6em}
\noindent
\begin{tcolorbox}[breakable,enhanced,title=\bf Instrucción en \verb"Mathematica", top=2pt,bottom=2pt,boxsep=0pt,before skip=2pt,after skip=0pt]
\begin{Verbatim}[formatcom=\letraverbatim]
ParametricPlot3D[{(1+u)*(1-Cos[t])*Cos[t],
(1+u)*(1-Cos[t])*Sin[t],2-(1+u)*(1-Cos[t])*Cos[t]/2},
{u,0,1},{t,0,2*Pi},Mesh->False,
PlotStyle->RGBColor[0.6,0.3,0.3],PlotPoints->80]
\end{Verbatim}
\end{tcolorbox}
\vspace{0.7em}
\noindent



La figura del sólido se consigue utilizando coordenadas cilíndricas y ejecutando las siguientes tres instrucciones:



\vspace{0.6em}
\noindent
\begin{tcolorbox}[breakable,enhanced,title=\bf Instrucción en \verb"Mathematica", top=2pt,bottom=2pt,boxsep=0pt,before skip=2pt,after skip=0pt]
\begin{Verbatim}[formatcom=\letraverbatim]
{r,t}={Sqrt[x^2+y^2],ArcTan[x,y]};
RegionPlot3D[1-Cos[t]<r<2-2*Cos[t]&&0<z<2-x/2,
{x,-4,2},{y,-4,4},{z,0,4},PlotPoints->90,Mesh->False]
\end{Verbatim}
\end{tcolorbox}
\vspace{0.7em}
\noindent



En el plano $xy,$ los cilindros cardióidicos se pueden parametrizar, haciendo
$x = (1 -\cos(t))\cos(t),$ $y = (1 - \cos(t)) \sin(t),\,t \in[0, 2\pi];$ esto para el cilindro menor. Mientras que $x = (2 - 2\cos(t))\cos(t),$ $y = (2 - 2\cos(t))\sin(t)$ se usarán para el cilindro mayor.
Lo anterior, nos permite definir las funciones vectoriales que representan las curvas de intersección de los cilindros con el plano $xy$ ($z = 0$), que son
%\vskip0.2cm
\begin{eqnarray*}
\overrightarrow{\alpha} (t) &=& \langle(1-\cos(t))\cos(t), (1-\cos(t))\sin(t), 0\rangle ,\,\, t \in[0, 2\pi] \\
\overrightarrow{\beta} (t)& = &\langle(2-2\cos(t))\cos(t), (2-2\cos(t))\sin(t), 0\rangle ,\,\, t \in[0, 2\pi.]
\end{eqnarray*}
%\vskip0.2cm
Con estas funciones formamos la combinación convexa $(1-u)\overrightarrow{\alpha}(t)+u\overrightarrow{\beta}(t)$que determina la región contenida entre las dos curvas (en el plano $z=0$). Esto es:
%\vskip0.2cm
{\small
\begin{equation}\label{36a}
\overrightarrow{r}_1(t,u)=(1+u)\left\langle 2 \sin ^2\left(\frac{t}{2}\right) \cos (t),\sin (t) (1-\cos (t)),0 \right\rangle,
\end{equation}
}
%\vskip0.2cm
para $u\in[0, 1]$ y $t \in[0, 2\pi]$; proyectando esta región sobre el plano $z=2-x/2$, obtenemos la función
{\small
\begin{eqnarray}
\overrightarrow{r}_2(t,u)=\bigg\langle 2 (u+1) \sin ^2\left(\frac{t}{2}\right) \cos(t),(u+1) \sin (t) (1-\cos (t)),\nonumber\\
2-(u+1) \sin^2\left(\frac{t}{2}\right) \cos (t) \bigg\rangle \label{36b}
\end{eqnarray}
}

\vskip0.2cm
\noindent
\begin{minipage}{\textwidth}
	\centering
	\captionof{figure}{Sólido definido por $1 - \cos t \leq r \leq 2 - 2\cos t$, \quad $0 \leq z \leq 2 - \frac{x}{2}$.}
	\label{fig:solido_cardioide_inclinado}
	\begin{minipage}{0.35\linewidth}
		\centering
		\vspace{-15mm}
		\includegraphics[width=\linewidth]{Fig/459.pdf}
	\end{minipage}
	\hspace{1em}
	\begin{minipage}{0.35\linewidth}
		\centering
		\vspace{-15mm}
		\includegraphics[width=\linewidth]{Fig/460.pdf}
	\end{minipage}
	
	\vspace{-8mm}
	\fuentepropia
\end{minipage}


\vskip0.2cm
Un punto $Q,$ sobre la base del sólido, tendrá la forma de la expresión \ref{36a}, mientras que un punto $P$ del sólido sobre el plano $z=2-x/2$,
tiene la forma \ref{36b}. Con el fin de parametrizar la región en el espacio, formamos una combinación convexa de los puntos $P$ y $Q$, de la forma
$(1-w)\overrightarrow{r}_1(t,u)+w\overrightarrow{r}_2(t,u).$ Esto define la función
\begin{multline*}
{\bf r}(t,u,w) = \bigg \langle 2 (u+1) \sin ^2\left(\frac{t}{2}\right)\cos(t),-(u+1) \sin (t) (\cos (t)-1),\\
w \left(2-(u+1)\sin ^2\left(\frac{t}{2}\right) \cos(t)\right)\bigg \rangle,
\end{multline*}
con $t\in [0,2\pi]$, $u\in [0,1]$, $w\in [0,1].$
Ahora obtenemos los vectores
{\small
	\[
	\mathbf{r}_{t} = (1+u)\left\langle \sin (2 t)-\sin (t),2\sin^2\left(\tfrac{t}{2}\right) (2\cos(t)+1),
	\frac{1}{2} w (\sin (t)-\sin (2 t)) \right\rangle,
	\]
	\[
	\mathbf{r}_{u} = \left\langle2 \sin ^2\left(\tfrac{t}{2}\right) \cos (t),-\sin(t)(\cos (t)+1),-w \sin^2\left(\tfrac{t}{2}\right) \cos (t) \right\rangle,
	\]
	\[
	\mathbf{r}_{w} = \left\langle 0,0,2-(u+1) \sin ^2\left(\tfrac{t}{2}\right) \cos(t)\right\rangle.
	\]
}
Además,
{\small
\[\left| {\bf r}_{t}\cdot \left( {\bf r}_{u}\times {\bf r}_{w}\right) \right| =
\left|(u+1) \sin^4\left(\frac{t}{2}\right) ((u+1)(\cos (2 t)-2 \cos (t))+u+9) \right|.\]
}
Por último, el volumen está dado por:
\vskip0.2cm
\[
\int_{0}^{1} \int_{0}^{1} \int_{0}^{2\pi}
\left| \mathbf{r}_{t} \cdot \left( \mathbf{r}_{u} \times \mathbf{r}_{w} \right) \right|
\, dt\, du\, dw = \frac{107}{8}\pi.
\]
\vskip0.2cm
La parametrización y el cálculo anterior, se pueden hacer en \verb"Mathematica" con ayuda de las siguientes instrucciones


\vspace{0.6em}
\noindent
\begin{tcolorbox}[breakable,enhanced,title=\bf Instrucción en \verb"Mathematica", top=2pt,bottom=2pt,boxsep=0pt,before skip=2pt,after skip=0pt]
\begin{Verbatim}[formatcom=\letraverbatim]
{x1,y1}={Cos[t]*(1-Cos[t]),Sin[t]*(1-Cos[t])};
r=(1-u)*{x1,y1,0}+u*{x2,y2,0};
s=(1-u)*{x1,y1,2-x1/2}+u*{2*x1,2*y1,2-x1};
v=(1-w)*s+w*r;
A=Dot[D[v,t],Cross[D[v,u],D[v,w]]];
Integrate[Abs[A],{t,0,2*Pi},{u,0,1},{w,0,1}]
\end{Verbatim}
\end{tcolorbox}
\vspace{0.7em}
\noindent


Otra manera de realizar dicho cálculo es el uso de coordenadas cilíndricas, de esta manera el volumen está dado por:
\[\int_{0}^{2\pi }\int_{1-\cos\theta}^{2-2\cos\theta}\int_{0}^{2-\frac{1}{2}r\cos\theta}rdzdrd\theta = \frac{107}{8}\pi.\]

\subsection*{Cilindros cardioidicos y el plano z=2-y/2}

Las ecuaciones $r = 1 - \cos(t)$ y $r = 2 -2 \cos(t)$ representan dos cardioides en el plano, mientras que en el espacio las mismas ecuaciones representan dos cilindros cardioides.
\vskip0.2cm

\noindent
\begin{minipage}{\textwidth}
	\centering
	\captionof{figure}{Región delimitada por $r = 1 - \cos(t)$, $r = 2 - 2\cos(t)$, \quad y \quad $z = 2 - \frac{y}{2}$.}
	\label{fig:cardioide_y_inclinado}
	\begin{minipage}{0.26\linewidth}
		\centering
		\vspace{-10mm}
		\includegraphics[width=\linewidth]{Fig/461.pdf}
	\end{minipage}
	\hspace{1em}
	\begin{minipage}{0.24\linewidth}
		\centering
		\vspace{-10mm}
		\includegraphics[width=\linewidth]{Img/35e.pdf}
	\end{minipage}
	
	\vspace{-5mm}
	\parbox{\linewidth}{\centering\fontsize{9pt}{10pt}\selectfont\textbf{Fuente:} elaboración propia.}
\end{minipage}
\vskip0.2cm


Consideremos el sólido acotado por las figuras de las expresiones $z = 2 -y/2,$
$z=0$ y los cilindros $r=1-\cos(t)$ y $r=2- 2 \cos(t).$

\begin{figure}[H]
\centering
\caption{Sólido definido por $1 - \cos(t) \leq r \leq 2 - 2\cos(t)$ \quad y \quad $0 \leq z \leq 2 - \frac{y}{2}$.}
\label{fig:cardioide_y_inclinado_2} \includegraphics[height=4.5cm]{Img/35a.pdf}
\hspace{0.3cm}
\includegraphics[height=4.5cm]{Img/35b.pdf}
\hspace{0.3cm}
\includegraphics[height=4.5cm]{Img/35c.pdf}
\parbox{\linewidth}{\centering\fontsize{8pt}{10pt}\selectfont\textbf{Fuente:} elaboración propia.}
\end{figure}

Las figuras se obtienen en coordenadas cilíndricas, definiendo $r$ y $t=\theta$.


\vspace{0.6em}
\noindent
\begin{tcolorbox}[breakable,enhanced,title=\bf Instrucción en \verb"Mathematica", top=2pt,bottom=2pt,boxsep=0pt,before skip=2pt,after skip=0pt]
\begin{Verbatim}[formatcom=\letraverbatim]
{r,t}={Sqrt[x^2+y^2],ArcTan[x,y]};
RegionPlot3D[1-Cos[t]<r<2-2*Cos[t]&&0<z<2-y/2,
{x,-4,2},{y,-4,4},{z,0,4},PlotPoints->90,Mesh->False]
\end{Verbatim}
\end{tcolorbox}
\vspace{0.7em}
\noindent


De igual manera, se pueden usar las siguientes instrucciones:


\vspace{0.6em}
\noindent
\begin{tcolorbox}[breakable,enhanced,title=\bf Instrucción en \verb"Mathematica", top=2pt,bottom=2pt,boxsep=0pt,before skip=2pt,after skip=0pt]
\begin{Verbatim}[formatcom=\letraverbatim]
a=ParametricPlot3D[{{(1-Cos[t])*Cos[t],
(1-Cos[t])*Sin[t],u},{(2-2*Cos[t])*Cos[t],
(2-2*Cos[t])*Sin[t],u}},{u,0,3.5},{t,0,2*Pi},
PlotStyle->{Orange,RGBColor[0.7,0.1,0.1]},
Mesh->False,PlotPoints->80];
b=ContourPlot3D[{z==2-y/2},{x,-5,2},{y,-4,4},{z,-2,
3.5},Mesh->False,ContourStyle->RGBColor[0.2,0.4,0.3]];
Show[a,b]
\end{Verbatim}
\end{tcolorbox}
\vspace{0.7em}
\noindent


El cilindro menor se parametriza por $x = (1-\cos(t)) \cos(t),$ $y = (1-\cos(t))\sin(t)$
y para el cilindro mayor se hace $x = (2- 2\cos(t)) \cos(t),$
$y = (2 -2\cos(t))\sin(t)$, con $t \in [0, 2\pi]$. Las siguientes son las funciones vectoriales que representan las curvas de intersección de los cilindros con el plano $xy$ $(z=0)$:
%\vskip0.2cm
{\small
\begin{eqnarray*}
\overrightarrow{\alpha }(t)&=&\langle ( 1-\cos ( t)) \cos(t),( 1-\cos(t)) \sin(t),0\rangle \\
\overrightarrow{\beta }(t)&=&\langle ( 2-2\cos(t))\cos(t),( 2-2\cos(t)) \sin(t),0\rangle,
\end{eqnarray*}
}
%\vskip0.2cm
con $t\in [0,2\pi ].$ Posteriormente, formamos una combinación convexa de ellas que parametriza la región contenida entre las dos curvas.
%\vskip0.2cm
{\small
\begin{eqnarray}
\overrightarrow{r}_1(t,u) &=& ( 1-u) \overrightarrow{\alpha }(t)+u\overrightarrow{\beta } (t)\nonumber \\
&=& (u+1)\langle 2 \sin ^2\left(\frac{t}{2}\right)\cos(t),\sin(t)(1-\cos(t)),0\rangle, \label{625x}
\end{eqnarray}
}
%\vskip0.2cm
para $u\in [0,1]$ y $t\in [0,2\pi ].$
Proyectando esta región sobre el plano $z=2-y/2$, esta superficie se parametriza como
%\vskip0.2cm
{\small
\begin{eqnarray}
\overrightarrow{r}_2(t,u) &=& \big \langle 2(u+1)\sin^2\left(\frac{t}{2}\right)\cos(t),(u+1)\sin(t)(1-\cos(t)), \nonumber\\
&& \hskip3.25 cm \frac{1}{2}(u+1)\sin(t)(\cos(t)-1)+2 \big \rangle \label{626x}
\end{eqnarray}
}
\vskip0.2cm
Cuya figura es la siguiente
\vskip0.2cm
\noindent
\begin{minipage}{\textwidth}
	\centering
	\captionof{figure}{Superficie $r_2(t,u)$ definida sobre región cardioidal con altura $z = 2 - \frac{r \sin(t)}{2}$.}
	\label{fig:superficie_cardioide_inclinada}
	\begin{minipage}{0.26\linewidth}
		\centering
		\vspace{-12mm}
		\includegraphics[width=\linewidth]{Img/35d.pdf}
	\end{minipage}
	\hspace{1em}
	\begin{minipage}{0.26\linewidth}
		\centering
		\vspace{-12mm}
		\includegraphics[width=\linewidth]{Fig/462.pdf}
	\end{minipage}
	
	\vspace{-8mm}
	\fuentepropia
\end{minipage}

\vskip0.2cm
Un punto $Q$ sobre la base del sólido tendrá la forma de la expresión (\ref{625x}), mientras que un punto $P$ del sólido sobre el plano $z=2-y/2$, tiene la forma \ref{626x}.
\vskip0.2cm

\noindent
\begin{minipage}{\textwidth}
	\centering
	\captionof{figure}{Sólido definido por $1 - \cos(t) \leq r \leq 2 - 2\cos(t)$ \quad y \quad $0 \leq z \leq 2 - \frac{y}{2}$.}
	\label{fig:solido_cardioide_y_inclinado}
	\begin{minipage}{0.35\linewidth}
		\centering
		\vspace{-12mm}
		\includegraphics[width=\linewidth]{Fig/463.pdf}
	\end{minipage}
	\hspace{1em}
	\begin{minipage}{0.35\linewidth}
		\centering
		\vspace{-12mm}
		\includegraphics[width=\linewidth]{Fig/464.pdf}
	\end{minipage}
	
	\vspace{-8mm}
	\fuentepropia
\end{minipage}


\vskip0.2cm

Con el fin de parametrizar la región en el espacio formamos una combinación convexa de los puntos $P$ y $Q$ de la forma $(1-w)\overrightarrow{r}_1(t,u)+w \overrightarrow{r}_2(t,u),$ que al simplificarse define la función

\vspace{-8mm}
\begin{multline}
{\bf r}( t,u,w) =\bigg\langle 2 (u+1) \sin ^2\left(\frac{t}{2}\right)\cos(t),(u+1)\sin (t) (1-\cos(t)), \\
w \left(\frac{1}{2} (u+1) \sin (t) (\cos (t)-1)+2\right)\bigg \rangle
\end{multline}
De la cual se obtienen los siguientes vectores
\begin{equation*}
	\begin{aligned}
		{\bf r}_{t} &= (u+1)\Big \langle (\sin(2t)-\sin(t)), \\
		&\qquad 2\sin^2\left(\frac{t}{2}\right)(2\cos (t)+1), w \sin ^2\left(\frac{t}{2}\right) (2 \cos (t)+1) \Big \rangle \\
		{\bf r}_{u} &= \Big\langle 2\sin^2\left(\frac{t}{2}\right)\cos(t),\sin(t)(-(\cos(t)-1)),\\
		&\qquad \frac{1}{2}w\sin(t)(\cos(t)-1)\Big\rangle \\
		{\bf r}_{w} &= \Big\langle 0,0,\frac{1}{2} (u+1) \sin (t) (\cos (t)-1)+2 \Big\rangle.
	\end{aligned}
\end{equation*}
Con los cuales se obtiene el elemento diferencial de volumen
$|{\bf r}_{t}\cdot({\bf r}_{u}\times{\bf r}_{w})|=\big|-(u+1)\sin^4\left(\frac{t}{2}\right)(8-2(u+1)\sin(t)+(u+1)\sin(2t)) \big| .$
Finalmente, el volumen del sólido está dado por
\vskip0.2cm
{\small
\[\int_{0}^{1}\int_{0}^{1}\int_{0}^{2\pi}\bigg|-(u+1)\sin^4\left(\frac{t}{2}\right)(8-(u+1)(2\sin(t)-\sin(2t)))\bigg| dtdudw= 9\pi\]
}
\vskip0.2cm
Este valor también se puede obtener utilizando coordenadas cilíndricas.
\vskip0.2cm
\[\int_{0}^{2\pi } \int_{1-\cos \theta }^{2-2\cos \theta }\int_{0}^{2-\frac{r}{2}\sin \theta }rdzdrd\theta = 9\pi.\]
\vskip0.2cm
Los cálculos anteriores se pueden realizar en \verb"Mathematica" ejecutando las siguientes sentencias.



\vspace{0.6em}
\noindent
\begin{tcolorbox}[breakable,enhanced,title=\bf Instrucción en \verb"Mathematica", top=2pt,bottom=2pt,boxsep=0pt,before skip=2pt,after skip=0pt]
\begin{Verbatim}[formatcom=\letraverbatim]
r1={Cos[t]*(1-Cos[t]),Sin[t]*(1-Cos[t])};
r2={Cos[t]*2*(1-Cos[t]),2*Sin[t]*(1-Cos[t])};
r={x,y}=Simplify[(1-u)*r1+u*r2];
R=Simplify[(1-w)*{x,y,0}+w*{x,y,2-y/2}];
A=Simplify[Dot[D[R,t],Cross[D[R,u],D[R,w]]]]
Integrate[Abs[A],{t,0,2*Pi},{u,0,1},{w,0,1}]
\end{Verbatim}
\end{tcolorbox}
\vspace{0.7em}
\noindent


\subsection*{Dos cilindros cardioidicos y un paraboloide}
El paraboloide con ecuación $4z = x^2 + y^2$, el plano $xy$ y los cilindros cuya base son las cardioiodes con ecuaciones $r=1-\cos t$ y $r=2 - 2 \cos t,$ acotan un sólido.




\vskip0.2cm
\noindent
\begin{minipage}{\textwidth}
	\centering
	\captionof{figure}{Sólido definido por $r = 1 - \cos t$, \quad $r = 2 - 2\cos t$, y $4z = x^2 + y^2$}
	\label{fig:solido_cardioide_paraboloide}
	\begin{minipage}{0.26\linewidth}
		\centering
		\includegraphics[width=\linewidth]{Img/37a.pdf}
	\end{minipage}
	\begin{minipage}{0.26\linewidth}
		\centering
		\includegraphics[width=\linewidth]{Img/37b.pdf}
	\end{minipage}
	\begin{minipage}{0.26\linewidth}
		\centering
		\includegraphics[width=\linewidth]{Img/37c.pdf}
	\end{minipage}
	
	\vspace{0.3em}
	\fuentepropia
\end{minipage}
\vskip0.2cm


Con las siguientes indicaciones se consiguen las figuras presentadas



\vspace{0.6em}
\noindent
\begin{tcolorbox}[breakable,enhanced,title=\bf Instrucción en \verb"Mathematica", top=2pt,bottom=2pt,boxsep=0pt,before skip=2pt,after skip=0pt]
\begin{Verbatim}[formatcom=\letraverbatim]
a=ParametricPlot3D[{{(1-Cos[t])*Cos[t],
(1-Cos[t])*Sin[t],u},{(2-2*Cos[t])*Cos[t],
(2-2*Cos[t])*Sin[t],u}},{u,0,4},{t,0,2*Pi},
PlotStyle->{Orange,RGBColor[0.7,0.1,0.1]},
Mesh->False,PlotPoints->80];
b=ContourPlot3D[{4*z==x^2+y^2},{x,-4,4},{y,-4,4},
{z,0,4},ContourStyle->RGBColor[0.4,0.5,0.3],
Mesh->False];
Show[b,a]
\end{Verbatim}
\end{tcolorbox}
\vspace{0.7em}
\noindent




\vskip0.2cm
\noindent
\begin{minipage}{\textwidth}
	\centering
	\captionof{figure}{Sólido definido por $1 - \cos(t) \leq r \leq 2 - 2\cos(t)$ \quad y \quad $0 \leq 4z \leq r^2$}
	\label{fig:cardioide_paraboloide}
	\begin{minipage}{0.25\linewidth}
		\centering
		\includegraphics[width=\linewidth]{Img/37d.pdf}
	\end{minipage}
	\begin{minipage}{0.25\linewidth}
		\centering
		\includegraphics[width=\linewidth]{Img/37f.pdf}
	\end{minipage}
	\begin{minipage}{0.25\linewidth}
		\centering
		\includegraphics[width=\linewidth]{Img/37e.pdf}
	\end{minipage}
	
	\vspace{0.3em}
	\fuentepropia
\end{minipage}
\vskip0.2cm



El sólido se consigue definiendo las coordenadas polares y ejecutando las siguientes instrucciones:


\vspace{0.6em}
\noindent
\begin{tcolorbox}[breakable,enhanced,title=\bf Instrucción en \verb"Mathematica", top=2pt,bottom=2pt,boxsep=0pt,before skip=2pt,after skip=0pt]
\begin{Verbatim}[formatcom=\letraverbatim]
{r,t}={Sqrt[x^2+y^2],ArcTan[x,y]};
RegionPlot3D[1-Cos[t]<r<2-2*Cos[t]&&0<z<(x^2+y^2)/4,
{x,-4,2},{y,-4,4},{z,0,4.2},PlotPoints->90,Mesh->False]
\end{Verbatim}
\end{tcolorbox}
\vspace{0.7em}
\noindent


La base del sólido es igual a la región definida por la fórmula (\ref{625x}). Sobre el paraboloide la coordenada $z$ del sólido es $z = (u+1)^2 \sin ^4\left(\frac{t}{2}\right),$ esto permite definir la función
{\small
\begin{multline} \label{628x}
\overrightarrow{r}_2(t,u)=(u+1)\bigg \langle 2\sin^2\left(\frac{t}{2}\right)\cos(t),\sin(t)(1-\cos(t)),\\ (u+1)\sin^4\left(\frac{t}{2}\right) \bigg\rangle
\end{multline}
}



\noindent
\begin{minipage}{\textwidth}
	\centering
	\captionof{figure}{Sólido definido por $1 - \cos(t) \leq r \leq 2 - 2\cos(t)$ \quad y \quad $0 \leq 4z \leq r^2$.}
	\label{fig:solido_cardioide_paraboloide_2}
	\begin{minipage}{0.33\linewidth}
		\centering
		\vspace{-18mm}
		\includegraphics[width=\linewidth]{Fig/464.pdf}
	\end{minipage}
	\hspace{1em}
	\begin{minipage}{0.33\linewidth}
		\centering
		\vspace{-18mm}
		\includegraphics[width=\linewidth]{Fig/465.pdf}
	\end{minipage}
	
	\vspace{-10mm}
	\fuentepropia
\end{minipage}
\vskip0.3cm


Un punto $Q$ sobre la base del sólido tendrá coordenadas determinadas por la expresión \ref{625x}
y un punto $P$ sobre el paraboloide tendrá la forma dada por la expresión \ref{628x};
con ellos se puede formar una combinación convexa que parametriza el sólido. Esto permite definir la función
\begin{equation*}
	\mathbf{r}(t,u,w) = (u+1)\left\langle 2\sin^2\left(\tfrac{t}{2}\right)\cos(t),\sin(t)(1-\cos(t)),(u+1)w\sin^4\left(\tfrac{t}{2}\right)\right\rangle,
\end{equation*}
en donde $t\in [0,2\pi],$ $u \in [0,1]$, $w \in [0,1].$
Con esta función se pueden obtener los siguientes vectores
{\footnotesize
\begin{align*}
	{\bf r}_{t} &= (u+1)\left \langle\sin(2t)-\sin(t), 2\sin^2\left(\frac{t}{2}\right)(2\cos (t)+1), w(u+1) \sin ^2\left(\frac{t}{2}\right) \sin (t) \right \rangle \\
	{\bf r}_{u} &=\left \langle 2 \sin ^2\left(\frac{t}{2}\right) \cos (t), -\sin (t)(\cos (t)+1), 2(u+1) w \sin ^4\left(\frac{t}{2}\right) \right \rangle \\
	{\bf r}_{w} &=\left \langle 0,0,(u+1)^2 \sin ^4\left(\frac{t}{2}\right) \right \rangle.
\end{align*}
}
Finalmente, el volumen está dado por
\vskip0.2cm
\[\int_{0}^{1}\int_{0}^{1}\int_{0}^{2\pi }\left | -4 (u+1)^3 \sin ^8\left(\frac{t}{2}\right)\right|dtdudw= \frac{525}{64}\pi.\]
\vskip0.2cm
Utilizando coordenadas cilíndricas, también se puede obtener este valor
\vskip0.2cm
\[\int_{0}^{2\pi}\int_{1-\cos\theta}^{2-2\cos\theta}\int_{0}^{\frac{1}{4}r^{2}}rdzdrd\theta=\frac{525}{64}\pi.\]
\vskip0.2cm
Las siguientes instrucciones permiten construir la parametrización del sólido, hallar el elemento diferencial de volumen y evaluar la integral que proporciona su volumen:

\vspace{0.6em}
\noindent
\begin{tcolorbox}[breakable,enhanced,title=\bf Instrucción en \verb"Mathematica", top=2pt,bottom=2pt,boxsep=0pt,before skip=2pt,after skip=0pt]
\begin{Verbatim}[formatcom=\letraverbatim]
r1={Cos[t]*(1-Cos[t]),Sin[t]*(1-Cos[t])};
r2={Cos[t]*2*(1-Cos[t]),2*Sin[t]*(1-Cos[t])};
{x,y}=Simplify[(1-u)*r1+u*r2];
R=Simplify[(1-w)*{x,y,0}+w*{x,y,(x^2+y^2)/4}];
A=Simplify[Dot[D[R,t],Cross[D[R,u],D[R,w]]]]
Integrate[Abs[A],{t,0,2*Pi},{u,0,1},{w,0,1}]
\end{Verbatim}
\end{tcolorbox}
\vspace{0.7em}
\noindent


\subsection*{Un cilindro cardióidico y dos planos}

Consideremos el sólido acotado por cilindro cardióidico con ecuación $r=2-2\cos (t)$ y los planos con ecuaciones $z=0$ y $z=1-y/3.$




\vskip0.2cm
\noindent
\begin{minipage}{\textwidth}
	\centering
	\captionof{figure}{Sólido definido por $r = 2 - 2\cos(t)$ \quad y \quad $z = 1 - \frac{y}{3}$}
	\label{fig:cardioide_plano_inclinado}
	\begin{minipage}{0.25\linewidth}
		\centering
		\includegraphics[width=\linewidth]{Img/38a.pdf}
	\end{minipage}
	\begin{minipage}{0.25\linewidth}
		\centering
		\includegraphics[width=\linewidth]{Img/38b.pdf}
	\end{minipage}
	\begin{minipage}{0.25\linewidth}
		\centering
		\includegraphics[width=\linewidth]{Img/38c.pdf}
	\end{minipage}
	
	\vspace{0.3em}
	\fuentepropia
\end{minipage}
\vskip0.2cm


Con las siguientes instrucciones se consiguen las figuras presentadas:


\vspace{0.6em}
\noindent
\begin{tcolorbox}[breakable,enhanced,title=\bf Instrucción en \verb"Mathematica", top=2pt,bottom=2pt,boxsep=0pt,before skip=2pt,after skip=0pt]
\begin{Verbatim}[formatcom=\letraverbatim]
a=ParametricPlot3D[{(2-2*Cos[t])*Cos[t],
(2-2*Cos[t])*Sin[t],u},{u,0,2},{t,0,2*Pi},
Mesh->False,PlotStyle->Orange,PlotPoints->90];
b=ContourPlot3D[z==1-y/3,{x,-4,2},{y,-3,3},{z,0,2},
Mesh->False,ContourStyle->RGBColor[0.4,0.5,0.3],
PlotPoints->80];
Show[b,a]
{r,t}={Sqrt[x^2+y^2],ArcTan[x,y]};
RegionPlot3D[r<2-2*Cos[t]&&0<z<1-y/3,{x,-4,2},
{y,-3,3},{z,0,2},PlotPoints->90,Mesh->False]
\end{Verbatim}
\end{tcolorbox}
\vspace{0.7em}
\noindent



A partir de la ecuación $r=2-2\cos(t)$ se sigue que, en el plano $xy$, la cardioide se parametriza
$x=2u\cos(t)(1-\cos(t))$ y $y=2u\sin(t)(1-\cos(t)).$
Consideremos un punto $Q(2u\cos(t)(1-\cos(t)),2u\sin(t)(1-\cos(t)),0),$ sobre la base del sólido y un punto
$P\big(2u\cos(t)(1-\cos(t)), \linebreak 2u\sin(t)(1-\cos(t)),1+\frac{2}{3}u\sin(t)(\cos(t)-1)\big),$
sobre el plano $z=1-y/3,$ para $u\in [0,1]$, $t\in [0,2\pi].$
Una combinación convexa de la forma $(1-w)P+wQ$
parametriza el sólido, esto es
\begin{eqnarray}
{\bf r}\left(t,u,w\right)=\big\langle 2u(1-\cos(t))\cos(t),2u\sin(t)(1-\cos(t)),\nonumber\\
\frac{1}{3}(1-w)(3-2u\sin(t)+u\sin(2t)) \big\rangle, \label{38a}
\end{eqnarray}
para $t \in [0,2\pi ]$, $u\in [0,1]$, $w\in [0,1]$.

\vskip0.2cm
\noindent
\begin{minipage}{\textwidth}
	\centering
	\captionof{figure}{Sólido definido por $0 \leq r \leq 2(1 - \cos t)$ \quad y \quad $0 \leq 3z \leq 3 - y$.}
	\label{fig:cardioide_plano_inclinado_final}
	\begin{minipage}{0.35\linewidth}
		\centering
		\vspace{-18mm}
		\includegraphics[width=\linewidth]{Fig/467.pdf}
	\end{minipage}
	\hspace{1em}
	\begin{minipage}{0.35\linewidth}
		\centering
		\vspace{-18mm}
		\includegraphics[width=\linewidth]{Fig/468.pdf}
	\end{minipage}
	
	\vspace{-10mm}
	\fuentepropia
\end{minipage}


\vskip0.2cm
Derivando parcialmente la función \ref{38a} se obtienen los siguientes vectores
{\small
\begin{flalign*}
	&{\bf r}_t = \Big \langle 2u\sin(t)(2\cos(t)-1), &\\
	&\qquad 4u\sin^2\left(\frac{t}{2}\right)(2\cos(t)+1), \frac{2}{3}u(w-1)(\cos(t)-\cos(2t)) \Big \rangle ,&\\
	& {\bf r}_{u} = \Big\langle 2(1-\cos(t))\cos(t), 2\sin(t)(1-\cos(t)),&\\
	&\qquad \frac{1}{3}(1-w)(\sin(2t)-2\sin(t)) \Big\rangle \quad \text{y}&\\
	&{\bf r}_{w}=\Big\langle 0,0,\frac{1}{3}(2u\sin(t)-2u\sin(t)\cos(t)-3) \Big\rangle .&
\end{flalign*}
}
\vskip0.2cm
Esto implica que {\small
$\left| {\bf r}_t \cdot \left( {\bf r}_{u}\times {\bf r}_{w}\right) \right| =
\left | \frac{16}{3}u\sin^4\left(\frac{t}{2}\right)(3-2u\sin(t)+u\sin(2t)) \right|.$
}
Por tanto, el volumen del sólido está dado por
{\small\[\int_{0}^{1}\int_{0}^{1}\int_{0}^{2\pi }\left( \frac{4}{3}u\left( \cos (t) -1\right) ^{2}\left( 2u\sin (t) \left( \cos (t) -1\right) +3\right)\right) dtdudw= 6\pi.\]
}
En \verb"Mathematica" el proceso anterior se ejecuta como sigue.


\vspace{0.6em}
\noindent
\begin{tcolorbox}[breakable,enhanced,title=\bf Instrucción en \verb"Mathematica", top=2pt,bottom=2pt,boxsep=0pt,before skip=2pt,after skip=0pt]
\begin{Verbatim}[formatcom=\letraverbatim]
P={u*(2-2*Cos[t])*Cos[t],u*(2-2*Cos[t])*Sin[t],
1-(u*(2-2*Cos[t])*Sin[t])/3};
Q={u*(2-2*Cos[t])*Cos[t],u*(2-2*Cos[t])*Sin[t],0};
r=Simplify[(1-w)*P+w*Q];
A=Simplify[Dot[D[r,t],Cross[D[r,u],D[r,w]]]];
Integrate[A,{w,0,1},{u,0,1},{t,0,2*Pi}]
\end{Verbatim}
\end{tcolorbox}
\vspace{0.7em}
\noindent


Para ver el resultado en cada instrucción, el usuario debe, además de ejecutar cada sentencia, eliminar el punto y coma (;) en algunas líneas.
En coordenadas cilíndricas, el volumen está dado por la siguiente integral
\vskip0.2cm
\[\int_{0}^{2\pi }\int_{0}^{2-2\cos \theta }\int_{0}^{1-\frac{1}{3} r\sin \theta }rdzdrd\theta = 6\pi.\]


\subsection*{Un cilindro cardióidico y un paraboloide}

Consideremos el parabóloide con ecuación $3z=9-(x+2)^{2}-(y+1)^{2}$ y el cilindro cardióidico
$r=1-\cos (t)$.
Esta región se presenta en la sección \secref{cilcardiopar}. \vskip0.2cm
En dos dimensiones, para $t\in [0,2\pi], $ la cardioide se parametriza como
{\small
\[\vec{r}(t)=\left\langle (1-\cos(t))\cos(t),(1-\cos(t))\sin(t)\right\rangle.\]
}

\begin{figure}[H]
\centering
\caption{Sólido definido por $r = 1 - \cos t$ \quad y \quad $3z = 4 - 4x - x^2 - 2y - y^2$.}
\label{fig:cardioide_superficie_parabolica}
\includegraphics[width=4cm,height=3cm]{Img/58a.pdf}\hskip0.5cm
\includegraphics[width=4cm,height=3cm]{Img/58b.pdf}
\fuentepropia
\end{figure}

La región interior a la cardioide en el plano \textit{xy} se parametriza como
\begin{align}
x(u,t)&= u(1-\cos (t))\cos (t) \label{38b} \\
y(u,t)&= u(1-\cos (t))\sin (t) \label{38c}
\end{align}
% ${\bf r}_1(t,u)=\left\langle u(1-\cos (t))\cos (t),u(1-\cos (t))\sin (t),0 \right\rangle ,$
La coordenada $z$ sobre el paraboloide satisface que,
\begin{equation}
	\resizebox{0.85\textwidth}{!}{%
		$z(u,t)= \frac{1}{3}\left(9-(u(1-\cos (t))\cos (t)+2)^{2} -(u(1-\cos (t))\sin(t)+1)^{2}\right)$
	}
	\label{38d}
\end{equation}

Este término en la ecuación \ref{38d} se reescribe como
{\footnotesize
\[\frac{1}{6}\left( (4u-u^2)\cos(2t) + (4u^2-8u)\cos(t) -4u\sin(t)+2u\sin(2t)-3u^2+4u+8 \right).\]
}


\noindent
\begin{minipage}{\textwidth}
	\centering
	\captionof{figure}{Región sobre el paraboloide interior al cilindro.}
	\label{fig:paraboloide_interior_cilindro}
	\begin{minipage}{0.35\linewidth}
		\centering
		\vspace{-15mm}
		\includegraphics[width=\linewidth]{Fig/469.pdf}
	\end{minipage}
	\hspace{1em} 
	\begin{minipage}{0.35\linewidth}
		\centering
		\vspace{-15mm}
		\includegraphics[width=\linewidth]{Fig/470.pdf}
	\end{minipage}
	
	\vspace{-10mm}
	\fuentepropia
\end{minipage}


\vskip0.2cm
La base del sólido se describe por medio de la función 
\vskip0.2cm
\[{\bf r}_1(u,t)=\langle x(u,t),y(u,t),0\rangle.\]
\vskip0.2cm
% \begin{equation} {\bf r}_1(u,t)=\langle x(u,t),y(u,t),0\rangle. \end{equation}
Utilizando las expresiones \ref{38b}, \ref{38c} y \ref{38d} se usará la función
%\vskip0.2cm
\begin{equation}
{\bf r}_2(u,t)=\langle x(u,t),y(u,t),z(u,t)\rangle,
\end{equation}
%\vskip0.2cm
con $u\in [0,1],$ $t\in [0,2\pi ],$ como una parametrización de la superficie del paraboloide dentro del cilindro cardióidico.
El sólido se parametriza con $w\in [0,1],$ mediante una combinación convexa entre un punto sobre el paraboloide y un punto sobre el plano $xy$, esto es $(1-w){\bf r}_1(u,t)+w{\bf r}_2(u,t)$.
Con lo cual, el sólido queda determinado por medio de la función
\begin{multline*}
{\bf r} (t,u,w) =\bigg\langle u(1-\cos(t))\cos(t),u\sin(t)(1-\cos(t)),\frac{1}{6}w\bigg(2u\sin(2t) \\
4u\sin(t)+4(u-2) u\cos(t)-(u-4)u\cos(2t)-3u^2+4u+8\bigg) \bigg\rangle.
\end{multline*}
Puede mostrarse que
\begin{multline*}
\left| {\bf r}_{t}\cdot \left( {\bf r}_{u}\times {\bf r}_{w}\right) \right|
=\bigg | \frac{1}{6}u(\cos(t)-1)^2\bigg(8+4u-3u^2-u((8-4u)\cos(t)+\\ (u-4)\cos(2t)+ 4\sin(t)-2\sin(2t))\bigg)\bigg |,
\end{multline*}
entonces,
\[\int_{0}^{1}\int_{0}^{1}\int_{0}^{2\pi}\left| {\bf r}_{t}\cdot \left( {\bf r}_{u}\times {\bf r}_{w}\right) \right|dtdudw= \frac{47}{16}\pi.\]

%Al expresar la ecuación del parabolide en coordenadas polares,
%$z = \frac{1}{3}\left( 9-(r\cos \theta +2)^{2}-(r\sin \theta +1)^{2}\right)=\frac{4}{3}-\frac{4}{3}r\cos \theta -\frac{1}{3}r^{2}-\frac{2}{3}r\sin \theta ,$
El volumen también puede hallarse evaluando la siguiente integral:
\[
\bigintss_{0}^{2\pi}
\bigintss_{0}^{1 - \cos \theta}
\bigintss_{0}^{\frac{4}{3} - \frac{4}{3}r\cos \theta - \frac{1}{3}r^{2} - \frac{2}{3}r\sin \theta}
r\, dz\, dr\, d\theta = \frac{47}{16}\pi
\]

Los cálculos anteriores se pueden obtener en \verb"Mathematica", definiendo $x$ e
$y,$ como la parametrización de la cardioide. Luego definir la parametrización $r$ del sólido, hallar el producto vectorial triple se denomina $A$ y se integra con los límites descritos anteriormente.


\vspace{0.6em}
\noindent
\begin{tcolorbox}[breakable,enhanced,title=\bf Instrucción en \verb"Mathematica", top=2pt,bottom=2pt,boxsep=0pt,before skip=2pt,after skip=0pt]
\begin{Verbatim}[formatcom=\letraverbatim]
{x,y}={u*(1-Cos[t])*Cos[t],u*(1-Cos[t])*Sin[t]};
r=(1-w)*{x,y,0}+w*{x,y,(9-(x+2)^2-(y+1)^2)/3};
A=Simplify[Dot[D[r,t],Cross[D[r,u],D[r,w]]]];
Integrate[Abs[A],{t,0,2*Pi},{u,0,1},{w,0,1}]
Integrate[r,{\[Theta],0,2*Pi},{r,0,1-Cos[\[Theta]]},
{z,0,4/3-4/3*r*Cos[\[Theta]]
r^2/3-2/3*r*Sin[\[Theta]]}]
\end{Verbatim}
\end{tcolorbox}
\vspace{0.7em}
\noindent


\section{Sólidos acotados por otras superficies}
%\vskip0.9cm
\subsection*{Superficie exponencial y  paraboloide}

Las superficies definidas por las ecuaciones $z=x^{2}+y^{2}$ y $z=e^{1-x^{2}-y^{2}}$ el sólido que se muestra a continuación:

\begin{figure}[H]
\centering
\caption{Sólido definido por $x^{2}+y^{2} \leq z \leq e^{1-x^{2}-y^{2}}$.}
\label{fig:paraboloide_exponencial}
\includegraphics[height=4.5cm]{Img/54a.pdf} \hspace{0.2cm}
\includegraphics[height=4.5cm]{Img/54b.pdf} \hspace{0.2cm}
\includegraphics[height=4.5cm]{Img/54c.pdf}
\fuentepropia
\end{figure}

Estos gráficos se obtienen por medio de las siguientes instrucciones.


\vspace{0.6em}
\noindent
\begin{tcolorbox}[breakable,enhanced,title=\bf Instrucción en \verb"Mathematica", top=2pt,bottom=2pt,boxsep=0pt,before skip=2pt,after skip=0pt]
\begin{Verbatim}[formatcom=\letraverbatim]
RegionPlot3D[x^2+y^2<z<Exp[1-x^2-y^2]&&-Sqrt[1-x^2]<y<
Sqrt[1-x^2]&&-1<x<1,{x,-1,1},{y,-1,1},{z,0,3},
PlotPoints->90,PlotStyle->{Blend[{Orange,Red},1/4]},
Mesh->False]
ContourPlot3D[{z==x^2+y^2,z==Exp[1-x^2-y^2]},{x,-2,2},
{y,-2,2},{z,0,2.8},ContourStyle->{Directive[Orange,
Opacity[0.7]],Directive[RGBColor[0.6,0.1,0.1],
Opacity[0.6]]},Mesh->False,BoxRatios->{2,2,1}]
\end{Verbatim}
\end{tcolorbox}
\vspace{0.7em}
\noindent


La intersección de las superficies se da en $z=1,$
por lo tanto, en el plano $xy,$ la intersección es el conjunto de puntos que satisface la ecuación $x^{2}+y^{2}=1.$


\noindent
\begin{minipage}{\textwidth}
	\centering
	\captionof{figure}{Proyección del sólido en los planos $xz$ y $xy$.}
	\label{fig:proyecciones_xz_xy}
	\begin{minipage}{0.30\linewidth}
		\centering
		\vspace{-15mm}
		\includegraphics[width=\linewidth]{Fig/471.pdf}
	\end{minipage}
	\hspace{1em}
	\begin{minipage}{0.30\linewidth}
		\centering
		\vspace{-15mm}
		\includegraphics[width=\linewidth]{Fig/472.pdf}
	\end{minipage}
	
	\vspace{-10mm}
	\fuentepropia
\end{minipage}


\vskip0.2cm
Ahora presentamos las integrales que se requieren evaluar para hallar el volumen del mencionado sólido.
%\vskip0.2cm
{\small
\[\int_{-1}^{1}\int_{-\sqrt{1-x^{2}}}^{\sqrt{1-x^{2}}}\int_{x^{2}+y^{2}}^{e^{1-x^{2}-y^{2}}}dzdydx= \int_{-1}^{1}\int_{-\sqrt{1-y^{2}}}^{\sqrt{1-y^{2}}}\int_{x^{2}+y^{2}}^{e^{1-x^{2}-y^{2}}}dzdxdy \approx 3.8273.\]
}
%\vskip0.2cm

Al despejar $y,$ de las expresiones que determinan las superficies e integrar primero con respecto a $y,$ se obtiene por un lado que {\small $y=\pm \sqrt{1-x^{2}-\ln z}$} y {\small $y=\pm \sqrt{z-x^{2}},$} por el otro. Entonces, el volumen está dado por
%\vskip0.2cm
{\small
\[\int_{0}^{1}\int_{-\sqrt{z}}^{\sqrt{z}}\int_{-\sqrt{z-x^{2}}}^{\sqrt{z-x^{2}}}dydxdz+\int_{1}^{e}\int_{-\sqrt{1-\ln z}}^{\sqrt{1-\ln z}}\int_{-\sqrt{1-x^{2}-\ln z}}^{\sqrt{1-x^{2}-\ln z}}dydxdz=\frac{1}{2}\pi ( 2e-3)\]
}
\[\int_{-1}^{1}\int_{x^{2}}^{1}\int_{-\sqrt{z-x^{2}}}^{\sqrt{z-x^{2}}}dydzdx+\int_{-1}^{1}\int_{1}^{e^{1-x^{2}}}\int_{-\sqrt{1-x^{2}-\ln z}}^{\sqrt{1-x^{2}-\ln z}}dydzdx \approx 3.8273.\]
%\vskip0.2cm
Por la simetría de la figura, al intercambiar \(x\) con \(y\) en las ecuaciones que determinan las superficies, es evidente que se obtiene:
%\vskip0.2cm

{\small
\[
\int_{0}^{1} \int_{-\sqrt{z}}^{\sqrt{z}} \int_{-\sqrt{z - y^{2}}}^{\sqrt{z - y^{2}}} dx\, dy\, dz
+ \int_{1}^{e} \int_{-\sqrt{1 - \ln z}}^{\sqrt{1 - \ln z}} \int_{-\sqrt{1 - y^{2} - \ln z}}^{\sqrt{1 - y^{2} - \ln z}} dx\, dy\, dz
= \pi e - \frac{3}{2}\pi
\]
}
\[
\int_{-1}^{1} \int_{y^{2}}^{1} \int_{-\sqrt{z - y^{2}}}^{\sqrt{z - y^{2}}} dx\, dz\, dy
+ \int_{-1}^{1} \int_{1}^{e^{1 - y^{2}}} \int_{-\sqrt{1 - y^{2} - \ln z}}^{\sqrt{1 - y^{2} - \ln z}} dx\, dz\, dy
\approx 3.8273.
\]
\vskip0.2cm
Algunas de estas integrales se evaluaron usando \verb|Mathematica|, las indicaciones son las siguientes:


\vspace{0.6em}
\noindent
\begin{tcolorbox}[breakable,enhanced,title=\bf Instrucción en \verb"Mathematica", top=2pt,bottom=2pt,boxsep=0pt,before skip=2pt,after skip=0pt]
\begin{Verbatim}[formatcom=\letraverbatim]
NIntegrate[1,{x,-1,1},{y,-Sqrt[1-x^2],Sqrt[1-x^2]},
{z,x^2+y^2,Exp[1-x^2-y^2]}]
Integrate[1,{y,-1,1},{x,-Sqrt[1-y^2],Sqrt[1-y^2]},
{z,x^2+y^2,Exp[1-x^2-y^2]}]
Integrate[1,{z,0,1},{x,-Sqrt[z],Sqrt[z]},
{y,-Sqrt[z-x^2],Sqrt[z-x^2]}]+Integrate[1,
{z,1,Exp[1]},{x,-Sqrt[1-Log[z]],Sqrt[1-Log[z]]},
{y,-Sqrt[1-x^2-Log[z]],Sqrt[1-x^2-Log[z]]}]
NIntegrate[1,{y,-1,1},{z,y^2,1},{x,-Sqrt[z-y^2],
Sqrt[z-y^2]}]+NIntegrate[1,{y,-1,1},{z,1,Exp[1-y^2]},
{x,-Sqrt[1-Log[z]-y^2],Sqrt[1-Log[z]-y^2]}]
\end{Verbatim}
\end{tcolorbox}
\vspace{0.7em}
\noindent
\vskip0.2cm


Utilizando el siguiente proceso también aporta el mismo valor.

%\begin{tcolorbox}[title=\bfseries Parametrizando el sólido]
	
La región de integración en el plano $xy$ es el círculo unitario, un punto en este conjunto se determina por {\small $(u\cos t,u\sin t),$} para $u\in [0,1]$ y {\small $t\in [0,2\pi].$}
Un punto sobre el paraboloide posee coordenadas {\small $P( u\cos t,u\sin t,u^{2}),$}
y un punto sobre la superficie exponencial tiene la forma
{\small $Q(u\cos t,u\sin t,e^{1-u^{2}}).$} Con una combinación convexa de estos puntos se parametriza el sólido en mención.
\begin{multline*}
{\bf r}( u,t,w) = (1-w) Q +wP = \langle u\cos t,u\sin t,e^{1-u^{2}}( 1-w) +wu^{2}\rangle,
\end{multline*}
con $u\in [0,1]$, $w\in [0,1]$ y $t\in [0,2\pi ].$

%\vskip0.2cm
De esta función se obtiene
\begin{eqnarray*}
{\bf r}_{t} &=& \langle -u\sin t,u\cos t,0\rangle \\
{\bf r}_{u} &=&\langle \cos t,\sin t,-2( 1-w)
ue^{1-u^{2}}+2wu\rangle \\
{\bf r}_{w}&=& \langle 0,0,u^{2}-e^{1-u^{2}}\rangle \\
|{\bf r}_{t}\cdot ({\bf r}_{u}\times {\bf r}_{w})| &=& |u\exp(1-u^{2})-u^{3}|=ue^{1-u^{2}}-u^{3}.
\end{eqnarray*}
%\vskip0.2cm
El volumen del sólido se halla evaluando la siguiente integral
%\vskip0.2cm
\[\int_{0}^{1}\int_{0}^{1}\int_{0}^{2\pi }( ue^{1-u^{2}}-u^{3})
dtdwdu= \pi e-\frac{3}{2}\pi.\]
%\vskip0.2cm
%\end{tcolorbox}

Con el cambio a coordenadas cilíndricas, la integral respectiva es
\[\int_{0}^{2\pi }\int_{0}^{1}\int_{r^{2}}^{e^{1-r^{2}}}rdzdrd\theta
= \frac{1}{2}\pi ( 2e-3).\]
%\vskip0.2cm
Con las siguientes instrucciones se obtuvieron los cálculos anteriores.


\vspace{0.6em}
\noindent
\begin{tcolorbox}[breakable,enhanced,title=\bf Instrucción en \verb"Mathematica", top=2pt,bottom=2pt,boxsep=0pt,before skip=2pt,after skip=0pt]
\begin{Verbatim}[formatcom=\letraverbatim]
{x,y}={u*Cos[t],u*Sin[t]}
r=Simplify[(1-w)*{x,y,x^2+y^2}+w*{x,y,Exp[1-x^2-y^2]}]
A=Simplify[Dot[D[r,w],Cross[D[r,u],D[r,t]]]]
Integrate[A,{u,0,1},{t,0,2*Pi},{w,0,1}]
Integrate[r,{\[Theta],0,2*Pi},{r,0,1},
{z,r^2,Exp[1-r^2]}]
\end{Verbatim}
\end{tcolorbox}
\vspace{0.7em}
\noindent


\subsection*{Una semiesfera y un cilindro} \index{Semiesfera}
Hallar el volumen del sólido determinado por las siguientes expresiones.
$x^2 +y^2 +z^2 \leq 9,$ $ y^2 +z^2 \geq 1,$ $z \leq 0.$ La figura se forma cuando una esfera de radio $3$ se perfora por un cilindro de radio $1$ cuyo eje de simetría es una recta que pasa por el centro de la esfera; luego, un plano que contiene al eje de simetría del cilindro divide esta región del espacio en dos porciones iguales y se toma una de estas.

\vskip0.2cm
\noindent
\begin{minipage}{\textwidth}
	\centering
	\captionof{figure}{Sólido definido por $1 \leq y^2 + z^2 \leq 9 - x^2$ \quad y \quad $z \geq 0$.}
	\label{fig:cilindro_parabolico_superior}
	\begin{minipage}{0.25\linewidth}
		\centering
		\includegraphics[width=\linewidth]{Img/56c.pdf}
	\end{minipage}
	\hspace{1em}
	\begin{minipage}{0.30\linewidth}
		\centering
		\includegraphics[width=\linewidth]{Img/56d.pdf}
	\end{minipage}
	
	\vspace{0.3em}
	\fuentepropia
\end{minipage}

\vskip0.2cm
Las siguientes instrucciones permiten obtener el gráfico del sólido.


\vspace{0.6em}
\noindent
\begin{tcolorbox}[breakable,enhanced,title=\bf Instrucción en \verb"Mathematica", top=2pt,bottom=2pt,boxsep=0pt,before skip=2pt,after skip=0pt]
\begin{Verbatim}[formatcom=\letraverbatim]
RegionPlot3D[x^2+y^2+z^2<9&&y^2+z^2>1&&-3<z<0,
{x,-3,3},{y,-3,3},{z,-3,2},PlotPoints->90,Mesh->False,
PlotStyle->Orange,Ticks->{{-4,-2,0,2,4},{-3,-1,1,3},
{-2,0,2}},AxesLabel->{x,y,z}]
\end{Verbatim}
\end{tcolorbox}
\vspace{0.7em}
\noindent


Con el fin de plantear las integrales que determinan el volumen del sólido, se presentan las siguientes figuras.

\vskip0.2cm
\noindent
\begin{minipage}{\textwidth}
	\centering
	\captionof{figure}{Proyección del sólido en los planos coordenados.}
	\label{figsemicirc}
	\begin{minipage}{0.28\linewidth}
		\centering
		\vspace{-10mm}
		\includegraphics[width=\linewidth]{Fig/473.pdf}
	\end{minipage}
	\hspace{1em}
	\begin{minipage}{0.28\linewidth}
		\centering
		\vspace{-10mm}
		\includegraphics[width=\linewidth]{Fig/474.pdf}
	\end{minipage}
	\hspace{1em}
	\begin{minipage}{0.28\linewidth}
		\centering
		\vspace{-10mm}
		\includegraphics[width=\linewidth]{Fig/475.pdf}
	\end{minipage}
	
	\vspace{-4mm}
	\fuentepropia
\end{minipage}


\vskip0.2cm
El volumen del sólido que se considera está dado por las siguientes integrales.
En el orden $dxdzdy$ se debe observar que $x$ recorre puntos desde la parte negativa de la esfera hasta su parte positiva.
%\vskip0.2cm
\vspace{-2mm}
\begin{multline*}
\int_{-3}^{-1}\int_{-\sqrt{9-y^2 }}^{0}\int_{-\sqrt{9-y^2 -z^2 }}^{\sqrt{9-y^2 -z^2 }}dxdzdy+\int_{-1}^{1}\int_{-\sqrt{9-y^2 }}^{-\sqrt{1-y^2 }}\int_{-\sqrt{9-y^2 -z^2 }}^{\sqrt{9-y^2 -z^2 }}dxdzdy\\
+\int_{1}^{3}\int_{-\sqrt{9-y^2 }}^{0}\int_{-\sqrt{9-y^2 -z^2 }}^{\sqrt{9-y^2 -z^2 }}dxdzdy = \frac{32}{3}\sqrt{2}\pi
\end{multline*}
%\vskip0.2cm
\vspace{-7mm}
\begin{multline*}
\int_{-1}^{0}\int_{-\sqrt{9-z^2 }}^{-\sqrt{1-z^2 }}\int_{-\sqrt{9-y^2 -z^2 }}^{\sqrt{9-y^2 -z^2 }}dxdydz+\int_{-3}^{-1}\int_{-\sqrt{9-z^2 }}^{\sqrt{9-z^2 }}\int_{-\sqrt{9-y^2 -z^2 }}^{\sqrt{9-y^2 -z^2}}dxdydz\\
+\int_{-1}^{0}\int_{\sqrt{1-z^2 }}^{\sqrt{9-z^2 }}\int_{-\sqrt{9-y^2 -z^2 }}^{\sqrt{9-y^2 -z^2 }}dxdydz= \frac{32}{3}\sqrt{2}\pi
\end{multline*}

%\vskip0.2cm

En el plano $xz$ hay una porción en la región de integración que corresponde a un rectángulo y existe una región en el plano $xz$ que no contiene puntos $x$ en $[-3,\sqrt{8}]$ y $[-3,\sqrt{8}],$ entre otros.
%\vskip0.2cm
\begin{multline*}
\int_{-1}^{0}\int_{-\sqrt{8}}^{\sqrt{8}}\int_{-\sqrt{9-x^{2}-z^{2}}}^{-\sqrt{1-z^{2}}}dydxdz
+\int_{-1}^{0}\int_{-\sqrt{8}}^{\sqrt{8}}\int_{\sqrt{1-z^{2}}}^{\sqrt{9-x^{2}-z^{2}}}dydxdz\\
+\int_{-3}^{-1}\int_{-\sqrt{9-z^{2}}}^{\sqrt{9-z^{2}}}\int_{-\sqrt{9-x^{2}-z^{2}}}^{\sqrt{9-x^{2}-z^{2}}}dydxdz = \frac{32}{3}\sqrt{2}\pi
\end{multline*}
%\vskip0.2cm
\vspace{-7mm}
\begin{multline*}
\int_{-\sqrt{8}}^{\sqrt{8}} \int_{-1}^{0}\int_{-\sqrt{9-x^{2}-z^{2}}}^{-\sqrt{1-z^{2}}}dydzdx
+\int_{-\sqrt{8}}^{\sqrt{8}}\int_{-1}^{0}\int_{\sqrt{1-z^{2}}}^{\sqrt{9-x^{2}-z^{2}}}dydzdx \\
+\int_{-\sqrt{8}}^{\sqrt{8}}\int_{-\sqrt{9-x^{2}}}^{-1}\int_{-\sqrt{9-x^{2}-z^{2}}}^{\sqrt{9-x^{2}-z^{2}}}dydzdx = \frac{32}{3}\sqrt{2}\pi.
\end{multline*}
%\vskip0.2cm
Integrando primero, con respecto a $z,$ existe en el plano $xy$ una región donde $z$ recorre los puntos desde la parte negativa de la esfera al plano
$z=0$ y otra región donde $z$ recorre los puntos de la parte negativa de la esfera hasta la parte negativa del cilindro con ecuación $y^2+z^2=1.$
\vspace{-2mm}
{\small
\begin{multline*}
\int_{1}^{3}\int_{-\sqrt{9-y^2}}^{\sqrt{9-y^2}}\int_{-\sqrt{9-x^2 -y^2 }}^{0}dzdxdy+\int_{-3}^{-1}\int_{-\sqrt{9-y^2 }}^{\sqrt{9-y^2 }}\int_{-\sqrt{9-x^2 -y^2 }}^{0}dzdxdy\\
+\int_{-1}^{1}\int_{-\sqrt{8}}^{\sqrt{8}}\int_{-\sqrt{9-x^2 -y^2 }}^{-\sqrt{1-y^2 }}dzdxdy= \frac{32}{3}\sqrt{2}\pi
\end{multline*}
}
\vspace{-9mm}
{\small
\begin{multline*}
\int_{-\sqrt{8}}^{\sqrt{8}}\int_{-\sqrt{9-x^2 }}^{-1}\int_{-\sqrt{9-x^2 -y^2 }}^{0}dzdydx+\int_{-\sqrt{8}}^{\sqrt{8}}\int_{1}^{\sqrt{9-x^2 }}\int_{-\sqrt{9-x^2 -y^2 }}^{0}dzdydx\\
+\int_{-\sqrt{8}}^{\sqrt{8}}\int_{-1}^{1}\int_{-\sqrt{9-x^2 -y^2 }}^{-\sqrt{1-y^2 }}dzdydx= \frac{32}{3}\sqrt{2}\pi.
\end{multline*}
}
%\vskip0.2cm
Al hacer un cambio a coordenadas cilíndricas
$y = r\cos \theta,\, z = r\sin \theta,$, el volumen está dado por.
%\vskip0.2cm
{\small
\[\int_{\pi}^{2\pi}\int_{1}^{3}\int_{-\sqrt{9-r^2}}^{\sqrt{9-r^2}}rdxdrd\theta=\frac{32}{3}\sqrt{2}\pi.\]
}
%\vskip0.2cm
Algunos de los cálculos anteriores se pueden evaluar usando \verb"Mathematica"
mediante la ejecución de las siguientes instrucciones.



\vspace{0.6em}
\noindent
\begin{tcolorbox}[breakable,enhanced,title=\bf Instrucción en \verb"Mathematica", top=2pt,bottom=2pt,boxsep=0pt,before skip=2pt,after skip=0pt]
\begin{Verbatim}[formatcom=\letraverbatim]
Integrate[r,{t,Pi,2*Pi},{r,1,3},{x,-Sqrt[9-r^2],
Sqrt[9-r^2]}]
Integrate[1,{y,-3,-1},{z,-Sqrt[9-y^2],0},
{x,-Sqrt[9-y^2-z^2],Sqrt[9-y^2-z^2]}]+Integrate[1,
{y,-1,1},{z,-Sqrt[9-y^2],-Sqrt[1-y^2]},{x,
Sqrt[9-y^2-z^2],Sqrt[9-y^2-z^2]}]+Integrate[1,
{y,1,3},{z,-Sqrt[9-y^2],0},{x,-Sqrt[9-y^2-z^2],
Sqrt[9-y^2-z^2]}]

Integrate[1,{z,-1,0},{x,-Sqrt[8],Sqrt[8]},
{y,-Sqrt[9-x^2-z^2],-Sqrt[1-z^2]}]+Integrate[1,
{z,-1,0},{x,-Sqrt[8],Sqrt[8]},{y,Sqrt[1-z^2],
Sqrt[9-x^2-z^2]}]+Integrate[1,{z,-3,-1},{x,
Sqrt[9-z^2],Sqrt[9-z^2]},{y,-Sqrt[9-x^2-z^2],
Sqrt[9-x^2-z^2]}]
NIntegrate[1,{x,-Sqrt[8],Sqrt[8]},{y,-Sqrt[9-x^2],-1},
{z,-Sqrt[9-x^2-z^2],0}]+NIntegrate[1,{x,-Sqrt[8],
Sqrt[8]},{y,1,Sqrt[9-x^2]},{z,-Sqrt[9-x^2-z^2],0}]
+NIntegrate[1,{x,-Sqrt[8],Sqrt[8]},{y,-1,1},
{z,-Sqrt[9-x^2-z^2],-Sqrt[1-y^2]}]
\end{Verbatim}
\end{tcolorbox}
\vspace{0.7em}
\noindent


El volumen también se obtiene mediante el siguiente procedimiento.

%\begin{tcolorbox}[title=\bfseries Parametrizando el sólido]
Consideremos la proyección del sólido en el plano $yz,$ (véase la figura \ref{figsemicirc}). Un punto sobre la semicircunferencia de radio $1,$ posee coordenadas
$( \cos( t) ,\sin ( t) ) $ y un punto sobre la semicircunferencia de radio $3,$ es de la forma
$( 3\cos ( t),3\sin ( t) ) ,$ con $t \in [ \pi ,2\pi ].$
Por lo tanto, una combinación convexa de estos dos puntos parametriza la región en mención:
$( 1-u) ( \cos ( t) ,\sin ( t) )+u( 3\cos ( t) ,3\sin ( t) ) $
que se simplifica como $ ( ( 1+2u) \cos t,( 1+2u) \sin t) .$
\vskip0.2cm
Sobre el sólido, $x$ recorre los puntos que van desde la parte negativa de la esfera hasta la parte positiva de la misma. Dos puntos arbitrarios
$P$ y $Q$ en estos dos lugares tienen coordenadas
$P=(-2\sqrt{( u+2)( 1-u) },(1+2u)\cos t,(1+2u)\sin t), \,\, \text{y} $
$Q=(2\sqrt{( u+2) ( 1-u) },(1+2u)\cos t,(1+2u)\sin t).$
\vskip0.2cm
Es posible parametrizar el sólido mediante una combinación convexa de la forma $(1-w)P+wQ$, esto define la función ${\bf r}.$
%\vskip0.2cm
\vspace{-2mm}
{\small
\begin{multline*}
{\bf r}(t,u,w) = \big \langle 2( 2w-1) \sqrt{( u+2)( 1-u) },\cos t+2u\cos t, \\ \sin t+2u\sin t \big \rangle,
\end{multline*}
}
en donde $t \in [0,2\pi],$ $u \in [0,1]$ y $w \in [0,1].$
Al derivar parcialmente con respecto a $t$, $u$ y $w$ hallamos respectivamente, los siguientes vectores:
%\vskip0.2cm
\vspace{-2mm}
{\small
\begin{eqnarray*}
{\bf r}_{t} &=& \langle 0,-\sin t-2u\sin t,\cos t+2u\cos t\rangle \\
{\bf r}_{u} &=& \left \langle \frac{( 1-2w)( 1+2u)}{\sqrt{( u+2) ( 1-u) }},2\cos t,2\sin t\right \rangle \\
{\bf r}_{w} &=& ( 4\sqrt{( u+2) ( 1-u) },0,0).
\end{eqnarray*}
}
%\vskip0.2cm
\vspace{-2mm}
Su triple producto vectorial triple es
{\small
\[|{\bf r}_{t}\cdot({\bf r}_{u}\times{\bf r}_{w})|=8(1+2u)\sqrt{(u+2)(1-u)};\]
}
%\vskip0.2cm
al integrarlo con los límites definidos, nos lleva al volumen del
sólido
%\vskip0.2cm
{\small
\[\int_{0}^{1}\int_{0}^{1}\int_{\pi}^{2\pi}8(1+2u)\sqrt{(u+2)(1-u)}\,dtdudw=\frac{32}{3}\sqrt{2}\pi.\]
}
%\end{tcolorbox}

\subsection*{Sólidos cardióidicos} \index{Cardioide} \index{Sólidos!cardioidicos}

En esta sección consideraremos el sólido generado por la rotación de una cardioide alrededor de una recta exterior a ella.
Compararemos los valores obtenidos parametrizando el sólido y comparándolos con el valor obtenido aplicando el teorema de Pappus o la conocida formulación para superficies de revolución. \index{Teorema de Pappus!área superficial}

\index{Toro cardioidico!volumen}
\subsection*{Toro cardióidico generado por $\bm{\mathsf{r=1-\sin t}}$}

Para $0\leq t \leq 2\pi,$, la cardioide con ecuación $r=1-\sin t$ su área está dada por:
\[\int_{0}^{2\pi }\int_{0}^{1-\sin \theta}rdrd\theta = \frac{3}{2}\pi.\]
El volumen del sólido generado cuando la curva parametrizada en la forma {\small $\langle x(t),y(t)\rangle,$} con {\small $t \in [a,b]$}, gira alrededor del eje $y,$
está dado por {\small $\ds \pi \int_{a}^{b}\left( x(t)\right) ^{2}y^{\prime }(t)dt.$}
El volumen está dado por {\small $\ds\pi \int_{a}^{b}\left( y(t)\right) ^{2}x^{\prime }(t)dt,$}
si el giro se hace alrededor del eje $x.$
En el caso de la cardioide $r=1-\sin \theta$ que se desplazada tres unidades sobre el eje polar, esta curva se puede parametrizar como {\small $x(t)=3+\left( 1-\sin (t) \right) \cos (t),$}
{\small $y(t)=\left( 1-\sin (t) \right) \sin (t) ,$} {\small $0\leq t \leq 2\pi.$} Entonces el volumen $V$ del sólido que se obtiene al girar la cardiode sobre la recta $\theta = \pi/2,$ está dado por
{\small \[V=\pi \int_{0}^{2\pi }\left( 3+\left( 1-\sin (t) \right) \cos(t) \right) ^{2}(\cos t-2\cos t\sin t)dt= 9\pi ^{2}.\]
}
\vspace{-6mm}
%\vskip0.2cm
\begin{figure}[H]
\centering
\caption{Toro cardióidico generado por $r = 1 - \sin \theta$.}
\label{fig:toro_cardioidico}
\includegraphics[height=4.5cm]{Img/63d.pdf} \hspace{0.4cm}
\includegraphics[height=4.5cm]{Img/63e.pdf}
\fuentepropia
\end{figure}

Las coordenadas del centro de masa, $(\overline{x},\overline{y})$ están determinadas por
\[\frac{2}{3\pi }\int_{0}^{2\pi }\int_{0}^{1-\sin \theta }\left( 3+r\cos \theta \right) rdrd\theta = 3,\]
%\vskip0.2cm
\[\frac{2}{3\pi }\int_{0}^{2\pi }\int_{0}^{1-\sin \theta }\left( r\sin \theta \right) rdrd\theta = -\frac{5}{6}.\]
El volumen de este sólido, usando el teorema de Pappus, es \index{Teorema de Pappus!volumen}
$2\pi \left( \frac{3}{2}\pi \right) \left( 3\right) = 9\pi ^{2}.$

% \begin{eqnarray*} \index{Centro de masa} \overline{x} &=& \frac{2}{3\pi }\int_{0}^{2\pi }\int_{0}^{1-\sin \theta }\left( 3+r\cos \theta \right) rdrd\theta = 3 \\ \overline{y}&=& \frac{2}{3\pi }\int_{0}^{2\pi }\int_{0}^{1-\sin \theta }\left( r\sin \theta \right) rdrd\theta = -\frac{5}{6}
% \end{eqnarray*}

%\begin{tcolorbox}[title=\bfseries Parametrizando el sólido]
%\footnotesize 
%\parskip=0pt 
%\setlength{\belowdisplayskip}{0pt} 
%\setlength{\abovedisplayskip}{0pt} 
\vskip0.2cm
La cardioide se rellena usando la parametrización
%\vskip0.2cm
\[\langle 3+u\left( 1-\sin (t) \right) \cos (t),u\left( 1-\sin (t) \right) \sin (t) \rangle,\]
%\vskip0.2cm
con $t\in [0,2\pi]$ y $u\in [0,1]$.
Agregando a la primera componente $\cos(v)$ y $\sin(v)$, se obtiene una parametrización de la región rotada,
\begin{multline*}
{\bf r}(t,u,v)=\bigg \langle\left ( 3+u\left( 1-\sin (t) \right) \cos (t) \right) \cos (v) , \\ \left( 3+u\left( 1-\sin \left(
t\right) \right) \cos (t) \right) \sin (v) ,u\left(1-\sin (t) \right) \sin (t) \bigg \rangle,
\end{multline*}
\vskip0.2cm
con $t\in [0,2\pi ],\,\, u\in [0,1]$ y $v\in [0,2\pi].$
Derivando parcialmente se obtienen los vectores: 
\vskip0.2cm
${\bf r}_{u} = ( 1-\sin t) \langle \cos t\cos v, \cos t\sin v, \sin t\rangle ,$ \\
${\bf r}_{v} = (3+u(1-\sin t) \cos t) \langle -
\sin v,\cos t) \cos v,0\rangle ,$ \\
${\bf r}_{t} = \bigg \langle -u( \cos \left( 2t\right) +\sin t) \cos
v,-u\left( \cos \left( 2t\right) +\sin t\right) \sin v, \\
u\left( 1-2\sin t\right) \cos t \bigg \rangle.$
\vskip0.2cm
Con ellos, se obtiene el elemento diferencial de volumen
$|{\bf r}_{t}\cdot \left( {\bf r}_{u}\times {\bf r}_{v}\right)| =
|u\left( 6-3\cos ^{2}t+u\left(\sin t-3\right) \cos ^{3}t+4u(1-\sin t)\cos t-6\sin t\right)|,$
cuya integración nos proporciona el volumen del sólido
%\vskip0.2cm
\[\int_{0}^{2\pi }\int_{0}^{1}\int_{0}^{2\pi} |{\bf r}_{t}\cdot \left( {\bf r}_{u}\times {\bf r}_{v}\right)|
dtdudv= 9\pi ^{2}.\]
%\end{tcolorbox}

Si la densidad del solido es $\delta \left( x,y,z\right) =z^{2},$ entonces la masa del sólido es:
%\vskip0.2cm
\begin{equation*}
\int_{0}^{2\pi }\int_{0}^{1}\int_{0}^{2\pi}\left(u\left(1-\sin(t)\right)\sin(t)\right)^{2}|{\bf r}_{t}\cdot\left({\bf r}_{u}\times {\bf r}_{v}\right)|dtdudv= \frac{147}{16}\pi ^{2}.
\end{equation*}
\vskip0.2cm
Los cálculos anteriores se evaluan con las siguientes instrucciones: %en \verb"Mathematica" con ayuda de la siguiente rutina


\vspace{0.6em}
\noindent
\begin{tcolorbox}[breakable,enhanced,title=\bf Instrucción \verb"Mathematica", top=2pt,bottom=2pt,boxsep=0pt,before skip=2pt,after skip=0pt]
\begin{Verbatim}[formatcom=\letraverbatim]
{x,y,z}={3+u*(1-Sin[t])*Cos[t],(1-u*Sin[t])*Sin[t],y};
r={x*Cos[v],x*Sin[v],z};
A=Simplify[Dot[D[r,t],Cross[D[r,u],D[r,v]]]]
Integrate[A,{t,0,2*Pi},{u,0,1},{v,0,2*Pi}]
Integrate[z^2*A,{t,0,2*Pi},{u,0,1},{v,0,2*Pi}]
\end{Verbatim}
\end{tcolorbox}
\vspace{0.7em}
\noindent



\subsection*{Toro cardióidico generado por $\bm{\mathsf{r=1-\cos t}}$}

Consideraremos la cardioide con ecuación $r=1-\cos(t),$ $0\leq t \leq 2\pi.$
En el plano polar, si esta curva se desplaza tres unidades sobre el eje polar, y luego rota alrededor de una recta paralela a
$\theta=\pi/2$, se obtiene el sólido $E$.

\begin{figure}[H]
\centering
\caption{Toro cardioidico generado por $r = 1 - \cos \theta$.}
\label{fig:toro_cardioidico_cos}
\includegraphics[height=4.5cm]{Img/59f.pdf}
\hspace{0.4cm}
\includegraphics[height=4.5cm]{Img/59e.pdf}
\fuentepropia
\end{figure}

El área de la cardioide es
%\vskip0.2cm
\[\int_{0}^{2\pi }\int_{0}^{1-\cos \theta }rdrd\theta = \frac{3}{2}\pi,\]
%\vskip0.2cm
y las coordenadas del centro de masa están dadas por
%\vskip0.2cm
\begin{eqnarray*}
\overline{x} &=& \frac{2}{3\pi}\int_{0}^{2\pi}\int_{0}^{1-\cos \theta }\left( 3+r\cos \theta \right) rdrd\theta =\frac{13}{6} \\
\overline{y} &=& \frac{2}{3\pi}\int_0^{2\pi}\int_0^{1-\cos\theta}(r\sin\theta)rdrd\theta = 0.
\end{eqnarray*}
\vskip0.2cm
De acuerdo con el teorema de Pappus, el volumen del toro cardioidico generado por
$r=1-\cos t$, es:
\[2\pi \left( \frac{13}{6}\right) \frac{3}{2}\pi =\frac{13}{2}\pi ^{2}.\]
\vskip0.2cm
La cardioide se parametriza en el plano $xy,$ como
\[x(t)=3+\left( 1-\cos (t) \right) \cos (t),\quad y(t)=\left( 1-\cos (t) \right) \sin (t),\quad t\in[0,2\pi].\]

El volumen $V,$ del sólido generado por la revolución de esta región sobre la recta $\theta = \pi/2,$ está dado por
\[V=\pi \int_{0}^{2\pi }\left( 3+\left( 1-\cos (t) \right) \cos (t) \right) ^{2}(\cos t-2\cos ^{2}t+1)dt =
\frac{13}{2}\pi ^{2}.\]

%\begin{tcolorbox}[title=\bfseries Parametrizando el sólido]
En el plano $xy$, para rellenar la cardioide se usa la parametrización
${\bf r}_1(t,u)=\langle 3+u\left( 1-\cos (t) \right) \cos (t),u\left( 1-\cos (t) \right) \sin (t) \rangle ,$ con $t\in [0,2\pi ],\,\, u\in [0,1].$
Si esta región gira sobre la recta $x=3$, se genera un toro cardioidico, cualquier punto sobre este sólido tiene la forma determinada por la si\-guiente función
\begin{multline*}
{\bf r}(t,u,v)=\bigg \langle \left( 3+u\left( 1-\cos (t) \right) \cos \left(
t\right) \right) \cos (v) , \\
\left( 3+u\left( 1-\cos \left(t\right) \right) \cos (t) \right) \sin (v) , u\left(1-\cos (t) \right) \sin (t) \bigg \rangle,
\end{multline*}
con $t\in [0,2\pi ],\,\,v\in [0,2\pi ],\,\, u\in [0,1].$
Al derivar parcialmente con respecto a $t$, $u$ y $v$ se obtiene el elemento diferencial de volumen:
\begin{multline*}
\left|{\bf r}_{t}\cdot \left( {\bf r}_{u}\times {\bf r}_{v}\right)\right| 
= \Big|(1+u)\cos^2(t)\left(6 - u + 2u\cos(t) - u\cos(2t)\right) \\
\times \sin^2\left( \frac{t}{2} \right)\Big|
\end{multline*}
Esto se integra y nos proporciona el volumen del toro cardióidico generado por la curva $r = 1 - \cos t$:
\[
\int_{0}^{2\pi }\int_{0}^{2\pi }\int_{0}^{1}
\left|{\bf r}_{t}\cdot \left( {\bf r}_{u}\times {\bf r}_{v}\right)\right| \, dudvdt
= \frac{13}{2}\pi ^{2}.
\]
%\end{tcolorbox}

En \verb"Mathematica" se pueden efectuar estos cálculos ejecutando las si\-guientes instrucciones:



\vspace{0.6em}
\noindent
\begin{tcolorbox}[breakable,enhanced,title=\bf Instrucción \verb"Mathematica", top=2pt,bottom=2pt,boxsep=0pt,before skip=2pt,after skip=0pt]
\begin{Verbatim}[formatcom=\letraverbatim]
R=(3+u*(1-Cos[t])*Cos[t]);
r={x,y,z}={R*Cos[v],R*Sin[v],u*(1-Cos[t])*Sin[t]};
A=Simplify[Dot[D[r,t],Cross[D[r,u],D[r,v]]]];
Integrate[Abs[A],{t,0,2*Pi},{v,0,2*Pi},{u,0,1}]
\end{Verbatim}
\end{tcolorbox}
\vspace{0.7em}
\noindent



\subsection*{Un sólido cardióidico} \index{Cardioide}

En el plano $xy$, el interior de la cardioide $r=1-\cos(t)$
se parametriza por medio de las expresiones
$x(t,u)=u\left( 1-\cos (t)\right)\cos(t)$, $y(t,u)=u\left( 1-\cos (t) \right) \sin (t), $
para $u \in [0,1]$ y $t \in [0,2\pi].$
\vskip0.2cm
Construiremos una superficie que en el plano $xy$ coincida con la cardioide mencionada.
La ecuación $r=1-\cos(t)$ corresponde en coordenadas rectangulares a la expresión
$(x^2+y^2)=(x^2+y^2+x)^2$; es por esto que al hacer
$z=(x^2+y^2)-(x^2+y^2+x)^2,$ se satisface la igualdad en $z=0$.
Sustituyéndose $x$ e $y$ en la expresión anterior se obtiene:
%\vskip0.2cm
\[z(t,u)= 8 (u-1) u^2 \sin ^6\left(\frac{t}{2}\right)((u-1)\cos (t)-u-1).\]
%\vskip0.2cm
La superficie parametrizada por las ecuaciones $x(t,u),$ $y(t,u)$ $z(t,u)$ se muestra a continuación:


\vskip0.2cm
\noindent
\begin{minipage}{\textwidth}
	\centering
	\captionof{figure}{Superficie definida por $z = (x^2 + y^2) - (x^2 + y^2 + x)^2$}
	\label{fig:superficie_compuesta}
	\begin{minipage}{0.26\linewidth}
		\centering
		\includegraphics[width=\linewidth]{Img/65a.pdf}
	\end{minipage}
	\begin{minipage}{0.26\linewidth}
		\centering
		\includegraphics[width=\linewidth]{Img/65b.pdf}
	\end{minipage}
	
	\vspace{0.3em}
	\fuentepropia
\end{minipage}
\vskip0.2cm


El volumen del sólido acotado por esta superficie y el plano $z=0$
se hallan definiendo la función ${\bf r}\left( t,u,v\right),$ que es una combinación convexa de dos puntos de la forma $(x(t,u),y(t,u),0)$ y $(x(t,u),y(t,u),z(t,u)).$ Esto es:
\begin{multline*}
{\bf r}\left( t,u,v\right) = \bigg \langle u(1-\cos(t))\cos(t),u\sin(t)(1-\cos(t)), \\
8(u-1)u^2v\sin ^6\left(\frac{t}{2}\right)((u-1)\cos(t)-u-1 \bigg \rangle,
\end{multline*}
\vskip0.2cm
en donde \( t \in [0,2\pi] \), \( u \in [0,1] \) y \( v \in [0,1] \).
\vskip0.2cm
Con esta función, hallamos el producto vectorial triple:
%\vskip0.2cm
\[
\left| \mathbf{r}_{t} \cdot \left( \mathbf{r}_{u} \times \mathbf{r}_{v} \right) \right| =
\left| 32 (1 - u) u^3 \sin^{10}\left( \frac{t}{2} \right) \left( (u - 1) \cos(t) - u - 1 \right) \right|,
\]
\vskip0.2cm
que se integra y nos proporciona el volumen del sólido en mención:
{\small
\[
\int_{0}^{1} \int_{0}^{1} \int_{0}^{2\pi}
\left| 32 (1 - u) u^3 \sin^{10}\left( \frac{t}{2} \right) \left( (u - 1) \cos(t) - u - 1 \right) \right|
\, dt\, du\, dv = \frac{35}{32}\pi.
\]
}

El volumen se puede plantear en coordenadas cilíndricas como sigue:
%\vskip0.2cm
\[
\int_{0}^{2\pi} \int_{0}^{1 - \cos \theta}
\int_{0}^{r^{2} \sin^{2}(\theta) - r^{4} - 2r^{3} \cos(\theta)}
r\, dz\, dr\, d\theta = \frac{35}{32}\pi.
\]
%\vskip0.2cm
Las figuras y las integrales se consiguen con las siguientes instrucciones:


\vspace{0.6em}
\noindent
\begin{tcolorbox}[breakable,enhanced,title=\bf Instrucción \verb"Mathematica", top=2pt,bottom=2pt,boxsep=0pt,before skip=2pt,after skip=0pt]
\begin{Verbatim}[formatcom=\letraverbatim]
{x,y}={u*(1-Cos[t])*Cos[t],u*(1-Cos[t])*Sin[t]};
z=(x^2+y^2)-(x^2+y^2+x)^2;
R=(1-v)*{x,y,0}+v*{x,y,z};
ParametricPlot3D[{x,y,z},{t,0,2*Pi},{u,0,1},
Ticks->{{-2,-1,0},{-1,0,1},{0,2}},MeshStyle->Gray,
PlotPoints->80,PlotStyle->Orange,PlotRange->{{-2,1/2},
{-3/2,3/2},{0,2}}]
A=Simplify[Abs[Dot[D[R,t],Cross[D[R,u],D[R,v]]]]];
Integrate[A,{t,0,2*Pi},{u,0,1},{v,0,1}]
Clear[r,z,t]
Integrate[r,{t,0,2*Pi},{r,0,1-Cos[t]},
{z,0,r^2-(r^2+r*Cos[t])^2}]
\end{Verbatim}
\end{tcolorbox}
\vspace{0.7em}
\noindent



\subsection*{Un cardioide revolucionado
} \index{Cardioide}


La cardioide escrita en forma polar como $r=1-\cos \theta,$ se parametriza como $x=\left( 1-\cos (t) \right) \cos (t) $
y $y=\left( 1-\cos (t) \right) \sin (t) ,$ para $t\in [0,2\pi].$
Sea $E$ el sólido obtenido cuando esta curva gira alrededor del eje polar.

\vskip0.2cm
\noindent
\begin{minipage}{\textwidth}
	\centering
	\captionof{figure}{Sólido $E$ generado por la rotación de $r = 1 - \cos \theta$.}
	\label{fig:toro_cardioidico_E}
	\begin{minipage}{0.25\linewidth}
		\centering
		\includegraphics[width=\linewidth]{Img/66b.pdf}
	\end{minipage}
	\hspace{1em}
	\begin{minipage}{0.25\linewidth}
		\centering
		\includegraphics[width=\linewidth]{Img/66a.pdf}
	\end{minipage}
	
	\vspace{0.3em}
	\fuentepropia
\end{minipage}

\vskip0.2cm
Esta figura se genera en \verb"Mathematica" usando los siguientes comandos:


\vspace{0.6em}
\noindent
\begin{tcolorbox}[breakable,enhanced,title=\bf Instrucción \verb"Mathematica", top=2pt,bottom=2pt,boxsep=0pt,before skip=2pt,after skip=0pt]
\begin{Verbatim}[formatcom=\letraverbatim]
{x,y}={(1-Cos[t])*Cos[t],(1-Cos[t])*Sin[t]};
r={x,y*Cos[v],y*Sin[v]};
ParametricPlot3D[r,{t,0,Pi},{v,0,2*Pi},MeshStyle->Gray,
PlotStyle->Orange,PlotPoints->60,Mesh->{20,15}]
\end{Verbatim}
\end{tcolorbox}
\vspace{0.7em}
\noindent


El volumen de $E$, usando el teorema de Pappus, está dado por
%\vskip0.2cm
{\small
\begin{eqnarray*}
V&=&\pi \int_{0}^{\pi }\left(- \left( 1-\cos (t) \right) \sin
(t) \right) ^{2}\frac{d}{dt}\left( \left( 1-\cos (t)
\right) \cos (t) \right) dt \\
&=& \pi \int_{0}^{\pi }\left( \sin (t) \left( \cos (t)
+1\right) \left( 1-2\cos (t) \right) \left( \cos (t)
1\right) ^{3} \right) dt = \frac{8}{3}\pi
\end{eqnarray*}
}
%\vskip0.2cm
Otra manera de obtener este valor es la siguiente.

%\begin{tcolorbox}[title=\bfseries Parametrizando el sólido]
%\vskip0.2cm
Una parametrización de la región interior a la cardioide es
\[x=u\left( 1-\cos (t) \right) \cos (t),\]
%\vskip0.2cm
\[y=u\left( 1-\cos (t) \right) \sin (t),\]
en donde $t\in [0,\pi]$ y $u\in [0,1].$
Si esta región gira sobre el eje polar, agregando $\cos(v)$ y $\sin(v)$ a $y$, con $v \in [0,2\pi]$, se obtiene una parametrización de $E$, esto es:
%\vskip0.2cm
\[{\bf r}(t,u,v)=u(1-\cos (t))\left\langle \cos (t) ,\sin (t) \cos (v) ,\sin(t) \sin (v) \right \rangle.\]
%\vskip0.2cm
El producto vectorial triple es
$\left| {\bf r}_{t}\cdot \left( {\bf r}_{u}\times {\bf r}_{v}\right) \right| = \left|
8 u^2 \sin ^6\left(\frac{t}{2}\right) \sin (t)\right| ,$
Entonces, el volumen de $E$ está dado por
%\vskip0.2cm
\[\int_{0}^{2\pi }\int_{0}^{1}\int_{0}^{\pi }\left| u^{2}\sin (t)
\left( \cos (t) -1\right) ^{3}\right| dtdudv= \frac{8}{3}\pi.\]
%\vskip0.2cm
%\end{tcolorbox}
	
	
Estos cálculos se pueden hacer en \verb"Mathematica" con ayuda de las siguientes instrucciones:


\vspace{0.6em}
\noindent
\begin{tcolorbox}[breakable,enhanced,title=\bf Instrucción \verb"Mathematica", top=2pt,bottom=2pt,boxsep=0pt,before skip=2pt,after skip=0pt]
\begin{Verbatim}[formatcom=\letraverbatim]
r=u*(1-Cos[t])*{Cos[t],Sin[t]*Cos[v],Sin[t]*Sin[v]};
A=Simplify[Dot[D[r,t],Cross[D[r,u],D[r,v]]]]
Integrate[Abs[A],{t,0,Pi},{v,0,2Pi},{u,0,1}]
\end{Verbatim}
\end{tcolorbox}
\vspace{0.7em}
\noindent
%\clearpage


\subsection*{Ejercicios propuestos para terminar}
En esta sección se muestran varios sólidos con las instrucciones requeridas para obtener su gráfico y
al menos una posibilidad de obtener su volumen mediante integrales múltiples.
El lector deberá, como ejercicio, obtener el mismo valor a través del planteamiento de integrales triples en otros 
ordenes de integración o planteando cambios de variables.
\vskip0.2cm
\begin{enumerate}[leftmargin=0cm,itemindent=0.5cm,labelwidth=\itemindent,labelsep=0cm,align=left,label={{\bf \arabic*)}}]
	
\item En \cite{stewart_calculo_2018}, sección 12.1, se plantea el siguiente ejercicio:
Describa y trace un sólido con las siguientes propiedades: cuando es iluminado por rayos paralelos al eje $z,$ su sombra es un disco circular. Si los rayos son paralelos al eje $y,$ su sombra es un cuadrado. Si los rayos son paralelos al eje $x,$
su sombra es un triángulo isósceles. Presentamos un sólido con estas características.

\vskip0.2cm

La proyección en el plano $xy$ corresponde a una círculo, cuya frontera se parametriza en la forma
$x= \cos(t)$, $y= \sin(t)$ para $t\in [0,2\pi].$ El diámetro en la base es $2,$ por tanto consideraremos un segmento paralelo al eje $y$ que sea perpendicular y simétrico al eje $z$. Un punto en la base tendrá la forma $P(\cos(t),\sin(t),0),$ mientras que un punto en la parte superior posee la forma $Q(0,\sin(t),2)$. Formando segmentos rectilíneos entre los puntos anteriores, es decir simplificando $(1-u)P+uQ,$
se obtiene una función que parametriza dicha superficie:
\vspace{-2mm}
\begin{eqnarray*}
{\bf r}(t,u) &=& \langle (1-u)\cos(t),\sin(t),2u \rangle, \,\, t\in [0,2\pi],\, u\in [0,1].
\end{eqnarray*}

\vskip0.2cm
\noindent
\begin{minipage}{\textwidth}
	\centering
	\captionof{figure}{\small Sólido cuyas proyecciones son un rectángulo, un círculo o un triángulo}
	\begin{minipage}{0.21\linewidth}
		\centering
		\includegraphics[width=\linewidth]{Imgn/T_20a.pdf}
	\end{minipage}
	\hspace{1em}
	\begin{minipage}{0.21\linewidth}
		\centering
		\includegraphics[width=\linewidth]{Imgn/T_20b.pdf}
	\end{minipage}
	\hspace{1em}
	\begin{minipage}{0.21\linewidth}
		\centering
		\includegraphics[width=\linewidth]{Imgn/T_20c.pdf}
	\end{minipage}
	\vspace{0.3em}
	\fuentepropia
\end{minipage}

\vskip0.2cm		
Estas figuras se consiguen con las siguientes instrucciones:

\vspace{0.6em}
\noindent
\begin{tcolorbox}[breakable,enhanced,title=\bf Instrucción \verb"Mathematica", top=2pt,bottom=2pt,boxsep=0pt,before skip=2pt,after skip=0pt]
\begin{Verbatim}[formatcom=\letraverbatim]
R=(1-u)*{0,Sin[t],2}+u*{Cos[t],Sin[t],0};
ParametricPlot3D[R,{t,0,2*Pi},{u,0,1},PlotPoints->80,
Ticks->{{-1,0,1},{-1,0,1},{1,2}},MeshStyle->Black]
\end{Verbatim}
\end{tcolorbox}
	
\vskip0.4cm	
El sólido se parametriza rellenando la base de sólido que corresponde a un círculo y formando segmentos rectilíneos hasta el segmento superior. Es decir,
%\vskip0.2cm
\vspace{-2mm}
\begin{eqnarray*}
{\bf r}(t,u,w) &=&(1-w)(0,u\sin(t),2)+w(u\cos(t),u\sin(t),0) \\
&=& \langle uw\cos(t),u\sin(t),2(1-w) \rangle, 
\end{eqnarray*}
%\vskip0.2cm
con $t\in [0,2\pi],\, u\in [0,1],\, w\in [0,1].$
El elemento diferencial de volumen está dado por $|{\bf r}_w \cdot ({\bf r}_t \times {\bf r}_u) |=2uw,$ por lo tanto el volumen está dado por
%\vskip0.2cm

$$\int_0^1 \int_0^1 \int_0^{2\pi} 2 uw dt du dw = \pi.$$
%\vskip0.2cm

Otro sólido con los requerimientos hechos es el que satisface las expresiones
$0\leq z \leq 1-|x|$, $x^2+y^2\leq 1.$ Obviamente su volumen no es el mismo al obtenido anteriormente.
\vskip0.3cm
	
\item El sólido acotado por las superficies de los los paraboloides con ecuaciones $2z=6-2x^2-y^2$ y $z+3=x^2+y^2,$ se muestra a continuación.

\vskip0.3cm
\noindent
\begin{minipage}{\textwidth}
	\centering
	\captionof{figure}{\small $2x^2+2y^2-6\leq 2z\leq 6-2x^2-y^2$ y $,$}
	\begin{minipage}{0.23\linewidth}
		\centering
		\includegraphics[width=\linewidth]{Imgn/T_8a.pdf}
	\end{minipage}
	\hspace{1em}
	\begin{minipage}{0.25\linewidth}
		\centering
		\includegraphics[width=\linewidth]{Imgn/T_8e.pdf}
	\end{minipage}
	
	\vspace{0.3em}
	\fuentepropia
\end{minipage}

\vskip0.2cm
Ahora presentamos las instrucciones necesarias para conseguir esta figura.


\vspace{0.6em}
\noindent
\begin{tcolorbox}[breakable,enhanced,title=\bf Instrucción \verb"Mathematica", top=2pt,bottom=2pt,boxsep=0pt,before skip=2pt,after skip=0pt]
\begin{Verbatim}[formatcom=\letraverbatim]
RegionPlot3D[z<(6-2*x^2-y^2)/2&&z>x^2+y^2-3,{x,-2,2},
{y,-2,2},{z,-3,3},PlotPoints->90,Mesh->False];
\end{Verbatim}
\end{tcolorbox}

\vskip0.4cm	
Al despejar $z$ de las dos ecuaciones e igualar estas expresiones se obtiene $4x^{2}+3y^{2}=12$. Esta expresión corresponde a la intersección de las superficies consideradas y representa, en el plano
$xy$, a una elipse alargada en sentido de las $y$ centrada en el origen. Por tanto, el volumen del sólido está dado por
{\small
\begin{equation} \label{ejercicio1} 
\bigintss_{-\sqrt{3}}^{\sqrt{3}}\bigintss_{-\frac{2}{\sqrt{3}}\sqrt{3-x^{2}}}^{\frac{2}{\sqrt{3}}\sqrt{3-x^{2}}}\bigintss_{x^{2}+y^{2}-3}^{3-x^{2}-\frac{1}{2}y^{2}}dzdydx= 6\sqrt{3}\pi .
\end{equation}
}
\vskip0.2cm
Otra manera de hallar este valor es la siguiente. La elipse y su interior se parametriza como $x=\sqrt{3}u\cos \theta$, $y=2u\sin \theta ,$ con $u \in [0,1]$, $\theta \in [0,2\pi].$ Reemplazando estas expresiones en las ecuaciones de los paraboloides se sigue que

\vspace{-12mm}
\begin{eqnarray*}
z &=& 3-x^{2}-\frac{1}{2}y^{2}= 3-\frac{5}{2}u^{2}-\frac{1}{2}u^{2}\cos 2\theta \\
z&=&x^2+y^2-3= \frac{7}{2}u^{2}-\frac{1}{2}u^{2}\cos 2\theta -3.
\end{eqnarray*}

Sobre cada uno de los paraboloides consideremos dos puntos con coordenadas
$Q(\sqrt{3}u\cos \theta ,2u\sin \theta ,\frac{7}{2}u^{2}-\frac{1}{2}u^{2}\cos 2\theta-3)$
y $P(\sqrt{3}u\cos \theta,2u\sin \theta,\\ 3-\frac{5}{2}u^{2}-\frac{1}{2}u^{2}\cos 2\theta).$  
Una combinación convexa de estos puntos permite parametrizar el sólido, es decir, simplificando la expresión $(1-w)P+wQ$ se define la función
%\vskip0.2cm
{\small
\[{\bf r}(u,\theta ,w)=\left\langle\sqrt{3}u\cos\theta,2u\sin\theta,6u^{2}w-\frac{1}{2}u^{2}\cos2\theta -6w-\frac{5}{2}u^{2}+3 \right\rangle,\]}
%\vskip0.2cm

en donde $u \in [0,1]$, $\theta \in [0,2\pi] ,$ $w\in [0,1].$ 
Con esta función hallamos los siguientes vectores
\vspace{-4mm}
\begin{eqnarray*}
{\bf r}_{\theta}&=& \left\langle -\sqrt{3}u\sin \theta ,2u\cos \theta ,u^{2}\sin 2\theta \right\rangle\\
{\bf r}_{u} &=& \left\langle \sqrt{3}\cos \theta ,2\sin \theta ,12uw-u\cos 2\theta -5u\right\rangle \\
{\bf r}_{w} &=& \left\langle 0,0,6u^{2}-6\right\rangle \\
{\bf r}_{\theta }\times {\bf r}_{u}&=& \left\langle 12u^{2}\left( 2w-1\right)
\cos \theta ,4\sqrt{3}u^{2}\left( 3w-1\right) \sin \theta ,-2\sqrt{3}u \right\rangle.
\end{eqnarray*}

El elemento diferencial de volumen es $\left\vert {\bf r}_{w}\cdot \left({\bf r}_{\theta}\times {\bf r}_{u}\right) \right\vert= 12\sqrt{3}u\left( 1-u^{2}\right) ,$ cuya integración con los límites establecidos nos proporciona el volumen del sólido $$\int_{0}^{1}\int_{0}^{1}\int_{0}^{2\pi }12\sqrt{3}u\left( 1-u^{2}\right)
d\theta dudw= 6\sqrt{3}\pi.$$
\vskip0.2cm
	
\item Consideremos el sólido determinado por las expresiones $x^2+y^2+z^2\leq 9,$ $1+\cos\theta \leq r\leq 2+\cos\theta$. Esto es, la región del espacio acotada por dos cilindros cardióidicos e interior la esfera de radio $3$ centrada en el origen.
\vskip0.2cm
\noindent
\begin{minipage}{\textwidth}
	\centering
	\captionof{figure}{\small $x^2+y^2+z^2\leq 9$ \& $1+\cos\theta\leq r\leq 2+\cos\theta$}
	\begin{minipage}{0.22\linewidth}
		\centering
		\includegraphics[width=\linewidth]{Imgn/T_7a}
	\end{minipage}
	\hspace{1em}
	\begin{minipage}{0.22\linewidth}
		\centering
		\includegraphics[width=\linewidth]{Imgn/T_7b.pdf}
	\end{minipage}
	\hspace{1em}
	\begin{minipage}{0.22\linewidth}
		\centering
		\includegraphics[width=\linewidth]{Imgn/T_7c.pdf}
	\end{minipage}
	
	\vspace{0.3em}
	\fuentepropia
\end{minipage}

\vskip0.2cm
Esta figura y su proyección en el plano $xy$ se obtiene con estas instrucciones:
\vskip0.2cm	


\vspace{0.6em}
\noindent
\begin{tcolorbox}[breakable,enhanced,title=\bf Instrucción \verb"Mathematica", top=2pt,bottom=2pt,boxsep=0pt,before skip=2pt,after skip=0pt]
\begin{Verbatim}[formatcom=\letraverbatim]
PolarPlot[{1+Cos[t],2+Cos[t]},{t,0,2Pi}]
RegionPlot3D[x^2+y^2+z^2<9&&(x^2+y^2-x)^2>x^2+y^2&&
(x^2+y^2-x)^2<3*(x^2+y^2),{x,-1.2,3.2},{y,-2.5,2.5},
{z,-3,3},PlotPoints->90,Mesh->False]
\end{Verbatim}
\end{tcolorbox}

\vskip0.2cm
%\clearpage
La base del sólido se muestra a continuación:

\vskip0.2cm	
\noindent
\begin{minipage}{\textwidth}
	\centering
	\captionof{figure}{\small $1+\cos\theta\leq r\leq 2+\cos\theta$}
	\begin{minipage}{0.30\linewidth}
		\centering
		\vspace{-12mm}
		\includegraphics[width=\linewidth]{Imgn/T_600.pdf}
	\end{minipage}
	
	\vspace{-8mm}
	\fuentepropia
\end{minipage}
\vskip0.3cm

Sobre la proyección del sólido en el plano $xy$ consideremos dos puntos,
$P=((1+\cos(t))\cos(t),(1+\cos(t))\sin(t))$ y $Q=((2+\cos(t))\cos(t),(2+\cos(t))\sin(t)),$ 
uno sobre cada curva. La región se parametriza al formar una combinación convexa de la forma $(1-u)P+uQ$ entre estos puntos, es decir $(\cos(t)(1+u+\cos(t)),(1+u+\cos(t))\sin(t)).$

\vskip0.2cm 

Las siguientes instrucciones muestran en \verb"Mathematica" cómo hacer el procedimiento completo y además permiten hallar su área.
\vskip0.2cm	


\vspace{0.6em}
\noindent
\begin{tcolorbox}[breakable,enhanced,title=\bf Instrucción \verb"Mathematica", top=2pt,bottom=2pt,boxsep=0pt,before skip=2pt,after skip=0pt]	
\begin{Verbatim}[formatcom=\letraverbatim]
r={x,y}=Simplify[(1-u)*{(1+Cos[t])*Cos[t],
(1+Cos[t])*Sin[t]}+u*{(2+Cos[t])*Cos[t],
(2+Cos[t])*Sin[t]}];
w=(1-v)*{x,y,-Sqrt[9-x^2-y^2]}+v*{x,y,
Sqrt[9-x^2-y^2]};
A=Abs[Simplify[Dot[D[w,t],Cross[D[w,u],D[w,v]]]]];
NIntegrate[A,{t,0,2Pi},{u,0,1},{v,0,1}]
\end{Verbatim}
\end{tcolorbox}

\vskip0.4cm
Utilizando coordenadas cilíndricas, el volumen del mencionado sólido está dado por la siguiente integral:
%\vskip0.2cm
$$\int_0^{2\pi} \int_{1+\cos \theta}^{2+\cos \theta}\int_{-\sqrt{9-r^2}}^{\sqrt{9-r^2}}r dzdrd\theta \approx 40.7091.$$
\vskip0.2cm
La evaluación numérica de esta integral se consigue ejecutando la siguiente instrucción:


\vskip0.2cm
\vspace{0.6em}
\noindent
\begin{tcolorbox}[breakable,enhanced,title=\bf Instrucción \verb"Mathematica", top=2pt,bottom=2pt,boxsep=0pt,before skip=2pt,after skip=0pt]
\begin{Verbatim}[formatcom=\letraverbatim]
NIntegrate[r,{t,0,2Pi},{r,1+Cos[t],2+Cos[t]},
{z,-Sqrt[9-r^2],Sqrt[9-r^2]}]
\end{Verbatim}
\end{tcolorbox}
\vskip0.4cm	
	
\item El sólido acotado por las superficies del cilindro $x^2+y^2=4$ y la esfera $x^2+y^2+z^2=9$ se muestran a continuación.
\vskip0.2cm	

\noindent
\begin{minipage}{\textwidth}
	\centering
	\captionof{figure}{\small $x^2+y^2 \leq 4$ \& $6x^2+y^2+z^2\leq 9.$}
	\begin{minipage}{0.20\linewidth}
		\centering
		\includegraphics[width=\linewidth]{Imgn/T_9a.pdf}
	\end{minipage}
	\hspace{1em} 
	\begin{minipage}{0.18\linewidth}
		\centering
		\includegraphics[width=\linewidth]{Imgn/T_9d.pdf}
	\end{minipage}
	
	\vspace{0.3em}
	\fuentepropia
\end{minipage}

\vskip0.2cm
Las siguientes instrucciones permiten obtenerlo:


\vspace{0.6em}
\noindent
\begin{tcolorbox}[breakable,enhanced,title=\bf Instrucción \verb"Mathematica", top=2pt,bottom=2pt,boxsep=0pt,before skip=2pt,after skip=0pt]	
\begin{Verbatim}[formatcom=\letraverbatim]
RegionPlot3D[x^2+y^2<4&&x^2+y^2+z^2<=9,{x,-2,2},
{y,-2,2},{z,-3.5,3.5},PlotPoints->90,Mesh->False]
\end{Verbatim}
\end{tcolorbox}

\vskip0.4cm
Usando coordenadas rectangulares, se plantea la integral que representa el volumen de dicha región.
%\vskip0.2cm
$$\bigintsss_{-2}^2 \bigintsss_{-\sqrt{4-x^{2}}}^{\sqrt{4-x^{2}}}\bigintsss_{-\sqrt{9-x^{2}-y^{2}}}^{\sqrt{9-x^{2}-y^{2}}}dzdydx \approx 66.265. $$
%\vskip0.2cm
En el plano $xy$ la proyección del cilindro y su interior se parametriza como $x=2u\cos \theta, $
$y=2u\sin \theta $ con $u \in [0,1]$ y $\theta \in [0,2\pi];$ sobre la esfera la expresión correspondiente es
%\vskip0.2cm
\[z=\pm \sqrt{9-\left( 2u\cos \theta \right)^2 -\left( 2u\sin \theta \right) ^2 }= \pm \sqrt{9-4u^{2}}.\]
%\vskip0.2cm
De modo que el sólido se parametriza mediante una combinación convexa de dos puntos sobre la esfera dentro del cilindro, esto es:
%\vskip0.2cm
{\small
\[\left( 1-w\right) \left( 2u\cos \theta ,2u\sin
\theta,\sqrt{9-4u^{2}}\right) +w\left( 2u\cos \theta ,2u\sin \theta,-\sqrt{9-4u^{2}}\right).\]
}
%\vskip0.2cm

Lo que nos lleva a presentar la función
%\vskip0.2cm
\begin{eqnarray*}
{\bf r}\left( u,\theta ,w\right) &=& \left\langle 2u\cos \theta ,2u\sin \theta ,\left( 1-2w\right) \sqrt{9-4u^{2}}\right\rangle 
\end{eqnarray*}
%\vskip0.2cm

con $u \in [0,1]$, $\theta \in [0,2\pi],$ $w \in [0,2\pi];$ 
con esta se hallan los siguientes vectores

\vspace{-10mm}
\begin{eqnarray*}
{\bf r}_{\theta } &=& \left\langle -2u\sin \theta ,2u\cos \theta ,0\right\rangle \\
{\bf r}_{u} &=& \left\langle 2\cos \theta ,2\sin \theta ,\frac{4u\left( 2w-1\right) }{\sqrt{9-4u^{2}}}%
\right\rangle \\
{\bf r}_{w}&=& \left\langle 0,0,-2\sqrt{9-4u^{2}}\right\rangle. 
\end{eqnarray*}
El elemento diferencial de volumen es
%\vskip0.2cm
\begin{eqnarray*}
\left\vert {\bf r}_{w}\cdot \left({\bf r}_{u}\times {\bf r}_{\theta }\right) \right\vert 
=8u\sqrt{9-4u^{2}},
\end{eqnarray*}
%\vskip0.2cm
por lo tanto el volumen es
\vskip0.2cm
\begin{equation} \label{ejercicio2}
\int_{0}^{1}\int_{0}^{1}\int_{0}^{2\pi }8u\sqrt{9-4u^{2}}d\theta dudw= \frac{4}{3}\pi \left(27-5\sqrt{5}\right) 
\end{equation} 
\vskip0.2cm

\item El sólido que contiene a los puntos interiores, tanto a la esfera con ecuación $x^2+y^2+z^2=9$ como al cilindro cardióidico
$r=1+\cos \theta,$ se muestra a continuación.

\vskip0.2cm	
\noindent
\begin{minipage}{\textwidth}
	\centering
	\captionof{figure}{\small $x^2+y^2+z^2\leq 9$ \& $r\leq 1+\cos \theta.$}
	\begin{minipage}{0.20\linewidth}
		\centering
		\includegraphics[width=\linewidth]{Imgn/T_12a}
	\end{minipage}
	\hspace{1em}
	\begin{minipage}{0.20\linewidth}
		\centering
		\includegraphics[width=\linewidth]{Imgn/T_12b}
	\end{minipage}
	
	\vspace{0.3em}
	\fuentepropia
\end{minipage}

\vskip0.2cm	
Las instrucciones necesarias para conseguirlo son las siguientes:
\vskip0.2cm


\vspace{0.6em}
\noindent
\begin{tcolorbox}[breakable,enhanced,title=\bf Instrucción en \verb"Mathematica", top=2pt,bottom=2pt,boxsep=0pt,before skip=2pt,after skip=0pt]
\begin{Verbatim}[formatcom=\letraverbatim]
RegionPlot3D[(x^2+y^2-x)^2<x^2+y^2&&x^2+y^2+z^2<9,
{x,-1/2,2},{y,-3/2,3/2},{z,-3,3},Mesh->False,
PlotPoints->80]
\end{Verbatim}
\end{tcolorbox}

\vskip0.4cm
Usando coordenadas cilíndricas se puede plantear la integral que determina el volumen de esta región.
%\vskip0.2cm
\begin{multline} \label{ejercicio3}
\int_{0}^{2\pi }\int_{0}^{1+\cos \theta}\int_{-\sqrt{9-r^{2}}}^{\sqrt{9-r^{2}}}rdzdrd\theta \\ 
= 2\int_{0}^{2\pi }\int_{0}^{1+\cos \theta }r\sqrt{9-r^{2}}\,drd\theta \approx 25.8179 
\end{multline}
\vskip0.2cm
	
\item El siguiente es la figura del sólido acotado por el paraboloide $x^2+y^2=3-z$ y el paraboloide hiperbólico $x^2-y^2=3z+6.$

\vskip0.2cm
\noindent
\begin{minipage}{\textwidth}
	\centering
	\captionof{figure}{\small $x^2+y^2\geq 3-z$ \& $x^2-y^2\leq 3z+6$}
	%{\label{ejercicio4a}
	\begin{minipage}{0.20\linewidth}
		\centering
		\includegraphics[width=\linewidth]{Imgn/T_10a.pdf}
	\end{minipage}
	\hspace{1em}
	\begin{minipage}{0.20\linewidth}
		\centering
		\includegraphics[width=\linewidth]{Imgn/T_10b.pdf}
	\end{minipage}
	
	\vspace{0.3em}
	\fuentepropia
\end{minipage}

\vskip0.2cm
Con las siguientes instrucciones se consigue esta figura:


\vspace{0.6em}
\noindent
\begin{tcolorbox}[breakable,enhanced,title=\bf Instrucción \verb"Mathematica", top=2pt,bottom=2pt,boxsep=0pt,before skip=2pt,after skip=0pt]	
\begin{Verbatim}[formatcom=\letraverbatim]
RegionPlot3D[x^2+y^2<3-z&&x^2-y^2<3*z+6,{x,-2,2},
{y,-3,3},{z,-5,3},PlotPoints->90,Mesh->False]
\end{Verbatim}
\end{tcolorbox}

\vskip0.4cm	
	
La intersección de estas dos superficies se obtiene al igualar las expresiones para $z$, esto es, $$3-x^{2}-y^{2}=\frac{x^{2}-y^{2}-6}{3},$$
la cual se simplifica como $4x^2+2y^2=15$. Esta ecuación corresponde a una elipse en el plano $xy$ centrada en el origen y alargada en sentido de las $y.$ Por tanto, el volumen del sólido está dado por
\vskip0.2cm
{\small
\begin{equation} \label{ejercicio4}
\bigintss_{-\frac{\sqrt{15}}{2}}^{\frac{\sqrt{15}}{2}}\bigintss_{-\frac{1}{2}\sqrt{2} \sqrt{15-4x^{2}}}^{\frac{1}{2}\sqrt{2}\sqrt{15-4x^{2}}}\bigintss_{\frac{x^{2}-y^{2}-6}{3}}^{3-x^{2}-y^{2}}dzdydx= \frac{75}{8}\sqrt{2}\pi
\end{equation}
}
\vskip0.2cm

Otra posibilidad para hallar este valor es parametrizar la elipse mediante las expresiones
{\small $x=\frac{\sqrt{15}}{2}u\cos \theta $}, {\small $y=\sqrt{\frac{15}{2}}u\sin \theta $} con {\small $u \in [0,1]$} y {\small $\theta \in [0,2\pi].$}

Luego, reemplazar en las ecuaciones del paraboloide y el paraboloide hiperbólico, lo cual produce, respectivamente:	
%\vskip0.2cm
{\small
\begin{multline*}
z = 3-\left( \frac{\sqrt{15}}{2}u\cos \theta \right) ^2 -\left( \sqrt{\frac{15}{2}}u\sin \theta \right) ^2  = \frac{15}{8}u^{2}(\cos 2\theta -3)+3 \\
z =  \frac{1}{3}\left( \left( \frac{\sqrt{15}}{2}u\cos \theta \right)^2-\left( \sqrt{\frac{15}{2}}u\sin \theta \right) ^2 -6\right) =  \frac{15}{8}u^{2}(\cos 2\theta -1)-2.
\end{multline*}
}
%\vskip0.2cm
Con estas expresiones tendremos dos puntos de la forma 

{\small
\begin{eqnarray*}
P& =&\left( \frac{\sqrt{15}}{2}u\cos \theta,\sqrt{\frac{15}{2}}u\sin \theta , \frac{15}{8}u^{2}(\cos 2\theta -3)+3\right) \\
Q&=&\left( \frac{\sqrt{15}}{2}u\cos \theta,\sqrt{\frac{15}{2}}u\sin \theta,\frac{15}{8}u^{2}\cos 2\theta-\frac{5}{8}u^{2}-2\right)\!.
\end{eqnarray*}
}

\vskip0.2cm
Una combinación convexa de ellos, $(1-w)P+wQ,$ determina una parametrización del sólido
%\vskip0.2cm
{\small
\[{\bf r}(u,\theta,w) = \left\langle \frac{\sqrt{15}}{2}u\cos v,\sqrt{\frac{15}{2}}u\sin v,5w\left( u^{2}-1\right) +\frac{15}{8}u^{2}\left(\cos2v-3\right) +3 \right\rangle.\]
}

De esta función se obtienen los siguientes vectores:
{\small
\begin{align*}
{\bf r}_{\theta } &= \left\langle -\frac{1}{2}\sqrt{15}u\sin \theta ,\frac{1}{2}%
\sqrt{2}\sqrt{15}u\cos \theta ,-\frac{15}{4}u^{2}\sin 2\theta \right\rangle \\ 
{\bf r}_{u} &= \left\langle \frac{1}{2}\sqrt{15}\cos \theta ,\frac{1}{2}\sqrt{2} \sqrt{15}\sin \theta ,\frac{15}{4}u\left( \cos 2\theta -3\right) +10uw \right\rangle \\
{\bf r}_{w} &= \left\langle 0,0,5u^{2}-5\right\rangle \\
{\bf r}_{\theta }\times {\bf r}_{u} &= \left\langle \frac{5\sqrt{30}}{4} u^{2} (4w-3) \cos \theta ,\frac{5\sqrt{15}}{2}u^{2}(2w-3) \sin \theta ,-\frac{15\sqrt{2}}{4}u \right\rangle.
\end{align*}
}
\vskip0.2cm

El elemento diferencial de volumen es
$\left\vert {\bf r}_{w}\cdot \left({\bf r}_{\theta }\times {\bf r}_{u}\right) \right\vert
=  \left\vert \frac{75}{4}\sqrt{2}\left( u-u^{3}\right) \right\vert ,$
cuya integración nos proporciona el volumen deseado
\vskip0.2cm
$$\int_{0}^{1}\int_{0}^{1}\int_{0}^{2\pi }\left\vert \frac{75}{4}\sqrt{2}\left( u-u^{3}\right) \right\vert d\theta dudw
= \frac{75}{8}\sqrt{2}\pi .$$
\vskip0.2cm
	
\item El sólido cuya región, descrita en coordenadas esféricas, comprende los puntos que satisfacen la expresión $\rho \leq 1+\cos \phi,$ se presenta a continua\-ción:
\vskip0.2cm
\noindent
\begin{minipage}{\textwidth}
	\centering
	\captionof{figure}{\small Cardioide de revolución}
	\begin{minipage}{0.20\linewidth}
		\centering
		\includegraphics[width=\linewidth]{Imgn/T_13a.pdf}
	\end{minipage}
	\hspace{1em}
	\begin{minipage}{0.20\linewidth}
		\centering
		\includegraphics[width=\linewidth]{Imgn/T_13d.pdf}
	\end{minipage}
	
	\vspace{0.3em}
	\fuentepropia
\end{minipage}
\vskip0.2cm

Las siguientes instrucciones permiten obtener la figura:


\vspace{0.6em}
\noindent
\begin{tcolorbox}[breakable,enhanced,title=\bf Instrucción en \verb"Mathematica", top=2pt,bottom=2pt,boxsep=0pt,before skip=2pt,after skip=0pt]	
\begin{Verbatim}[formatcom=\letraverbatim]
RegionPlot3D[(x^2+y^2+z^2-x)^2<x^2+y^2+z^2,{x,-1,2},
{y,-2,2},{z,-3/2,3/2},PlotPoints->30,Mesh->False,
BoxRatios->{1,1,1}]
\end{Verbatim}
\end{tcolorbox}
	
\vskip0.4cm	
En coordenadas esféricas el volumen del sólido está dado por

\begin{equation}
\label{ejercicio5}
\int_{0}^{\pi }\int_{0}^{2\pi }\int_{0}^{1+\cos \phi }\rho ^2 \sin \phi
d\rho d\theta d\phi =\frac{8}{3}\pi.
\end{equation}
	
Sustituyendo $\rho =1+\cos \phi,$ en las fórmulas que definen las coordenadas esféricas, se obtienen las siguientes expresiones: 
$x=u\left( 1+\cos \phi \right) \sin \phi \sin \theta ,$ $y=u\left( 1+\cos \phi \right) \sin \phi \cos \theta ,$  $z=u\left( 1+\cos \phi \right) \cos \phi .$
Por tanto, se define la función 
%\vskip0.2cm
$${\bf r}(u,\theta,\phi)=(1+\cos \phi) \left\langle u \sin
\phi \cos \theta ,u\sin \phi \sin \theta ,u \cos \phi \right\rangle $$
\vskip0.2cm
que parametriza el mencionado sólido con $u\in[0,1]$, $\phi \in[0,\pi]$, $\theta\in[0,2\pi]$. 
Con la función ${\bf r}$ se hallan los siguientes vectores
\noindent
{\small
	\begin{gather*}
		{\bf r}_{u} = (1+\cos \phi) \langle \cos \theta \sin \phi, \sin \theta \sin \phi, \cos \phi \rangle \\[6pt]
		{\bf r}_{\phi} = u \langle (\cos \phi + \cos (2\phi)) \cos \theta, (\cos \phi + \cos (2\phi)) \sin \theta, -(2\cos \phi + 1) \sin \phi \rangle \\[6pt]
		{\bf r}_{\theta} = \langle -u(\cos \phi + 1) \sin \theta \sin \phi, u(\cos \phi + 1) \cos \theta \sin \phi, 0 \rangle 
	\end{gather*}
}
\vskip0.2cm

Además {\footnotesize ${\bf r}_{u}\times {\bf r}_{\phi }= \left\langle  -u\left( \cos \phi +1\right) ^2 \sin \theta ,u\left( \cos \phi +1\right) ^2 \cos \theta ,0 \right\rangle.$}
%\linebreak
El elemento diferencial de volumen es {\footnotesize $\left\vert {\bf r}_{\theta }\cdot \left({\bf r}_{u}\times {\bf r}_{\phi }\right) \right\vert =u^{2}\left( \cos \phi +1\right) ^{3}\sin \phi ;$} con lo anterior, se plantea la integral que determina el volumen del sólido
%\vskip0.2cm 
{\small
$$\int_{0}^{1}\int_{0}^{2\pi }\int_{0}^{\pi }u^{2}\left( \cos \phi +1\right) ^{3}\sin \phi d\phi d\theta du= \frac{8}{3}\pi.$$
}
%\vskip0.2cm

\item El sólido acotado por la superficie del semicono $y=\sqrt{x^2+z^2}$ y la esfera $x^2+y^2+z^2=9,$
se muestran enseguida. \label{ejercicio7a}
\vskip0.2cm

\noindent
\begin{minipage}{\textwidth}
	\centering
	\captionof{figure}{\small $\sqrt{x^2+z^2}\leq y$ \& $x^2+y^2+z^2\leq 9$}
	\begin{minipage}{0.21\linewidth}
		\centering
		\includegraphics[width=\linewidth]{Imgn/T_14a.pdf}
	\end{minipage}
	\hspace{1em}
	\begin{minipage}{0.21\linewidth}
		\centering
		\includegraphics[width=\linewidth]{Imgn/T_14b.pdf}
	\end{minipage}
	
	\vspace{0.3em}
	\fuentepropia
\end{minipage}

\vskip0.2cm	
Estas figuras se obtienen con ayuda de las siguientes instrucciones:


\vspace{0.6em}
\noindent
\begin{tcolorbox}[breakable,enhanced,title=\bf Instrucción en \verb"Mathematica", top=2pt,bottom=2pt,boxsep=0pt,before skip=2pt,after skip=0pt]	
\begin{Verbatim}[formatcom=\letraverbatim]
RegionPlot3D[y>Sqrt[x^2+z^2]&&x^2+y^2+z^2<9,
{x,-2.5,2.5},{y,0,3.2},{z,-2.5,2.5},PlotPoints->90,
Mesh->False,BoxRatios->{1,1,1}]
\end{Verbatim}
\end{tcolorbox}

\vskip0.4cm
	
La intersección de las dos superficies está dada por la expresión $x^{2}+z^{2}=\frac{9}{2}.$ Además, sobre el sólido $y$ recorre los puntos que van desde el semicono hasta la esfera. Por tanto, el volumen se determina evaluando la siguiente integral:
$$\bigintsss_{-\frac{3}{\sqrt{2}}}^{\frac{3}{\sqrt{2}}}\bigintsss_{-\sqrt{\frac{9}{2}-x^{2}}}^{\sqrt{\frac{9}{2}-x^{2}}}\bigintsss_{\sqrt{x^{2}+z^{2}}}^{\sqrt{9-x^{2}-z^{2}}}dydzdx = 9 (2-\sqrt{2}) \pi.$$
%\vskip0.2cm

Utilizando coordenadas esféricas la integral correspondiente es
%\vskip0.2cm
\begin{equation} \label{ejercicio7}
\int_{0}^{2\pi }\int_{0}^{\frac{3}{\sqrt{2}}}\int_{r}^{\sqrt{9-r^{2}}}rdydrd\theta = 9\pi \left( 2-\sqrt{2}\right).
\end{equation}
%\vskip0.2cm

También se puede calcular el volumen considerando un punto sobre el cono {\small $P\left(u\cos v,u,u\sin v\right)$} y otro punto {\small $Q\left( u\cos v,\sqrt{9-u^{2}},u\sin v\right)$} sobre la esfera. De esta manera de forma una combinación convexa entre estos puntos y se parametriza el sólido. Esto es
% ${\bf r}(u,v,w) =(1-w) P +wQ  = \left\langle u\cos v,u+w\sqrt{9-u^{2}}-uw,u\sin v\right\rangle $
\vspace{-2mm}
{\small
\begin{eqnarray*} 
{\bf r}(u,v,w) &=&(1-w) P +wQ  = \left\langle u\cos v,u+w\sqrt{9-u^{2}}-uw,u\sin v\right\rangle,
\end{eqnarray*}
}
%\vskip0.2cm
con {\small $u\in \lbrack 0,\frac{3}{\sqrt{2}}],$} {\small $v\in \lbrack 0,2\pi ]$} y {\small $w\in \lbrack 0,1].$} De esta función se obtienen los siguien\-tes vectores: {\footnotesize ${\bf r}_{u} =\left\langle \cos v,1-w-\frac{uw}{\sqrt{9-u^{2}}},\sin v\right\rangle ,$} {\footnotesize ${\bf r}_{v} = \left\langle -u\sin v,0,u\cos v\right\rangle$} y {\footnotesize ${\bf r}_{w} = \left\langle 0,\sqrt{9-u^{2}}-u,0\right\rangle.$} El elemento diferencial de volumen es
{\footnotesize $\left\vert {\bf r}_{u}\cdot \left({\bf r}_{v}\times {\bf r}_{w}\right)\right\vert 
\\
=\left\vert u^{2}-u\sqrt{9-u^{2}}\right\vert.$} Con lo anterior, el volumen está dado por
\vskip0.2cm

{\small
$$\int_{0}^{1}\int_{0}^{2\pi }\int_{0}^{\frac{3}{\sqrt{2}}}\left\vert u^{2}-u \sqrt{9-u^{2}}\right\vert 
dudvdw= 9\pi \left( 2-\sqrt{2}\right).$$
}

\vskip0.2cm
	
\item El hiperboloide de un manto con ecuación $2x^2+3y^2=z^2+7$ y el elipsoide $2x^2+y^2+z^2=9$ acotan el sólido que se muestra a continuación. Esta región se presenta en detalle en la página \pageref{elipsoidehiperboloide}.

\vskip0.3cm
\noindent
\begin{minipage}{\textwidth}
	\centering
	\captionof{figure}{\small $2x^2+3y^2\leq z^2+7$ \& $2x^2+y^2+z^2\leq 9$}
	\begin{minipage}{0.23\linewidth}
		\centering
		\includegraphics[width=\linewidth]{Imgn/T_11a.pdf}
	\end{minipage}
	\hspace{1em}
	\begin{minipage}{0.23\linewidth}
		\centering
		\includegraphics[width=\linewidth]{Imgn/T_11b.pdf}
	\end{minipage}
	
	\vspace{0.3em}
	\fuentepropia
\end{minipage}

\vskip0.3cm
Esta figura se consigue ejecutando las siguientes instrucciones:
%\vskip0.2cm


\vspace{0.6em}
\noindent
\begin{tcolorbox}[breakable,enhanced,title=\bf Instrucción en \verb"Mathematica", top=2pt,bottom=2pt,boxsep=0pt,before skip=2pt,after skip=0pt]
\begin{Verbatim}[formatcom=\letraverbatim]
RegionPlot3D[2*x^2+3*y^2<z^2+7&&2*x^2+y^2+z^2<9,
{x,-2,2},{y,-2,2},{z,-3,3.2},PlotPoints->90,
Mesh->False];
\end{Verbatim}
\end{tcolorbox}

\vskip0.4cm	
En el plano $xy$, la proyección de la parte más delgada del hiperboloide posee ecuación $2x^{2}+3y^{2}=7.$ Eliminando $z$ de las ecuaciones iniciales se obtiene $x^{2}+y^{2}=4,$ que corresponde a la proyección en el plano $xy$ de la intersección de las superficies. 
\vskip0.4cm
Hay una región en la cual los l\'{\i}mites de integración para $z$ van desde el elipsoide hasta el mismo elipsoide, mientras que en la región restante, los límites para $z$ van desde el hiperboloide hasta el elipsoide. Por la simetr\'{\i}a de la figura se puede establecer que el volumen está dado por

\vspace{-4mm}
{\small
\begin{multline}
\bigintsss_{0}^{\sqrt{\frac{7}{2}}}\bigintsss_{0}^{\sqrt{\frac{7-2x^{2}}{3}}}\bigintsss_{0}^{\sqrt{9-2x^{2}-y^{2}}}8dzdydx \\ 
+ \bigintsss_{0}^{\sqrt{\frac{7}{2}}}\bigintsss_{\sqrt{\frac{7-2x^{2}}{3}}}^{\sqrt{4-x^{2}}}\bigintsss_{\sqrt{2x^{2}+3y^{2}-7}}^{\sqrt{9-2x^{2}-y^{2}}}8dzdydx\\
+\bigintsss_{\sqrt{\frac{7}{2}}}^2 \bigintsss_{0}^{\sqrt{4-x^{2}}}\bigintsss_{\sqrt{2x^{2}+3y^{2}-7}}^{\sqrt{9-2x^{2}-y^{2}}}8dzdydx \approx 51.8616 \label{ejercicio6}
\end{multline}
}
%\vskip0.2cm	
En el plano $xz$ la proyección del sólido genera una elipse y una hipérbola; mientras que la curva de intersección proyectada en el plano $xz$ se obtiene eliminando $y$ de las ecuaciones iniciales, esto se reduce a la fórmula $x^2+z^2=5$. El volumen del sólido se determina integrando primero con respecto a $y$ sobre dos regiones,
en una de ellas $y$ recorre puntos que van desde el elipsoide hasta la misma superficie y en la otra región $y$ se recorre los puntos desde el hiperboloide hasta el mismo. Con lo anterior, el planteamiento es el siguiente:
{\small
\begin{multline} 
\hskip-0.4cm \bigintsss_{0}^2 \bigintsss_{\sqrt{5-x^{2}}}^{\sqrt{9-2x^{2}}}\bigintsss_{0}^{\sqrt{9-2x^{2}-z^{2}}}8dydzdx + 
\bigintsss_{\sqrt{\frac{7}{2}}}^2 \bigintsss_{\sqrt{2x^{2}-7}}^{\sqrt{5-x^{2}}}\bigintsss_{0}^{\sqrt{\frac{z^{2}-2x^{2}+7}{3}}}8dydzdx  \\
+\bigintsss_{0}^{\sqrt{\frac{7}{2}}}\bigintsss_{0}^{\sqrt{5-x^{2}}}\bigintsss_{0}^{\sqrt{\frac{z^{2}-2x^{2}+7}{3}}}8dydzdx \approx 51.8616
\end{multline}
}
%\vskip0.2cm

La evaluación numérica de las anteriores integrales se obtuvo mediante la 
ejecución de las siguientes instrucciones:
%\vskip0.2cm


\vspace{0.6em}
\noindent
\begin{tcolorbox}[breakable,enhanced,title=\bf Instrucción en \verb"Mathematica", top=2pt,bottom=2pt,boxsep=0pt,before skip=2pt,after skip=0pt]
\begin{Verbatim}[formatcom=\letraverbatim]
NIntegrate[8,{x,0,Sqrt[7/2]},{y,0,Sqrt[(7-2*x^2)/3],
{z,0,Sqrt[9-2*x^2-y^2]}]+NIntegrate[8,{x,0,Sqrt[7/2]},
{y,Sqrt[(7-2*x^2)/3],Sqrt[4-x^2]},{z,Sqrt[2*x^2
+3*y^2-7],Sqrt[9-2*x^2-y^2]}]+NIntegrate[8,
{x,Sqrt[7/2],2},{y,0,Sqrt[4-x^2]},{z,Sqrt[2*x^2+
3*y^2-7],Sqrt[9-2*x^2-y^2]}]
\end{Verbatim}
\end{tcolorbox}
\vskip0.4cm	
%NIntegrate[8,{x,0,2},{z,Sqrt[5-x^2],Sqrt[9-2*x^2]},
%{y,0,Sqrt[9-2*x^2-z^2]}]+NIntegrate[8,{x,0,Sqrt[7/2]},	%{z,0,Sqrt[5-x^2]},{y,0,Sqrt[(z^2-2*x^2+7)/3]}]+
%NIntegrate[8,{x,Sqrt[7/2],2},{z,Sqrt[2*x^2-7],
%Sqrt[5-x^2]},{y,0,Sqrt[(z^2-2*x^2+7)/3]}]

\item El sólido acotado por las superficies del paraboloide $4y=x^2+z^2$ y la esfera $x^2+y^2+z^2=12$ 
se muestran a continuación: \label{ejercicio8a}

\vskip0.2cm
\noindent
\begin{minipage}{\textwidth}
	\centering
	\captionof{figure}{\small $4y\geq x^2+z^2$ \& $x^2+y^2+z^2\leq 12.$}
	\begin{minipage}{0.23\linewidth}
		\centering
		\includegraphics[width=\linewidth]{Imgn/T_15a.pdf}
	\end{minipage}
	\hspace{1em}
	\begin{minipage}{0.23\linewidth}
		\centering
		\includegraphics[width=\linewidth]{Imgn/T_15b.pdf}
	\end{minipage}
	
	\vspace{0.3em}
	\fuentepropia
\end{minipage}

\vskip0.2cm
Esta figura se consigue ejecutando las siguientes instrucciones:


\vspace{0.6em}
\noindent
\begin{tcolorbox}[breakable,enhanced,title=\bf Instrucción en \verb"Mathematica", top=2pt,bottom=2pt,boxsep=0pt,before skip=2pt,after skip=0pt]	
\begin{Verbatim}[formatcom=\letraverbatim]
RegionPlot3D[4*y>x^2+z^2&&x^2+y^2+z^2<12,{x,-3,3},
{y,0,3.7},{z,-3,3},PlotPoints->90,Mesh->False,
BoxRatios->{1,1,1}]
\end{Verbatim}
\end{tcolorbox}
\vskip0.4cm
	
Eliminando $x^2+z^2$ de las dos ecuaciones, obtenemos que $y^{2}+4y=12$, cuya solución es $y=2$ y $y=-6;$ por la ubicación de la región en el espacio consideramos el caso $y=2.$ Esto implica que en el plano $xz,$ la proyección del sólido está determinada por la ecuación $x^{2}+z^{2}=8;$ esto corresponde a un c\'{\i}rculo
centrado en el origen de radio $4$. Por tanto, utilizando coordenadas rectangulares e integrando primero con respecto a $y,$ la integral que determina el volumen es
\vskip0.2cm
$$\int _{-\sqrt{8}}^{\sqrt{8}}\int _{-\sqrt{8-x^2}}^{\sqrt{8-x^2}}\int _{\frac{1}{4}
\left(x^2+y^2\right)}^{\sqrt{12-x^2-y^2}}dzdydx= \frac{8}{3}\pi(6\sqrt{3}-5)$$
\vskip0.2cm
Esta integral se puede evaluar al ejecutar la siguiente instrucción:


\vspace{0.6em}
\noindent
\begin{tcolorbox}[breakable,enhanced,title=\bf Instrucción en \verb"Mathematica", top=2pt,bottom=2pt,boxsep=0pt,before skip=2pt,after skip=0pt]	
\begin{Verbatim}[formatcom=\letraverbatim]
Integrate[1,{x,-Sqrt[8],Sqrt[8]},{y,-Sqrt[8-x^2],
Sqrt[8-x^2]},{z,(x^2+y^2)/4,Sqrt[12-x^2-y^2]}]
\end{Verbatim}
\end{tcolorbox}

\vskip0.4cm	
Al hacer un cambio a coordenadas cilíndricas, el volumen está dado por
\vskip0.2cm
\begin{equation} \label{ejercicio8}
\int_0^{\sqrt{8}}\int_{0}^{2\pi}\int_{\frac{r^{2}}{4}}^{\sqrt{12-r^{2}}}rdyd\theta dr=\frac{8}{3}\pi\left(6\sqrt{3}-5\right)\!. 
\end{equation}
%\vskip0.2cm
Por otra parte, la ecuación $x^{2}+z^{2}=8$ sugiere usar la parametrización $x=\sqrt{8}u\cos v,$ $z=\sqrt{8}u\sin v$, en donde $u \in [0,1],$ $v \in [0,2\pi]$. Reemplazando $x$ y $z$ en la expresión de la esfera y del paraboloide se obtienen expresiones correspondientes para $y$,
es decir $y=\sqrt{12-x^{2}-z^{2}}= 2\sqrt{3-2u^{2}},$ $y=\frac{x^{2}+z^{2}}{4}= 2u^{2}.$ 
\vskip0.2cm
De esta manera, tendremos dos puntos sobre las superficies mencionadas. Al hacer una combinación convexa de estos puntos se parametriza el sólido
{\footnotesize
\[\left( 1-w\right) \left( \sqrt{8}u\cos v,2u^{2},\sqrt{8}u\sin v\right) \linebreak 	+w\left( \sqrt{8}u\cos v,2\sqrt{3-2u^{2}},\sqrt{8}u\sin v\right)\!;\]
}
por esta razón se considera la función
%\vskip0.2cm
$$ {\bf r}(u,v,w)= \langle 2\sqrt{2}u\cos v,2w\sqrt{3-2u^{2}}-2u^{2}w+2u^{2},2\sqrt{2}u\sin v \rangle,$$
\vskip0.2cm
con $u \in [0,1],$ $v \in [0,2\pi],$ $w \in [0,1].$ De aquí se obtienen los siguientes vectores:
{\small
\begin{eqnarray*}
{\bf r}_{u} &=& \left\langle 2\sqrt{2}\cos v,4u-4uw-\frac{4uw}{\sqrt{3-2u^{2}}},2\sqrt{2}\sin v \right\rangle  \\
{\bf r}_{v} &=& \left\langle  -2\sqrt{2}u\sin v,0,2\sqrt{2}u\cos v\right\rangle  \\
{\bf r}_{w} &=& \left\langle  0,2\sqrt{3-2u^{2}}-2u^{2},0\right\rangle.
\end{eqnarray*}
}
%\vskip0.2cm

Por lo descrito antes, el elemento diferencial de volumen está dado por $\left\vert {\bf r}_{w} \cdot({\bf r}_{u} \times {\bf r}_{v}) \right\vert = \left\vert 16u\left( u^{2}-\sqrt{3-2u^{2}}\right) \right\vert,$ que al integrarlo con los límites establecidos, produce el volumen del mencionado sólido:
%\vskip0.2cm 
$$\int_{0}^{1}\int_{0}^{1}\int_{0}^{2\pi }\left\vert 16u\left( u^{2}-\sqrt{3-2u^{2}}\right) \right\vert dvdudw= \frac{8}{3}\pi \left( 6\sqrt{3}-5\right)\!.$$
%\vskip0.3cm

\item Hallar el volumen del sólido acotado por el semicono con ecuación 
$x=\sqrt{y^2+z^2}$ y el paraboloide $y^2+z^2+4x=12.$

\vskip0.4cm
\noindent
\begin{minipage}{\textwidth}
	\centering
	\captionof{figure}{\small $r\leq x \leq 3-r^2/4$}
	\begin{minipage}{0.21\linewidth}
		\centering
		\includegraphics[width=\linewidth]{Imgn/T_17a.pdf}
	\end{minipage}
	\hspace{1em}
	\begin{minipage}{0.21\linewidth}
		\centering
		\includegraphics[width=\linewidth]{Imgn/T_17b.pdf}
	\end{minipage}
	
	\vspace{0.3em}
	\fuentepropia
\end{minipage}

\vskip0.2cm
Las siguientes instrucciones producen la figura:
%\vskip0.2cm	


\vspace{0.6em}
\noindent
\begin{tcolorbox}[breakable,enhanced,title=\bf Instrucción en \verb"Mathematica", top=2pt,bottom=2pt,boxsep=0pt,before skip=2pt,after skip=0pt]
\begin{Verbatim}[formatcom=\letraverbatim]
RegionPlot3D[x>Sqrt[y^2+z^2]&&y^2+z^2+4*x<12,
{x,0,3.25},{y,-2,2},{z,-2,2},PlotPoints->90,
Mesh->False]
\end{Verbatim}
\end{tcolorbox}


\vskip0.4cm	
Estas superficies se intersectan en una región del plano $yz$
determinada por la ecuación $y^{2}+z^{2}=4.$
Se debe tener presente que $x$ recorre los puntos desde el semicono hasta el paraboloide. También se presenta la integral correspondiente al cambio a coordenadas cilíndricas; por tanto, 
el volumen del sólido está dado por
%\vskip0.2cm 
\vspace{-2mm}
\begin{multline} \label{ejercicio10}
\int_{-2}^2 \int_{-\sqrt{4-y^2}}^{\sqrt{4-y^2}}\int_{\sqrt{y^2+z^{2}}}^{\frac{12-y^{2}-z^{2}}{4}}dxdzdy \\
=\int_{0}^2 \int_{0}^{2\pi }\int_{r}^{3-\frac{r^{2}}{4}}rdxd\theta dr= \frac{14}{3}\pi.
\end{multline}
%\vskip0.3cm
%$$\int_{0}^2 \int_{0}^{2\pi }\int_{r}^{3-\frac{r^{2}}{4}}rdxd\theta dr= \frac{14}{3}\pi .$$
\item Presentamos ahora el sólido acotado por el paraboloide $x+3=y^2+z^2$ 
y el paraboloide hiperbólico $y^2-z^2=9x-3.$

\vskip0.2cm
\noindent
\begin{minipage}{\textwidth}
	\centering
	\captionof{figure}{\small $3-y^2-z^2\leq x$ \& $y^2-z^2\leq 9x-3$}
	\begin{minipage}{0.20\linewidth}
		\centering
		\includegraphics[width=\linewidth]{Imgn/T_16a.pdf}
	\end{minipage}
	\hspace{1em}
	\begin{minipage}{0.20\linewidth}
		\centering
		\includegraphics[width=\linewidth]{Imgn/T_16b.pdf}
	\end{minipage}
	\hspace{1em}
	\begin{minipage}{0.20\linewidth}
		\centering
		\includegraphics[width=\linewidth]{Imgn/T_16c.pdf}
	\end{minipage}
	
	\vspace{0.3em}
	\fuentepropia
\end{minipage}

\vskip0.2cm
Ejecutando las siguientes instrucciones, se consigue la siguiente figura:	
\vskip0.2cm


%\vspace{0.6em}
\noindent
\begin{tcolorbox}[breakable,enhanced,title=\bf Instrucción en \verb"Mathematica", top=2pt,bottom=2pt,boxsep=0pt,before skip=2pt,after skip=0pt]
\begin{Verbatim}[formatcom=\letraverbatim]
RegionPlot3D[x+3>y^2+z^2&&y^2-z^2>9*x-3,{x,-3.5,1.5,
{y,-2,2},{z,-2.2,2.2},PlotPoints->90,Mesh->False,
BoxRatios->{1,1,1}]
\end{Verbatim}
\end{tcolorbox}

\vskip0.4cm
	
Eliminando $x$ de las ecuaciones del paraboloide y del paraboloide hiperbólico se obtiene $4y^{2}+5z^{2}=15,$ que corresponde a una elipse
en el plano $yz;$ esta ecuación define la región de integración en dicho plano. Con esta información se plantea la integral que representa el volumen del
sólido,
%\vskip0.2cm
$$\bigintss_{-\frac{\sqrt{15}}{2}}^{\frac{\sqrt{15}}{2}}\bigintss_{-\sqrt{3-\frac{4y^{2}}{5}}}^{\sqrt{3-\frac{4y^{2}}{5}}}\bigintss_{y^{2}+z^{2}-3}^{\frac{y^{2}-z^{2}+3}{9}}dxdzdy= \frac{5}{2}\sqrt{5}\pi.$$
\vskip0.2cm	
La región del plano $xz$ determinada por la ecuación $4y^{2}+5z^{2}=15,$ se puede representar mediante las expresiones $y=\frac{\sqrt{15}}{2}u\cos v,$ $z=\sqrt{3}u\sin v$, con $u \in [0,1]$ y $v \in [0,2\pi].$ La expresión correspondiente para $x$ en las dos superficies que acotan el sólido son: 
%\vskip0.2cm
\begin{eqnarray*}
x &=&y^{2}+z^{2}-3= \frac{3}{8}u^{2}\cos 2v+\frac{27}{8}u^{2}-3\\
x&=&\frac{y^{2}-z^{2}+3}{9}= \frac{3}{8}u^{2}\cos 2v+\frac{1}{24}u^{2}+\frac{1}{3}
\end{eqnarray*} 
%\vskip0.2cm
Con lo mencionado antes, 
se pueden considerar dos puntos: uno sobre el paraboloide {\small $Q\left( \frac{3}{8}u^{2}\cos 2v+\frac{27}{8}u^{2}-3,\frac{\sqrt{15}}{2}u\cos v,\sqrt{3}u\sin v\right)$} y otro sobre el paraboloide hiperbólico
{\small $P\left(\frac{3}{8}u^{2}\cos2v+\frac{1}{24}u^{2}+\frac{1}{3},\frac{\sqrt{15}}{2}u\cos v,\sqrt{3}u\sin v\right);$}
Con ellos se forma una combinación convexa {\small $(1-w) P +wQ,$ con $w \in [0,1],$} que al simplificarse define la función que parametriza el sólido,
%\vskip0.2cm 
{\small
\begin{multline*} 
{\bf r}(u,v,w)= \bigg\langle \frac{10}{3}w(u^{2}-1)+\frac{1}{8}u^{2}\left(3\cos 2v+\frac{1}{3}\right)+\frac{1}{3}, \\
\frac{\sqrt{15}}{2}u\cos v,\sqrt{3}u\sin v \bigg \rangle.
\end{multline*}
}
%\vskip0.2cm
Derivando parcialmente hallamos los siguientes vectores
%\vskip0.2cm
{\small
\begin{eqnarray*}
{\bf r}_{u} &=&\left\langle \frac{1}{12}u+\frac{3}{4}u\cos 2v+\frac{20}{3}uw,\frac{1}{2}\sqrt{15}\cos v,\sqrt{3}\sin v\right\rangle\\
{\bf r}_{v}&=&\left\langle -\frac{3}{4}u^{2}\sin 2v,-\frac{1}{2}\sqrt{15}u\sin v,\sqrt{3}u\cos v\right\rangle  \\
{\bf r}_{w}&=&\left\langle \frac{10}{3}u^{2}-\frac{10}{3},0,0\right\rangle  \\
{\bf r}_{w}\times {\bf r}_{u}&=& \frac{10}{3} \left\langle 0,-\sqrt{3}\left( u^{2}-1\right) \sin v,\frac{1}{2}\sqrt{15}\left( u^{2}-1\right) \cos v \right\rangle.
\end{eqnarray*}
}
\vskip0.2cm
Por último, se evalúa el producto vectorial triple, que determina el ele\-mento diferencial de volumen
$\left\vert {\bf r}_{v}\cdot \left({\bf r}_{w}\times {\bf r}_{u}\right) \right\vert= \left\vert 5\sqrt{5}u\left( u^{2}-1\right) \right\vert.$ Así, se integra con los límites establecidos y obtendremos el volumen deseado
\vskip0.2cm
\begin{equation}\label{ejercicio9}
\int_{0}^{1}\int_{0}^{1}\int_{0}^{2\pi }\left\vert 5\sqrt{5}u\left(u^{2}-1\right) \right\vert dvdudw= \frac{5}{2}\sqrt{5}\pi.
\end{equation}
\vskip0.2cm

\item Hallar el volumen del sólido que se obtiene de perforar la esfera centrada en el origen de radio $3$ con el cilindro con ecuación $x^2+y^2=2x+y.$

\vskip0.2cm
\noindent
\begin{minipage}{\textwidth}
	\centering
	\captionof{figure}{\small $2x+y \leq x^2+y^2$ \& $x^2+y^2+z^2\leq 9.$}
	\begin{minipage}{0.23\linewidth}
		\centering
		\includegraphics[width=\linewidth]{Imgn/T_19a.pdf}
	\end{minipage}
	\hspace{1em}
	\begin{minipage}{0.20\linewidth}
		\centering
		\includegraphics[width=\linewidth]{Imgn/T_19b.pdf}
	\end{minipage}
	
	\vspace{0.3em}
	\fuentepropia
\end{minipage}

\vskip0.2cm
Ejecutando las siguientes instrucciones se consigue la figura:
\vskip0.2cm	



%\vspace{0.6em}
\noindent
\begin{tcolorbox}[breakable,enhanced,title=\bf Instrucción en \verb"Mathematica", top=2pt,bottom=2pt,boxsep=0pt,before skip=2pt,after skip=0pt]
\begin{Verbatim}[formatcom=\letraverbatim]
RegionPlot3D[2*x+y<x^2+y^2<9&&x^2+y^2+z^2<=9,{x,-3,3},
{y,-3,3},{z,-3,3},PlotPoints->90,Mesh->False]
\end{Verbatim}
\end{tcolorbox}


\vskip0.2cm	
Al integrar en coordenadas rectangulares con el orden $dzdydx$ es necesario evaluar las cuatro integrales triples que a continuación se presentan:
%\vskip0.2cm
\begin{multline*} 
\bigintss_{-3}^{1-\frac{\sqrt{5}}{2}}\bigintss_{-\sqrt{9-x^{2}}}^{\sqrt{9-x^{2}}}\bigintss_{-\sqrt{9-x^{2}-y^{2}}}^{\sqrt{9-x^{2}-y^{2}}}dzdydx+ \\
\bigintsss_{1-\frac{\sqrt{5}}{2}}^{1+\frac{\sqrt{5}}{2}}\bigintsss_{-\sqrt{9-x^{2}}}^{\frac{1}{2}-\sqrt{\frac{5}{4}-\left( x-1\right) ^2 }}\bigintsss_{-\sqrt{9-x^{2}-y^{2}}}^{\sqrt{9-x^{2}-y^{2}}}dzdydx +\\
\bigintsss_{1-\frac{\sqrt{5}}{2}}^{1+\frac{\sqrt{5}}{2}}\bigintsss_{\frac{1}{2}+\sqrt{\frac{5}{4}-\left( x-1\right) ^2 }}^{\sqrt{9-x^{2}}}\bigintsss_{-\sqrt{9-x^{2}-y^{2}}}^{\sqrt{9-x^{2}-y^{2}}}dzdydx\\
+\bigintsss_{1+\frac{\sqrt{5}}{2}}^{3}\bigintsss_{-\sqrt{9-x^{2}}}^{\sqrt{9-x^{2}}}\bigintsss_{-\sqrt{9-x^{2}-y^{2}}}^{\sqrt{9-x^{2}-y^{2}}}dzdydx \approx 92.2262.
\end{multline*}
%\vskip0.2cm
	
Estas integrales se evaluaron numéricamente en \verb"Mathematica":
\vskip0.2cm


%\vspace{0.6em}
\noindent
\begin{tcolorbox}[breakable,enhanced,title=\bf Instrucción en \verb"Mathematica", top=2pt,bottom=2pt,boxsep=0pt,before skip=2pt,after skip=0pt]
\begin{Verbatim}[formatcom=\letraverbatim]
Integrate[1,{x,-3,1-Sqrt[5]/2},{y,-Sqrt[9-x^2],
Sqrt[9-x^2]},{z,-Sqrt[9-x^2-y^2],Sqrt[9-x^2-y^2]}]
+NIntegrate[1,{x,1-Sqrt[5]/2,1+Sqrt[5]/2},
{y,-Sqrt[9-x^2],1/2-Sqrt[5/4-(x-1)^2]},{z,
-Sqrt[9-x^2-y^2],Sqrt[9-x^2-y^2]}]+Integrate[1,
{x,1+Sqrt[5]/2,3},{y,-Sqrt[9-x^2],Sqrt[9-x^2]},
{z,-Sqrt[9-x^2-y^2],Sqrt[9-x^2-y^2]}]+NIntegrate[1,
{x,1-Sqrt[5]/2,1+Sqrt[5]/2},{y,1/2+Sqrt[5/4-(x-1)^2],
Sqrt[9-x^2]},{z,-Sqrt[9-x^2-y^2],Sqrt[9-x^2-y^2]}]
\end{Verbatim}
\end{tcolorbox}

\vskip0.4cm	
El volumen del sólido también puede calcularse como la diferencia entre el volumen de la esfera y
el volumen de la región acotada por la esfera e interior al cilindro, esto es,
%\vskip0.2cm
$$36\pi -\bigintss_{1-\frac{\sqrt{5}}{2}}^{1+\frac{\sqrt{5}}{2}}\bigintss_{\frac{1}{2}- \sqrt{\frac{5}{4}-(x-1)^2 }}^{\frac{1}{2}+\sqrt{\frac{5}{4}-(x-1) ^2 }}\bigintss_{-\sqrt{9-x^{2}-y^{2}}}^{\sqrt{9-x^{2}-y^{2}}}dzdydx \approx 92.2262.$$
%\vskip0.2cm
	
\item El sólido acotado por las figuras de las ecuaciones $z=4-y^{2},$ $z=4-x,$ $z=0,$ $x=0,$ 
se perfora por el cilindro $\left( z-2\right) ^{2}+y^{2}=1.$ Esta región se presentó en la sección \secref{solidohueco3} de la página \pageref{solidohueco3}; esta se muestra en las siguientes figuras.

\vskip0.2cm
\noindent
\begin{minipage}{\textwidth}
	\centering
	\captionof{figure}{\small $1\leq (z-2)^{2}+y^{2}$ \& $0\leq z\leq 4-y^2$}
	\begin{minipage}{0.18\linewidth}
		\centering
		\vspace{-4mm}
		\includegraphics[width=\linewidth]{Imgn/5b.pdf}
	\end{minipage}
	\hspace{1em}
	\begin{minipage}{0.18\linewidth}
		\centering
		\vspace{-4mm}
		\includegraphics[width=\linewidth]{Imgn/5a.pdf}
	\end{minipage}
	\hspace{1em}
	\begin{minipage}{0.19\linewidth}
		\centering
		\vspace{-4mm}
		\includegraphics[width=\linewidth]{Imgn/5c.pdf}
	\end{minipage}
	
	\vspace{0.3em}
	\fuentepropia
\end{minipage}
\vskip0.2cm

El volumen puede obtenerse como la diferencia entre el volumen del sólido sin perforar y el volumen de la parte cilíndrica. Esto es,
\vspace{-2mm}
$$ \int _{-2}^2\int _0^{4-y^2}\int _0^{4-z}dxdzdy-\int _{-1}^1\int _{2-\sqrt{1-y^2}}^{2+\sqrt{1-y^2}}\int _0^{4-z}dxdzdy = \frac{128}{5}-2\pi.$$
%\vskip0.2cm


%\vspace{0.6em}
\noindent
\begin{tcolorbox}[breakable,enhanced,title=\bf Instrucción en \verb"Mathematica", top=2pt,bottom=2pt,boxsep=0pt,before skip=2pt,after skip=0pt]
\begin{Verbatim}[formatcom=\letraverbatim]
Integrate[1,{y,-2,2},{z,0,4-y^2},{x,0,4-z}]
-Integrate[1,{y,-1,1},{z,2-Sqrt[1-y^2],2+
Sqrt[1-y^2]},{x,0,4-z}]
\end{Verbatim}
\end{tcolorbox}

\vskip0.4cm

El mencionado sólido se parametrizó mediante la expresión (\ref{solidohueco2}).
Con las siguientes instrucciones se define la parametrización, se halla el elemento diferencial de 
volumen y se evalúan las integrales con los límites establecidos:
\vskip0.2cm

%\vspace{0.6em}
\noindent
\begin{tcolorbox}[breakable,enhanced,title=\bf Instrucción en \verb"Mathematica", top=2pt,bottom=2pt,boxsep=0pt,before skip=2pt,after skip=0pt]	
\begin{Verbatim}[formatcom=\letraverbatim]
r1={0,(1-u)*Cos[t]+u*(2-4*t/Pi),(1-u)*(2+Sin[t])
+16*t*u/Pi*(1-t/Pi)};
r2={0,(1-u)*Cos[t]+u*(4*(t-Pi)/Pi-2),(1-u)*(2+Sin[t])};
r3={4-(1-u)*(2+Sin[t])-16*t*u/Pi*(1-t/Pi),(1-u)*Cos[t]
+u*(2-4*t/Pi),(1-u)*(2+Sin[t])+16*t*u/Pi*(1-t/Pi)};
r4={4-((1-u)*(2+Sin[t])),(1-u)*Cos[t]
+u*(4*(t-Pi)/Pi-2),(1-u)*(2+Sin[t])};
{R1,R2}={(1-w)*r1+w*r3,(1-w)*r2+w*r4};
Integrate[Dot[D[R1,w],Cross[D[R1,u],D[R1,t]]],
{t,0,Pi},{u,0,1},{w,0,1}]+Integrate[Dot[D[R2,w],
Cross[D[R2,u],D[R2,t]]],{t,Pi,2*Pi},{u,0,1},{w,0,1}]
\end{Verbatim}
\end{tcolorbox}

\end{enumerate}
\vskip0.4cm

\begin{ejer}
Determine el volumen del sólido acotado por las figuras de las expresiones señaladas.
Hágalo de diversas maneras, incluyendo la parametrización de la región. Se agrega la figura y las instrucciones para conseguirlo.
\vskip0.2cm

\begin{enumerate}[leftmargin=1.2em, itemindent=0pt,labelwidth=\labelwidth,labelsep=0.5em]
\item $x^2+y^2=4+z$, $4z=(4-x^2-y^2)e^{1-x^2-y^2}$. 
\normalfont
\vskip0.2cm


%\vspace{0.6em}
\noindent
\begin{tcolorbox}[breakable,enhanced,title=\bf Instrucción en \verb"Mathematica", top=2pt,bottom=2pt,boxsep=0pt,before skip=2pt,after skip=0pt]
\begin{Verbatim}[formatcom=\letraverbatim]
RegionPlot3D[x^2+y^2-4<=z&&z<=(4-x^2-y^2)*Exp[1
-x^2-y^2],{x,-2,2},{y,-2,2},{z,-4,12},
PlotPoints->90,Mesh->False]
\end{Verbatim}
\end{tcolorbox}

\vskip0.4cm

\noindent
\begin{minipage}{\textwidth}
	\centering
	\captionof{figure}{\small $x^2+y^2-4\leq z \leq (4-x^2-y^2)e^{1-x^2-y^2}$}
	\begin{minipage}{0.22\linewidth}
		\centering
		\includegraphics[width=\linewidth]{Imgn/T_42a.pdf}
	\end{minipage}
	\hspace{1em}
	\begin{minipage}{0.22\linewidth}
		\centering
		\includegraphics[width=\linewidth]{Imgn/T_42b.pdf}
	\end{minipage}
	\hspace{1em}
	\begin{minipage}{0.22\linewidth}
		\centering
		\includegraphics[width=\linewidth]{Imgn/T_42c.pdf}
	\end{minipage}
	
	\vspace{0.3em}
	\fuentepropia
\end{minipage}
\vskip0.2cm
		
\item $z=x^2+y^2,$ $4-z=3(x^2+y^2)^3.$
\vskip0.2cm 
\normalfont


%\vspace{0.6em}
\noindent
\begin{tcolorbox}[breakable,enhanced,title=\bf Instrucción en \verb"Mathematica", top=2pt,bottom=2pt,boxsep=0pt,before skip=2pt,after skip=0pt]
\begin{Verbatim}[formatcom=\letraverbatim]
RegionPlot3D[x^2+y^2<=z<=4-3*(x^2+y^2)^3&&x^2
+y^2<=1,{x,-1,1},{y,-1,1},{z,0,4},Mesh->False,
PlotPoints->90,PlotStyle->RGBColor[0.3,0.6,0.3]]
\end{Verbatim}
\end{tcolorbox}

\vskip0.2cm

\noindent
\begin{minipage}{\textwidth}
	\centering
	\captionof{figure}{\small $x^2+y^2\leq z\leq 4-3(x^2+y^2)^3$}
	\begin{minipage}{0.22\linewidth}
		\centering
		\includegraphics[width=\linewidth]{Imgn/T_40a.pdf}
	\end{minipage}
	\hspace{1em}
	\begin{minipage}{0.22\linewidth}
		\centering
		\includegraphics[width=\linewidth]{Imgn/T_40b.pdf}
	\end{minipage}
	\hspace{1em}
	\begin{minipage}{0.22\linewidth}
		\centering
		\includegraphics[width=\linewidth]{Imgn/T_40c.pdf}
	\end{minipage}
	
	\vspace{0.3em}
	\fuentepropia
\end{minipage}
\vskip0.2cm
		
\item $z=4-x^2-y^2,$ $z=3(x^2+y^2)^6.$ 

\vskip0.2cm
\noindent
\begin{minipage}{\textwidth}
	\centering
	\captionof{figure}{\small $3(x^2+y^2)^6 \leq z \leq 4-x^2-y^2$} 
	\begin{minipage}{0.20\linewidth}
		\centering
		\includegraphics[width=\linewidth]{Imgn/T_41a.pdf}
	\end{minipage}
	\hspace{1em}
	\begin{minipage}{0.20\linewidth}
		\centering
		\includegraphics[width=\linewidth]{Imgn/T_41b.pdf}
	\end{minipage}
	
	\vspace{0.3em}
	\fuentepropia
\end{minipage}
\vskip0.2cm
 
\normalfont


%\vspace{0.6em}
\noindent
\begin{tcolorbox}[breakable,enhanced,title=\bf Instrucción en \verb"Mathematica", top=2pt,bottom=2pt,boxsep=0pt,before skip=2pt,after skip=0pt]
\begin{Verbatim}[formatcom=\letraverbatim]
RegionPlot3D[3*(x^2+y^2)^6<=z<=4-x^2-y^2&&x^2
+y^2<=1,{x,-1,1},{y,-1,1},{z,0,4},PlotPoints->90,
Mesh->False,PlotStyle->RGBColor[0.6,0.2,0.2]]
\end{Verbatim}
\end{tcolorbox}

\vskip0.2cm
\end{enumerate}
\end{ejer}
